% આ ભાગનો પૂર્ણ સંપાદનીય પાઠ છે. શૈલી, ફોન્ટ, ચિત્રો અને પ્રકરણસંદર્ભ માટે સંપૂર્ણ સ્રોત ZIP તથા reader/full-reader.tex જુઓ. આ એકલી સ્વતંત્ર કમ્પાઇલ થતી મુખ્ય ફાઇલ નથી.

% BEGIN OLP-0407
\olpart{nml}{સામાન્ય મોડલ તર્કશાસ્ત્રો}
\phantomsection\label{guunit:OLP-0407}


\begin{editorial}
  આ ભાગ સામાન્ય મોડલ તર્કશાસ્ત્રોના અધિસિદ્ધાંતને આવરી લે છે.
  હાલમાં તેમાં મૂળભૂત મોડલ તર્કશાસ્ત્ર માટેના શાસ્ત્રીય
  અનુરૂપતા સિદ્ધાંત અંગે આલ્ડો એન્ટોનેલીની નોંધો છે.
\end{editorial}

\begingroup
\makeatletter\def\ol@sectioncs{section}\makeatother

% BEGIN OLP-0408
\olchapter{nml}{syn}{વાક્યરચના અને અર્થવિજ્ઞાન}
\olchapterid{nml}{syn}
\phantomsection\label{guunit:OLP-0408}


\begingroup
\makeatletter\def\ol@sectioncs{section}\makeatother

% BEGIN OLP-0409
\olfileid{nml}{syn}{int}

\olsection{પરિચય}
\phantomsection\label{guunit:OLP-0409}


મોડલ તર્કશાસ્ત્ર \emph{મોડલ વિધાનો} અને તેમની વચ્ચેના ફલિતતા
સંબંધોનો અભ્યાસ કરે છે. મોડલ વિધાનોના દાખલા આ છે:
\begin{enumerate}
\item $2+2=4$ હોવું અનિવાર્ય છે.
\item આવતી કાલે વરસાદ પડવો શક્ય છે તે અનિવાર્ય છે.
\item જો~$!A$ શક્ય હોવું અનિવાર્ય હોય, તો~$!A$
  શક્ય છે.
\end{enumerate}
શક્યતા અને અનિવાર્યતા જ એકમાત્ર મોડલતાઓ નથી:
બીજા એકસ્થાની સંયોજકોને પણ મોડલતાઓ ગણવામાં આવે છે.
દાખલા તરીકે, ``$!A$ હોવું જોઈએ,'' ``$!A$ થશે,''
``ડાના જાણે છે કે~$!A$,'' અથવા ``ડાના માને છે કે~$!A$.''

મોડલ તર્કશાસ્ત્ર સૌપ્રથમ એરિસ્ટોટલના \emph{De
  Interpretatione}માં દેખાય છે. અનિવાર્યતામાંથી શક્યતા
ફલિત થાય છે, પરંતુ ઊલટું નહીં; શક્યતા અને અનિવાર્યતા
એકબીજાની મદદથી વ્યાખ્યાયિત થઈ શકે છે; જો $!A \land !B$
શક્યપણે સત્ય હોય તો $!A$ અને $!B$ બંને શક્યપણે સત્ય છે,
પરંતુ ઊલટું નહીં; અને જો $!A \to !B$ અનિવાર્ય હોય,
તો $!A$ અનિવાર્ય હોવાથી~$!B$ પણ અનિવાર્ય છે---આ બાબતો
સૌપ્રથમ તેમણે નોંધેલી.

મોડલ તર્કશાસ્ત્ર પ્રત્યેનો પ્રથમ આધુનિક અભિગમ સી.~આઇ.~લુઇસના
કાર્યમાં જોવા મળે છે; તેનો પરિપાક લુઇસ અને લેંગફોર્ડના
\emph{Symbolic Logic} (1932)માં થયો. લુઇસ અને લેંગફોર્ડને
વસ્તુગત શરતી વડે પ્રેરણનું નિરૂપણ સંતોષકારક લાગતું ન હતું:
``$!A$માંથી $!B$ ફલિત થાય છે'' તેના માટે $!A \lif !B$
યોગ્ય વિકલ્પ નથી. તેથી તેમણે પ્રેરણને ``અનિવાર્યપણે,
જો $!A$ તો $!B$'' તરીકે લક્ષણિત કરવાનો પ્રસ્તાવ કર્યો,
જે $!A \strictif !B$ વડે દર્શાવાય છે. જુદા ગુણધર્મોને
વ્યવસ્થિત કરતાં લુઇસે પાંચ જુદી મોડલ પ્રણાલીઓ ઓળખી:
\Log{S1}, \ldots, \Log{S4}, \Log{S5}; છેલ્લી બેનો આજે પણ
ઉપયોગ થાય છે.

લુઇસ અને લેંગફોર્ડનો અભિગમ સંપૂર્ણપણે \emph{વાક્યરચનાત્મક}
હતો: તેમણે યોગ્ય સ્વયંસિદ્ધિઓ અને નિયમો ઓળખ્યાં અને
તેમના વડે શું સાબિત થઈ શકે તેની તપાસ કરી. અર્થવિજ્ઞાનલક્ષી
અભિગમ લાંબા સમય સુધી હાથ ન લાગ્યો. રુડોલ્ફ કાર્નાપે
\emph{Meaning and Necessity} (1947)માં \emph{અવસ્થા-વર્ણન}
વાપરી પ્રથમ પ્રયત્ન કર્યો; એટલે પરમાણુ વાક્યોનો એવો સંગ્રહ,
જે તે અવસ્થા-વર્ણનમાં ``સત્ય'' હોય. સત્યની વ્યાખ્યા કોઈપણ
વાક્ય~$!A$ સુધી વિસ્તારીને કાર્નાપે $!A$ને \emph{અનિવાર્યપણે
સત્ય} ત્યારે કહ્યું જ્યારે તે બધાં અવસ્થા-વર્ણનોમાં સત્ય હોય.
કાર્નાપનો અભિગમ \emph{પુનરાવર્તિત} મોડલતાઓને સંભાળી શકતો
ન હતો: ``શક્યપણે અનિવાર્યપણે \ldots શક્યપણે $!A$''
સ્વરૂપનાં વાક્યો હંમેશા સૌથી અંદરની મોડલતામાં સંકોચાઈ જાય છે.

મોડલ અર્થવિજ્ઞાનમાં મોટો સીમાચિહ્ન સૉલ ક્રિપ્કેનો લેખ
``A Completeness Theorem in Modal Logic'' (JSL 1959) હતો.
ક્રિપ્કેએ લાઇબ્નિત્સના એ વિચારને આધાર બનાવ્યો કે કોઈ વિધાન
``બધાં શક્ય જગતોમાં'' સત્ય હોય તો તે અનિવાર્યપણે સત્ય છે.
પરંતુ આ વિચારને પણ કાર્નાપના અભિગમ જેવી જ મર્યાદા છે:
કોઈ જગત~$w$ (અથવા અવસ્થા-વર્ણન~$s$)માં વિધાનનું સત્ય
$w$ પર જરાય આધાર રાખતું નથી. તેથી ક્રિપ્કેએ માન્યું કે
જગતો વચ્ચે \emph{સુલભતા સંબંધ} $R$ છે, અને
``અનિવાર્યપણે $!A$'' સ્વરૂપનું વિધાન જગત~$w$માં ત્યારે અને
માત્ર ત્યારે જ સત્ય છે જ્યારે $w$માંથી \emph{સુલભ}
બધાં જગતો~$w'$માં $!A$ સત્ય હોય. આ અભિગમનું કોઈ સ્વરૂપ
આપતા અર્થવિજ્ઞાનને ક્રિપ્કે અર્થવિજ્ઞાન કહે છે. તેનાથી
મોડલ તર્કશાસ્ત્રોના ઝડપી વિકાસને વેગ મળ્યો.

ક્રિપ્કે અર્થવિજ્ઞાન મુજબ અર્થઘટન કરતાં મોડલ તર્કશાસ્ત્ર
\emph{સંબંધાત્મક સંરચનાઓ} ``અંદરથી'' કેવી દેખાય છે તે બતાવે છે.
સંબંધાત્મક સંરચના એટલે દ્વિસ્થાની સંબંધ ધરાવતો ગણ
(દાખલા તરીકે, વર્ગના વિદ્યાર્થીઓનો ગણ, જેને તેમના
સામાજિક સુરક્ષા ક્રમાંક મુજબ ક્રમબદ્ધ કર્યો હોય).
પરંતુ સંબંધાત્મક સંરચનાઓ અનેક ક્ષેત્રોમાં મળે છે:
જગતની અવસ્થાઓની સાપેક્ષ શક્યતા ઉપરાંત, કોઈ વ્યક્તિની
જ્ઞાનાત્મક અવસ્થાઓ જ્ઞાનાત્મક શક્યતા વડે જોડાયેલી હોય
અથવા કોઈ ગતિશીલ તંત્રની અવસ્થાઓ અવસ્થા-સંક્રમણ વડે
જોડાયેલી હોય, વગેરે. આ બધાંનું નિદર્શન મોડલ તર્કશાસ્ત્રથી
થઈ શકે છે: પ્રથમથી સામાન્ય સત્યલક્ષી મોડલ તર્કશાસ્ત્ર,
અને બીજાંથી જ્ઞાનમીમાંસક તર્કશાસ્ત્ર, ગતિશીલ તર્કશાસ્ત્ર
વગેરે મળે છે.

અહીં અમે એક વિશિષ્ટ દૃષ્ટિકોણ પર ધ્યાન આપીએ છીએ, જેને
મોડલ તર્કશાસ્ત્રીઓ ``અનુરૂપતા સિદ્ધાંત'' કહે છે. ક્રિપ્કેની
આરંભની મહત્વની શોધોમાંની એક એ છે કે સુલભતા સંબંધ~$R$ના
ઘણા ગુણધર્મો (જેમ કે પરંપરિતતા, સમમિતિ વગેરે) યોગ્ય
``મોડલ યોજનાઓ'' વડે \emph{મોડલ ભાષામાં જ} લક્ષણિત કરી
શકાય છે. દાખલા તરીકે, મોડલ તર્કશાસ્ત્રીઓ કહે છે કે
$R$ની સ્વવાચકતા ``જો અનિવાર્યપણે~$!A$, તો~$!A$''
યોજનાને ``અનુરૂપ'' છે. અમે મુખ્યત્વે \Ax{D}, \Ax{T},
\Ax{B}, \Ax{4} અને~\Ax{5} યોજનાઓના સંયોજનથી મળતી
મોડલ તર્કશાસ્ત્રની કેટલીક શાસ્ત્રીય પ્રણાલીઓ (જેમ કે
\Log{S4} અને \Log{S5})ના અનુરૂપતા સિદ્ધાંતની તપાસ કરીએ છીએ.
% END OLP-0409
\endgroup

\begingroup
\makeatletter\def\ol@sectioncs{section}\makeatother

% BEGIN OLP-0410
\olfileid{nml}{syn}{lan}

\olsection{મૂળભૂત મોડલ તર્કશાસ્ત્રની ભાષા}
\phantomsection\label{guunit:OLP-0410}


\begin{defn}
મોડલ તર્કશાસ્ત્રની મૂળભૂત ભાષામાં નીચેની વસ્તુઓ હોય છે:
\begin{enumerate}
  \item અસત્ય માટેનો વિધાનાત્મક અચલ~$\lfalse$.
  \item સત્ય માટેનો વિધાનાત્મક અચલ~$\ltrue$.
  \item પ્રતિજ્ઞાપન ચલોનો ગણનીય ગણ: $\Obj
    p_0$, $\Obj p_1$, $\Obj p_2$, \dots
  \item વિધાનાત્મક સંયોજકો: \startycommalist
  \ycomma $\lnot$ (નિષેધ)%
  \ycomma $\land$ (સંયોજન)%
  \ycomma $\lor$ (વિયોજન)%
  \ycomma $\lif$ (પ્રેરણ)%
  \ycomma $\liff$ (દ્વિમુખી પ્રેરણ).
  \item મોડલ કારક $\Box$.
  \item મોડલ કારક $\Diamond$.
\end{enumerate}
\end{defn}

\begin{defn}
મૂળભૂત મોડલ ભાષાનાં \emph{સૂત્રો} નીચે પ્રમાણે
આગમનાત્મક રીતે વ્યાખ્યાયિત થાય છે:
\begin{enumerate}
\item $\lfalse$ એ પરમાણુ સૂત્ર છે.

\item $\ltrue$ એ પરમાણુ સૂત્ર છે.

\item દરેક પ્રતિજ્ઞાપન ચલ $\Obj p_i$ એ (પરમાણુ) સૂત્ર છે.

\item જો $!A$ એ સૂત્ર હોય, તો $\lnot !A$ પણ
  સૂત્ર છે.

\item જો $!A$ અને $!B$ એ સૂત્રો હોય, તો $(!A \land
  !B)$ એ સૂત્ર છે.

\item જો $!A$ અને $!B$ એ સૂત્રો હોય, તો $(!A \lor !B)$
  એ સૂત્ર છે.

\item જો $!A$ અને $!B$ એ સૂત્રો હોય, તો $(!A \lif !B)$
  એ સૂત્ર છે.

\item જો $!A$ અને $!B$ એ સૂત્રો હોય, તો $(!A \liff !B)$
  એ સૂત્ર છે.

\item જો $!A$ એ સૂત્ર હોય, તો $\Box !A$ પણ
  સૂત્ર છે.

\item જો $!A$ એ સૂત્ર હોય, તો $\Diamond !A$ પણ
  સૂત્ર છે.

\item બીજું કંઈ સૂત્ર નથી.
\end{enumerate}
\end{defn}


\sourcecorrection{OLMD-001}{સ્થિર અંગ્રેજી મૂળમાં
\texttt{defIf}ની પ્રથમ શાખામાં $\lnot !A \lor !B)$
લખતાં બંધ કૌંસ માટે ખુલતું કૌંસ નથી. અહીં વધારાનું
બંધ કૌંસ દૂર કરીને $\lnot !A \lor !B$ રાખ્યું છે.}

જો સૂત્ર~$!A$માં $\Box$ કે $\Diamond$ ન હોય,
તો તેને \emph{મોડલ-મુક્ત} કહીએ છીએ.
% END OLP-0410
\endgroup

\begingroup
\makeatletter\def\ol@sectioncs{section}\makeatother

% BEGIN OLP-0411
\olfileid{nml}{syn}{sub}

\olsection{એકસાથે પ્રતિસ્થાપન}
\phantomsection\label{guunit:OLP-0411}


સૂત્ર~$!A$માં પ્રતિજ્ઞાપન ચલનાં બધાં આવર્તનોને
કોઈ બીજા સૂત્રથી બદલીને મળતું પરિણામ એ~$!A$નો
\emph{દૃષ્ટાંત} છે. પ્રામાણ્ય અને નિષ્પન્નક્ષમતા બંનેની ચર્ચામાં
સૂત્રોના આવા દૃષ્ટાંતોનો વારંવાર ઉલ્લેખ આવશે. તેથી આ ખ્યાલની
ચોક્કસ વ્યાખ્યા આપવી ઉપયોગી છે.

\begin{defn}\ollabel{def:subst-inst}
  ધારો કે $!A$ એ મોડલ સૂત્ર છે જેમાંનાં બધાં પ્રતિજ્ઞાપન ચલો $p_1$, \dots, $p_n$ પૈકીનાં છે, અને $!D_1$, \dots,
  $!D_n$ પણ મોડલ સૂત્રો છે. ત્યારે
  $\SSubst{!A}{\subst{!D_1}{p_1}, \dots, \subst{!D_n}{p_n}}$ને
  $!A$માં દરેક $p_i$ માટે એકસાથે $!D_i$નું પ્રતિસ્થાપન
  કરવાથી મળતું પરિણામ વ્યાખ્યાયિત કરીએ છીએ. ચોક્કસ રીતે,
  આ વ્યાખ્યા~$!A$ પર અનુમાન વડે અપાય છે:
  \begin{enumerate}
    \item \indcase{!A}{\lfalse}{%
        $\SSubst{\indfrm}{\subst{!D_1}{p_1}, \dots,
          \subst{!D_n}{p_n}}$ એ $\lfalse$ છે.}
    \item \indcase{!A}{\ltrue}{%
        $\SSubst{\indfrm}{\subst{!D_1}{p_1}, \dots,
          \subst{!D_n}{p_n}}$ એ $\ltrue$ છે.}
    \item \indcase{!A}{q}{$\SSubst{\indfrm}{\subst{!D_1}{p_1}, \dots,
        \subst{!D_n}{p_n}}$ એ $q$ છે, જો $i
      = 1$, \dots,~$n$ માટે $q \not\ident p_i$ હોય.}
    \item \indcase{!A}{p_i}{$\SSubst{\indfrm}{\subst{!D_1}{p_1},
        \dots, \subst{!D_n}{p_n}}$ એ $!D_i$ છે.}
    \item \indcase{!A}{\lnot !B}{%
        $\SSubst{\indfrm}{\subst{!D_1}{p_1}, \dots, \subst{!D_n}{p_n}}$ એ
        $\lnot \SSubst{!B}{\subst{!D_1}{p_1}, \dots, \subst{!D_n}{p_n}}$ છે.}
    \item \indcase{!A}{(!B \land
        !C)}{$\SSubst{\indfrm}{\subst{!D_1}{p_1}, \dots,
          \subst{!D_n}{p_n}}$ એ
        \[(\SSubst{!B}{\subst{!D_1}{p_1}, \dots, \subst{!D_n}{p_n}} \land
          \SSubst{!C}{\subst{!D_1}{p_1}, \dots, \subst{!D_n}{p_n}}).\] છે.}
    \item \indcase{!A}{(!B \lor
          !C)}{$\SSubst{\indfrm}{\subst{!D_1}{p_1}, \dots,
          \subst{!D_n}{p_n}}$ એ
        \[(\SSubst{!B}{\subst{!D_1}{p_1}, \dots, \subst{!D_n}{p_n}} \lor
          \SSubst{!C}{\subst{!D_1}{p_1}, \dots, \subst{!D_n}{p_n}}).\] છે.}
    \item \indcase{!A}{(!B \lif
          !C)}{$\SSubst{\indfrm}{\subst{!D_1}{p_1}, \dots,
          \subst{!D_n}{p_n}}$ એ
        \[(\SSubst{!B}{\subst{!D_1}{p_1}, \dots, \subst{!D_n}{p_n}} \lif
          \SSubst{!C}{\subst{!D_1}{p_1}, \dots, \subst{!D_n}{p_n}}).\] છે.}
    \item \indcase{!A}{(!B \liff
          !C)}{$\SSubst{\indfrm}{\subst{!D_1}{p_1}, \dots,
          \subst{!D_n}{p_n}}$ એ
        \[(\SSubst{!B}{\subst{!D_1}{p_1}, \dots, \subst{!D_n}{p_n}} \liff
          \SSubst{!C}{\subst{!D_1}{p_1}, \dots, \subst{!D_n}{p_n}}).\] છે.}
    \item \indcase{!A}{\Box !B}{%
        $\SSubst{\indfrm}{\subst{!D_1}{p_1}, \dots, \subst{!D_n}{p_n}}$ એ
        $\Box
        \SSubst{!B}{\subst{!D_1}{p_1}, \dots, \subst{!D_n}{p_n}}.$ છે}
    \item \indcase{!A}{\Diamond !B}{%
        $\SSubst{\indfrm}{\subst{!D_1}{p_1}, \dots, \subst{!D_n}{p_n}}$ એ
        $\Diamond
        \SSubst{!B}{\subst{!D_1}{p_1}, \dots, \subst{!D_n}{p_n}}.$ છે}
  \end{enumerate}
  સૂત્ર $\SSubst{!A}{\subst{!D_1}{p_1}, \dots,
    \subst{!D_n}{p_n}}$ને $!A$નો \emph{પ્રતિસ્થાપન દૃષ્ટાંત}
  કહે છે.
\end{defn}
\sourcecorrection{OLMD-002}{સ્થિર અંગ્રેજી મૂળ
દ્વિમુખી પ્રેરણના આગમન-કિસ્સાને ભૂલથી
\texttt{prvIf} ટૅગ આપે છે, જ્યારે તેની સૂત્રરચના
$!B \liff !C$ છે. મૂળ ભાષાની રચના-વ્યાખ્યા
સાથે સુસંગત \texttt{prvIff} ટૅગ અહીં રાખ્યો છે.}
\sourcecorrection{OLMD-003}{સ્થિર અંગ્રેજી મૂળ
$\Box !B$ના આગમન-કિસ્સાને અશરતી \texttt{item}
રૂપે રાખે છે, જ્યારે મૂળ ભાષામાં $\Box$નો પ્રાથમિક
કારક \texttt{prvBox} ટૅગથી પસંદ થાય છે. તેની સાથે
સુસંગત રીતે આ કિસ્સો \texttt{prvBox}થી શરતિત કર્યો છે.}

\begin{ex}
  માનો કે $!A$ એ $p_1 \lif \Box(p_1 \land p_2)$ છે,
  $!D_1$ એ $\Diamond(p_2 \lif p_3)$ છે અને $!D_2$
  એ $\lnot\Box p_1$ છે. ત્યારે
  $\SSubst{!A}{\subst{!D_1}{p_1}, \subst{!D_2}{p_2}}$ આ છે:
  \begin{align*}
    \Diamond(p_2 \lif p_3) &
    \lif \Box(\Diamond(p_2 \lif p_3) \land \lnot\Box p_1)
    \intertext{જ્યારે $\SSubst{!A}{\subst{!D_2}{p_1}, \subst{!D_1}{p_2}}$ આ છે:}
    \lnot\Box p_1 &
    \lif \Box(\lnot\Box p_1 \land \Diamond(p_2 \lif p_3))
    \intertext{નોંધો કે એકસાથે પ્રતિસ્થાપન સામાન્ય રીતે
      ક્રમિક પ્રતિસ્થાપન જેવું નથી. દાખલા તરીકે, ઉપરના
      $\SSubst{!A}{\subst{!D_1}{p_1}, \subst{!D_2}{p_2}}$ને
      $\Subst{(\Subst{!A}{!D_1}{p_1})}{!D_2}{p_2}$ સાથે સરખાવો; તે આ છે:}
    \Diamond(p_2 \lif p_3) & \Subst{\lif \Box(\Diamond(p_2 \lif p_3) \land p_2)}{\lnot\Box p_1}{p_2}, \text{ અર્થાત્,}\\
    \Diamond(\lnot\Box p_1 \lif p_3) &
    \lif \Box(\Diamond(\lnot\Box p_1
    \lif p_3) \land \lnot\Box p_1)\\
    \intertext{અને $\Subst{(\Subst{!A}{!D_2}{p_2})}{!D_1}{p_1}$ સાથે પણ સરખાવો:}
    p_1 & \lif \Subst{\Box(p_1 \land \lnot\Box p_1)}{\Diamond(p_2 \lif p_3)}{p_1}, \text{ અર્થાત્,}\\
    \Diamond(p_2 \lif p_3) &
    \lif \Box(\Diamond(p_2
    \lif p_3) \land \lnot\Box\Diamond(p_2 \lif p_3)).
  \end{align*}
\end{ex}
% END OLP-0411
\endgroup

\begingroup
\makeatletter\def\ol@sectioncs{section}\makeatother

% BEGIN OLP-0412
\olfileid{nml}{syn}{rel}

\olsection{સંબંધાત્મક નિદર્શનો}
\phantomsection\label{guunit:OLP-0412}


સામાન્ય મોડલ તર્કશાસ્ત્રોના અર્થવિજ્ઞાનનો મૂળભૂત ખ્યાલ
\emph{સંબંધાત્મક નિદર્શન} છે. તેમાં જગતોનો એક ગણ હોય છે,
જેનાં જગતો દ્વિસ્થાની ``સુલભતા સંબંધ'' વડે જોડાયેલાં હોય છે;
સાથે એવી નિમણૂક હોય છે જે કયા જગતમાં કયા
પ્રતિજ્ઞાપન ચલોને ``સત્ય'' ગણવા તે નક્કી કરે છે.

\begin{defn}
  મૂળભૂત મોડલ ભાષાનું \emph{નિદર્શન} એ ત્રિક
  $\mModel{M} = \tuple{W, R, V}$ છે, જ્યાં
  \begin{enumerate}
  \item $W$ એ ``જગતો''નો અખાલી ગણ છે,
  \item $R$ એ~$W$ પરનો દ્વિસ્થાની સુલભતા સંબંધ છે, અને
  \item $V$ એવું વિધેય છે જે દરેક પ્રતિજ્ઞાપન ચલ~$p$ને શક્ય જગતોનો ગણ $V(p)$ આપે છે.
  \end{enumerate}
  $Rww'$ હોય ત્યારે કહીએ કે $w'$ એ~$w$માંથી \emph{સુલભ}
  છે. $w \in V(p)$ હોય ત્યારે કહીએ કે $p$ એ~$w$માં
  \emph{સત્ય} છે.
\end{defn}

સંબંધાત્મક અર્થવિજ્ઞાનનો મોટો લાભ એ છે કે નિદર્શનોને
\olref{fig:simple} જેવી સરળ આકૃતિઓથી દર્શાવી શકાય છે.
જગતો ગાંઠો વડે દર્શાવાય છે; જગત~$w'$ એ~$w$માંથી
ત્યારે અને માત્ર ત્યારે જ સુલભ છે જ્યારે~$w$થી~$w'$
સુધી તીર હોય. વધુમાં, $w \in V(p)$ હોય ત્યારે ગાંઠને
(જગતને)~$\mTrue{p}$થી અને અન્યથા~\mFalse{p}થી
ચિહ્નિત કરીએ છીએ. \olref{fig:simple}માં
$W = \{w_1, w_2, w_3\}$, $R = \{\tuple{w_1, w_2},\tuple{w_1,
  w_3}\}$, $V(p) = \{w_1, w_2\}$ અને $V(q) = \{w_2\}$
ધરાવતું નિદર્શન દર્શાવ્યું છે.

\begin{figure}
  \begin{center}
    \begin{tikzpicture}[modal]
      \node[world] (w1) [label={[align=right]right:\mTrue{p}\\ \mFalse{q}}]{$w_1$}; 
      \node[world] (w2) [label={[align=right]right:\mTrue{p}\\ \mTrue{q}},
        above right=of w1]{$w_2$}; 
      \node[world] (w3) [label={[align=right]right:\mFalse{p}\\ \mFalse{q}},
        below right=of w1] {$w_3$};
      \draw[->] (w1) to (w2);
      \draw[->] (w1) to (w3);
    \end{tikzpicture}
  \end{center}
  \caption{સરળ નિદર્શન.}
  \ollabel{fig:simple}
\end{figure}
% END OLP-0412
\endgroup

\begingroup
\makeatletter\def\ol@sectioncs{section}\makeatother

% BEGIN OLP-0413
\olfileid{nml}{syn}{trw}

\olsection{જગતમાં સત્યતા}
\phantomsection\label{guunit:OLP-0413}


દરેક મોડલ નિદર્શન તેમાં કયા જગતમાં કયાં મોડલ સૂત્રો
સત્ય ગણાય છે તે નક્કી કરે છે. ``નિદર્શન $\mModel{M}$
સૂત્ર~$!A$ને જગત~$w$માં સત્ય બનાવે છે'' એ સંબંધ
સંબંધાત્મક અર્થવિજ્ઞાનનો મૂળભૂત ખ્યાલ છે. આ સંબંધ
આગમનાત્મક રીતે વ્યાખ્યાયિત થાય છે અને બિનમોડલ
કારકો માટે સત્યતાસારણીઓથી મળતી સામાન્ય વ્યાખ્યા
સાથે મેળ ખાય છે.

\begin{defn}\ollabel{defn:mmodels}
  નિદર્શન~$\mModel M$માં \emph{સૂત્ર~$!A$ની $w$માં સત્યતા},
  સંકેતમાં $\mSat{M}{!A}[w]$, નીચે પ્રમાણે આગમનાત્મક રીતે
  વ્યાખ્યાયિત થાય છે:
  \begin{enumerate}
  \item \indcase{!A}{\lfalse}{$\mSat{M}{\lfalse}[w]$ ક્યારેય નહીં.}
  \item \indcase{!A}{\ltrue}{$\mSat{M}{\ltrue}[w]$ હંમેશા.}
  \item $\mSat{M}{p}[w]$ ત્યારે અને માત્ર ત્યારે જ જ્યારે $w \in V(p)$.
  \item \indcase{!A}{\lnot !B}{$\mSat{M}{\indfrm}[w]$ ત્યારે અને માત્ર ત્યારે જ જ્યારે
    $\mSat/{M}{!B}[w]$}.
  \item \indcase{!A}{(!B \land !C)}{$\mSat{M}{\indfrm}[w]$ ત્યારે અને માત્ર ત્યારે જ જ્યારે
    $\mSat{M}{!B}[w]$ અને $\mSat{M}{!C}[w]$}.
  \item \indcase{!A}{(!B \lor !C)}{$\mSat{M}{\indfrm}[w]$ ત્યારે અને માત્ર ત્યારે જ જ્યારે
    $\mSat{M}{!B}[w]$ અથવા $\mSat{M}{!C}[w]$} (અથવા બંને).
  \item \indcase{!A}{(!B \lif !C)}{$\mSat{M}{\indfrm}[w]$ ત્યારે અને માત્ર ત્યારે જ જ્યારે
    $\mSat/{M}{!B}[w]$ અથવા $\mSat{M}{!C}[w]$}.
  \item \indcase{!A}{(!B \liff !C)}{$\mSat{M}{\indfrm}[w]$ ત્યારે અને માત્ર ત્યારે જ જ્યારે
    $\mSat{M}{!B}[w]$ અને $\mSat{M}{!C}[w]$ બંને હોય અથવા
    $\mSat{M}{!B}[w]$ અને $\mSat{M}{!C}[w]$ બેમાંથી એકેય ન હોય}.
  \item \ollabel{defn:sub:mmodels-box}
    \indcase{!A}{\Box !B}{$\mSat{M}{\indfrm}[w]$ ત્યારે અને માત્ર ત્યારે જ જ્યારે
    $Rww'$ ધરાવતા બધા $w' \in W$ માટે $\mSat{M}{!B}[w']$ હોય.}
  \item \ollabel{defn:sub:mmodels-diamond}
    \indcase{!A}{\Diamond !B}{$\mSat{M}{\indfrm}[w]$ ત્યારે અને માત્ર ત્યારે જ જ્યારે
    $Rww'$ ધરાવતા ઓછામાં ઓછા એક $w' \in W$ માટે $\mSat{M}{!B}[w']$ હોય.}
  \end{enumerate}
\end{defn}
\sourcecorrection{OLMD-004}{સ્થિર અંગ્રેજી મૂળમાં
$\Box$નો સત્યતા-કિસ્સો અશરતી \texttt{item} છે,
જ્યારે ભાષામાં આ પ્રાથમિક કારક \texttt{prvBox}
ટૅગથી પસંદ થાય છે. તે કિસ્સાને એ જ ટૅગથી
શરતિત કર્યો છે; તેના સત્ય-સૂત્રમાં ફેરફાર નથી.}
\sourcecorrection{OLMD-005}{સ્થિર અંગ્રેજી મૂળમાં
$\Diamond$નો સત્યતા-કિસ્સો પણ અશરતી
\texttt{item} છે, જ્યારે ભાષામાં તે \texttt{prvDiamond}
ટૅગથી પસંદ થાય છે. અહીં સત્ય-સૂત્ર જાળવી
તે કિસ્સાને સંબંધિત ટૅગથી શરતિત કર્યો છે.}

\olref{defn:sub:mmodels-box} ખંડ મુજબ, $Rww'$ ધરાવતું
કોઈ~$w'$ ન હોય ત્યારે સૂત્ર~$\Box
!B$ એ~$w$માં સત્ય છે. આ કિસ્સામાં $\Box !B$
એ~$w$માં \emph{તુચ્છ રીતે} સત્ય છે. વળી,
$!B$ સત્ય ન હોય છતાં $\Box !B$ એ~$w$માં સંતોષાઈ શકે છે.
$!B$ એ~$w$માં સત્ય હોય તેથી $\Diamond !B$ પણ~$w$માં
સત્ય હશે એવું જરૂરી નથી. જોકે $Rww$ હોય, જેમ કે
$R$ સ્વવાચક હોય, ત્યારે એ ફલિત થાય છે. જો $Rww'$
ધરાવતું કોઈ $w'$ ન હોય, તો કોઈપણ~$!A$ માટે
$\mSat/{M}{\Diamond !A}[w]$ હોય છે.

\begin{prob}
  \olref[nml][syn][rel]{fig:simple}ના નિદર્શનને વિચારો.
  નીચેનામાંથી કયાં સાચાં છે?
    \begin{enumerate}
    \item $\mSat{M}{q}[w_1]$;
    \item $\mSat{M}{\lnot q}[w_3]$;
    \item $\mSat{M}{p \lor q}[w_1]$;
    \item $\mSat{M}{\Box (p \lor q)}[w_1]$;
    \item $\mSat{M}{\Box q}[w_3]$;
    \item $\mSat{M}{\Box \bot}[w_3]$;
    \item $\mSat{M}{\Diamond q}[w_1]$;
    \item $\mSat{M}{\Box q}[w_1]$;
    \item $\mSat{M}{\lnot \Box\Box \lnot q}[w_1]$.
    \end{enumerate}
\end{prob}

\begin{prop}\ollabel{prop:dual}
  \begin{enumerate}
  \item $\mSat{M}{\Box !A}[w]$ ત્યારે અને માત્ર ત્યારે જ જ્યારે $\mSat{M}{\lnot\Diamond\lnot !A}[w]$.
  \item $\mSat{M}{\Diamond !A}[w]$ ત્યારે અને માત્ર ત્યારે જ જ્યારે $\mSat{M}{\lnot\Box\lnot !A}[w]$.
  \end{enumerate}
\end{prop}

\begin{proof}
  \begin{enumerate}
    \item $\mSat{M}{\lnot\Diamond\lnot !A}[w]$ ત્યારે અને માત્ર ત્યારે જ જ્યારે $\mSat/{M}{\Diamond\lnot
      !A}[w]$; આ $\mSat{M}{}[w]$ની વ્યાખ્યા મુજબ છે. $\mSat{M}{\Diamond\lnot
      !A}[w]$ ત્યારે અને માત્ર ત્યારે જ જ્યારે $Rww'$ ધરાવતા કોઈ $w'$ માટે
      $\mSat{M}{\lnot !A}[w']$ હોય. તેથી $\mSat/{M}{\Diamond\lnot !A}[w]$
      ત્યારે અને માત્ર ત્યારે જ જ્યારે $Rww'$ ધરાવતા બધા $w'$ માટે
      $\mSat/{M}{\lnot !A}[w']$ હોય. વળી, $\mSat/{M}{\lnot !A}[w']$
      ત્યારે અને માત્ર ત્યારે જ જ્યારે $\mSat{M}{!A}[w']$ હોય. આ બધું
      સાથે લેતાં, $\mSat{M}{\lnot\Diamond\lnot !A}[w]$ ત્યારે અને માત્ર
      ત્યારે જ જ્યારે $Rww'$ ધરાવતા બધા $w'$ માટે $\mSat{M}{!A}[w']$ હોય.
      ફરી $\mSat{M}{}[w]$ની વ્યાખ્યા મુજબ, આ ત્યારે અને માત્ર ત્યારે જ
      જ્યારે $\mSat{M}{\Box !A}[w]$ હોય.
    \item $\mSat{M}{\lnot\Box\lnot !A}[w]$
      ત્યારે અને માત્ર ત્યારે જ જ્યારે $\mSat/{M}{\Box\lnot !A}[w]$.
      $\mSat{M}{\Box\lnot !A}[w]$ ત્યારે અને માત્ર ત્યારે જ જ્યારે
      $Rww'$ ધરાવતા બધા $w'$ માટે $\mSat{M}{\lnot !A}[w']$ હોય. તેથી
      $\mSat/{M}{\Box\lnot !A}[w]$ ત્યારે અને માત્ર ત્યારે જ જ્યારે
      $Rww'$ ધરાવતા કોઈ $w'$ માટે $\mSat/{M}{\lnot !A}[w']$ હોય.
      વળી, $\mSat/{M}{\lnot !A}[w']$ ત્યારે અને માત્ર ત્યારે જ જ્યારે
      $\mSat{M}{!A}[w']$ હોય. તેથી $\mSat{M}{\lnot\Box\lnot !A}[w]$
      ત્યારે અને માત્ર ત્યારે જ જ્યારે $Rww'$ ધરાવતા કોઈ~$w'$ માટે
      $\mSat{M}{!A}[w']$ હોય. ફરી $\mSat{M}{}[w]$ની વ્યાખ્યા મુજબ,
      આ ત્યારે અને માત્ર ત્યારે જ જ્યારે $\mSat{M}{\Diamond
        !A}[w]$ હોય.
  \end{enumerate}
\end{proof}
\sourcecorrection{OLMD-006}{સ્થિર અંગ્રેજી મૂળ
દ્વિતીય દ્વૈતતાના પ્રમાણમાં
$\mSat/{M}{\Box\lnot !A}$ પાસે જગત-દલીલ $[w]$
છોડી દે છે. આસપાસનાં બધા સત્યતા-સંબંધો અને
વ્યાખ્યાની જેમ એ જ $[w]$ અહીં પૂરું કર્યું છે.}

\begin{probtag}{probDiamond}
  \olref[nml][syn][trw]{prop:dual}નું પ્રમાણ પૂર્ણ કરો.
\end{probtag}

\begin{prob}
  માનો કે $\mModel{M} = \tuple{W, R, V}$ નિદર્શન છે
  અને $w_1, w_2 \in W$ માટે નીચેની શરતો છે:
  \begin{enumerate}
  \item $w_1 \in V(p)$ જો અને માત્ર જો $w_2 \in V(p)$
    (દરેક પ્રતિજ્ઞાપન ચલ $p$ માટે); અને
  \item બધા $w \in W$ માટે: $Rw_1w$ જો અને માત્ર જો $Rw_2w$.
  \end{enumerate}
  સૂત્રો પર અનુમાન વડે બતાવો કે બધા સૂત્રો $!A$ માટે:
  $\mSat{M}{!A}[w_1]$ જો અને માત્ર જો $\mSat{M}{!A}[w_2]$.
\end{prob}

\begin{probtag}{prvDiamond}
  માનો કે $\mModel{M} = \tuple{W, R, V}$. બતાવો કે
  $\mSat{M}{\lnot\Diamond !A}[w]$ જો અને માત્ર જો
  $\mSat{M}{\Box\lnot!A}[w]$.
\end{probtag}
% END OLP-0413
\endgroup

\begingroup
\makeatletter\def\ol@sectioncs{section}\makeatother

% BEGIN OLP-0414
\olfileid{nml}{syn}{tru}

\olsection{નિદર્શનમાં સત્યતા}
\phantomsection\label{guunit:OLP-0414}


ક્યારેક કોઈ નિદર્શનના દરેક જગતમાં કયાં સૂત્રો સત્ય છે
તે જાણવામાં આપણને રસ હોય છે. તે માટે સંકેત રજૂ કરીએ.

\begin{defn}
  સૂત્ર~$!A$ એ નિદર્શન $M = \tuple{W, R,
    V}$માં \emph{સત્ય} છે, જે $\mSat{M}{!A}$થી દર્શાવાય છે,
  ત્યારે અને માત્ર ત્યારે જ જ્યારે દરેક $w \in W$ માટે
  $\mSat{M}{!A}[w]$ હોય.
\end{defn}

\begin{prop}\ollabel{prop:truthfacts}
  \begin{enumerate}
  \item જો $\mSat{M}{!A}$ હોય તો $\mSat/{M}{\lnot !A}$ હોય,
    પણ તેનું વિપરીત વિધાન \emph{સાચું નથી}.
  \item જો $\mSat{M}{!A \lif !B}$ હોય તો $\mSat{M}{!A}$ હોય
    ત્યારે $\mSat{M}{!B}$ પણ હોય; પરંતુ તેનું વિપરીત વિધાન
    \emph{સાચું નથી}.
  \end{enumerate}
\end{prop}

\begin{proof}
  \begin{enumerate}
  \item જો $\mSat{M}{!A}$ હોય તો $!A$ એ $W$ના બધા જગતોમાં
    સત્ય છે. $W \neq \emptyset$ હોવાથી $\mSat{M}{\lnot!A}$
    શક્ય નથી; નહીંતર કોઈ એક જગતમાં $!A$ સત્ય પણ અને અસત્ય પણ
    હોવું પડે.

    બીજી બાજુ, જો $\mSat/{M}{\lnot !A}$ હોય તો $!A$ એ કોઈ
    $w \in W$માં સત્ય છે. તે પરથી \emph{દરેક}~$w \in W$ માટે
    $\mSat{M}{!A}[w]$ ફલિત થતું નથી. ઉદાહરણ તરીકે,
    \olref[rel]{fig:simple}ના નિદર્શનમાં $\mSat/{M}{\lnot p}$
    અને $\mSat/{M}{p}$ બંને છે.
  \item $\mSat{M}{!A \lif !B}$ અને $\mSat{M}{!A}$ માનો.
    $\mSat{M}{!B}$ બતાવવા માટે $w \in W$ કોઈપણ જગત લો.
    ત્યારે $\mSat{M}{!A \lif !B}[w]$ અને $\mSat{M}{!A}[w]$
    હોવાથી $\mSat{M}{!B}[w]$ છે. $w$ ગમે તે હતું, તેથી
    $\mSat{M}{!B}$.

    વિપરીત વિધાન ખોટું છે તે બતાવવા એવું નિદર્શન
    $\mModel{M}$ જોઈએ જેમાં $\mSat{M}{!A}$ હોય તો
    $\mSat{M}{!B}$ હોય, છતાં $\mSat/{M}{!A \lif !B}$ હોય.
    ફરી \olref[rel]{fig:simple}નું નિદર્શન લો:
    $\mSat/{M}{p}$ હોવાથી (પૂર્વપક્ષ ખોટો હોવાથી)
    $\mSat{M}{p}$ હોય તો $\mSat{M}{q}$ હોય એવું વિધાન સાચું છે.
    તેમ છતાં $\mSat/{M}{p \lif
      q}$ છે, કેમ કે~$w_1$માં $p$ સત્ય છે અને $q$ અસત્ય છે.
  \end{enumerate}
\end{proof}

\begin{prob}
  માત્ર $p_1$, $p_2$, $p_3$ને પ્રતિજ્ઞાપન ચલો તરીકે
  ધરાવતી ભાષા માટે નીચેનું નિદર્શન $\mModel{M}$ વિચારો:
  \begin{center}
    \begin{tikzpicture}[modal]
      \node[world] (w1) [label={[align=right]left:\mTrue{p_1}\\\mFalse{p_2}\\\mFalse{p_3}}]
            {$w_1$} ; 
      \node[world] (w2) [label={[align=right]right:\mTrue{p_1}\\\mTrue{p_2}\\\mFalse{p_3}},
        below right=of w1]  {$w_2$}; 
      \node[world] (w3) [label={[align=right]right:\mTrue{p_1}\\\mTrue{p_2}\\\mTrue{p_3}},
        above right=of w2] {$w_3$};
      \draw[reflexive above] (w3) to (w3); 
      \draw[->] (w1) to (w2);
      \draw[->] (w2) to (w3);
      \draw[->] (w1) to (w3); 
    \end{tikzpicture}
  \end{center}
  નીચેના સૂત્રો અને સૂત્ર-ઢાંચા નિદર્શન
  $\mModel{M}$માં, એટલે કે $\mModel{M}$ના દરેક જગતમાં, સત્ય છે?
  સમજાવો.
  \begin{enumerate}
  \item $p\lif \Diamond p$ ($p$ પરમાણુ હોય ત્યારે);
  \item $!A\lif \Diamond !A$ ($!A$ કોઈપણ સૂત્ર હોય ત્યારે);
  \item $\Box p \lif p$ ($p$ પરમાણુ હોય ત્યારે);
  \item $\lnot p \lif \Diamond \Box p$ ($p$ પરમાણુ હોય ત્યારે);
  \item $\Diamond \Box !A$ ($!A$ કોઈપણ સૂત્ર હોય ત્યારે);
  \item $\Box \Diamond p$ ($p$ પરમાણુ હોય ત્યારે).
  \end{enumerate}
\end{prob}
% END OLP-0414
\endgroup

\begingroup
\makeatletter\def\ol@sectioncs{section}\makeatother

% BEGIN OLP-0415
\olfileid{nml}{syn}{val}

\olsection{માન્યતા}
\phantomsection\label{guunit:OLP-0415}


\begin{explain}
  બધાં નિદર્શનોમાં, એટલે કે દરેક નિદર્શનના દરેક જગતમાં,
  સત્ય હોય તેવાં સૂત્રો ખાસ રસપ્રદ છે. શક્ય જગતો
  પરના કોઈ સુલભતા સંબંધથી ઉદ્ભવતા અર્થમાં અર્થઘટન
  ``સામાન્ય'' હોય ત્યાં સુધી $\Box$ અને $\Diamond$નું
  અર્થઘટન ગમે તેમ કરવામાં આવે, આ મોડલ વિધાનો
  સત્ય રહે છે. આવાં સૂત્રોને આપણે \emph{માન્ય} કહીએ.
  ઉદાહરણ તરીકે, $\Box(p \land q) \lif \Box p$ માન્ય છે.
  પરંતુ $\Box$ના સત્યલક્ષી અર્થઘટનને આધારે માન્ય હોવાની
  અપેક્ષા રાખી શકાય એવાં કેટલાંક સૂત્રો, જેમ કે
  $\Box p \lif p$, માન્ય નથી. સંબંધાત્મક નિદર્શનો
  રસપ્રદ હોવાનું એક કારણ એ છે કે વિવિધ પ્રકારના
  સુલભતા સંબંધો વડે $\Box$ અને $\Diamond$નાં વિવિધ
  અર્થઘટનો વ્યક્ત કરી શકાય છે. તેથી માન્યતાની વ્યાખ્યા
  માત્ર \emph{બધાં} નિદર્શનો સાપેક્ષે જ નહીં, પરંતુ
  \emph{ચોક્કસ પ્રકારનાં} બધાં નિદર્શનો સાપેક્ષે પણ
  આપવી જોઈએ. દાખલા તરીકે, જે નિદર્શનોમાં દરેક
  જગત પોતામાંથી સુલભ છે, એટલે કે $R$ સ્વવાચક છે,
  તેમાં $\Box p \lif p$ સત્ય છે. નિદર્શનોના વર્ગ
  સાપેક્ષે માન્યતા વ્યાખ્યાયિત કરવાથી આ વાત ટૂંકમાં
  કહી શકાય: $\Box p \lif p$ સ્વવાચક નિદર્શનોના
  વર્ગમાં માન્ય છે.
\end{explain}

\begin{defn}
  સૂત્ર~$!A$ નિદર્શનોના વર્ગ $\mClass{C}$માં
  \emph{માન્ય} છે, જો તે $\mClass{C}$ના દરેક નિદર્શનમાં
  સત્ય હોય (એટલે કે $\mClass{C}$ના દરેક નિદર્શનના
  દરેક જગતમાં સત્ય હોય). જો $!A$ એ~$\mClass{C}$માં
  માન્ય હોય, તો $\mClass{C} \Entails !A$ લખીએ;
  અને $!A$ \emph{બધાં} નિદર્શનોના વર્ગમાં માન્ય હોય,
  તો $\Entails !A$ લખીએ.
\end{defn}

\begin{prop}\ollabel{prop:subset-class}
  જો $!A$ વર્ગ $\mClass{C}$માં માન્ય હોય, તો તે દરેક
  ઉપવર્ગ $\mClass{C}' \subseteq \mClass{C}$માં પણ માન્ય છે.
\end{prop}

\begin{prop}\ollabel{prop:Nec-rule}
  જો $!A$ માન્ય હોય, તો $\Box!A$ પણ માન્ય છે.
\end{prop}

\begin{proof}
  $\Entails !A$ માનો. $\Entails \Box!A$ બતાવવા માટે
  કોઈપણ નિદર્શન $\mModel{M} =
  \tuple{W, R, V}$ અને $w \in W$ લો. જો $Rww'$ હોય,
  તો $!A$ માન્ય હોવાથી $\mSat{M}{!A}[w']$ છે;
  તેથી $\mSat{M}{\Box!A}[w]$ પણ છે. નિદર્શન
  $\mModel{M}$ અને જગત $w$ ગમે તે હતાં, તેથી
  $\Entails \Box!A$.
\end{proof}

\begin{prob}
  નીચેના સૂત્રો માન્ય છે તે બતાવો:
  \begin{enumerate}
  \item $\Box p \lif \Box (q \lif p)$;
  \item $\Box \lnot \lfalse$;
  \item $\Box p \lif (\Box q \lif \Box p)$.
  \end{enumerate}
\end{prob}

\begin{prob}
  $W = \{w\}$ ધરાવતા નિદર્શનો $\mModel{M} =
  \tuple{W, R, V}$ના વર્ગ $\mClass{C}$માં
  $!A \lif \Box!A$ માન્ય છે તે બતાવો. એ જ રીતે,
  $R = \emptyset$ ધરાવતા નિદર્શનો $\mModel{M} =
  \tuple{W, R, V}$ના વર્ગમાં $!B \lif \Box !A$
  અને $\Diamond !A \lif !B$ માન્ય છે તે બતાવો.
\end{prob}
% END OLP-0415
\endgroup

\begingroup
\makeatletter\def\ol@sectioncs{section}\makeatother

% BEGIN OLP-0416
\olfileid{nml}{syn}{tau}

\olsection{પુનરુક્તિના પ્રતિસ્થાપન દૃષ્ટાંતો}
\phantomsection\label{guunit:OLP-0416}


\begin{explain}
મોડલ-મુક્ત સૂત્ર દરેક સત્યમૂલ્ય-નિમણૂક હેઠળ સત્ય હોય
તો તે પુનરુક્તિ છે. સ્પષ્ટ છે કે દરેક પુનરુક્તિ દરેક
નિદર્શનના દરેક જગતમાં સત્ય છે. પરંતુ $\Box$
અને~$\Diamond$ ધરાવતા સૂત્રો માટે પુનરુક્તિનો
ખ્યાલ વ્યાખ્યાયિત નથી. તો શું, દાખલા તરીકે,
$\Box p \lor \lnot \Box p$---મધ્યના નિષેધના સિદ્ધાંતનો
એક દૃષ્ટાંત---માન્ય છે? \emph{પુનરુક્તિનો પ્રતિસ્થાપન દૃષ્ટાંત}
એ ખ્યાલ ઉપયોગી છે: તે (બિનમોડલ) પુનરુક્તિનો
પ્રતિસ્થાપન દૃષ્ટાંત હોય એવું સૂત્ર. દરેક
આવો દૃષ્ટાંત માન્ય છે તે નવાઈની વાત નથી, છતાં
તેનું પ્રમાણ જરૂરી છે.
\end{explain}

\begin{defn}
  મોડલ સૂત્ર~$!B$ \emph{પુનરુક્તિનો પ્રતિસ્થાપન દૃષ્ટાંત}
  ત્યારે અને માત્ર ત્યારે જ છે જ્યારે પ્રતિજ્ઞાપન ચલો $p_1$, \dots,~$p_n$ ધરાવતી મોડલ-મુક્ત
  પુનરુક્તિ~$!A$ અને સૂત્રો $!D_1$,
  \dots,~$!D_n$ એવાં હોય કે $!B \ident \SSubst{!A}{\subst{!D_1}{p_1},
    \dots, \subst{!D_n}{p_n}}$.
\end{defn}

\begin{lem}\ollabel{lem:valid-taut}
  માનો કે $!A$ એ મોડલ-મુક્ત સૂત્ર છે જેનાં પ્રતિજ્ઞાપન ચલો $p_1$, \dots, $p_n$ છે, અને $!D_1$, \dots,
  $!D_n$ મોડલ સૂત્રો છે. તો કોઈપણ નિમણૂક $\pAssign{v}$,
  કોઈપણ નિદર્શન $\mModel{M} = \tuple{W, R, V}$ અને એવા કોઈપણ
  $w \in W$ માટે કે જ્યાં $\pAssign{v}(p_i) = \True$ ત્યારે અને માત્ર
  ત્યારે જ જ્યારે $\mSat{M}{!D_i}[w]$ હોય, નીચેનું સાચું છે:
  $\pSat{v}{!A}$ ત્યારે અને માત્ર ત્યારે જ જ્યારે
  $\mSat{M}{\SSubst{!A}{\subst{!D_1}{p_1}, \dots,
      \subst{!D_n}{p_n}}}[w]$.
\end{lem}

\begin{proof}
  $!A$ પર આગમનથી.
  \begin{enumerate}
  \item \indcase{!A}{\lfalse}{$\pSat/{v}{\lfalse}$
      અને $\mSat/{M}{\lfalse}[w]$ બંને ખોટાં છે.}
  \item \indcase{!A}{\ltrue}{$\pSat{v}{\ltrue}$
      અને $\mSat{M}{\ltrue}[w]$ બંને સાચાં છે.}
  \item \indcase{!A}{p_i}{%
    \begin{align*}
      \pSat{v}{p_i} \Leftrightarrow {} & \pAssign{v}(p_i) = \True \\
      & \qquad \text{$\pSat{v}{p_i}$ની વ્યાખ્યા મુજબ}\\
      \Leftrightarrow {} & \mSat{M}{!D_i}[w] \\
      &\qquad \text{ધારણા મુજબ}\\
      \Leftrightarrow {} & \mSat{M}{\SSubst{p_i}{\subst{!D_1}{p_1}, \dots,
          \subst{!D_n}{p_n}}}[w]\\
      &\qquad \text{કેમ કે $\SSubst{p_i}{\subst{!D_1}{p_1}, \dots,
          \subst{!D_n}{p_n}} \ident !D_i$}.
      \end{align*}}
  \item \indcase{!A}{\lnot !B}{%
      \begin{align*}
        \pSat{v}{\lnot !B} \Leftrightarrow {} & \pSat/{v}{!B}\\
        &\qquad \text{$\pSat{v}{}$ની વ્યાખ્યા મુજબ};\\
        \Leftrightarrow {} & \mSat/{M}{\SSubst{!B}{\subst{!D_1}{p_1}, \dots,
          \subst{!D_n}{p_n}}}[w]\\
        &\qquad \text{આગમન-ધારણા મુજબ}\\
        \Leftrightarrow {} &
        \mSat{M}{\SSubst{\lnot !B}{\subst{!D_1}{p_1}, \dots,
          \subst{!D_n}{p_n}}}[w]\\
        &\qquad \text{$\mSat{M}{}[w]$ની વ્યાખ્યા મુજબ}.
      \end{align*}}
  \item \indcase{!A}{(!B \land !C)}{%
      \begin{align*}
        \pSat{v}{!B \land !C} \Leftrightarrow {} &
        \pSat{v}{!B} \text{ અને } \pSat{v}{!C}\\
        &\qquad \text{$\pSat{v}{}$ની વ્યાખ્યા મુજબ}\\
        \Leftrightarrow {} &
        \mSat{M}{\SSubst{!B}{\subst{!D_1}{p_1}, \dots,
            \subst{!D_n}{p_n}}}[w] \text{ અને } \\
        & \mSat{M}{\SSubst{!C}{\subst{!D_1}{p_1}, \dots,
          \subst{!D_n}{p_n}}}[w]\\
        &\qquad \text{આગમન-ધારણા મુજબ}\\
        \Leftrightarrow{} &
        \mSat{M}{\SSubst{(!B \land !C)}{\subst{!D_1}{p_1}, \dots,
          \subst{!D_n}{p_n}}}[w]\\
        &\qquad \text{$\mSat{M}{}[w]$ની વ્યાખ્યા મુજબ}.
      \end{align*}}
  \item \indcase{!A}{(!B \lor !C)}{%
      \begin{align*}
        \pSat{v}{!B \lor !C} \Leftrightarrow{} &
        \pSat{v}{!B} \text{ અથવા } \pSat{v}{!C}\\
        &\qquad \text{$\pSat{v}{}$ની વ્યાખ્યા મુજબ};\\
        \Leftrightarrow{} & \mSat{M}{\SSubst{!B}{\subst{!D_1}{p_1}, \dots,
            \subst{!D_n}{p_n}}}[w] \text{ અથવા }\\
        & \mSat{M}{\SSubst{!C}{\subst{!D_1}{p_1}, \dots,
          \subst{!D_n}{p_n}}}[w]\\
        &\qquad \text{આગમન-ધારણા મુજબ}\\
        \Leftrightarrow{} &
        \mSat{M}{\SSubst{(!B \lor !C)}{\subst{!D_1}{p_1}, \dots,
          \subst{!D_n}{p_n}}}[w]\\
        &\qquad \text{$\mSat{M}{}[w]$ની વ્યાખ્યા મુજબ}.
      \end{align*}}
  \item \indcase{!A}{(!B \lif !C)}{%
      \begin{align*}
        \pSat{v}{!B \lif !C} \Leftrightarrow{} &
        \pSat/{v}{!B} \text{ અથવા } \pSat{v}{!C}\\
        &\qquad \text{$\pSat{v}{}$ની વ્યાખ્યા મુજબ}\\
        \Leftrightarrow{} &
        \mSat/{M}{\SSubst{!B}{\subst{!D_1}{p_1}, \dots,
            \subst{!D_n}{p_n}}}[w] \text{ અથવા }\\
        & \mSat{M}{\SSubst{!C}{\subst{!D_1}{p_1}, \dots,
          \subst{!D_n}{p_n}}}[w]\\
        &\qquad \text{આગમન-ધારણા મુજબ}\\
        \Leftrightarrow{} &
        \mSat{M}{\SSubst{(!B \lif !C)}{\subst{!D_1}{p_1}, \dots,
          \subst{!D_n}{p_n}}}[w]\\
        &\qquad \text{$\mSat{M}{}[w]$ની વ્યાખ્યા મુજબ}.
      \end{align*}}
  \item \indcase{!A}{(!B \liff !C)}{%
      \begin{align*}
        \pSat{v}{!B \liff !C} \Leftrightarrow {} &
        \text{ક્યાં તો } \pSat{v}{!B} \text{ અને } \pSat{v}{!C}\\
        & \text{અથવા } \pSat/{v}{!B} \text{ અને } \pSat/{v}{!C}\\
        &\qquad \text{$\pSat{v}{}$ની વ્યાખ્યા મુજબ}\\
        \Leftrightarrow {} &
        \text{ક્યાં તો } \mSat{M}{\SSubst{!B}{\subst{!D_1}{p_1}, \dots,
            \subst{!D_n}{p_n}}}[w] \text{ અને }\\
        & \qquad \mSat{M}{\SSubst{!C}{\subst{!D_1}{p_1}, \dots,
          \subst{!D_n}{p_n}}}[w]\\ 
        &  \text{ અથવા } \mSat/{M}{\SSubst{!B}{\subst{!D_1}{p_1}, \dots,
            \subst{!D_n}{p_n}}}[w] \text{ અને }\\
        & \qquad \mSat/{M}{\SSubst{!C}{\subst{!D_1}{p_1}, \dots,
          \subst{!D_n}{p_n}}}[w]\\
        &\qquad\text{આગમન-ધારણા મુજબ}\\
        \Leftrightarrow {} &
        \mSat{M}{\SSubst{(!B \liff !C)}{\subst{!D_1}{p_1}, \dots,
          \subst{!D_n}{p_n}}}[w]\\
        &\qquad \text{$\mSat{M}{}[w]$ની વ્યાખ્યા મુજબ}.
      \end{align*}}
  \end{enumerate}
  \end{proof}
\sourcecorrection{OLMD-007}{સ્થિર અંગ્રેજી મૂળમાં
નકારનો આગમન-કિસ્સો \texttt{prvFalse} હેઠળ મૂકાયો છે;
અહીં એ જ પ્રમાણ \texttt{prvNot} હેઠળ મૂક્યું છે.}
\sourcecorrection{OLMD-008}{દ્વિપક્ષીયના આગમન-કિસ્સાની
પ્રથમ સમતામાં મૂળે \texttt{lif} લખાયેલું છે; કિસ્સાના
સૂત્ર અને અંતિમ સમતા પ્રમાણે તેને \texttt{liff} કર્યું છે.}
\sourcecorrection{OLMD-009}{નકારના કિસ્સાની અંતિમ
સમતાનું કારણ મૂળે \texttt{pSat}ની વ્યાખ્યા કહે છે;
તે ખરેખર નિદર્શનની \texttt{mSat} સત્યતા-વ્યાખ્યા છે.}

\begin{prop}\ollabel{prop:valid-taut}
  પુનરુક્તિના બધા પ્રતિસ્થાપન દૃષ્ટાંતો માન્ય છે.
\end{prop}

\begin{proof}
  પ્રતિધનાત્મક રીતે, માનો કે કોઈ નિદર્શન $\mModel{M}$ અને
  જગત~$w$ માટે એવું $!A$ છે કે
  $\mSat/{M}{\SSubst{!A}{\subst{!D_1}{p_1}, \dots,
  \subst{!D_n}{p_n}}}[w]$.
  એવી નિમણૂક $\pAssign{v}$ વ્યાખ્યાયિત કરો કે $\pAssign{v}(p_i) =
  \True$ ત્યારે અને માત્ર ત્યારે જ જ્યારે $\mSat{M}{!D_i}[w]$
  હોય (અને $q \notin \{p_1, \dots, p_n \}$ માટે $\pAssign{v}$
  ગમે તે સત્યમૂલ્યો આપે). તો \olref{lem:valid-taut} મુજબ
  $\pSat/{v}{!A}$, તેથી $!A$ પુનરુક્તિ નથી.
\end{proof}
% END OLP-0416
\endgroup

\begingroup
\makeatletter\def\ol@sectioncs{section}\makeatother

% BEGIN OLP-0417
\olfileid{nml}{syn}{sch}

\olsection{સૂત્ર-ઢાંચા અને માન્યતા}
\phantomsection\label{guunit:OLP-0417}


\begin{defn}
  \emph{સૂત્ર-ઢાંચો} એ સૂત્રોનો એવો ગણ છે જેમાં
  કોઈ મોડલ સૂત્ર~$!C$ના બધા અને માત્ર બધા
  પ્રતિસ્થાપન દૃષ્ટાંતો હોય; એટલે કે,
  \[
  \Setabs{!B}{\lexists[!D_1], \dots, \lexists[!D_n] \left(!B =
  \SSubst{!C}{\subst{!D_1}{p_1}, \dots, \subst{!D_n}{p_n}} \right) }.
  \]
  સૂત્ર~$!C$ને આ સૂત્ર-ઢાંચાનું \emph{લક્ષણરૂપ}
  સૂત્ર કહે છે; પ્રતિજ્ઞાપન ચલોનાં નામ બદલવા સિવાય તે અનન્ય છે.
  સૂત્ર~$!A$ આ ગણનું સભ્ય હોય તો તેને
  સૂત્ર-ઢાંચાનો \emph{દૃષ્ટાંત} કહે છે.
\end{defn}

અણુ-ઘટકોને બદલે `$!A$', `$!B$',~\dots{} મૂકી
$!C$માંથી મળતી અધિભાષાકીય અભિવ્યક્તિ વડે
સૂત્ર-ઢાંચો દર્શાવવો અનુકૂળ છે. તેથી, દાખલા તરીકે,
`$!A$', `$!A \lif \Box!A$', `$!A \lif (!B \lif !A)$'
નીચેના સૂત્ર-ઢાંચા દર્શાવે છે. તેમનાં લક્ષણરૂપ
સૂત્રો અનુક્રમે $p$, $p \lif \Box p$, $p \lif (q
\lif p)$ છે. `$!A$' સૂત્ર-ઢાંચો \emph{બધાં}
સૂત્રોનો ગણ દર્શાવે છે.

\begin{defn}
  સૂત્ર-ઢાંચો તેના બધા દૃષ્ટાંતો નિદર્શનમાં સત્ય હોય
  ત્યારે અને માત્ર ત્યારે જ તે નિદર્શનમાં \emph{સત્ય} છે;
  અને તે દરેક નિદર્શનમાં સત્ય હોય ત્યારે અને માત્ર
  ત્યારે જ \emph{માન્ય} છે.
\end{defn}

\begin{prop}\ollabel{prop:Kvalid}
  નીચેનો સૂત્ર-ઢાંચો \Ax{K} માન્ય છે:
  \begin{equation*}
  \Box(!A \lif !B) \lif (\Box !A \lif \Box !B). \tag{\Ax{K}}
  \end{equation*}
\end{prop}

\begin{proof}
  સૂત્ર-ઢાંચાના બધા દૃષ્ટાંતો દરેક નિદર્શનના દરેક
  જગતમાં સત્ય છે તે બતાવવું છે. તેથી કોઈપણ
  $\mModel{M} = \tuple{W,R,V}$ અને $w \in
  W$ લો. શરતી વિધાન જગતમાં સત્ય છે તે બતાવવા,
  તેનો પૂર્વપક્ષ સત્ય માની તેનો ઉત્તરપક્ષ પણ સત્ય
  છે તે બતાવીએ. અહીં $\mSat{M}{\Box(!A \lif !B)}[w]$
  અને $\mSat{M}{\Box !A}[w]$ માનો. હવે
  $\mSat{M}{\Box !B}[w]$ બતાવવું છે. તો $Rww'$
  ધરાવતું કોઈપણ $w'$ લો. પ્રથમ ધારણા મુજબ
  $\mSat{M}{!A \lif !B}[w']$ અને બીજી મુજબ
  $\mSat{M}{!A}[w']$. તેથી $\mSat{M}{!B}[w']$.
  $w'$ ગમે તે હતું, તેથી $\mSat{M}{\Box !B}[w]$.
\end{proof}

\begin{prop}\ollabel{prop:Dual-valid}
  નીચેનો સૂત્ર-ઢાંચો \Dual{} માન્ય છે:
  \begin{equation*}
  \Diamond !A \liff \lnot\Box\lnot !A. \tag{\Dual}
  \end{equation*}
\end{prop}

\begin{proof}
  અભ્યાસપ્રશ્ન.
\end{proof}

\begin{prob}
  \olref[nml][syn][sch]{prop:Dual-valid} સાબિત કરો.
\end{prob}

\begin{prop}\ollabel{prop:soundMP}
  કોઈ નિદર્શનના એક જગતમાં $!A$ અને $!A \lif !B$
  સત્ય હોય તો $!B$ પણ ત્યાં સત્ય છે. તેથી માન્ય
  સૂત્રો મોડસ પોનેન્સ હેઠળ બંધ છે.
\end{prop}

\begin{prop}\ollabel{prop:valid-instances}
  સૂત્ર~$!A$ ત્યારે અને માત્ર ત્યારે જ માન્ય છે
  જ્યારે તેના બધા પ્રતિસ્થાપન દૃષ્ટાંતો માન્ય હોય.
  બીજા શબ્દોમાં, સૂત્ર-ઢાંચો ત્યારે અને માત્ર ત્યારે જ
  માન્ય છે જ્યારે તેનું લક્ષણરૂપ સૂત્ર માન્ય હોય.
\end{prop}

\begin{proof}
  ``જો'' દિશા સ્પષ્ટ છે, કારણ કે $!A$ પોતાનો જ
  પ્રતિસ્થાપન દૃષ્ટાંત છે.

  ``માત્ર જો'' દિશા સાબિત કરવા નીચેનું બતાવીએ. માનો કે
  $\mModel{M} = \tuple{W, R, V}$ મોડલ નિદર્શન છે,
  અને $!B \ident
  \SSubst{!A}{\subst{!D_1}{p_1},\dots,\subst{!D_n}{p_n}}$
  એ~$!A$નો પ્રતિસ્થાપન દૃષ્ટાંત છે. $V'(p_i) =
  \Setabs{w}{\mSat{M}{!D_i}[w]}$ મૂકી $\mModel{M'} = \tuple{W, R,
    V'}$ વ્યાખ્યાયિત કરો. ત્યારે કોઈપણ~$w \in W$ માટે
  $\mSat{M}{!B}[w]$ જો અને માત્ર જો $\mSat{M'}{!A}[w]$.
  (આ દાવાનું પ્રમાણ અભ્યાસપ્રશ્ન તરીકે છોડીએ છીએ.)
  હવે માનો કે $!A$ માન્ય છે, પણ તેનો કોઈ પ્રતિસ્થાપન
  દૃષ્ટાંત $!B$ એ~$!A$નો હોવા છતાં માન્ય નથી. તો કોઈ $\mModel{M} =
  \tuple{W, R, V}$ અને કોઈ $w \in W$ માટે
  $\mSat/{M}{!B}[w]$. પરંતુ દાવા મુજબ ત્યારે
  $\mSat/{M'}{!A}[w]$, એટલે $!A$ માન્ય નથી;
  આ વિરોધાભાસ છે.
\end{proof}

\begin{prob}
  \olref[nml][syn][sch]{prop:valid-instances}ના પ્રમાણની
  ``માત્ર જો'' દિશામાંનો દાવો સાબિત કરો. (સૂચન:
  $!A$ પર આગમન વાપરો.)
\end{prob}

જોકે, સૂત્ર-ઢાંચો કોઈ નિદર્શનમાં સત્ય હોય ત્યારે
અને માત્ર ત્યારે જ તેનું લક્ષણરૂપ સૂત્ર ત્યાં સત્ય
હોય એવું સાચું નથી. ``માત્ર જો'' દિશા અલબત્ત સાચી છે:
જો $!A$નો દરેક દૃષ્ટાંત~$\mModel{M}$માં સત્ય હોય,
તો $!A$~પોતે પણ~$\mModel{M}$માં સત્ય છે. પરંતુ
$!A$ એ~$\mModel{M}$માં સત્ય હોય છતાં $!A$નો કોઈ દૃષ્ટાંત
$\mModel{M}$ના કોઈ જગતમાં અસત્ય હોઈ શકે છે. સરળ
વિરોધી દૃષ્ટાંત માટે~$p$ લો અને માત્ર એક જગત~$w$ અને
$V(p) = \{w\}$ ધરાવતું નિદર્શન લો; તેમાં $p$
એ~$w$માં સત્ય છે. પરંતુ $\lfalse$ એ~$p$નો
દૃષ્ટાંત છે અને~$w$માં સત્ય નથી.

\begin{prob}
  નીચેના સૂત્રોમાંથી એકેય માન્ય નથી તે બતાવો:
  \begin{enumerate}
    \item[\Ax{D}:] \quad $\Box p \lif \Diamond p$;
    \item[\Ax{T}:] \quad $\Box p \lif p$;
    \item[\Ax{B}:] \quad $p \lif \Box\Diamond p$;
    \item[\Ax{4}:] \quad $\Box p \lif \Box \Box p$;
    \item[\Ax{5}:] \quad $\Diamond p \lif \Box \Diamond
      p$.
  \end{enumerate}
\end{prob}

\begin{table}[t]
    \centering
    \begin{tabular}{| l || l |}
      \hline
      {\emph{માન્ય સૂત્ર-ઢાંચા}} & {\emph{અમાન્ય સૂત્ર-ઢાંચા}} \\
      \hline\hline
      $\Box(!A \lif !B) \lif (\Diamond !A \lif \Diamond !B)$
      & $\Box (!A \lor !B) \lif (\Box !A \lor \Box !B)$ \\
      $\Diamond (!A \lif !B) \lif (\Box !A \lif \Diamond
      !B)$
      & $(\Diamond !A \land \Diamond !B) \lif \Diamond (!A
      \land !B)$\\
      $\Box (!A \land !B) \liff (\Box !A \land \Box !B)$
      & $!A \lif \Box !A$ \\
      $\Box !A \lif \Box (!B \lif !A)$
      & $\Box \Diamond !A \lif !B$ \\
      $\lnot \Diamond !A \lif \Box (!A \lif !B)$
      & $\Box \Box !A \lif \Box !A$ \\
      $\Diamond (!A \lor !B) \liff (\Diamond !A \lor
      \Diamond !B)$
      & $\Box \Diamond !A \lif \Diamond \Box !A$. \\
      \hline
    \end{tabular}
    \caption{માન્ય અને અમાન્ય સૂત્ર-ઢાંચા.}
    \ollabel{tab:valid-invalidSchemas}
\end{table}
\sourcecorrection{OLMD-010}{સ્થિર અંગ્રેજી મૂળના કોષ્ટક-શીર્ષકમાં
``and (or?) invalid'' એવો અનિર્ણીત સંપાદકીય કૌંસ છે.
બંને સ્તંભની માન્યતા અને અમાન્યતા સાન્ત વિરોધી
નિદર્શનો સામે તપાસીને ગુજરાતી શીર્ષકમાં તે કૌંસ દૂર કર્યો છે.}

\begin{prob}%\ollabel{ex:in/validSchemas}
  \olref[nml][syn][sch]{tab:valid-invalidSchemas}ના
  પ્રથમ સ્તંભના સૂત્ર-ઢાંચા માન્ય છે અને બીજા સ્તંભના
  સૂત્ર-ઢાંચા માન્ય નથી તે સાબિત કરો.
\end{prob}

\begin{prob}
  નીચેના સૂત્ર-ઢાંચા માન્ય છે કે અમાન્ય, તે નક્કી કરો:
  \begin{enumerate}
  \item $(\Diamond !A \lif \Box !B) \lif (\Box !A \lif \Box
    !B)$;
  \item $\Diamond(!A \lif !B) \lor \Box(!B \lif !A)$.
  \end{enumerate}
\end{prob}

\begin{prob}
  નીચેના દરેક સૂત્ર-ઢાંચા માટે એવું નિદર્શન $\mModel{M}$
  શોધો કે $\mModel{M}$માં તે સૂત્રનો દરેક દૃષ્ટાંત સત્ય હોય:
  \begin{enumerate}
  \item $p \lif \Diamond\Diamond p$;
  \item $\Diamond p \lif \Box p$.
  \end{enumerate}
\end{prob}
% END OLP-0417
\endgroup

\begingroup
\makeatletter\def\ol@sectioncs{section}\makeatother

% BEGIN OLP-0418
\olfileid{nml}{syn}{ent}

\olsection{ફલિતતા}
\phantomsection\label{guunit:OLP-0418}


\begin{explain}
  જગતમાં સત્યતાની વ્યાખ્યાના આધારે આપણે
  સૂત્રો વચ્ચે ફલિતતાનો સંબંધ વ્યાખ્યાયિત કરી
  શકીએ. સૂત્ર~$!B$માંથી~$!A$ ત્યારે અને માત્ર
  ત્યારે જ ફલિત થાય છે જ્યારે $!B$ સત્ય હોય તે દરેક
  સંજોગમાં $!A$ પણ સત્ય હોય. અહીં ``દરેક સંજોગમાં''નો
  અર્થ ``આપણે ગમે તે નિદર્શન લઈએ'' અને ``તે
  નિદર્શનનું ગમે તે જગત લઈએ'' બંને છે.
\end{explain}

\begin{defn}
  જો $\Gamma$ એ સૂત્રોનો ગણ હોય અને $!A$
  સૂત્ર હોય, તો $\Gamma$માંથી~$!A$
  \emph{ફલિત થાય છે}, સંકેતમાં $\Gamma \Entails !A$,
  ત્યારે અને માત્ર ત્યારે જ જ્યારે દરેક નિદર્શન
  $\mModel{M} = \tuple{W, R, V}$ અને જગત $w \in W$
  માટે, દરેક $!B \in \Gamma$ માટે $\mSat{M}{!B}[w]$
  હોય તો $\mSat{M}{!A}[w]$ હોય. જો $\Gamma$માં એક જ
  સૂત્ર~$!B$ હોય, તો $!B \Entails !A$ લખીએ.
\end{defn}

\begin{ex}
  એક સૂત્રમાંથી બીજું ફલિત થાય છે તે બતાવવા
  $\mSat{M}{}[w]$ની વ્યાખ્યા વાપરી બધાં નિદર્શનો
  વિશે વિચારવું પડે. દાખલા તરીકે,
  $p \lif \Diamond p \Entails \Box\lnot p \lif \lnot p$
  બતાવવા આ રીતે દલીલ કરી શકાય: નિદર્શન
  $\mModel{M} = \tuple{W, R,
    V}$ અને $w \in W$ લો, અને $\mSat{M}{p \lif \Diamond
    p}[w]$ માનો. $\mSat{M}{\Box\lnot p \lif \lnot
    p}[w]$ બતાવવું છે. તે ખોટું છે એમ માનો. તો
  $\mSat{M}{\Box\lnot p}[w]$ અને
  $\mSat/{M}{\lnot p}[w]$. $\mSat/{M}{\lnot p}[w]$
  હોવાથી $\mSat{M}{p}[w]$. ધારણા મુજબ
  $\mSat{M}{p \lif \Diamond p}[w]$ હોવાથી
  $\mSat{M}{\Diamond p}[w]$. $\mSat{M}{\Diamond
    p}[w]$ની વ્યાખ્યા મુજબ $Rww'$ ધરાવતું એવું કોઈ
  $w'$ છે કે $\mSat{M}{p}[w']$. પરંતુ
  $\mSat{M}{\Box \lnot p}[w]$ પણ હોવાથી
  $\mSat{M}{\lnot p}[w']$; આ વિરોધાભાસ છે.

  કોઈ સૂત્ર~$!B$માંથી બીજું~$!A$ ફલિત થતું
  નથી તે બતાવવા વિરોધી દૃષ્ટાંત આપવો પડે: એવું
  નિદર્શન $\mModel{M} = \tuple{W, R,
    V}$, જેમાં કોઈ જગત~$w \in W$ પર $\mSat{M}{!B}[w]$
  હોય પણ $\mSat/{M}{!A}[w]$ હોય. હવે બતાવીએ કે
  $p \lif \Diamond p \Entails/
  \Box p \lif p$. \olref{fig:counterex}નું નિદર્શન લો.
  તેમાં $\mSat{M}{\Diamond p}[w_1]$ અને તેથી
  $\mSat{M}{p \lif
    \Diamond p}[w_1]$. પરંતુ $\mSat{M}{\Box p}[w_1]$
  હોવા છતાં $\mSat/{M}{p}[w_1]$, તેથી
  $\mSat/{M}{\Box p \lif p}[w_1]$.
  \begin{figure}
  \begin{center}
    \begin{tikzpicture}[modal]
      \node[world] (w1) [label=right:\mFalse{p}]{$w_1$}; 
      \node[world] (w2) [label=right:\mTrue{p}, above left=of w1]{$w_2$}; 
      \node[world] (w3) [label=right:\mTrue{p}, above right=of w1] {$w_3$};
      \draw[->] (w1) to (w2);
      \draw[->] (w1) to (w3);
    \end{tikzpicture}
  \end{center}
  \caption{$p \lif \Diamond p \Entails \Box p \lif p$નો વિરોધી દૃષ્ટાંત.}
  \ollabel{fig:counterex}
  \end{figure}

  ઘણીવાર બહુ સરળ વિરોધી દૃષ્ટાંત પૂરતો હોય છે.
  $W' = \{w\}$, $R' = \emptyset$ અને $V'(p) =
  \emptyset$ ધરાવતું નિદર્શન $\mModel{M'} =
  \tuple{W', R', V'}$ પણ વિરોધી દૃષ્ટાંત છે:
  $\mSat/{M'}{p}[w]$ હોવાથી
  $\mSat{M'}{p \lif \Diamond p}[w]$. જગત~$w$માંથી
  કોઈ જગત સુલભ નથી, તેથી $\mSat{M'}{\Box p}[w]$;
  અને આમ $\mSat/{M'}{\Box p
    \lif p}[w]$.
\end{ex}
\sourcecorrection{OLMD-011}{સ્થિર અંગ્રેજી મૂળમાં
એક-જગત નિદર્શનને ત્રિકને બદલે ગણ-કૌંસમાં
લખાયું છે; નિદર્શનની વ્યાખ્યા પ્રમાણે ક્રમિત
ત્રિકનો સંકેત મૂક્યો છે.}
\sourcecorrection{OLMD-012}{અગાઉ મૂળ શીર્ષકમાં
ફલિતતાનું નિષેધ-ચિહ્ન ખૂટે છે એવું ખોટું માન્યું હતું.
મૂળનું ``ફલિતતાનો વિરોધી દૃષ્ટાંત'' શીર્ષક સાચું છે;
અહીં મૂળનું નિષેધ-ચિહ્ન વિનાનું સૂત્ર ફરી મૂક્યું છે.
આ મૂળની ભૂલ નહિ, અગાઉની સંપાદકીય ભૂલની પાછી ખેંચણી છે.}

\begin{prob}
  $\Box (!A \land !B) \Entails \Box !A$ બતાવો.
\end{prob}

\begin{prob}
  $\Box (p \lif q) \Entails/ p \lif \Box q$ અને $p \lif \Box
  q \Entails/\Box (p \lif q)$ બતાવો.
\end{prob}
% END OLP-0418
\endgroup


\OLEndChapterHook


\begin{editorial}
આગળનો સંબંધિત અંશ મૂળના સક્રિય આયાતક્રમમાં નથી. તેને અહીં સ્પષ્ટ વૈકલ્પિક સંદર્ભમાં જાળવ્યો છે.
\end{editorial}
\begingroup
\makeatletter\def\ol@sectioncs{section}\makeatother

% BEGIN OLP-0654
\olfileid{nml}{syn}{nor}

\olsection{સામાન્ય મોડલ તર્કો}
\phantomsection\label{guunit:OLP-0654}


મોડલ તર્કને યોગ્ય રીતે સુવ્યવસ્થિત
સૂત્રોનો ગણ માની શકાય.
કેટલાક મોડલ તર્કો ખાસ રસના
છે: તેમને રસપ્રદ અર્થઘટનો અને
ઉપયોગો છે, પરંતુ મુખ્યત્વે તેઓ
સૈદ્ધાંતિક રીતે સારી રીતે
સમજાયેલા છે. તેમને કહેવાતા
સામાન્ય મોડલ તર્કો કહે છે.

\begin{defn}\ollabel{defn:modal-logic}
\emph{મોડલ તર્ક}~$\Log L$
એ મૂળભૂત મોડલ ભાષાનાં
સૂત્રોનો એવો ગણ છે જે
\begin{enumerate}
\item બધી વિધાનાત્મક પુનરુક્તિઓ
  ધરાવે છે,
\item એકરૂપ પ્રતિસ્થાપન હેઠળ
  બંધ છે, અને
\item મોડસ પોનેન્સ હેઠળ બંધ છે:
  જો $!A$ અને $(!A \lif
  !B) \in \Log L$ હોય,
  તો $!B \in \Log L$ પણ.
\end{enumerate}
\end{defn}

\begin{ex}
ખૂબ રસપ્રદ ન હોવા છતાં
આ વ્યાખ્યા સંતોષતા,
અને તેથી મોડલ તર્કો
ગણાતા, સૂત્રોના
કેટલાક ગણ આ છે:
\begin{enumerate}
\item બધી પુનરુક્તિઓનો ગણ;
\item બધા સૂત્રોનો ગણ
  (અસુસંગત મોડલ તર્ક).
\end{enumerate}
\end{ex}

\begin{defn}\label{defn:normal-modal-logic}
મોડલ તર્ક~$\Log L$
\emph{સામાન્ય} ત્યારે અને
માત્ર ત્યારે જ જ્યારે
\begin{enumerate}
\item તે નીચેનાં બંને
  સૂત્રો ધરાવે:
\begin{align}
\tag{K} \Box(p \lif q) & \lif (\Box p \lif \Box q) \\
\tag{Dual} \Diamond p & \liff \lnot \Box \lnot p
\end{align}
\item અને તે આવશ્યકીકરણ
  હેઠળ બંધ હોય: જો
  $!A \in
  \Log L$ હોય,
  તો $\Box !A \in \Log L$
  પણ હોય.
\end{enumerate}
\end{defn}

મોડલ તર્કો વ્યાખ્યાયિત કરવાની
બે મુખ્ય રીતો આપણે જોઈશું.
એક રીત સાબિતી-સૈદ્ધાંતિક
છે: અમુક નિષ્પત્તિ
પ્રણાલીઓમાં નિષ્પન્ન થઈ
શકતાં સૂત્રોનો ગણ
લેવો. બીજી રીત
નિદર્શ-સૈદ્ધાંતિક છે:
નિદર્શોના અમુક વર્ગોમાં
પ્રામાણ્ય ધરાવતાં
સૂત્રોના ગણ લેવા.
આ સામાન્ય વિધાનાત્મક
તર્કની, અને તે બાબતે
પ્રથમ-ક્રમ તર્કની પણ,
પરિસ્થિતિને મળતું આવે
છે. ત્યાં પણ
સૂત્રોનો ગણ
બે રીતે અલગ તારવી
શકીએ: બધી સત્યમૂલ્ય
નિયુક્તિઓમાં સાચી
પુનરુક્તિઓના ગણ તરીકે,
અથવા કોઈ નિષ્પત્તિ
પ્રણાલીમાં બધાં
નિષ્પન્ન કરી શકાય એવું
સૂત્રોના ગણ
તરીકે. સામાન્ય
મોડલ તર્કો માટે
આ બે રીતો
સંબંધાત્મક નિદર્શો
અને સ્વયંસિદ્ધિમૂલક
(હિલ્બર્ટ-પ્રકારની)
સાબિતી પ્રણાલીઓ
વડે શક્ય બને છે.
% END OLP-0654
\endgroup
% END OLP-0408
\endgroup


\begingroup
\makeatletter\def\ol@sectioncs{section}\makeatother

% BEGIN OLP-0419
\olchapter{nml}{frd}{ફ્રેમ-વ્યાખ્યાયિતતા}
\olchapterid{nml}{frd}
\phantomsection\label{guunit:OLP-0419}


\begingroup
\makeatletter\def\ol@sectioncs{section}\makeatother

% BEGIN OLP-0420
\olfileid{nml}{frd}{int}

\olsection{પરિચય}
\phantomsection\label{guunit:OLP-0420}


મોડલ તર્કશાસ્ત્રીઓને રસ પડે એવો એક પ્રશ્ન સુલભતા
સંબંધ અને તે સંબંધ ધરાવતા નિદર્શનોમાં ચોક્કસ
સૂત્રોની સત્યતા વચ્ચેના સંબંધનો છે. ઉદાહરણ
તરીકે, માનો કે સુલભતા સંબંધ સ્વવાચક છે, એટલે કે
દરેક $w
\in W$ માટે $Rww$. બીજા શબ્દોમાં, દરેક જગત
પોતામાંથી સુલભ છે. તેનો અર્થ એ કે કોઈ જગત~$w$માં
$\Box !A$ સત્ય હોય ત્યારે $w$ પોતે પણ એવા સુલભ
જગતોમાંનું એક છે જ્યાં~$!A$ સત્ય હોવું પડે. તેથી
$\mModel{M}$નો સુલભતા સંબંધ~$R$ સ્વવાચક હોય, તો
કોઈપણ જગત $w$ અને સૂત્ર $!A$ લઈએ, ત્યાં $\Box !A
\lif !A$ સત્ય છે (બીજા શબ્દોમાં, સૂત્ર-ઢાંચો $\Box p \lif
p$ અને તેના બધા પ્રતિસ્થાપન દૃષ્ટાંતો~$\mModel{M}$માં
સત્ય છે).

પરંતુ વિપરીત વિધાન ખોટું છે. દાખલા તરીકે,
$\Box p \lif p$ એ~$\mModel{M}$માં સત્ય હોય તો $R$
સ્વવાચક છે એવું નથી. એક જ જગત~$w$ ધરાવતું એવું
અસ્વવાચક નિદર્શન~$\mModel{M}$ સરળતાથી મળે જેમાં
$\Box p \lif
p$ દરેક જગતમાં સત્ય છે: તે પોતામાંથી સુલભ
ન હોય, પણ $w \in V(p)$ હોય એવું નિદર્શન લો.
$p$નું સત્યમૂલ્ય યોગ્ય રીતે પસંદ કરીને અસ્વવાચક
નિદર્શનમાં $\Box !A \lif !A$ સત્ય બનાવી શકીએ છીએ.

ઉપાય એ છે કે સમીકરણમાંથી ચલ-નિમણૂક~$V$ને દૂર કરીએ.
જો~$\mModel{M}$ના બધા જગતોમાં $\Box p \lif p$ સત્ય
હોવાની માગણી~$V(p)$માં કયાં જગતો છે તેનાથી સ્વતંત્ર
રીતે કરીએ, તો $R$ સ્વવાચક હોવું જરૂરી છે. કેમ કે
કોઈપણ અસ્વવાચક નિદર્શનમાં ઓછામાં ઓછું એક જગત~$w$
એવું હશે કે $Rww$ ન હોય. જો $V(p) =
W \setminus \{w\}$ મૂકીએ, તો $p$ એ~$w$ સિવાયનાં
બધાં જગતોમાં, તેથી~$w$માંથી સુલભ બધા જગતોમાં,
સત્ય હશે ($w$ એ~$w$માંથી સુલભ નથી અને $p$ માત્ર
$w$માં જ અસત્ય છે). બીજી બાજુ, $p$ એ $w$માં
અસત્ય છે, તેથી $\Box
p \lif p$ એ~$w$માં અસત્ય છે.

આથી નિમણૂક વિનાની નિદર્શન-સંરચનાઓ માટે સંકેત
રજૂ કરવો યોગ્ય લાગે છે: તેમને \emph{ફ્રેમ} કહીએ.
ફ્રેમ $\mModel{F}$ એ માત્ર જગતોના ગણ અને તેમના
સુલભતા સંબંધની જોડ $\tuple{W, R}$ છે. ત્યારે
દરેક નિદર્શન $\tuple{W, R, V}$ આ ફ્રેમ
$\tuple{W, R}$ \emph{પર આધારિત} છે એમ કહીએ.
વિપરીત રીતે, ફ્રેમ તેના પર આધારિત નિદર્શનોનો
વર્ગ નક્કી કરે છે; અને ફ્રેમોના વર્ગથી તે વર્ગની
કોઈપણ ફ્રેમ પર આધારિત નિદર્શનોનો વર્ગ નક્કી થાય છે.
વળી, ફ્રેમમાં સૂત્રની \emph{માન્યતા} માટે
$\mModel{F} \Entails !A$ની વ્યાખ્યા આપી શકીએ: આનો
અર્થ~$\mModel{F}$ પર આધારિત દરેક $\mModel{M}$ માટે
$\mSat{M}{!A}$ છે.

આ સંકેતથી સૂત્રો અને ફ્રેમોના વર્ગો વચ્ચે
અનુરૂપતા સંબંધો સ્થાપી શકીએ: દાખલા તરીકે,
$\mModel{F} \Entails \Box p \lif p$ ત્યારે અને માત્ર
ત્યારે જ જ્યારે $\mModel{F}$ સ્વવાચક હોય.
% END OLP-0420
\endgroup

\begingroup
\makeatletter\def\ol@sectioncs{section}\makeatother

% BEGIN OLP-0421
\olfileid{nml}{frd}{acc}

\olsection{સુલભતા સંબંધોના ગુણધર્મો}
\phantomsection\label{guunit:OLP-0421}


ઘણાં મોડલ સૂત્રો સુલભતા સંબંધના સરળ, અને
પરિચિત પણ, ગુણધર્મોને લક્ષણરૂપ નીવડે છે. એક
દિશામાં તેનો અર્થ એ છે કે જે નિદર્શનનો સંબંધ
આપેલો ગુણધર્મ ધરાવે છે તેમાં અનુરૂપ સૂત્ર
(અને તેના બધા પ્રતિસ્થાપન દૃષ્ટાંતો) સત્ય હોય છે.
આપણે સુલભતા સંબંધોના પાંચ શાસ્ત્રીય પ્રકારો
અને તેઓ જે સૂત્રોની સત્યતા સુનિશ્ચિત કરે છે
તેનાથી શરૂઆત કરીએ.

\begin{thm}\ollabel{thm:soundschemas}
  $\mModel{M} = \tuple{W, R, V}$ નિદર્શન લો.
  જો $R$ એ \olref{tab:five}ની ડાબી બાજુનો ગુણધર્મ
  ધરાવે, તો જમણી બાજુના સૂત્રનો દરેક
  દૃષ્ટાંત~$\mModel{M}$માં સત્ય છે.
\end{thm}

\begin{table}[t]
    \begin{tabular}{| l || l |}
      \hline
      {\emph{જો $R$ \dots હોય}} & {\emph{તો \dots એ~$\mModel{M}$માં સત્ય છે:}} \\
      \hline \hline
      \emph{સિરિયલ}: $\forall u \exists v Ruv$ & \hbox to.43\textwidth
           {$\Box p \lif \Diamond p$ \hfill (\Ax{D})} \\
      \hline
      \emph{સ્વવાચક}: $\forall w Rww$  
      & \hbox to.43\textwidth
          {$\Box p \lif p$ \hfill (\Ax{T})} \\
      \hline
      \emph{સંમિત}: &  \hbox to .43\textwidth
           {$p \lif \Box\Diamond p$ \hfill (\Ax{B})} \\
           $\forall u\forall v(Ruv \lif Rvu)$ & \\
      \hline
      \emph{પરંપરિત}: & \hbox to.43\textwidth
           {$\Box p \lif \Box \Box p$ \hfill (\Ax{4})} \\
      $\forall u \forall v \forall w ((Ruv \land Rvw) \lif Ruw)$ & \\
      \hline 
      \emph{યુક્લિડિયન}: & \hbox to .43\textwidth
        {$\Diamond p \lif \Box\Diamond p$ \hfill (\Ax{5})} \\
      $\forall w \forall u \forall v ((Rwu \land Rwv) \lif Ruv)$ &\\
      \hline
    \end{tabular}
    \caption{અનુરૂપતાનાં પાંચ તથ્યો.}
    \ollabel{tab:five}
  \end{table} 


\begin{proof}
  અહીં \Ax{B}નો કિસ્સો જોઈએ: સૂત્ર-ઢાંચો નિદર્શનમાં
  સત્ય છે તે બતાવવા તેના બધા દૃષ્ટાંતો નિદર્શનના બધા
  જગતોમાં સત્ય છે તે બતાવવું પડે. તેથી \Ax{B}નો
  કોઈ દૃષ્ટાંત $!A \lif \Box\Diamond !A$ લો અને
  $w \in W$ કોઈપણ જગત લો. $\Box \Diamond
  !A$ એ $w$માં સત્ય છે તે બતાવવા $!A$ પૂર્વપક્ષ
  $w$માં સત્ય છે એમ માનો. એટલે $w$માંથી સુલભ
  દરેક $w'$માં $\Diamond !A$ સત્ય છે તે બતાવવું છે.
  હવે $Rww'$ ધરાવતા કોઈપણ $w'$ માટે સંમિતતાની
  ધારણાથી $Rw'w$ પણ છે (જુઓ \olref{fig:Bsymm}).
  $\mSat{M}{!A}[w]$ હોવાથી
  $\mSat{M}{\Diamond !A}[w']$ છે. $Rww'$ ધરાવતું
  $w'$ ગમે તે હતું, તેથી
  $\mSat{M}{\Box\Diamond!A}[w]$.

  બાકીના કિસ્સા અભ્યાસપ્રશ્ન તરીકે છોડીએ છીએ.
\end{proof}

\begin{prob}
  \olref[nml][frd][acc]{thm:soundschemas}નું પ્રમાણ
  પૂરું કરો.
\end{prob}

\begin{figure}
  \begin{center}
    \begin{tikzpicture}[modal]
      \node[world] (w1) [label={below:$\mSat{{}}{\formula{A}}$\\
          $\mSat{{}}{\Box\Diamond \formula{A}}$}] {$w$}; 
      \node[world] (w2) [label=below:{$\mSat{{}}{\Diamond \formula{A}}$},
        right=of w1] {$w'$}; 
      \draw[->, bend left] (w1) to (w2); 
      \draw[->, bend left] (w2) to (w1); 
    \end{tikzpicture}
  \end{center}
\caption{સંમિતતાથી થતી દલીલ.}
\ollabel{fig:Bsymm}
\end{figure}

ધ્યાન આપો કે \olref{thm:soundschemas}નાં વિપરીત
વિધાનો સાચાં નથી: કોઈ નિદર્શનમાં સૂત્ર-ઢાંચો
સત્ય હોય તો તે નિદર્શનનો સુલભતા સંબંધ અનુરૂપ
ગુણધર્મ ધરાવે જ એવું નથી. \Ax{T} અને સ્વવાચક
નિદર્શનોના કિસ્સામાં \Ax{T} પોતે નિદર્શનમાં ખોટું
પડે એવો દાખલો સરળતાથી મળે: $W = \{w\}$ અને
$V(p) = \emptyset$ લો, તથા સુલભતા સંબંધ ખાલી લો.
ત્યારે $R$ સ્વવાચક નથી, પણ
$\mSat{M}{\Box
  p}[w]$ અને $\mSat/{M}{p}[w]$. જોકે અહીં
\Ax{T}નો માત્ર એક દૃષ્ટાંત~$\mModel{M}$માં
ખોટો છે; $\Box \lnot p \lif \lnot p$ જેવા
બીજા દૃષ્ટાંતો સાચા છે. \Ax{T}નો
\emph{દરેક પ્રતિસ્થાપન દૃષ્ટાંત}~$\mModel{M}$માં
સત્ય હોય અને $\mModel{M}$ સ્વવાચક ન હોય એવાં
નિદર્શનો આપવાં વધુ મુશ્કેલ છે. પરંતુ એવાં
નિદર્શનો પણ છે:
\sourcecorrection{OLMD-013}{સ્થિર અંગ્રેજી મૂળના
એક-જગત વિરોધી નિદર્શનમાં સુલભતા સંબંધનો
ઉલ્લેખ કર્યા વિના તે અસ્વવાચક હોવાનું કહે છે.
અહીં ખાલી સંબંધ સ્પષ્ટ કર્યો છે; નહિતર દાવો
આપેલી માહિતીમાંથી ફલિત થતો નથી.}

\begin{prop}\ollabel{prop:reflexive}
  $\mModel{M} = \tuple{W, R, V}$ એવું નિદર્શન
  લો કે $W = \{u, v
  \}$ અને જગતો $u$ તથા $v$ સંબંધ $R$થી
  જોડાયેલાં છે: એટલે કે $Ruv$ અને $Rvu$ બંને છે;
  આ બે સિવાય સંબંધમાં બીજી કોઈ જોડી નથી.
  વળી, બધા $p$ માટે $u \in V(p) \Leftrightarrow v
  \in V(p)$ માનો. ત્યારે:
  \begin{enumerate}
  \item બધા $!A$ માટે: $\mSat{M}{!A}[u]$ ત્યારે અને માત્ર ત્યારે
    જ જ્યારે $\mSat{M}{!A}[v]$ (અહીં $!A$ પર આગમન વાપરો).
  \item \Ax{T}નો દરેક દૃષ્ટાંત $\mModel{M}$માં સત્ય છે.
  \end{enumerate}
  $\mModel{M}$ સ્વવાચક નથી (હકીકતમાં તે
  \emph{અસ્વવાચક} છે), તેથી \Ax{T}ના કિસ્સામાં
  \olref{thm:soundschemas}નું વિપરીત વિધાન ખોટું છે
  (\olref{thm:soundschemas}માં જણાવેલા અન્ય કેટલાંક---જોકે બધાં નહીં---
  સૂત્ર-ઢાંચાઓ માટે પણ આવી જ દલીલો આપી શકાય).
\end{prop}
\sourcecorrection{OLMD-014}{સ્થિર અંગ્રેજી મૂળમાં
માત્ર બંને આડાં તીરોની ધારણા છે, છતાં અંતે
નિદર્શનને અસ્વવાચક કહે છે. સ્વ-તીરોને બહાર રાખતી
શરત સ્પષ્ટ કરી છે; તે વિના અંતિમ નિષ્કર્ષ આવતો નથી.}

\begin{prob}
  \olref[nml][frd][acc]{prop:reflexive}ના દાવા સાબિત કરો.
\end{prob}

આપણે પાંચ શાસ્ત્રીય સૂત્રો \Ax{D},
\Ax{T}, \Ax{B}, \Ax{4} અને \Ax{5} પર ધ્યાન કેન્દ્રિત
કરીશું; છતાં \olref{tab:anotherfive}માં સુલભતા
સંબંધોના થોડા વધુ ગુણધર્મો નોંધીએ. સુલભતા
સંબંધ~$R$ \emph{આંશિક વિધેયાત્મક} છે જો દરેક
જગતમાંથી વધુમાં વધુ એક જગત સુલભ હોય. દરેક
જગતમાંથી બરાબર એક જગત સુલભ હોય તો તેને
\emph{વિધેયાત્મક} કહીએ. (આમ, વિધેયાત્મક સંબંધો
બરાબર સિરિયલ અને આંશિક વિધેયાત્મક બંને હોય
તેવા સંબંધો છે.) તેમને ``વિધેયાત્મક'' કહે છે કારણ કે
સુલભતા સંબંધ (આંશિક) વિધેયની જેમ વર્તે છે.
સંબંધ \emph{નબળે સઘન} છે જો $Ruv$ હોય ત્યારે
$u$ અને~$v$ની ``વચ્ચે'' કોઈ~$w$ હોય. તેથી
નબળે સઘન સંબંધો એક અર્થમાં પરંપરિત સંબંધોથી
વિપરીત છે: પરંપરિત સંબંધમાં~$u$માંથી~$w$ થઈને
$v$ સુધી પહોંચી શકાય તો~$u$માંથી સીધા~$v$
સુધી પણ પહોંચી શકાય; નબળે સઘન સંબંધમાં
~$u$માંથી સીધા~$v$ સુધી પહોંચી શકાય તો
કોઈ~$w$ થઈને પણ ત્યાં પહોંચી શકાય.
સંબંધ \emph{નબળે દિશિત} છે જો કોઈ જગત~$w$માંથી
$u$ અને~$v$ બંને સુધી પહોંચી શકાય ત્યારે~$u$
અને~$v$ બંનેમાંથી એક જ જગત~$t$ સુધી પહોંચી
શકાય---આને ક્યારેક ``હીરા ગુણધર્મ'' અથવા
``અભિસરણ'' કહે છે.

\begin{table}[t]
    \begin{tabular}{| p{.50\textwidth} || p{.43\textwidth} |}
      \hline
      {\emph{જો $R$ \dots હોય}} & {\emph{તો \dots એ~$\mModel{M}$માં સત્ય છે:}} \\
      \hline \hline
      \emph{આંશિક વિધેયાત્મક}: & 
      \multirow{2}{*}{$\Diamond p \lif \Box p$} \\
      $ \forall w \forall u \forall v ((Rwu \land Rwv) \lif u =v )$ & \\
      \hline
      \multirow{2}{*}{\emph{વિધેયાત્મક}: $\forall w \exists v \forall u(Rwu \liff \eq[u][v]) $} 
      & \multirow{2}{*}{$\Diamond p \liff \Box p$} \\
      & \\
      \hline
      \emph{નબળે સઘન}: &  \multirow{2}{*}{$\Box \Box p \lif
        \Box p$} \\
      $\forall u\forall v(Ruv \lif \exists w(Ruw \land Rwv))$ & \\
      \hline
      \emph{નબળે જોડાયેલો}: &
      \multirow{3}{*}{\hbox to.43\textwidth{$
        \begin{array}{@{}l@{}}
          \Box ((p \land \Box p) \lif q) \lor {} \\
          \qquad \Box ((q \land \Box q) \lif p)
        \end{array}$ \hfill (\Ax{L})}
      } \\
      $\forall w \forall u \forall v ((Rwu \land Rwv) \lif {}$ &\\
      \qquad $(Ruv \lor u=v \lor Rvu))$ & \\
      \hline 
      \emph{નબળે દિશિત}: & \multirow{3}{*}{\hbox to .43\textwidth
        {$\Diamond\Box p \lif \Box\Diamond p$\hfill (\Ax{G})}} \\
      $\forall w \forall u \forall v ((Rwu \land Rwv) \lif {}$ & \\
      \qquad $\exists t (Rut \land Rvt))$ & \\
      \hline
    \end{tabular}
    \caption{અનુરૂપતાનાં વધુ પાંચ તથ્યો.}
    \ollabel{tab:anotherfive}
  \end{table} 

\begin{prob}
  $\mModel{M} = \tuple{W, R, V}$ નિદર્શન લો.
  જો $R$ એ \olref[nml][frd][acc]{tab:anotherfive}ની
  ડાબી બાજુના ગુણધર્મો ધરાવતું હોય, તો અનુરૂપ
  જમણી બાજુના સૂત્રનો દરેક દૃષ્ટાંત~$\mModel{M}$માં
  સત્ય છે તે બતાવો.
\end{prob}
% END OLP-0421
\endgroup

\begingroup
\makeatletter\def\ol@sectioncs{section}\makeatother

% BEGIN OLP-0422
\olfileid{nml}{frd}{fra}

\olsection{ફ્રેમો}
\phantomsection\label{guunit:OLP-0422}


\begin{defn}
  \emph{ફ્રેમ} એ જોડ $\mModel{F} = \tuple{W,R}$ છે, જ્યાં
  $W$ જગતોનો અખાલી ગણ અને $R$ એ~$W$ પરનો
  દ્વિસ્થાની સંબંધ છે. નિદર્શન $\mModel{M}$ ફ્રેમ
  $\mModel{F} = \tuple{W,R}$ \emph{પર આધારિત} છે
  ત્યારે અને માત્ર ત્યારે જ જ્યારે કોઈ નિમણૂક $V$
  માટે $\mModel{M} = \tuple{W, R, V}$ હોય.
\end{defn}

\begin{defn}
  $\mModel{F}$ ફ્રેમ હોય, તો $!A$ એ
  \emph{$\mModel{F}$માં માન્ય છે}, સંકેતમાં
  $\mModel{F} \Entails !A$, એમ કહીએ જો
  ફ્રેમ~$\mModel{F}$ પર આધારિત દરેક નિદર્શન~$\mModel{M}$
  માટે $\mSat{M}{!A}$ હોય.
  
  $\mClass{F}$ ફ્રેમોનો વર્ગ હોય, તો $!A$ એ
  \emph{$\mClass{F}$માં માન્ય છે}, સંકેતમાં
  $\mClass{F} \Entails !A$, ત્યારે અને માત્ર ત્યારે જ
  જ્યારે દરેક ફ્રેમ~$\mModel{F} \in \mClass{F}$ માટે
  $\mModel{F} \Entails !A$ હોય.
\end{defn}

ફ્રેમો રસપ્રદ હોવાનું કારણ એ છે કે સૂત્ર-ઢાંચા અને
સુલભતા સંબંધ~$R$ના ગુણધર્મો વચ્ચેની અનુરૂપતા
ફ્રેમોના સ્તરે છે, \emph{નિદર્શનોના સ્તરે નથી}.
ઉદાહરણ તરીકે, \Ax{T} બધાં સ્વવાચક નિદર્શનોમાં
સત્ય હોવા છતાં જેમાં \Ax{T} સત્ય હોય તે દરેક
નિદર્શન સ્વવાચક નથી. પરંતુ \Ax{T} બધાં સ્વવાચક
\emph{ફ્રેમો}માં \emph{માન્ય} છે અને, વધુમાં, જેમાં
\Ax{T} માન્ય હોય તે દરેક ફ્રેમ સ્વવાચક છે.

\begin{rem}
ફ્રેમોના વર્ગમાં માન્યતા એ નિદર્શનોના વર્ગમાં
માન્યતાનો વિશેષ કિસ્સો છે: $\mClass{F} \Entails !A$
ત્યારે અને માત્ર ત્યારે જ જ્યારે
$\mClass{C} \Entails !A$, જ્યાં $\mClass{C}$
એ~$\mClass{F}$ની કોઈ ફ્રેમ પર આધારિત બધા
નિદર્શનોનો વર્ગ છે.

સ્પષ્ટ છે કે કોઈ સૂત્ર અથવા સૂત્ર-ઢાંચો
માન્ય હોય, એટલે કે \emph{બધાં} નિદર્શનોના વર્ગ
સાપેક્ષે માન્ય હોય, તો તે ફ્રેમોના કોઈપણ
વર્ગ~$\mClass{F}$ સાપેક્ષે પણ માન્ય છે.
\end{rem}
% END OLP-0422
\endgroup

\begingroup
\makeatletter\def\ol@sectioncs{section}\makeatother

% BEGIN OLP-0423
\olfileid{nml}{frd}{def}

\olsection{ફ્રેમ-વ્યાખ્યાયિતતા}
\phantomsection\label{guunit:OLP-0423}


\olref[acc]{thm:soundschemas}નાં વિપરીત વિધાનો
નિદર્શનો માટે સાચાં નથી, છતાં ``નિદર્શન''ને બદલે
``ફ્રેમ'' મૂકીએ તો તે સાચાં છે: \olref[acc]{thm:soundschemas}માં
વિચારેલા ગુણધર્મો માટે સૂત્ર કોઈ \emph{ફ્રેમ}માં
માન્ય હોય તો તે ફ્રેમનો સુલભતા સંબંધ અનુરૂપ
ગુણધર્મ ધરાવે છે. તેથી વિચારેલાં સૂત્રો તે
ગુણધર્મ ધરાવતી ફ્રેમોના વર્ગોને \emph{વ્યાખ્યાયિત}
કરે છે.

\begin{defn}
  $\mClass{F}$ ફ્રેમોનો વર્ગ હોય, તો $!A$ એ
  $\mClass{F}$ને \emph{વ્યાખ્યાયિત કરે છે} એમ કહીએ
  ત્યારે અને માત્ર ત્યારે જ જ્યારે બધી અને માત્ર
  બધી ફ્રેમો~$\mModel{F} \in \mClass{F}$ માટે
  $\mModel{F} \Entails !A$ હોય.
\end{defn}

હવે ફ્રેમો માટે વ્યાખ્યાયિતતાના પૂર્ણ પરિણામો
સ્થાપીએ.

\begin{thm}\ollabel{thm:fullCorrespondence}
\olref[acc]{tab:five}ની જમણી બાજુનું સૂત્ર
ફ્રેમ~$\mModel{F}$માં માન્ય હોય, તો $\mModel{F}$
ડાબી બાજુનો ગુણધર્મ ધરાવે છે.
\end{thm}

\begin{proof}
  \begin{enumerate}
  \item માનો કે \Ax{D} એ $\mModel{F} = \tuple{W, R}$માં
    માન્ય છે, એટલે કે
    $\mModel{F} \Entails \Box p \lif \Diamond p$.
    ફ્રેમ~$\mModel{F}$ પર આધારિત નિદર્શન $\mModel{M} =
    \tuple{W, R, V}$ અને $w \in
    W$ લો. એવું કોઈ $v$ છે કે~$Rwv$ તે બતાવવું છે.
    નથી એમ માનો: તો કોઈપણ~$!A$ માટે
    $\mSat{M}{\Box !A}[w]$ અને
    $\mSat/{M}{\Diamond !A}[w]$ બંને છે, ખાસ કરીને~$p$
    માટે. પરંતુ ત્યારે $\mSat/{M}{\Box p \lif
      \Diamond p}[w]$, જે $\mModel{F}
    \Entails \Box p \lif \Diamond p$ની ધારણાથી વિરોધાભાસ છે.
  \item માનો કે \Ax{T} એ $\mModel{F}$માં માન્ય છે,
    એટલે કે $\mModel{F} \Entails \Box
    p \lif p$. $w \in W$ કોઈપણ જગત લો;
    $Rww$ બતાવવું છે. $u \in V(p)$ ત્યારે અને માત્ર
    ત્યારે જ મૂકો જ્યારે $Rwu$ હોય ($q$ એ $p$ સિવાય
    બીજું ચલ હોય ત્યારે $V(q)$ ગમે તેમ લો, જેમ કે
    $V(q) = \emptyset$). હવે
    $\mModel{M} = \tuple{W, R, V}$ લો. રચના મુજબ
    $Rwu$ ધરાવતા બધા $u$ માટે $\mSat{M}{p}[u]$,
    તેથી $\mSat{M}{\Box p}[w]$. પરંતુ ધારણા મુજબ
    $\Box p \lif p$ એ~$w$માં સત્ય છે, તેથી
    $\mSat{M}{p}[w]$; અને $V$ની વ્યાખ્યા મુજબ આ
    માત્ર~$Rww$ હોય તો જ શક્ય છે.
  \item પ્રતિધનાત્મક રીતે સાબિત કરીએ: માનો કે
    $\mModel{F}$ સંમિત નથી; તો \Ax{B}, એટલે કે
    $p \lif \Box\Diamond p$, એ $\mModel{F}= \tuple{W, R}$માં
    માન્ય નથી તે બતાવીએ. $\mModel{F}$ સંમિત ન હોય,
    તો એવાં $u$, $v \in W$ છે કે $Ruv$ છે પણ
    $Rvu$ નથી. $V$ એવું વ્યાખ્યાયિત કરો કે
    $w \in V(p)$ ત્યારે અને માત્ર ત્યારે જ જ્યારે
    $Rvw$ ન હોય (બાકીનાં માટે $V$ ગમે તેમ હોય).
    $\mModel{M} =\tuple{W, R,
      V}$ લો. હવે $V$ની વ્યાખ્યા મુજબ $Rvw$ ન હોય
    એવા બધા $w$ માટે $\mSat{M}{p}[w]$; ખાસ કરીને,
    $Rvu$ ન હોવાથી $\mSat{M}{p}[u]$. વળી,
    $Rvw$ ત્યારે અને માત્ર ત્યારે જ જ્યારે
    $w \notin V(p)$, તેથી $Rvw$ અને
    $\mSat{M}{p}[w]$ બંને હોય એવું કોઈ~$w$ નથી;
    પરિણામે $\mSat/{M}{\Diamond p}[v]$. $Ruv$
    પણ હોવાથી $\mSat/{M}{\Box\Diamond p}[u]$.
    તેથી $\mSat/{M}{p \lif
      \Box\Diamond p}[u]$, અને આમ $\Ax{B}$
    ફ્રેમ~$\mModel{F}$માં માન્ય નથી.
  \item માનો કે \Ax{4} એ $\mModel{F} = \tuple{W,R}$માં
    માન્ય છે, એટલે કે
    $\mModel{F} \Entails \Box p \lif \Box\Box p$;
    અને $u$, $v$, $w \in W$ એવાં કોઈપણ જગતો લો
    કે $Ruv$ અને $Rvw$. $Ruw$ બતાવવું છે.
    $V$ એવું વ્યાખ્યાયિત કરો કે $z \in V(p)$ ત્યારે
    અને માત્ર ત્યારે જ જ્યારે $Ruz$ હોય (બાકીનાં
    માટે $V$ ગમે તેમ હોય). હવે
    $\mModel{M} =\tuple{W, R, V}$ લો. $V$ની વ્યાખ્યા
    મુજબ $Ruz$ ધરાવતા બધા $z$ માટે
    $\mSat{M}{p}[z]$, તેથી $\mSat{M}{\Box p}[u]$.
    પરંતુ ધારણા મુજબ \Ax{4}, એટલે કે $\Box p
    \lif \Box \Box p$, એ $u$માં સત્ય છે, તેથી
    $\mSat{M}{\Box \Box
      p}[u]$. $Ruv$ અને $Rvw$ હોવાથી
    $\mSat{M}{p}[w]$; $V$ની વ્યાખ્યા પ્રમાણે આ
    માત્ર $Ruw$ હોય તો જ શક્ય છે, જેમ જોઈએ તેમ.
  \item પ્રતિધનાત્મક રીતે આગળ વધીએ: ફ્રેમ
    $\mModel{F} = \tuple{W, R}$ યુક્લિડિયન નથી એમ
    માની તે \Ax{5}ને ખોટું પાડે છે, એટલે કે
    $\mModel{F} \Entails/ \Diamond p \lif
      \Box\Diamond p$, તે બતાવીએ. એવાં જગતો
    $u$, $v$, $w \in W$ માનો કે $Rwu$ અને $Rwv$
    છે, પણ $Ruv$ નથી. $V$ એવું વ્યાખ્યાયિત કરો
    કે બધાં જગતો $z$ માટે $z \in V(p)$ ત્યારે અને
    માત્ર ત્યારે જ જ્યારે~$Ruz$ \emph{ન હોય}.
    $\mModel{M} =\tuple{W, R, V}$ લો. ત્યારે ધારણા
    મુજબ $\mSat{M}{p}[v]$ અને $Rwv$ હોવાથી
    $\mSat{M}{\Diamond p}[w]$ પણ છે. જોકે
    $Ruy$ અને $\mSat{M}{p}[y]$ બંને હોય એવું કોઈ
    જગત $y$ નથી, તેથી
    $\mSat/{M}{\Diamond
      p}[u]$. $Rwu$ હોવાથી પછી
    $\mSat/{M}{\Box\Diamond
      p}[w]$, એટલે \Ax{5}, $\Diamond p \lif \Box\Diamond p$,
    એ~$w$માં ખોટું પડે છે.
  \end{enumerate}
\end{proof}
\sourcecorrection{OLMD-015}{સ્થિર અંગ્રેજી મૂળમાં
\Ax{D}ના પ્રમાણની એક સંતોષ-અભિવ્યક્તિમાં
જગતનો નિર્દેશ છૂટી ગયો છે; અન્ય જગત-સાપેક્ષ
અભિવ્યક્તિઓની જેમ ત્યાં પણ નિર્દેશ પૂર્યો છે.}
\sourcecorrection{OLMD-016}{\Ax{T}ના પ્રમાણમાં
સ્થિર અંગ્રેજી મૂળે $V(q)$ના ઉદાહરણનું બંધ
કૌંસ ગણિત-અભિવ્યક્તિની અંદર મૂકે છે;
તેને ગણિત-અભિવ્યક્તિની બહાર મૂક્યું છે.}

\Ax{D}ના પ્રમાણ અને અન્ય કિસ્સાઓ વચ્ચેનો એક
ભેદ ધ્યાનમાં આવશે: નિમણૂક~$V$નો તેમાં કોઈ
ઉલ્લેખ થયો નથી. હકીકતમાં, આપણે સાબિત કર્યું
કે $\mSat{M}{\Ax{D}}$ હોય તો $\mModel{M}$ સિરિયલ
છે. તેથી \Ax{D} માત્ર ફ્રેમોનો નહીં, પરંતુ સિરિયલ
\emph{નિદર્શનોનો} વર્ગ વ્યાખ્યાયિત કરે છે.

\begin{cor}\ollabel{prop:D-serial}
  જે નિદર્શનમાં \Ax{D} સત્ય હોય તે સિરિયલ છે.
\end{cor}

\begin{cor}
\olref[acc]{tab:five}ની જમણી બાજુનું દરેક
સૂત્ર ડાબી બાજુનો ગુણધર્મ ધરાવતી
ફ્રેમોનો વર્ગ વ્યાખ્યાયિત કરે છે.
\end{cor}

\begin{proof}
  \olref[acc]{thm:soundschemas}માં આપણે સાબિત કર્યું
  કે નિદર્શન ડાબી બાજુનો ગુણધર્મ ધરાવે તો જમણી
  બાજુનું સૂત્ર તેમાં સત્ય છે. તેથી ફ્રેમ~$\mModel{F}$
ડાબી બાજુનો ગુણધર્મ ધરાવે તો જમણી બાજુનું
  સૂત્ર એ~$\mModel{F}$માં માન્ય છે.
  \olref{thm:fullCorrespondence}માં આપણે વિપરીત
  વિધાનો સાબિત કર્યાં: જમણી બાજુનું સૂત્ર
  એ~$\mModel{F}$માં માન્ય હોય તો $\mModel{F}$ ડાબી
  બાજુનો ગુણધર્મ ધરાવે છે.
\end{proof}

\begin{prob}
જો \olref[nml][frd][acc]{tab:anotherfive}ની જમણી
બાજુનું સૂત્ર ફ્રેમ~$\mModel{F}$માં માન્ય હોય,
તો $\mModel{F}$ ડાબી બાજુનો ગુણધર્મ ધરાવે છે તે
બતાવો. તે માટે ડાબી બાજુનો ગુણધર્મ \emph{ન} ધરાવતી
ફ્રેમ લો અને યોગ્ય~$V$ વ્યાખ્યાયિત કરો જેથી
જમણી બાજુનું સૂત્ર કોઈ જગતમાં અસત્ય થાય.
\end{prob}

\olref{thm:fullCorrespondence} એ પણ બતાવે છે કે
ગુણધર્મો ભેગા કરી શકાય: દાખલા તરીકે, \Ax{B}
અને \Ax{4} બંને~$\mModel{F}$માં માન્ય હોય તો ફ્રેમ
સંમિત અને પરંપરિત બંને છે, વગેરે. કેટલાંક
ફ્રેમ-ગુણધર્મો ભેગાં ધરાવતી બધી ફ્રેમોમાં માન્ય
સૂત્રોના ગણ તરીકે ઘણાં મહત્ત્વનાં મોડલ
તર્કશાસ્ત્રો લક્ષણિત થાય છે; તેથી અનુરૂપ વ્યાખ્યાયિત
કરતાં સૂત્રો માન્ય હોય તેવી બધી ફ્રેમોમાં
માન્ય સૂત્રોના ગણ તરીકે પણ તેમને લક્ષણિત
કરી શકીએ. દાખલા તરીકે, શાસ્ત્રીય પ્રણાલી
\Log{S4} એ બધી સ્વવાચક અને પરંપરિત ફ્રેમોમાં
માન્ય સૂત્રોનો ગણ છે, એટલે કે \Ax{T}
અને~\Ax{4} બંને માન્ય હોય એવી બધી ફ્રેમોમાં
માન્ય સૂત્રોનો ગણ. \Log{S5} એ બધી સ્વવાચક,
સંમિત અને યુક્લિડિયન ફ્રેમોમાં માન્ય સૂત્રોનો ગણ છે,
એટલે કે \Ax{T}, \Ax{B} અને \Ax{5} ત્રણેય
માન્ય હોય એવી બધી ફ્રેમોમાં માન્ય સૂત્રોનો ગણ.

$R$ના ગુણધર્મો વચ્ચેના તાર્કિક સંબંધો સામાન્ય રીતે
તેમને વ્યાખ્યાયિત કરતાં સૂત્રો વચ્ચેના સંબંધોને
અનુરૂપ હોય છે. દાખલા તરીકે, દરેક સ્વવાચક સંબંધ
સિરિયલ છે; તેથી ફ્રેમમાં \Ax{T} માન્ય હોય તો~\Ax{D}
પણ માન્ય છે. (ધ્યાન આપો કે આ સંબંધ
\emph{ફલિતતાનો નથી}. $\mSat{M}{\Ax{T}}[w]$ હોય
ત્યારે હંમેશાં $\mSat{M}{\Ax{D}}[w]$ હોય એવું નથી.)
આવા કેટલાક સંબંધો નોંધીએ.

\begin{prop}\ollabel{prop:relation-facts}
  $R$ એ ગણ $W$ પરનો દ્વિસ્થાની સંબંધ હોય; તો:
  \begin{enumerate}
  \item $R$ સ્વવાચક હોય તો તે સિરિયલ છે.
  \item $R$ સંમિત હોય, તો તે પરંપરિત ત્યારે અને માત્ર
    ત્યારે જ છે જ્યારે તે યુક્લિડિયન હોય.
  \item $R$ સંમિત અથવા યુક્લિડિયન હોય તો તે
    નબળે દિશિત છે (તેમાં ``હીરા ગુણધર્મ'' છે).
  \item $R$ યુક્લિડિયન હોય તો તે નબળે જોડાયેલો છે.
  \item $R$ વિધેયાત્મક હોય તો તે સિરિયલ છે.
  \end{enumerate}
\end{prop}

\begin{prob}
  \olref[nml][frd][def]{prop:relation-facts} સાબિત કરો.
\end{prob}
% END OLP-0423
\endgroup

\begingroup
\makeatletter\def\ol@sectioncs{section}\makeatother

% BEGIN OLP-0424
\olfileid{nml}{frd}{fol}

\olsection{પ્રથમ-ક્રમ વ્યાખ્યાયિતતા}
\phantomsection\label{guunit:OLP-0424}


આપણે જોયું છે કે ફ્રેમોના સુલભતા સંબંધના કેટલાંક
ગુણધર્મો મોડલ સૂત્રોથી વ્યાખ્યાયિત કરી શકાય છે.
દાખલા તરીકે, ફ્રેમોની સંમિતતા સૂત્ર~\Ax{B},
$p \lif \Box\Diamond
p$થી વ્યાખ્યાયિત થાય છે. અત્યાર સુધી મળેલી બધી
શરતો એક દ્વિસ્થાની વિધેય ધરાવતી
પ્રથમ-ક્રમની ભાષામાંનાં સૂત્રોથી
વ્યક્ત કરી શકાય છે. દાખલા તરીકે, સંમિતતા
$\lforall[x][\lforall[y][(\Atom{Q}{x,y} \lif \Atom{Q}{y, x})]]$થી
એ અર્થમાં વ્યાખ્યાયિત થાય છે કે $\Domain{M} =
W$ અને $\Assign{Q}{M} = R$ ધરાવતી પ્રથમ-ક્રમની
સંરચના~$\Struct{M}$ એ ઉપરનું સૂત્ર ત્યારે અને
માત્ર ત્યારે જ સંતોષે છે જ્યારે $R$ સંમિત હોય.
આથી નીચેની વ્યાખ્યા સૂચાય છે:

\begin{defn}
  ફ્રેમોનો વર્ગ~$\mClass{F}$ \emph{પ્રથમ-ક્રમમાં વ્યાખ્યાયિત}
  છે, જો એક દ્વિસ્થાની વિધેય~$Q$ ધરાવતી
  પ્રથમ-ક્રમની ભાષામાં એવું વાક્ય~$!A$
  હોય કે $\mModel{F} =
  \tuple{W, R} \in \mClass{F}$ ત્યારે અને માત્ર ત્યારે જ
  જ્યારે $\Domain{M} = W$ અને $\Assign{Q}{M}
  = R$ ધરાવતી પ્રથમ-ક્રમની સંરચના~$\Struct{M}$માં
  $\Sat{M}{!A}$ હોય.
\end{defn}

અત્યાર સુધી વિચારેલા ગુણધર્મો અને તેમને વ્યાખ્યાયિત
કરતાં મોડલ સૂત્રો અપવાદરૂપ છે. દરેક
સૂત્ર ફ્રેમોના પ્રથમ-ક્રમમાં વ્યાખ્યાયિત વર્ગને
વ્યાખ્યાયિત કરતું નથી; અને ફ્રેમોના દરેક પ્રથમ-ક્રમમાં
વ્યાખ્યાયિત વર્ગને કોઈ મોડલ સૂત્રથી વ્યાખ્યાયિત
કરી શકાતો નથી.

પહેલા દાવાનું વિરોધી દૃષ્ટાંત લોબનું સૂત્ર છે:
\begin{equation}
\Box(\Box p \lif p) \lif \Box p. \tag{\Ax{W}}
\end{equation}
\Ax{W} એ પરંપરિત અને પ્રતિલોમ-સુસ્થાપિત ફ્રેમોનો
વર્ગ વ્યાખ્યાયિત કરે છે. સંબંધ સુસ્થાપિત હોય છે જો
$Rw_2w_1$, $Rw_3w_2$, \dots{} ધરાવતી અનંત શ્રેણી
$w_1$, $w_2$, \dots{} ન હોય. દાખલા તરીકે,
~$\Nat$ પરનો $<$ સંબંધ સુસ્થાપિત છે, જ્યારે
~$\Int$ પરનો $<$ સંબંધ એવો નથી. સંબંધ
પ્રતિલોમ-સુસ્થાપિત ત્યારે અને માત્ર ત્યારે જ છે
જ્યારે તેનો પ્રતિલોમ સુસ્થાપિત હોય. તેથી
પ્રતિલોમ-સુસ્થાપિત સંબંધોમાં $Rw_1w_2$,
$Rw_2w_3$, \dots{} ધરાવતી અનંત શ્રેણી
$w_1$, $w_2$, \dots{} હોતી નથી.

પરંતુ પરંપરિત પ્રતિલોમ-સુસ્થાપિત સંબંધોને
વ્યાખ્યાયિત કરતું કોઈ પ્રથમ-ક્રમનું સૂત્ર
નથી. કેમ કે માનો કે~$\Sat{M}{!F}$ ત્યારે અને માત્ર
ત્યારે જ જ્યારે $R =
\Assign{Q}{M}$ પરંપરિત અને પ્રતિલોમ-સુસ્થાપિત હોય.
સઘનતાની દલીલ માટે ભાષામાં નવા અચલ સંકેતો
$a_1$, $a_2$, \dots{} ઉમેરો. હવે $!A_n$ નીચેનું
સૂત્ર હોય, જ્યાં સૂચક ઓછામાં ઓછો બે છે:
\[
(\Atom{Q}{a_1,a_2} \land \dots \land
    \Atom{Q}{a_{n-1},a_{n}})
\]
હવે સૂત્રોનો ગણ લો:
\[
\Gamma = \{!F, !A_2, !A_3, \dots\}.
\]
$\Gamma$નો દરેક સાન્ત ઉપગણ સંતોષ્ય છે: ઉપગણમાં
કોઈ $!A_k$ હોય તો તેમાંનું સૌથી મોટું સૂચક $k$
લો; એવું એકેય ન હોય તો બે લો. પછી
$\Domain{M_k} = \{1, \dots, k\}$,
$\Assign{a_i}{M_k} = i$ અને $\Assign{Q}{M_k} = <$ મૂકો.
$\{1,
\dots, k\}$ પરનો $<$ સંબંધ પરંપરિત અને
પ્રતિલોમ-સુસ્થાપિત હોવાથી $\Sat{M_k}{!F}$.
રચના મુજબ બે કે વધુ હોય એવા બધા $i \le k$ માટે
$\Sat{M_k}{!A_i}$ પણ છે.
પ્રથમ-ક્રમ તર્કશાસ્ત્રના સઘનતા પ્રમેયથી $\Gamma$
કોઈ સંરચના~$\Struct{M}$માં સંતોષ્ય છે.
ધારણા મુજબ $\Sat{M}{!F}$ હોવાથી
$\Assign{Q}{M}$ પ્રતિલોમ-સુસ્થાપિત છે. પરંતુ
$\Assign{a_1}{M}$, $\Assign{a_2}{M}$,
\dots{} સ્પષ્ટ રીતે એવી અનંત શ્રેણી બનાવે છે
જે પ્રતિલોમ-સુસ્થાપિતતા પ્રતિબંધિત કરે છે.
\sourcecorrection{OLMD-017}{સ્થિર અંગ્રેજી મૂળમાં
શ્રેણી-સૂત્રનું $!A_1$ સૂચક અનિર્ધારિત છે: તેનું
પ્રથમ પગલું બીજા અચલથી શરૂ થાય છે. ગણને
બીજા સૂચકથી શરૂ કર્યો અને કોઈ શ્રેણી-સૂત્ર
વિના સાન્ત ઉપગણનો કિસ્સો પણ સ્પષ્ટ કર્યો છે.}
\sourcecorrection{OLMD-018}{સ્થિર અંગ્રેજી મૂળની
વ્યાખ્યામાં માત્ર દ્વિસ્થાની વિધેયચિહ્નવાળી ભાષા છે,
પરંતુ સઘનતાના પ્રમાણમાં નવા અચલ સંકેતો વપરાય છે.
પ્રમાણમાં ભાષાનો આ અચલ-વિસ્તાર સ્પષ્ટ કર્યો છે.}

બીજા દાવાનું વિરોધી દૃષ્ટાંત સાર્વત્રિકતાનો ગુણધર્મ
છે: દરેક $u$ અને $v$ માટે $Ruv$. સાર્વત્રિક
ફ્રેમો પ્રથમ-ક્રમના સૂત્ર
$\lforall[x][\lforall[y][\Atom{Q}{x,y}]]$થી
વ્યાખ્યાયિત થાય છે. પરંતુ એવો કોઈ મોડલ સૂત્ર
નથી જે બધી અને માત્ર સાર્વત્રિક ફ્રેમોમાં માન્ય હોય.
આ સ્વતંત્ર રીતે રસપ્રદ પરિણામનું ફળ છે: સાર્વત્રિક
ફ્રેમોમાં માન્ય સૂત્રો બરાબર સ્વવાચક, સંમિત અને
પરંપરિત ફ્રેમોમાં માન્ય સૂત્રો જ છે. સ્વવાચક,
સંમિત અને પરંપરિત છતાં અસાર્વત્રિક ફ્રેમો છે;
તેથી બધી સાર્વત્રિક ફ્રેમોમાં માન્ય દરેક સૂત્ર
કેટલીક અસાર્વત્રિક ફ્રેમોમાં પણ માન્ય હોય છે.
% END OLP-0424
\endgroup

\begingroup
\makeatletter\def\ol@sectioncs{section}\makeatother

% BEGIN OLP-0425
\olfileid{nml}{frd}{es5}

\olsection{સામ્ય સંબંધો અને \Log{S5}}
\phantomsection\label{guunit:OLP-0425}


મોડલ તર્કશાસ્ત્ર \Log{S5} એ બધી સાર્વત્રિક ફ્રેમોમાં
માન્ય સૂત્રોના ગણ તરીકે લક્ષણિત થાય છે;
એટલે કે દરેક જગત પોતાને સહિત દરેક જગતમાંથી
સુલભ છે. આવી સ્થિતિમાં $\Box$ અનિવાર્યતાને
અને $\Diamond$ શક્યતાને અનુરૂપ છે: $!A$ \emph{દરેક}
જગતમાં સત્ય હોય તો $\Box !A$ સત્ય છે, અને $!A$
\emph{કોઈક} જગતમાં સત્ય હોય તો $\Diamond !A$
સત્ય છે. \Log{S5}ને બધી સ્વવાચક, સંમિત અને
પરંપરિત ફ્રેમોમાં, એટલે કે બધા \emph{સામ્ય સંબંધો}
પર, માન્ય સૂત્રો તરીકે પણ લક્ષણિત કરી શકાય છે.

\begin{defn}
  $W$ પરનો દ્વિસ્થાની સંબંધ $R$ \emph{સામ્ય સંબંધ}
  ત્યારે અને માત્ર ત્યારે જ છે જ્યારે તે સ્વવાચક,
  સંમિત અને પરંપરિત હોય. $W$ પરનો સંબંધ $R$
  \emph{સાર્વત્રિક} ત્યારે અને માત્ર ત્યારે જ છે જ્યારે
  બધા $u,v \in
  W$ માટે $Ruv$ હોય.
\end{defn}

\Ax{T}, \Ax{B} અને \Ax{4} અનુક્રમે સ્વવાચક,
સંમિત અને પરંપરિત ફ્રેમોને લક્ષણિત કરતાં હોવાથી,
સુલભતા સંબંધ સામ્ય સંબંધ હોય તેવી ફ્રેમો બરાબર
એ છે જેમાં આ ત્રણેય સૂત્રો માન્ય છે.
સામ્ય સંબંધોને સૂત્રોનાં અન્ય સંયોજનો વડે
પણ લક્ષણિત કરી શકાય છે, કેમ કે તેમની વ્યાખ્યામાં
વાપરેલી શરતો $R$ પરની અન્ય પરિચિત શરતોનાં
સંયોજનોને સમકક્ષ છે.

\begin{prop}\ollabel{prop:equivalences}
  નીચેની શરતો સમકક્ષ છે:
  \begin{enumerate}
  \item $R$ સામ્ય સંબંધ છે;
  \item $R$ સ્વવાચક અને યુક્લિડિયન છે;
  \item $R$ સિરિયલ, સંમિત અને યુક્લિડિયન છે;
  \item $R$ સિરિયલ, સંમિત અને પરંપરિત છે.
  \end{enumerate}
\end{prop}

\begin{proof}
  અભ્યાસપ્રશ્ન.
\end{proof}

\begin{prob}
  નીચેનું બતાવી \olref[nml][frd][es5]{prop:equivalences}
  સાબિત કરો:
  \begin{enumerate}
  \item $R$ સંમિત અને પરંપરિત હોય તો યુક્લિડિયન છે.
  \item $R$ સ્વવાચક હોય તો સિરિયલ છે.
  \item $R$ સ્વવાચક અને યુક્લિડિયન હોય તો સંમિત છે.
  \item $R$ સંમિત અને યુક્લિડિયન હોય તો પરંપરિત છે.
  \item $R$ સિરિયલ, સંમિત અને પરંપરિત હોય તો સ્વવાચક છે.
  \end{enumerate}
  આ શરતો સમકક્ષ છે તે સાબિત કરવા આ કેમ પૂરતું છે
  તે સમજાવો.
\end{prob}

\olref{prop:equivalences} એ \olref[prf][prs]{prop:S5}નો
અર્થવિજ્ઞાનલક્ષી સમકક્ષ છે: તે એવી ફ્રેમોના મોડલ
તર્કશાસ્ત્રનું સમકક્ષ લક્ષણન આપે છે જેમનો $R$
સામ્ય સંબંધ છે (જેને પરંપરાગત રીતે \Log{S5} કહે છે).

સાર્વત્રિક અને સામ્ય સંબંધો વચ્ચે શું સંબંધ છે?
દરેક સાર્વત્રિક સંબંધ સામ્ય સંબંધ છે, છતાં દરેક
સામ્ય સંબંધ સાર્વત્રિક નથી. પરંતુ બધી સાર્વત્રિક
ફ્રેમોમાં માન્ય સૂત્રો બરાબર બધી સામ્ય-સંબંધવાળી
ફ્રેમોમાં માન્ય સૂત્રો જ છે.

\begin{prop}
  $R$ સામ્ય સંબંધ માનો, અને દરેક $w \in W$ માટે
  $w$નો \emph{સામ્ય વર્ગ} એ ગણ $[w] = \{w'\in W :
  Rww'\}$ વ્યાખ્યાયિત કરો. ત્યારે:
  \begin{enumerate}
  \item $w \in [w]$;
  \item દરેક સામ્ય વર્ગ $[w]$ પર $R$ સાર્વત્રિક છે;
  \item સામ્ય વર્ગોનો સંગ્રહ $W$ને પરસ્પર વિયુક્ત
    અને મળીને સંપૂર્ણ એવા ઉપગણોમાં વિભાજિત કરે છે.
  \end{enumerate}
\end{prop}

\begin{prop}\ollabel{prop:S5=univ}
  સૂત્ર~$!A$ એ એવી બધી ફ્રેમો
  $\mModel{F} =\tuple{W,R}$માં માન્ય છે જ્યાં $R$
  સામ્ય સંબંધ છે, ત્યારે અને માત્ર ત્યારે જ જ્યારે
  તે એવી બધી ફ્રેમો $\mModel{F} =\tuple{W,R}$માં
  માન્ય છે જ્યાં $R$ સાર્વત્રિક છે. તેથી સાર્વત્રિક
  ફ્રેમોનું તર્કશાસ્ત્ર બરાબર \Log{S5} છે.
\end{prop}

\begin{proof}
  $W$ પરનો સાર્વત્રિક સંબંધ~$R$ સામ્ય સંબંધ છે
  તે તરત ચકાસી શકાય છે. તેથી $!A$ એ $R$ સામ્ય
  સંબંધ હોય એવી બધી ફ્રેમોમાં માન્ય હોય, તો તે
  બધી સાર્વત્રિક ફ્રેમોમાં પણ માન્ય છે. બીજી
  દિશા માટે પ્રતિધનાત્મક દલીલ કરીએ: માનો કે
  સૂત્ર~$!B$ એ જગત $w$માં અસત્ય છે,
  જ્યાં નિદર્શન $\mModel{M} = \tuple{W, R, V}$
  એવી ફ્રેમ $\tuple{W,R}$ પર આધારિત છે જેમાં $R$
  એ $W$ પરનો સામ્ય સંબંધ છે. એટલે
  $\mSat/{M}{!B}[w]$. હવે નિદર્શન
  $\mModel{M}' = \tuple{W',
    R', V'}$ આ રીતે વ્યાખ્યાયિત કરો:
  \begin{enumerate}
  \item $W' = [w]$;
  \item $R'$ એ $W'$ પર સાર્વત્રિક છે;
  \item $V'(p) = V(p) \cap W'$.
  \end{enumerate}
  (અર્થાત્ \olref{fig:partition}ના ઘેરા વિસ્તારમાં
  નિદર્શન $\mModel{M}'$ના જગતોનો ગણ $W'$
  દર્શાવ્યો છે.) $W'$ પર $R$ અને $R'$ એકસરખા
  છે તે સરળતાથી જોઈ શકાય છે. પછી સૂત્રો
  પર આગમનથી બતાવી શકાય કે બધા $w' \in W'$ માટે:
  દરેક $!A$ માટે $\mSat{M'}{!A}[w']$ ત્યારે અને
  માત્ર ત્યારે જ જ્યારે $\mSat{M}{!A}[w']$ હોય
  ($W'
  \subseteq W$ હોવાથી આ વિધાન અર્થપૂર્ણ છે).
  ખાસ કરીને $\mSat/{M'}{!B}[w]$, એટલે $!B$
  સાર્વત્રિક ફ્રેમ પર આધારિત નિદર્શનમાં અસત્ય છે.
\end{proof}

\begin{figure}[t]
  \centering
  \begin{tikzpicture}[node distance=2cm, auto, thick]
    \clip [rounded corners] (0,0) -- (6,0) -- (6,4) -- (0,4) --  cycle;
    \begin{scope}
      \clip (0,0) -- (2,0)  .. controls (1.5,1.5) and (4.5,1.5)
      .. (4,4) -- (0,4) -- (0,0) -- cycle;
      \filldraw[fill=gray!40] (0,4) -- (0,2) .. controls (1.5,1.5)
      and (4.5,1.5) .. (4.5,0) -- (6,0) -- (6,4) -- (0,4) --  cycle;
    \end{scope}
    \draw [name path=line1] (2,0) .. controls (1.5,1.5) and (4.5,1.5) .. (4,4);
    \draw [name path=line2] (0,2)  .. controls (1.5,1.5) and (4.5,1.5) .. (4.5,0);
    \path [name intersections={of=line1 and line2}];
    \draw [rounded corners, use as bounding box, very thick] (0,0)
    -- (6,0) -- (6,4) -- (0,4) --  cycle; 
    \node at (2,3) {$[w]$} ;
    \node at (1,1) {$[u]$} ;
    \node at (3,0.75) {$[v]$} ;
    \node at (4.75,2) {$[z]$} ;
  \end{tikzpicture}
  \caption{સામ્ય વર્ગોમાં $W$નું વિભાજન.}
  \ollabel{fig:partition}
\end{figure}
% END OLP-0425
\endgroup

\begingroup
\makeatletter\def\ol@sectioncs{section}\makeatother

% BEGIN OLP-0426
\olfileid{nml}{frd}{st}

\olsection{દ્વિતીય-ક્રમ વ્યાખ્યાયિતતા}
\phantomsection\label{guunit:OLP-0426}


મોડલ સૂત્રોથી વ્યાખ્યાયિત કરી શકાય એવો દરેક
ફ્રેમ-ગુણધર્મ પ્રથમ-ક્રમમાં વ્યાખ્યાયિત થઈ શકતો
નથી. પરંતુ એકસ્થાની વિધેયધર્મો પર પરિમાણીકરણ
(એટલે કે એકસ્થાની દ્વિતીય-ક્રમ પરિમાણીકરણ)
મંજૂર કરીએ, તો બધા મોડલ રીતે વ્યાખ્યાયિત
ફ્રેમ-ગુણધર્મો વ્યાખ્યાયિત કરી શકાય છે. યુક્તિ
એ છે કે કોઈ મોડલ સૂત્ર જગતમાં કઈ શરતો હેઠળ
સત્ય છે અને પ્રથમ-ક્રમનાં સૂત્રો વચ્ચેના
ક્રમબદ્ધ સંબંધનો ઉપયોગ કરવો. આ મોડલ સૂત્રોનો
કહેવાતો \emph{પ્રમાણભૂત અનુવાદ} છે: તે માત્ર
સુલભતા સંબંધ માટેનું દ્વિસ્થાની વિધેય~$Q$
જ નહીં, પરંતુ~$!A$માં આવતાં દરેક
પ્રતિજ્ઞાપન ચલો~$\Obj p_i$ માટે
એકસ્થાની વિધેય~$P_i$ પણ ધરાવતી
પ્રથમ-ક્રમની ભાષામાં સૂત્રો આપે છે.

\begin{defn}
  \emph{પ્રમાણભૂત અનુવાદ}~$\ST_x(!A)$ નીચે પ્રમાણે
  આગમનાત્મક રીતે વ્યાખ્યાયિત થાય છે:
  \begin{enumerate}
  \item \indcase{!A}{\lfalse}{$\ST_x(\indfrm) = \lfalse$.}
  \item \indcase{!A}{\ltrue}{$\ST_x(\indfrm) = \ltrue$.}
  \item \indcase{!A}{\Obj p_i}{$\ST_x(\indfrm) = \Atom{P_i}{x}$.}
  \item \indcase{!A}{\lnot !B}{$\ST_x(\indfrm) = \lnot \ST_x(!B)$.}
  \item \indcase{!A}{(!B \land !C)}{$\ST_x(\indfrm) = (\ST_x(!B)
      \land \ST_x(!C))$.}
  \item \indcase{!A}{(!B \lor !C)}{$\ST_x(\indfrm) = (\ST_x(!B)
      \lor \ST_x(!C))$.}
  \item \indcase{!A}{(!B \lif !C)}{$\ST_x(\indfrm) = (\ST_x(!B)
      \lif \ST_x(!C))$.}
  \item \indcase{!A}{(!B \liff !C)}{$\ST_x(\indfrm) = (\ST_x(!B)
      \liff \ST_x(!C))}$.
  \item \indcase{!A}{\Box !B}{$\ST_x(\indfrm) =
      \lforall[y][(\Atom{Q}{x,y} \lif \ST_y(!B))]$.}
  \item \indcase{!A}{\Diamond !B}{$\ST_x(\indfrm) =
      \lexists[y][(\Atom{Q}{x,y} \land \ST_y(!B))]$.}
  \end{enumerate}
\end{defn}
\sourcecorrection{OLMD-019}{સ્થિર અંગ્રેજી મૂળના
\texttt{prvTrue} કિસ્સામાં આગમનનું સૂત્ર ભૂલથી
અસત્ય સ્થિરાંક છે, જ્યારે અનુવાદનું પરિણામ સત્ય
સ્થિરાંક છે. કિસ્સાનું સૂત્ર પરિણામ સાથે સુસંગત કર્યું છે.}

દાખલા તરીકે, $\ST_x(\Box p \lif p)$ એ
$\lforall[y][(\Atom{Q}{x,y}
  \lif \Atom{P}{y})] \lif \Atom{P}{x}$ છે.
~$\ST_x(!A)$ની ભાષાનું કોઈપણ સંરચના
માટે પ્રદેશ,~$Q$ને આપેલો દ્વિસ્થાની સંબંધ અને
એકસ્થાની વિધેયો~$P_i$ને આપેલા પ્રદેશના
ઉપગણો જોઈએ. બીજા શબ્દોમાં, આવા સંરચનાના
ઘટકો બરાબર~$!A$ના નિદર્શનના ઘટકો જ છે:
પ્રદેશ જગતોનો ગણ છે,~$Q$ને આપેલો દ્વિસ્થાની
સંબંધ સુલભતા સંબંધ છે, અને~$P_i$ને આપેલા
ઉપગણો બરાબર નિમણૂકો~$V(\Obj p_i)$ છે.
મોડલ નિદર્શનમાં $!A$નો સંતોષ અને અનુરૂપ
સંરચનામાં $\ST_x(!A)$નો સંતોષ મેળ ખાતા
હોય તે આશ્ચર્યની વાત નથી:

\begin{prop}\ollabel{prop:st}
  $\mModel{M} = \tuple{W, R, V}$ લો. $\Struct{M'}$ એવું
  પ્રથમ-ક્રમનું સંરચના હોય કે
  $\Domain{M'} = W$, $\Assign{Q}{M'} = R$ અને
  $\Assign{P_i}{M'} = V(\Obj p_i)$; તથા $s(x) = w$.
  ત્યારે
  \[
  \mSat{M}{!A}[w] \text{ ત્યારે અને માત્ર ત્યારે જ જ્યારે } \Sat{M'}{\ST_x(!A)}[s]
  \]
\end{prop}

\begin{proof}
  $!A$ પર આગમનથી.
\end{proof}

\begin{prop}
  માનો કે $!A$ મોડલ સૂત્ર છે અને
  $\mModel{F} = \tuple{W, R}$ ફ્રેમ છે.
  $\Struct{F'}$ એવું પ્રથમ-ક્રમનું સંરચના લો કે
  $\Domain{F'} = W$ અને $\Assign{Q}{F'} = R$;
  અને $!A'$ નીચેનું દ્વિતીય-ક્રમ સૂત્ર હોય:
  \[
  \lforall[X_1][\dots\lforall[X_n][\lforall[x][
        \SSubst{\ST_x(!A)}{\subst{X_1}{P_1}, \dots,
          \subst{X_n}{P_n}}]]],
  \]
  જ્યાં $P_1$, \dots, $P_n$ એ~$\ST_x(!A)$માંનાં
  બધાં એકસ્થાની વિધેયો છે. ત્યારે
  \[
  \mModel{F} \Entails !A \text{ ત્યારે અને માત્ર ત્યારે જ જ્યારે } \Sat{F'}{!A'}
  \]
\end{prop}

\begin{proof}
  $\Sat{F'}{!A'}$ ત્યારે અને માત્ર ત્યારે જ જ્યારે
  $\Domain{M'}=W$ અને $\Assign{Q}{M'}=R$ ધરાવતાં
  દરેક સંરચના~$\mModel{M'}$માં
  $i = 1$, \dots,~$n$ માટે
  $\Assign{P_i}{M'} \subseteq W$ હોય, અને
  $s(x) \in W$ ધરાવતી દરેક $s$ માટે
  $\Sat{M'}{\ST_x(!A)}[s]$ હોય. \olref{prop:st}
  મુજબ આ ત્યારે અને માત્ર ત્યારે જ છે જ્યારે
  ફ્રેમ~$\mModel{F}$ પર આધારિત બધા નિદર્શનો~$\mModel{M}$
  અને દરેક જગત~$w \in W$ માટે $\mSat{M}{!A}[w]$
  હોય, એટલે કે $\mModel{F} \Entails !A$.
\end{proof}
\sourcecorrection{OLMD-020}{સ્થિર અંગ્રેજી મૂળના
બીજા પ્રસ્તાવના પ્રમાણમાં રચના માટે નિર્ધારિત
પ્રદેશ અને દ્વિસ્થાની સંબંધને મધ્યવર્તી નિદર્શનમાં
સ્પષ્ટ રાખવામાં આવ્યા નથી. તે $W$ અને $R$
શરતો લખી છે; નહીં તો ફ્રેમ પર આધારિત
નિદર્શનો સાથે સમકક્ષતા ફલિત થતી નથી.}

\begin{defn}
  ફ્રેમોનો વર્ગ~$\mClass{F}$ \emph{દ્વિતીય-ક્રમમાં
  વ્યાખ્યાયિત} છે જો એક દ્વિસ્થાની વિધેય~$P$
  અને માત્ર એકસ્થાની ગણ-ચલો પરના પરિમાણકો
  ધરાવતી દ્વિતીય-ક્રમ ભાષામાં એવું વાક્ય~$!A$
  હોય કે $\mModel{F} = \tuple{W, R} \in \mClass{F}$
  ત્યારે અને માત્ર ત્યારે જ જ્યારે $\Domain{M} = W$
  અને $\Assign{P}{M} = R$ ધરાવતી
  સંરચના~$\Struct{M}$માં $\Sat{M}{!A}$ હોય.
\end{defn}

\begin{cor}
  ફ્રેમોનો વર્ગ કોઈ સૂત્ર~$!A$થી
  વ્યાખ્યાયિત થતો હોય, તો અનુરૂપ સુલભતા
  સંબંધોનો વર્ગ એકસ્થાની દ્વિતીય-ક્રમ વાક્યથી
  વ્યાખ્યાયિત થાય છે.
\end{cor}

\begin{proof}
  અગાઉના પ્રમાણનું એકસ્થાની દ્વિતીય-ક્રમ
  વાક્ય~$!A'$ જરૂરી ગુણધર્મ ધરાવે છે.
\end{proof}

ઉદાહરણ તરીકે ફરી સૂત્ર~$\Box p \lif p$ લો.
તે સ્વવાચકતા વ્યાખ્યાયિત કરે છે. સ્વવાચકતા
અલબત્ત પ્રથમ-ક્રમના વાક્ય~$\lforall[x][\Atom{Q}{x,x}]$થી
વ્યાખ્યાયિત થાય છે. પરંતુ એકસ્થાની દ્વિતીય-ક્રમના
વાક્યથી પણ તે વ્યાખ્યાયિત થાય છે:
\[
\lforall[X][\lforall[x][(\lforall[y][(\Atom{Q}{x,y} \lif \Atom{X}{y})]
    \lif \Atom{X}{x})]].
\]
એટલે આ બંને વાક્યો સમકક્ષ છે. આ વાત સીધી
કેવી રીતે સમજાય તે જોઈએ. પહેલાં માનો કે
દ્વિતીય-ક્રમ વાક્ય સંરચના~$\Struct{M}$માં
સત્ય છે. $x$ અને $X$ બંને પર સર્વપરિમાણક
હોવાથી બાકીનું કોઈપણ $x \in W$ અને
ગણ~$X \subseteq W$ માટે સાચું હોવું જોઈએ;
દાખલા તરીકે, $R =
\Assign{Q}{M}$ હોય ત્યારે ગણ
$\Setabs{z}{Rxz}$ લો. તો $s(x) \in W$ અને
$s(X) =
\Setabs{z}{Rxz}$ ધરાવતી કોઈપણ~$s$ માટે
$\Sat{M}{\lforall[y][(\Atom{Q}{x,y} \lif
    \Atom{X}{y})] \lif \Atom{X}{x}}[s]$ હોય.
પરંતુ $s(X)$ની પસંદગીથી આનો અર્થ
$\Sat{M}{\lforall[y][(\Atom{Q}{x,y} \lif
    \Atom{Q}{x,y})] \lif \Atom{Q}{x,x}}[s]$ હોય,
જેનો પૂર્વપક્ષ માન્ય હોવાથી તે
$\Atom{Q}{x,x}$ને સમકક્ષ છે. $s(x)$ ગમે તે
હોવાથી $\Sat{M}{\lforall[x][\Atom{Q}{x,x}]}$ મળે.
\sourcecorrection{OLMD-021}{સ્થિર અંગ્રેજી મૂળની
સીધી દલીલમાં મુક્ત ચલ ધરાવતી એક સંતોષ-અભિવ્યક્તિ
સાથે નિમણૂક દર્શાવાઈ નથી. આસપાસની દલીલ મુજબ
તેમાં નિમણૂકનો નિર્દેશ પૂર્યો છે.}

હવે $\Sat{M}{\lforall[x][\Atom{Q}{x,x}]}$ માનો
અને $\Sat{M}{\lforall[X][\lforall[x][(\lforall[y][(\Atom{Q}{x,y} \lif
        \Atom{X}{y})] \lif \Atom{X}{x})]]}$ બતાવો.
કોઈપણ નિમણૂક~$s$ લો અને
$\Sat{M}{\lforall[y][(\Atom{Q}{x,y} \lif
    \Atom{X}{y})]}[s]$ માનો. $s'$ એ~$s$નું એવું
$y$-રૂપાંતર હોય કે $s'(y)
= s(x)$; તો $\Sat{M}{\Atom{Q}{x,y} \lif \Atom{X}{y}}[s']$,
એટલે કે $\Sat{M}{\Atom{Q}{x,x} \lif \Atom{X}{x}}[s]$.
$\Sat{M}{\lforall[x][\Atom{Q}{x,x}]}$ હોવાથી
પૂર્વપક્ષ સત્ય છે અને $\Sat{M}{\Atom{X}{x}}[s]$
મળે છે; આ જ બતાવવાનું હતું.

ફ્રેમોના કેટલાક વ્યાખ્યાયિત વર્ગો પ્રથમ-ક્રમમાં
વ્યાખ્યાયિત થતા નથી; તેથી~$!A'$ સ્વરૂપના દરેક
એકસ્થાની દ્વિતીય-ક્રમ વાક્યને પ્રથમ-ક્રમના
વાક્યને સમકક્ષ બનાવી શકાય એવું નથી.
કયાં બનાવી શકાય તે નક્કી કરવાની કોઈ અસરકારક
પદ્ધતિ નથી.
% END OLP-0426
\endgroup


\OLEndChapterHook
% END OLP-0419
\endgroup


\begingroup
\makeatletter\def\ol@sectioncs{section}\makeatother

% BEGIN OLP-0427
\olchapter{nml}{prf}{સ્વયંસિદ્ધિમૂલક નિષ્પત્તિઓ}
\olchapterid{nml}{prf}
\phantomsection\label{guunit:OLP-0427}


\begingroup
\makeatletter\def\ol@sectioncs{section}\makeatother

% BEGIN OLP-0428
\olfileid{nml}{prf}{int}\sourcecorrection{OLMD-022}{મૂળમાં આ પરિચયનું પ્રકરણ-ઓળખચિહ્ન axs છે, જ્યારે તેને આયાત કરતું પ્રકરણ prf ઓળખચિહ્ન વાપરે છે; અહીં બંને એકસરખાં કર્યાં છે.}

\olsection{પરિચય}
\phantomsection\label{guunit:OLP-0428}


આપણી પાસે મૂળભૂત મોડલ ભાષા માટે મોડલ નિદર્શનો પર આધારિત અર્થવિજ્ઞાન છે.
સૂત્રની પ્રામાણ્યતા એટલે બધા નિદર્શનોનાં બધાં જગતોમાં તેની સત્યતા;
નિદર્શનોના કોઈ વર્ગ કે ફ્રેમને અનુલક્ષીને પ્રામાણ્યતા એટલે તે વર્ગનાં બધાં
નિદર્શનોનાં, અથવા તે ફ્રેમ પર આધારિત બધાં નિદર્શનોનાં, બધાં જગતોમાં તેની
સત્યતા. તર્કશાસ્ત્ર સામાન્ય રીતે પ્રામાણ્યતાનાં આવાં અર્થાનુસારી વર્ણનોને
નિષ્પન્નક્ષમતાના પ્રમાણસિદ્ધાંતલક્ષી ખ્યાલ સાથે જોડે છે. આપણો હેતુ
કોઈ તંત્રમાં નિષ્પન્નક્ષમતા એવી રીતે વ્યાખ્યાયિત કરવાનો છે કે સૂત્ર
પ્રામાણ્ય હોય તો અને માત્ર તો જ તે નિષ્પન્ન કરી શકાય એવું હોય.

સૌથી સરળ અને ઐતિહાસિક રીતે સૌથી જૂનાં નિષ્પત્તિ તંત્રો હિલ્બર્ટ-પ્રકારનાં,
અથવા સ્વયંસિદ્ધિમૂલક, નિષ્પત્તિ તંત્રો છે. અનેક મોડલ તર્કો માટે આવાં
નિષ્પત્તિ તંત્રો રચવાં પ્રમાણમાં સહેલાં છે; અધિસૈદ્ધાંતિક અભ્યાસમાં પણ તે સરળ છે
(દાખલા તરીકે, યથાર્થતા અને પૂર્ણતા સાબિત કરવા). જોકે, કુદરતી નિષ્કર્ષણનાં
તંત્રોની સરખામણીમાં તેમાં સૂત્રો સાબિત કરવાં ઘણાં અઘરાં છે.

હિલ્બર્ટ-પ્રકારનાં નિષ્પત્તિ તંત્રોમાં સૂત્રની નિષ્પત્તિ
એ ચોક્કસ સ્વયંસિદ્ધિઓથી શરૂ થઈને થોડા અનુમાન-નિયમો વડે સંબંધિત સૂત્ર
સુધી પહોંચતી સૂત્રોની શ્રેણી છે. નિષ્પત્તિ તંત્ર અર્થવિજ્ઞાન સાથે
મેળ ખાય તે માટે આપણે ખાતરી કરવી પડે કે નિષ્પન્ન કરી શકાય એવું સૂત્રો બધાં
નિદર્શનોમાં સત્ય છે (અથવા બધાં સ્વયંસિદ્ધિઓ સત્ય હોય એવાં બધાં
નિદર્શનોમાં સત્ય છે). પહેલાં આપણે આ માટે જરૂરી મોડલ તર્કોના કેટલાક
ગુણધર્મો અલગ તારવીશું: આ ``સામાન્ય'' મોડલ તર્કો છે. સામાન્ય મોડલ તર્કો
માટે અનુમાનના માત્ર બે નિયમો ધારવા પડે છે: મોડસ પોનેન્સ અને આવશ્યકીકરણ.
સ્વયંસિદ્ધિઓ તરીકે આપણે પુનરુક્તિઓના બધા પ્રતિસ્થાપન દૃષ્ટાંતો અને,
જે મોડલ તર્કની ચર્ચા કરીએ તેના આધારે, કેટલીક મોડલ સ્વયંસિદ્ધિઓ લઈએ
છીએ. જો આપણને બધાં નિદર્શનોના વર્ગમાં જ રસ હોય, તોપણ~$\Ax{K}$ અને~$\Ax{Dual}$ના
બધા પ્રતિસ્થાપન દૃષ્ટાંતોને સ્વયંસિદ્ધિઓ ગણવાં પડે છે. માત્ર આટલું
ન્યૂનતમ સામાન્ય મોડલ તર્ક~$\Log K$ પેદા કરે છે.

\begin{defn}
\emph{મોડસ પોનેન્સ}નો નિયમ આ અનુમાન-ઢાંચો છે:
\begin{prooftree}
\AxiomC{$!A$}
\AxiomC{$!A \lif !B$}
\RightLabel{\MP}
\BinaryInfC{$!B$}
\end{prooftree}
જો $!C \ident !A \lif !B$ હોય, તો આપણે કહીએ છીએ કે સૂત્રો $!A$, $!C$
માંથી મોડસ પોનેન્સ વડે સૂત્ર~$!B$ \emph{ફલિત થાય છે}.
\end{defn}

\begin{defn}
\emph{આવશ્યકીકરણ}નો નિયમ આ અનુમાન-ઢાંચો છે:
\begin{prooftree}
\AxiomC{$!A$}
\RightLabel{\Nec}
\UnaryInfC{$\Box !A$}
\end{prooftree}
જો $!B \ident \Box !A$ હોય, તો આપણે કહીએ છીએ કે સૂત્ર~$!A$માંથી
આવશ્યકીકરણ વડે સૂત્ર~$!B$ ફલિત થાય છે.
\end{defn}

\begin{defn}
સ્વયંસિદ્ધિઓના ગણ~$\Sigma$માંથી \emph{નિષ્પત્તિ} એ સૂત્રો
$!B_1$, $!B_2$, \dots, $!B_n$ની એવી શ્રેણી છે કે દરેક $!B_i$ આમાંથી
એક હોય:
\begin{enumerate}
\item પુનરુક્તિનો પ્રતિસ્થાપન દૃષ્ટાંત, અથવા
\item $\Sigma$માંના સૂત્રનો પ્રતિસ્થાપન દૃષ્ટાંત, અથવા
\item $j$, $k < i$ હોય એવાં બે સૂત્રો $!B_j$, $!B_k$માંથી
  મોડસ પોનેન્સ વડે ફલિત થતું સૂત્ર, અથવા
\item $j < i$ હોય એવા સૂત્ર~$!B_j$માંથી આવશ્યકીકરણ વડે
  ફલિત થતું સૂત્ર.
\end{enumerate}
જો $!B_n \ident !A$ હોય એવી નિષ્પત્તિ હોય, તો આપણે કહીએ છીએ કે
$!A$ એ~$\Sigma$માંથી \emph{નિષ્પન્ન કરી શકાય એવું} છે; સંકેતમાં $\Sigma \Proves !A$.
\end{defn}

આ વ્યાખ્યા અનુસાર નિષ્પન્ન કરી શકાય એવું સૂત્રોનો ગણ સામાન્ય મોડલ તર્ક બને છે,
અને દરેક નિષ્પન્ન કરી શકાય એવું સૂત્ર એવા દરેક નિદર્શનમાં સત્ય છે જેમાં
દરેક સ્વયંસિદ્ધિ સત્ય છે. નિષ્પત્તિઓના આ ગુણધર્મને \emph{યથાર્થતા}
કહે છે. તેના વિપરીત ગુણધર્મ, \emph{પૂર્ણતા}, સાબિત કરવો વધુ અઘરો છે.
% END OLP-0428
\endgroup

\begingroup
\makeatletter\def\ol@sectioncs{section}\makeatother

% BEGIN OLP-0429
\olfileid{nml}{prf}{nor}

\olsection{સામાન્ય મોડલ તર્કશાસ્ત્રો}
\phantomsection\label{guunit:OLP-0429}


મોડલ સૂત્રોના દરેક ગણને સ્વયંસિદ્ધિઓના કોઈ ગણમાંથી નિષ્પન્ન થતા
સૂત્રોના ગણ તરીકે સહેલાઈથી લક્ષણિત કરી શકાતો નથી. આપણે મોડલ તર્કો
સુવ્યવસ્થિત હોય એવું ઇચ્છીએ છીએ. સૌપ્રથમ, પરંપરાગત વિધાનાત્મક તર્કમાં
જે કંઈ આપણે નિષ્પન્ન કરી શકીએ તે હજી પણ નિષ્પન્ન કરી શકાય એવું હોવું જોઈએ;
અલબત્ત, સૂત્રોમાં હવે $\Box$
    અને~$\Diamond$ પણ હોઈ શકે છે. આ માટે
અમે માગીએ છીએ કે મોડલ તર્કમાં પુનરુક્તિઓના બધા દૃષ્ટાંતો હોય અને તે
મોડસ પોનેન્સ હેઠળ બંધ હોય.

\begin{defn}
  \emph{મોડલ તર્કશાસ્ત્ર} એ મોડલ સૂત્રોનો એવો ગણ~$\Sigma$ છે જે
  \begin{enumerate}
  \item બધી પુનરુક્તિઓ ધરાવે છે, અને
  \item પ્રતિસ્થાપન હેઠળ બંધ છે, એટલે કે જો $!A \in \Sigma$ હોય અને
    $!D_1$, \dots, $!D_n$ સૂત્રો હોય, તો
    \[
    \SSubst{!A}{\subst{!D_1}{p_1}, \dots, \subst{!D_n}{p_n}} \in \Sigma,
    \]
    \item \emph{મોડસ પોનેન્સ} હેઠળ બંધ છે, એટલે કે જો $!A$ અને $!A
      \lif !B \in \Sigma$ હોય, તો $!B \in \Sigma$ હોય છે.
  \end{enumerate}
\end{defn}

મોડલ તર્કો માટે સંબંધાત્મક અર્થવિજ્ઞાન વાપરવું હોય, તો બધાં મોડલ
નિદર્શનોમાં પ્રામાણ્ય હોય એવાં બધાં સૂત્રો પણ તેમાં સામેલ હોય તેવી
માગ રાખવી પડે. \Ax{K} અને~\Dual{}ના બધા દૃષ્ટાંતો
નિષ્પન્ન કરી શકાય એવું હોય, અને સૂત્ર~$!A$ નિષ્પન્ન કરી શકાય એવું હોય ત્યારે
$\Box !A$ પણ નિષ્પન્ન હોય, તો આ માગ પૂરી થાય છે. આ શરતો પૂરી કરતો
મોડલ તર્ક \emph{સામાન્ય} કહેવાય છે. (નિશ્ચિતપણે, અસામાન્ય મોડલ તર્કો પણ
છે, પરંતુ સામાન્ય સંબંધાત્મક નિદર્શનો તેમને માટે પર્યાપ્ત નથી.)

\begin{defn}
  મોડલ તર્ક $\Sigma$ \emph{સામાન્ય} છે, જો તે આ સૂત્રો ધરાવે:
  \begin{align*}
    \tag{\Ax{K}} & \Box(p \lif q) \lif (\Box p \lif \Box q),
    \\
      \tag{\Dual} & \Diamond p \liff \lnot\Box\lnot p
  \end{align*}
  અને તે \emph{આવશ્યકીકરણ} હેઠળ બંધ હોય, એટલે કે $!A \in
  \Sigma$ હોય, તો $\Box !A \in \Sigma$ પણ હોય.
\end{defn}

નોંધો કે પુનરુક્તિમૂલક પ્રેરણ ``સૂક્ષ્મ'' અર્થમાં કોઈ \emph{જગતમાં
સત્યતા} જાળવવા જેટલું પૂરતું છે, જ્યારે નિયમ \Nec{} માત્ર \emph{નિદર્શનમાં
સત્યતા} જાળવે છે (અને તેથી ફ્રેમમાં કે ફ્રેમોના વર્ગમાં પ્રામાણ્યતા પણ
જાળવે છે).

\begin{prop}\ollabel{prop:rk}
  દરેક સામાન્ય મોડલ તર્કશાસ્ત્ર નિયમ \RK હેઠળ બંધ છે:
  \begin{prooftree}
    \AxiomC{$!A_1 \lif (!A_2 \lif \cdots (!A_{n-1} \lif !A_n)\cdots)$}
    \RightLabel{\RK}
    \UnaryInfC{$\Box!A_1 \lif (\Box!A_2 \lif \cdots (\Box!A_{n-1}
      \lif \Box!A_n)\cdots).$}
  \end{prooftree}
\end{prop}

\begin{proof}
  $n$ પર આગમનથી: $n = 1$ હોય, તો આ નિયમ માત્ર \Nec છે, અને દરેક
  સામાન્ય મોડલ તર્ક \Nec હેઠળ બંધ છે.

  હવે માનો કે પરિણામ $n-1$ માટે સાચું છે; આપણે તે~$n$ માટે બતાવીએ.

  માનો કે
  \begin{align*}
  & !A_1 \lif (!A_2 \lif \cdots (!A_{n-1} \lif !A_n)\cdots) \in \Sigma
  \intertext{આગમન-પરિકલ્પના મુજબ આપણને મળે છે}
  & \Box!A_1 \lif (\Box!A_2 \lif \cdots \Box(!A_{n-1} \lif !A_n)\cdots)
  \in \Sigma
  \intertext{$\Sigma$ સામાન્ય મોડલ તર્ક હોવાથી તે~$\Ax{K}$ના બધા દૃષ્ટાંતો,
    ખાસ કરીને આ દૃષ્ટાંત, ધરાવે છે}
  & \Box(!A_{n-1} \lif !A_n) \lif (\Box!A_{n-1} \lif \Box!A_n) \in \Sigma
  \intertext{મોડસ પોનેન્સ અને પુનરુક્તિઓના યોગ્ય દૃષ્ટાંતો વાપરવાથી મળે છે}
  & \Box!A_1 \lif (\Box!A_2 \lif \cdots (\Box!A_{n-1}
  \lif \Box!A_n)\cdots) \in \Sigma. 
  \end{align*}
\end{proof}

\begin{prop}\ollabel{prop:notDiamondBot}
  દરેક સામાન્ય મોડલ તર્ક $\Sigma$માં~$\lnot\Diamond\lfalse$ હોય છે.
\end{prop}

\begin{prob}
  \olref[nml][prf][nor]{prop:notDiamondBot} સાબિત કરો.
\end{prob}

\begin{prop}
  માનો કે $!A_1$, \dots, $!A_n$ સૂત્રો છે. ત્યારે $!A_1$,
  \dots, $!A_n$ના બધા દૃષ્ટાંતો ધરાવતો સૌથી નાનો સામાન્ય મોડલ તર્ક
  $\Sigma$ છે.
\end{prop}\sourcecorrection{OLMD-023}{મૂળ પ્રસ્તાવમાં માત્ર ``મોડલ તર્ક'' લખાયું છે, પરંતુ તેની સાબિતી સામાન્ય મોડલ તર્કોના છેદથી સૌથી નાનો સામાન્ય મોડલ તર્ક રચે છે; આ વિશેષણ અહીં સ્પષ્ટ કર્યું છે.}

\begin{proof}
  $!A_1$, \dots, $!A_n$ આપેલાં હોય ત્યારે $\Sigma$ને એવાં બધાં સામાન્ય
  મોડલ તર્કોના છેદ તરીકે વ્યાખ્યાયિત કરો જે $!A_1$, \dots, $!A_n$ના
  બધા દૃષ્ટાંતો ધરાવે છે. આવો છેદ ખાલી નથી, કારણ કે બધાં સૂત્રોનો
  ગણ~$\Frm[L]$ પોતે એવો મોડલ તર્ક છે.
\end{proof}

\begin{defn}
$!A_1$, \dots, $!A_n$ ધરાવતો સૌથી નાનો સામાન્ય મોડલ તર્ક
\emph{મોડલ તંત્ર} કહેવાય છે અને તેને $\Log{K} !A_1 \dots
!A_n$ વડે દર્શાવાય છે. સૌથી નાનો સામાન્ય મોડલ તર્ક~\Log{K} વડે
દર્શાવાય છે.
\end{defn}
% END OLP-0429
\endgroup

\begingroup
\makeatletter\def\ol@sectioncs{section}\makeatother

% BEGIN OLP-0430
\olfileid{nml}{prf}{prf}

\olsection{નિષ્પત્તિઓ અને મોડલ તંત્રો}
\phantomsection\label{guunit:OLP-0430}


પહેલાં આપણે સામાન્ય મોડલ તર્કો માટે નિષ્પત્તિ શું છે તે
વ્યાખ્યાયિત કરીએ. અંદાજે, નિષ્પત્તિ એ સૂત્રોની એવી શ્રેણી છે
જેમાં દરેક ઘટક કાં તો કેટલાંક \emph{સ્વયંસિદ્ધિઓ}માંનું એક
(અથવા તેનો પ્રતિસ્થાપન દૃષ્ટાંત) હોય છે, અથવા અગાઉનાં ઘટકોમાંથી
થોડા અનુમાન-નિયમોમાંના કોઈ નિયમ વડે ફલિત થાય છે. સામાન્ય મોડલ તર્કો
માટે પુનરુક્તિઓના બધા દૃષ્ટાંતો, \Ax{K} અને~\Dual{}
સ્વયંસિદ્ધિઓ ગણાય છે. તેનાથી સૌથી નાનો સામાન્ય મોડલ તર્ક, મોડલ
તંત્ર~$\Log{K}$, મળે છે. બીજા તંત્રો મેળવવા આપણે વધારાનાં
સ્વયંસિદ્ધિઓ ઉમેરી શકીએ. નિયમો હંમેશા મોડસ પોનેન્સ~\MP{} અને
આવશ્યકીકરણ~\Nec છે.

\begin{defn}
  મોડલ તંત્ર $\Log{K} !A_1 \dots !A_n$ અને સૂત્ર~$!B$ આપેલાં
  હોય, ત્યારે $!B$ એ $\Log{K} !A_1 \dots
  !A_n$માં \emph{નિષ્પન્ન કરી શકાય એવું} છે, જે
  $\Log{K} !A_1 \dots !A_n \Proves !B$ લખાય છે, જો અને માત્ર જો
  એવાં સૂત્રો $!C_1$, \dots, $!C_k$ હોય કે $!C_k = !B$ અને દરેક
  $!C_i$ કાં તો પુનરુક્તિનો દૃષ્ટાંત હોય, અથવા
  $\Ax{K}$, $\Dual$, $!A_1$,
  \dots, $!A_n$ પૈકી કોઈ એકનો દૃષ્ટાંત હોય, અથવા અગાઉનાં
  સૂત્રોમાંથી \MP{} કે~\Nec નિયમ વડે ફલિત થાય.
\end{defn}

નીચેનો પ્રસ્તાવ~$!B$ની $\Sigma$-નિષ્પત્તિ દર્શાવીને
$!B \in \Sigma$ સાબિત કરવા દે છે.

\begin{prop}
  $\Log{K} !A_1 \dots !A_n = \Setabs{!B}{\Log{K} !A_1 \dots !A_n \Proves !B}$.
\end{prop}

\begin{proof}
  નિષ્પત્તિઓની લંબાઈ પર આગમન વાપરીને આપણે બતાવીએ છીએ કે
  $\Setabs{!B}{\Log{K} !A_1 \dots !A_n \Proves !B} \subseteq
  \Log{K} !A_1 \dots !A_n$.

  જો~$!B$ની નિષ્પત્તિની લંબાઈ~$1$ હોય, તો તેમાં એક જ સૂત્ર
  છે. તે સૂત્ર અગાઉનાં સૂત્રોમાંથી \MP{} કે \Nec વડે ફલિત થઈ
  શકે નહિ; એટલે તે પુનરુક્તિનો દૃષ્ટાંત, \Ax{K}નો દૃષ્ટાંત,
  $\Dual$નો દૃષ્ટાંત, અથવા $!A_1$, \dots,~$!A_n$ પૈકી કોઈ એકનો
  દૃષ્ટાંત હોવું જોઈએ. પણ $\Log{K}!A_1\dots!A_n$માં આ બધાં હોય છે;
  તેથી $!B \in \Log{K}!A_1 \dots !A_n$.

  જો $!B$ની નિષ્પત્તિની લંબાઈ~$> 1$ હોય, તો $!B$ એ નિષ્પત્તિની
  છેલ્લી પંક્તિમાં ન આવતાં અગાઉનાં સૂત્રોમાંથી \MP{} કે \Nec{}
  વડે પણ મેળવાયેલું હોઈ શકે છે. જો $!B$ એ $!C$ અને $!C \lif !B$માંથી
  (\MP વડે) ફલિત થતું હોય, તો આગમન-પરિકલ્પના અનુસાર $!C$ અને $!C \lif !B \in
  \Log{K}!A_1 \dots !A_n$. પણ દરેક મોડલ તર્ક મોડસ પોનેન્સ હેઠળ બંધ
  છે; તેથી $!B \in \Log{K}!A_1 \dots
  !A_n$. જો $!B \equiv \Box !C$ એ $!C$માંથી \Nec વડે ફલિત થતું હોય,
  તો આગમન-પરિકલ્પના અનુસાર $!C
  \in \Log{K}!A_1 \dots !A_n$. પણ દરેક સામાન્ય મોડલ તર્ક $\Nec$
  હેઠળ બંધ છે; તેથી $!B \in
  \Log{K}!A_1\dots!A_n$.

  વિપરીત સમાવેશ સાબિત કરવા આપણે બતાવીએ છીએ કે
  $\Sigma = \Setabs{!B}{\Log{K} !A_1 \dots !A_n \Proves !B }$ એ
  $!A_1$, \dots, $!A_n$ના બધા દૃષ્ટાંતો ધરાવતો સામાન્ય મોડલ તર્ક છે.
  પછી વ્યાખ્યા મુજબ $\Log{K} !A_1 \dots !A_n$ સૌથી નાનો એવો તર્ક
  છે, એ અવલોકન વાપરીએ છીએ.
  \begin{enumerate}
    \item દરેક પુનરુક્તિ~$!B$ પુનરુક્તિનો દૃષ્ટાંત હોવાથી
      $\Log{K}!A_1\dots!A_n \Proves !B$; તેથી $\Sigma$ બધી
      પુનરુક્તિઓ ધરાવે છે.
    \item જો $\Log{K}!A_1\dots!A_n \Proves !C$ અને
      $\Log{K}!A_1\dots!A_n \Proves !C \lif !B$ હોય, તો
      $\Log{K}!A_1\dots!A_n \Proves !B$: $!C$ની નિષ્પત્તિને
      $!C \lif !B$ની નિષ્પત્તિ સાથે જોડો અને અંતે~$!B$ની પંક્તિ
      ઉમેરો. છેલ્લી પંક્તિ \MP{}થી સમર્થિત છે. તેથી $\Sigma$ મોડસ
      પોનેન્સ હેઠળ બંધ છે.
    \item જો $!B$ની નિષ્પત્તિ હોય, તો $!B$ના દરેક પ્રતિસ્થાપન
      દૃષ્ટાંતની પણ નિષ્પત્તિ હોય છે: નિષ્પત્તિમાંના દરેક
      સૂત્ર પર એ પ્રતિસ્થાપન લાગુ કરો. (અભ્યાસપ્રશ્ન: નિષ્પત્તિઓની
      લંબાઈ પર આગમનથી સાબિત કરો કે પરિણામ પણ યોગ્ય નિષ્પત્તિ છે.)
      તેથી $\Sigma$ એકસાથે પ્રતિસ્થાપન હેઠળ બંધ છે. (હવે આપણે
      સ્થાપ્યું છે કે $\Sigma$ મોડલ તર્કની બધી શરતો પૂરી કરે છે.)
    \item આપણી પાસે $\Log{K}!A_1\dots!A_n \Proves \Ax{K}$ છે;
      તેથી $\Ax{K} \in \Sigma$.\sourcecorrection{OLMD-024}{મૂળ નિષ્કર્ષમાં $K \in \Sigma$ છે, પરંતુ અહીં \Sigma સૂત્રોનો ગણ છે અને પૂર્વપક્ષમાં \Ax{K} છે; નિષ્કર્ષમાં પણ એ જ સૂત્રચિહ્ન મૂક્યું છે.}
    \item આપણી પાસે
      $\Log{K}!A_1\dots!A_n \Proves \Dual$ છે; તેથી $\Dual \in \Sigma$.
    \item જો $\Log{K}!A_1\dots!A_n \Proves !C$ હોય, તો વધારાની
      પંક્તિ~$\Box !C$ને~\Nec વડે સમર્થન મળે છે. તેથી $\Sigma$
      એ~\Nec હેઠળ બંધ છે; એટલે $\Sigma$ સામાન્ય છે.
  \end{enumerate}
\end{proof}
% END OLP-0430
\endgroup

\begingroup
\makeatletter\def\ol@sectioncs{section}\makeatother

% BEGIN OLP-0431
\olfileid{nml}{prf}{prk}

\olsection{\Log{K}માં પ્રમાણો}
\phantomsection\label{guunit:OLP-0431}


સૌથી નાના મોડલ તંત્રમાં પ્રમાણોનો અભ્યાસ કરવા માટે આપણે બતાવીએ છીએ કે
\olref[syn][sch]{tab:valid-invalidSchemas}ની ડાબી બાજુનાં પ્રામાણ્ય
સૂત્રો માટે \Log{K}-પ્રમાણો આપી શકાય છે.

\begin{prop}
  $\Log{K} \Proves \Box!A \lif \Box (!B\lif !A)$
\end{prop}

\begin{proof}
  \begin{derivation}
    1. & $!A \lif (!B \lif !A)$ & \Taut \\
    2. & $\Box(!A \lif (!B \lif !A))$ & \Nec, 1 \\
    3. & $\Box(!A \lif (!B \lif !A)) \lif
    (\Box!A \lif \Box(!B \lif !A))$ & \Ax{K}\\
    4. & $\Box!A \lif \Box (!B\lif !A)$ & \MP, 2, 3
  \end{derivation}
\end{proof}

\begin{prop}
  $\Log{K} \Proves \Box(!A \land !B) \lif (\Box !A \land \Box!B)$
\end{prop}

\begin{proof}
\begin{derivation}
  1. & $(!A \land !B) \lif !A$ & \Taut \\
  2. & $\Box((!A \land !B) \lif !A)$ & \Nec \\
  3. & $\Box((!A \land !B) \lif !A) \lif (\Box(!A \land !B) \lif \Box!A)$
  & \Ax{K} \\
  4. & $\Box(!A \land !B) \lif \Box!A$ & \MP, 2, 3 \\
  5. & $(!A \land !B) \lif !B$ & \Taut \\
  6. & $\Box((!A \land !B) \lif !B)$ & \Nec \\
  7. & $\Box((!A \land !B) \lif !B) \lif (\Box(!A \land !B) \lif \Box!B)$
  & \Ax{K} \\
  8. & $\Box(!A \land !B) \lif \Box!B$ & \MP, 6, 7 \\
  9. & $(\Box(!A \land !B) \lif \Box!A) \lif{}$ \\
  & \qquad $((\Box(!A \land !B) \lif \Box!B) \lif{}$ \\
  & \qquad $(\Box(!A \land !B) \lif (\Box !A \land \Box!B)))$ & \Taut\\
  10. &  $(\Box(!A \land !B) \lif \Box!B) \lif{}$ \\
  & \qquad $(\Box(!A \land !B) \lif (\Box !A \land \Box!B))$ & \MP, 4, 9\\
  11. & $\Box(!A \land !B) \lif (\Box !A \land \Box!B)$ & \MP, 8, 10.
\end{derivation}
નોંધો કે પંક્તિ~$9$ પરનું સૂત્ર આ પુનરુક્તિનો દૃષ્ટાંત છે:
\[
(p \lif q) \lif ((p \lif r) \lif (p \lif (q \land r))).
\]
\end{proof}

\begin{prop}
  $\Log{K}\Proves (\Box!A \land \Box!B) \lif \Box (!A \land !B)$
\end{prop}

\begin{proof}
  \begin{derivation}
    1. & $!A \lif (!B \lif (!A \land !B))$ & \Taut \\
    2. & $\Box(!A \lif (!B \lif (!A \land !B)))$ & \Nec, 1 \\
    3. & $\Box(!A \lif (!B \lif (!A \land !B))) \lif (\Box!A \lif \Box(!B \lif (!A \land !B)))$ & \Ax{K}\\
    4. & $\Box!A \lif \Box(!B \lif (!A \land !B))$ & \MP, 2, 3\\
    5. & $\Box(!B \lif (!A \land !B)) \lif (\Box!B \lif \Box(!A \land !B))$ & \Ax{K} \\
    6. & $(\Box!A \lif \Box(!B \lif (!A \land !B))) \lif {}$ \\
    & \qquad $(\Box(!B \lif (!A \land !B)) \lif (\Box!B \lif \Box(!A \land !B))) \lif {}$\\
    & \qquad $(\Box!A \lif (\Box !B \lif \Box(!A \land !B))))$ & \Taut\\
    7. & $(\Box(!B \lif (!A \land !B)) \lif (\Box!B \lif \Box(!A \land !B))) \lif {}$\\
    & \qquad $(\Box!A \lif (\Box !B \lif \Box(!A \land !B)))$ & \MP, 4, 6\\
    8. & $\Box!A \lif (\Box !B \lif \Box(!A \land !B)))$ & \MP, 5, 7\\
    9. & $(\Box!A \lif (\Box !B \lif \Box(!A \land !B)))) \lif {}$\\
    & \qquad $((\Box!A \land \Box!B) \lif \Box(!A \land !B))$ & \Taut\\
    10. & $(\Box!A \land \Box!B) \lif \Box(!A \land !B)$ & \MP, 8, 9
  \end{derivation}
  પંક્તિઓ $6$ અને $9$ પરનાં સૂત્રો અનુક્રમે આ પુનરુક્તિઓના દૃષ્ટાંતો છે:
  \begin{align*}
    (p \lif q) & \lif ((q \lif r) \lif (p \lif r)) \\
    (p \lif (q \lif r)) & \lif ((p \land q) \lif r)
  \end{align*}
\end{proof}

\begin{prop}
  $\Log{K} \Proves \lnot\Box p \lif \Diamond \lnot
    p$
\end{prop}

\begin{proof}
  % Box અને Diamond મૂળભૂત છે
  \begin{derivation}
    1. & $\Diamond \lnot p \liff \lnot\Box\lnot\lnot p$ & \Dual\\
    2. & $(\Diamond \lnot p \liff \lnot\Box\lnot\lnot p) \lif {}$ \\
    & \qquad $(\lnot \Box \lnot\lnot p \lif \Diamond \lnot p)$& \Taut\\
    3. & $\lnot \Box \lnot\lnot p \lif \Diamond \lnot p$ & \MP, 1, 2\\
    4. & $\lnot\lnot p \lif p$ & \Taut\\
    5. & $\Box(\lnot\lnot p \lif p)$ & \Nec, 4\\
    6. & $\Box(\lnot\lnot p \lif p) \lif (\Box \lnot\lnot p \lif \Box p)$ & \Ax{K}\\
    7. & $(\Box \lnot\lnot p \lif \Box p)$ & \MP, 5, 6\\
    8. & $(\Box \lnot\lnot p \lif \Box p) \lif (\lnot \Box p \lif \lnot\Box\lnot\lnot p)$ & \Taut\\
    9. & $\lnot \Box p \lif \lnot\Box\lnot\lnot p$ & \MP, 7, 8\\
    10. & $(\lnot \Box p \lif \lnot\Box\lnot\lnot p) \lif {}$ \\
    & \qquad $((\lnot \Box \lnot\lnot p \lif \Diamond \lnot p) \lif (\lnot \Box p \lif \Diamond\lnot p))$ & \Taut\\
    11. & $(\lnot \Box \lnot\lnot p \lif \Diamond \lnot p) \lif (\lnot \Box p \lif \Diamond\lnot p)$ & \MP, 9, 10\\
    12. & $\lnot\Box p \lif \Diamond \lnot p$ & \MP, 3, 11\\
  \end{derivation}
  પંક્તિઓ $8$ અને $10$ પરનાં સૂત્રો આ પુનરુક્તિઓના દૃષ્ટાંતો છે:
  \begin{align*}
    & (p \lif q) \lif (\lnot q \lif \lnot p) \\
    & (p \lif q) \lif ((q \lif r) \lif (p \lif r)).
  \end{align*}
\end{proof}

\begin{prob}
  નીચેના સૂત્રો માટે~$\Log{K}$માં નિષ્પત્તિઓ શોધો:
  \begin{enumerate}
    \item $\Box \lnot p \lif \Box(p \lif q)$
    \item $(\Box p \lor \Box q) \lif \Box(p \lor q)$
    \item $\Diamond p \lif \Diamond(p \lor q)$
  \end{enumerate}
\end{prob}
% END OLP-0431
\endgroup

\begingroup
\makeatletter\def\ol@sectioncs{section}\makeatother

% BEGIN OLP-0432
\olfileid{nml}{prf}{der}

\olsection{નિષ્પન્ન નિયમો}
\phantomsection\label{guunit:OLP-0432}


નિષ્પત્તિઓ શોધવી અને લખવી સ્પષ્ટપણે મુશ્કેલ, કંટાળાજનક અને
પુનરાવર્તિત કામ છે. દાખલા તરીકે, ઘણી વાર આપણે $!A \lif !B$માંથી
$\Box !A \lif \Box !B$ મેળવવું હોય છે, એટલે કે નિયમ~\RK વાપરવો હોય
છે. તે માટે પહેલાં \Nec વાપરવું, પછી~\Ax{K}નો યોગ્ય દૃષ્ટાંત લખવો
અને છેલ્લે~\MP વાપરવું પડે. $!A \lif
!B$ અને $!B \lif !C$માંથી $!A \lif !C$ મેળવવા માટે આ લાંબો
પુનરુક્તિ-દૃષ્ટાંત લખવો પડે:
\[
(!A \lif !B) \lif ((!B \lif !C) \lif (!A \lif !C))
\]
અને \MP{} બે વાર વાપરવું પડે. ઘણી વાર કોઈ ઉપ-સૂત્રને તેની સાથે
સમકક્ષ હોવાનું આપણને જાણીતું હોય એવા સૂત્રથી બદલવું હોય છે; દાખલા
તરીકે, $\Diamond
  !A$ને $\lnot\Box\lnot !A$થી, અથવા
$\lnot\lnot !A$ને $!A$થી. તેથી આખી નિષ્પત્તિ લખવાને બદલે
વચ્ચેનાં પગલાં નિષ્પન્ન કરી શકાય એવું કેમ છે એટલું નોંધવું વધુ અનુકૂળ છે. આ માટે
નિષ્પન્નક્ષમતા વિશે કેટલીક હકીકતો એકત્ર કરીએ.

\begin{prop}
  જો $\Log{K} \Proves !A_1$, \dots, $\Log{K} \Proves !A_n$ હોય અને
  $!A_1$, \dots, $!A_n$માંથી વિધાનાત્મક તર્ક અનુસાર $!B$ ફલિત થતું હોય,
  તો $\Log{K} \Proves !B$.
\end{prop}

\begin{proof}
  જો $!A_1$, \dots,~$!A_n$માંથી વિધાનાત્મક તર્ક અનુસાર $!B$ ફલિત થતું હોય, તો
  \[
  !A_1 \lif (!A_2 \lif \cdots (!A_n \lif !B)\dots)
  \]
  પુનરુક્તિનો દૃષ્ટાંત છે. \MP{}ને $n$ વાર વાપરવાથી~$!B$ની
  નિષ્પત્તિ મળે છે.
\end{proof}

આ પ્રસ્તાવના ઉપયોગને આપણે~\PL વડે દર્શાવીશું.

\begin{prop}
  જો $\Log{K} \Proves !A_1 \lif (!A_2 \lif \cdots (!A_{n-1} \lif
  !A_n)\dots)$, તો $\Log{K} \Proves \Box !A_1 \lif (\Box !A_2 \lif
  \cdots (\Box !A_{n-1} \lif \Box !A_n)\dots)$.
\end{prop}

\begin{proof}
  \olref[nor]{prop:rk}ની સાબિતીની જેમ $n$ પર આગમનથી.
\end{proof}

આ પ્રસ્તાવના ઉપયોગને આપણે~\RK વડે દર્શાવીશું. આ પરિણામો વડે
નિષ્પન્નક્ષમતા વિશેના દાવા કેવી રીતે સહેલાઈથી સ્થાપી શકાય તે જોઈએ.

\begin{prop}
  $\Log{K} \Proves (\Box!A \land \Box!B) \lif \Box (!A \land !B)$
\end{prop}

\begin{proof}
  \begin{derivation}
    1. & $\Log{K} \Proves !A \lif (!B \lif (!A \land !B))$ & \Taut \\
    2. & $\Log{K} \Proves \Box!A \lif (\Box !B \lif \Box(!A \land !B)))$ & \RK, 1\\
    3. & $\Log{K} \Proves (\Box!A \land \Box!B) \lif \Box(!A \land !B)$ & \PL, 2
  \end{derivation}
\end{proof}

\begin{prop}\ollabel{prop:rewriting}
  જો $\Log{K} \Proves !A \liff !B$ અને $\Log{K} \Proves
  \Subst{!C}{!A}{q}$ હોય, તો $\Log{K} \Proves \Subst{!C}{!B}{q}$.
\end{prop}\sourcecorrection{OLMD-025}{મૂળ પ્રસ્તાવના અંતિમ પ્રતિસ્થાપનમાં B પહેલાં સૂત્રચિહ્ન ! છૂટી ગયું છે; પૂર્વપક્ષ અને પછીના પ્રશ્નની જેમ અહીં !B મૂક્યું છે.}

\begin{proof}
  અભ્યાસપ્રશ્ન.
\end{proof}

\begin{prob}
  $!C$ની જટિલતા પર આગમન વડે, જો $\Log{K} \Proves !A \liff !B$ હોય,
  તો $\Log{K} \Proves \Subst{!C}{!A}{q} \liff \Subst{!C}{!B}{q}$
  સાબિત કરો અને તેના વડે \olref[nml][prf][der]{prop:rewriting} સાબિત કરો.
\end{prob}

આ પ્રસ્તાવ ખાસ કરીને $\Diamond$ને $\Box$માં (અથવા વિપરીત રીતે)
ફેરવવો હોય, કે સૂત્રની અંદર બેવડો નિષેધ દૂર કરવો હોય, ત્યારે
ઉપયોગી છે. આગળ જ્યારે~$!C(!A)$ને સૂત્ર~$!C(!B)$ તરીકે ફરી
લખીશું, ત્યારે \olref{prop:rewriting}ના ઉપયોગને ``$!A$ના સ્થાને $!B$''
વડે દર્શાવીશું. બીજા શબ્દોમાં, ``$!A$ના સ્થાને $!B$'' આનું સંક્ષેપ છે:\sourcecorrection{OLMD-026}{મૂળ સામાન્ય સંક્ષેપમાં ``A for B'' છે, છતાં તેની નીચેની ત્રણ પંક્તિમાં C(A)માંથી C(B) મળે છે અને પછીના ઉદાહરણમાં p જૂના બેવડા નિષેધને બદલે છે; બદલીની દિશા અહીં સ્પષ્ટ કરી છે.}
  \begin{derivation}
    & $\Proves !C(!A)$ \\
    & $\Proves !A \liff !B$\\
    & $\Proves !C(!B)$ & \olref{prop:rewriting} વડે
  \end{derivation}
દાખલા તરીકે:

\begin{prop}
  $\Log{K} \Proves \lnot\Box p \lif \Diamond \lnot p$
\end{prop}

\begin{proof}
  % Diamond મૂળભૂત છે
  \begin{derivation}
    1. & $\Log{K} \Proves \Diamond \lnot p \liff \lnot\Box\lnot\lnot p$ &
    \Dual\\
    2. & $\Log{K} \Proves \lnot \Box \lnot\lnot p \lif \Diamond \lnot p$ & \PL, 1\\
    3. & $\Log{K} \Proves \lnot\Box p \lif \Diamond \lnot p$ & $\lnot\lnot p$ના સ્થાને $p$\\
  \end{derivation}
\end{proof}

ઉપરની નિષ્પત્તિમાં છેલ્લું પગલું ``$\lnot\lnot p$ના સ્થાને $p$''
આનો સંક્ષેપ છે:
  \begin{derivation}
    & $\Log{K} \Proves \lnot \Box \lnot\lnot p \lif \Diamond \lnot
    p$\\
    & $\Log{K} \Proves \lnot\lnot p \liff p$ & \Taut\\
    & $\Log{K} \Proves \lnot\Box p \lif \Diamond \lnot p$ & \olref{prop:rewriting} વડે
  \end{derivation}
અહીં \olref{prop:rewriting}ના $!C(q)$, $!A$ અને $!B$ની ભૂમિકામાં અનુક્રમે
$\lnot \Box q \lif \Diamond \lnot p$,
$\lnot\lnot p$ અને~$p$ આવે છે.

જ્યારે સૂત્રમાં ઉપ-સૂત્ર $\lnot\Diamond !A$ હોય, ત્યારે
\olref{prop:rewriting} વડે તેની જગ્યાએ $\Box\lnot !A$ મૂકી શકીએ,
કારણ કે $\Log{K} \Proves \lnot\Diamond !A \liff \Box\lnot !A$.
આવા અને સમાન પ્રતિસ્થાપનોને આપણે માત્ર ``$\lnot\Diamond$ના સ્થાને
$\Box\lnot$'' વડે દર્શાવીશું.

નીચેનો પ્રસ્તાવ નિષ્પન્નક્ષમતા વિશેનાં પરિણામો ઢાંચારૂપે સ્થાપવાનું
સમર્થન આપે છે. દાખલા તરીકે, અગાઉનો પ્રસ્તાવ માત્ર
$\Log{K} \Proves \lnot\Box p \lif
\Diamond \lnot p$ જ નહિ, પરંતુ મનસ્વી~$!A$ માટે
$\Log{K} \Proves \lnot\Box !A \lif \Diamond
\lnot !A$ પણ સ્થાપે છે.

\begin{prop}
  જો $!A$ એ $!B$નો પ્રતિસ્થાપન દૃષ્ટાંત હોય અને $\Log{K} \Proves !B$,
  તો $\Log{K} \Proves !A$.
\end{prop}

\begin{proof}
  આખી નિષ્પત્તિમાં પ્રતિસ્થાપન કરતાં પણ યોગ્ય નિષ્પત્તિ
  મળે છે તે ચકાસવું લાંબું, પરંતુ નિયમિત કામ છે ($!B$ની
  નિષ્પત્તિની લંબાઈ પર આગમન કરો). ખાસ કરીને, પુનરુક્તિના
  દૃષ્ટાંતોના પ્રતિસ્થાપન દૃષ્ટાંતો પોતે પુનરુક્તિના દૃષ્ટાંતો છે;
  \Dual{} અને~\Ax{K}ના દૃષ્ટાંતોના પ્રતિસ્થાપન
  દૃષ્ટાંતો પોતે~\Dual{}
    અને~\Ax{K}ના દૃષ્ટાંતો છે; અને પૂર્વપક્ષ(ઓ) તથા નિષ્કર્ષમાં
  પ્રતિજ્ઞાપન ચલોના સ્થાને સૂત્રો મૂકતાં \MP{} અને~\Nec{}
  નો ઉપયોગ યોગ્ય જ રહે છે.
\end{proof}
% END OLP-0432
\endgroup

\begingroup
\makeatletter\def\ol@sectioncs{section}\makeatother

% BEGIN OLP-0433
\olfileid{nml}{prf}{mpr}

\olsection{\Log{K}માં વધુ પ્રમાણો}
\phantomsection\label{guunit:OLP-0433}


હવે \olref[der]{sec}માં રજૂ કરેલી સરળ પદ્ધતિ વાપરીને~\Log{K}માં
નિષ્પન્નક્ષમતાનાં થોડાં વધુ ઉદાહરણો જોઈએ.

\begin{prop}
  $\Log{K} \Proves \Box (!A \lif !B) \lif (\Diamond !A \lif \Diamond
  !B)$
\end{prop}

\begin{proof}
\begin{derivation}
1. & $\Log{K} \Proves (!A \lif !B) \lif (\lnot !B \lif \lnot !A)$ & \PL \\
2. &  $\Log{K} \Proves \Box(!A \lif !B) \lif (\Box\lnot!B \lif \Box\lnot!A)$ & \RK, 1 \\
3. & $\Log{K} \Proves (\Box\lnot!B \lif \Box\lnot!A) \lif (\lnot
\Box\lnot!A \lif \lnot\Box\lnot!B)$ & \Taut\\
4. & $\Log{K} \Proves \Box(!A \lif !B) \lif (\lnot
\Box\lnot!A \lif \lnot\Box\lnot!B)$ & \PL, 2, 3 \\
5. &  $\Log{K} \Proves \Box(!A \lif !B) \lif (\Diamond!A \lif \Diamond!B)$ & $\lnot\Box\lnot$ના સ્થાને $\Diamond$.
\end{derivation}
\end{proof}
  
\begin{prop}
$\Log{K}\Proves \Box!A \lif (\Diamond(!A \lif !B) \lif
  \Diamond !B)$
\end{prop}

\begin{proof}
  \begin{derivation}
  1. & $\Log{K} \Proves !A \lif (\lnot!B \lif \lnot (!A \lif !B))$ & \Taut \\
  2. & $\Log{K} \Proves \Box!A \lif  (\Box\lnot!B \lif \Box\lnot (!A\lif !B))$ &
  \RK, 1 \\
  3. &  $\Log{K} \Proves \Box!A \lif (\lnot \Box\lnot (!A\lif !B) \lif \lnot
  \Box\lnot!B)$ & \PL, 2 \\
  4. & $\Log{K} \Proves \Box!A \lif (\Diamond(!A \lif !B) \lif
  \Diamond !B)$ & $\lnot\Box\lnot$ના સ્થાને $\Diamond$.
  \end{derivation}
\end{proof}

\begin{prop}
  $\Log{K} \Proves (\Diamond!A \lor \Diamond!B) \lif \Diamond(!A \lor !B)$
\end{prop}

\begin{proof}
  \begin{derivation}
    1. & $\Log{K} \Proves \lnot(!A \lor !B) \lif \lnot!A$ & \Taut \\
    2. & $\Log{K} \Proves \Box\lnot(!A \lor !B) \lif \Box\lnot!A$ & \RK, 1 \\
    3. & $\Log{K} \Proves \lnot\Box\lnot !A \lif \lnot\Box\lnot(!A \lor!B)$ &
    \PL, 2\\
    4. & $\Log{K} \Proves \Diamond!A \lif \Diamond(!A \lor !B)$ &
    $\lnot\Box\lnot$ના સ્થાને $\Diamond$\\
    5. & $\Log{K} \Proves \Diamond!B \lif \Diamond(!A \lor !B)$ & એ જ રીતે\\
    6. & $\Log{K} \Proves (\Diamond!A \lor\Diamond!B) \lif \Diamond(!A \lor !B)$
    & \PL, 4, 5.
  \end{derivation}
\end{proof}

\begin{prop}
  $\Log{K} \Proves \Diamond(!A \lor!B) \lif (\Diamond!A \lor \Diamond!B)$
\end{prop}

\begin{proof}
  \begin{derivation}
    1. & $\Log{K} \Proves \lnot !A \lif (\lnot !B \lif \lnot (!A \lor !B))$ & \Taut \\
    2. & $\Log{K} \Proves \Box\lnot !A \lif
    (\Box\lnot!B \lif \Box \lnot (!A \lor !B))$ & \RK\\
    3. & $\Log{K} \Proves \Box\lnot !A \lif (\lnot \Box \lnot (!A \lor!B)
    \lif \lnot\Box\lnot !B)$ & \PL, 2\\
    4. & $\Log{K} \Proves \lnot \Box \lnot(!A \lor !B) \lif (\Box \lnot !A \lif 
    \lnot\Box\lnot !B)$ & \PL, 3\\
    5. & $\Log{K} \Proves \lnot \Box \lnot(!A\lor!B) \lif (\lnot
    \lnot\Box\lnot!B \lif \lnot\Box\lnot!A)$ & \PL, 4\\
    6. & $\Log{K} \Proves \Diamond(!A \lor !B) \lif (\lnot
    \Diamond!B \lif \Diamond!A)$ & $\lnot\Box\lnot$ના સ્થાને $\Diamond$\\
    7. & $\Log{K} \Proves \Diamond(!A\lor!B) \lif (\Diamond!B \lor \Diamond!A)$ & \PL, 6. \\
  \end{derivation}
\end{proof}

\begin{prob}
  નીચેના નિષ્પન્નક્ષમતાના દાવા સાચા છે તે બતાવો:
  \begin{enumerate}
  \item $\Log{K} \Proves \Diamond \lnot \lfalse \lif (\Box !A \lif
    \Diamond !A)$;
  \item $\Log{K} \Proves \Box(!A \lor !B) \lif (\Diamond !A \lor \Box
    !B)$;
  \item $\Log{K} \Proves (\Diamond !A \lif \Box !B) \lif \Box(!A \lif
    !B)$.
  \end{enumerate}
\end{prob}
% END OLP-0433
\endgroup

\begingroup
\makeatletter\def\ol@sectioncs{section}\makeatother

% BEGIN OLP-0434
\olfileid{nml}{prf}{dua}

\olsection{દ્વૈત સૂત્રો}
\phantomsection\label{guunit:OLP-0434}


\begin{defn}\ollabel{def:duals}
  સૂત્રો \Ax{T}, \Ax{B}, \Ax{4} અને \Ax{5} દરેકનું
  \emph{દ્વૈત} છે; તેને નીચેની જેમ અધોલિખિત હીરાથી દર્શાવીએ છીએ:
  \begin{align}
    \tag{\Ax{T_\Diamond}} p & \lif \Diamond p\\
    \tag{\Ax{B_\Diamond}} \Diamond\Box p & \lif p\\
    \tag{\Ax{4_\Diamond}} \Diamond\Diamond p & \lif \Diamond p\\
    \tag{\Ax{5_\Diamond}} \Diamond\Box p & \lif \Box p
    \end{align}
\end{defn}

ઉપરના દરેક દ્વૈત સૂત્રને અનુરૂપ સૂત્રમાંથી મેળવવા પહેલાં
$p$ના સ્થાને $\lnot p$ મૂકો, પછી પ્રતિપક્ષી વિધાન લો, પછી
$\lnot\Box\lnot$ના સ્થાને $\Diamond$ અને $\lnot\Diamond\lnot$ના
સ્થાને~$\Box$ મૂકો. આ અર્થમાં \Ax{D}, એટલે કે
$\Box!A \lif \Diamond!A$, પોતાનું જ દ્વૈત છે.

\begin{prop}\ollabel{prop:dualsys}
  \olref[nml][prf][dua]{def:duals}માંના દરેક સૂત્ર~$!A$ માટે:
  $\Log{K}!A = \Log{K}!A_{\Diamond}$.
\end{prop}
\begin{proof}
  અભ્યાસપ્રશ્ન.
\end{proof}
\begin{prob}
  \olref[nml][prf][dua]{prop:dualsys} સાબિત કરો.
\end{prob}
% END OLP-0434
\endgroup

\begingroup
\makeatletter\def\ol@sectioncs{section}\makeatother

% BEGIN OLP-0435
\olfileid{nml}{prf}{prs}

\olsection{મોડલ તંત્રોમાં પ્રમાણો}
\phantomsection\label{guunit:OLP-0435}


હવે આપણે~\Log{K} સિવાયનાં મોડલ તર્કનાં તંત્રોમાં પ્રમાણો જોઈએ.

\begin{prop}\ollabel{prop:S5facts}
નીચેના નિષ્પન્નતાના પરિણામો સાચાં છે:
  \begin{enumerate}
  \item $\Log{KT5} \Proves \Ax{B}$;
  \item $\Log{KT5} \Proves \Ax{4}$;
  \item $\Log{KDB4} \Proves \Ax{T}$;
  \item $\Log{KB4} \Proves \Ax{5}$;
  \item $\Log{KB5} \Proves \Ax{4}$;
  \item \ollabel{prop:S5facts-KT-D}$\Log{KT} \Proves \Ax{D}$.
  \end{enumerate}
\end{prop}

\begin{proof}
  આપણે દરેક માટે પ્રમાણ આપીએ છીએ.
  \begin{enumerate}
    \item $\Log{KT5} \Proves \Ax{B}$:
      \begin{derivation}
        1. & $\Log{KT5} \Proves \Diamond!A \lif \Box\Diamond!A$ & \Ax{5}\\
        2. & $\Log{KT5} \Proves !A \lif \Diamond!A$ & $\Ax{T_\Diamond}$\\
        3. & $\Log{KT5} \Proves !A \lif \Box\Diamond!A$ & \PL, 2, 1.
      \end{derivation}
    \item $\Log{KT5} \Proves \Ax{4}$:
      \begin{derivation}
        1. &$\Log{KT5} \Proves \Diamond\Box!A \lif \Box\Diamond\Box!A$ & \Ax{5}
        જેમાં $p$ના સ્થાને $\Box!A$ મૂક્યું છે\\
        2. & $\Log{KT5} \Proves \Box!A \lif \Diamond\Box!A$ & $\Ax{T_\Diamond}$
        જેમાં $p$ના સ્થાને $\Box!A$ મૂક્યું છે\\
        3. & $\Log{KT5} \Proves \Box!A \lif \Box\Diamond\Box!A$ & \PL, 2, 1\\
        4. & $\Log{KT5} \Proves \Diamond\Box!A \lif \Box!A$ & $\Ax{5_\Diamond}$\\
        5. & $\Log{KT5} \Proves \Box\Diamond\Box!A \lif \Box\Box!A$ & \RK{}, 4 \\
        6. & $\Log{KT5} \Proves \Box!A \lif \Box\Box!A$ & \PL, 3, 5. \\
      \end{derivation}
    \item $\Log{KDB4} \Proves \Ax{T}$:
      \begin{derivation}
        1. & $\Log{KDB4} \Proves \Diamond\Box!A \lif !A $ & $\Ax{B_\Diamond}$ \\
        2. & $\Log{KDB4} \Proves \Box\Box!A \lif \Diamond\Box!A$ & $\Ax{D}$
        જેમાં $p$ના સ્થાને $\Box!A$ મૂક્યું છે\\
        3. & $\Log{KDB4} \Proves \Box\Box!A \lif !A$ & \PL, 1, 2\\
        4. & $\Log{KDB4} \Proves \Box!A \lif \Box\Box!A$ & \Ax{4} \\
        5. & $\Log{KDB4} \Proves \Box!A \lif !A$ & \PL, 4, 3. \\
      \end{derivation}
    \item $\Log{KB4} \Proves \Ax{5}$:
      \begin{derivation}
        1. & $\Log{KB4} \Proves \Diamond!A \lif \Box \Diamond\Diamond!A$ & \Ax{B}
        જેમાં $p$ના સ્થાને $\Diamond!A$ મૂક્યું છે\\
        2. & $\Log{KB4} \Proves \Diamond\Diamond !A \lif \Diamond !A$ &
        $\Ax{4_\Diamond}$ \\
        3. & $\Log{KB4} \Proves \Box\Diamond\Diamond !A \lif \Box \Diamond!A$ &
        \RK{}, 2 \\
        4. & $\Log{KB4} \Proves \Diamond!A \lif \Box\Diamond!A$ & \PL, 1, 3.
      \end{derivation}
    \item $\Log{KB5} \Proves \Ax{4}$:
      \begin{derivation}
        1. & $\Log{KB5} \Proves \Box!A \lif \Box\Diamond\Box !A$ & \Ax{B};
        $p$ના સ્થાને $\Box!A$ મૂક્યું છે \\
        2. & $\Log{KB5} \Proves  \Diamond\Box!A \lif \Box !A$ &
        $\Ax{5_\Diamond}$\\
        3. & $\Log{KB5} \Proves  \Box\Diamond\Box!A \lif \Box\Box !A$ & \RK{}, 2 \\
        4. & $\Log{KB5} \Proves  \Box!A \lif \Box\Box!A$ & \PL, 1, 3.
      \end{derivation}
    \item $\Log{KT} \Proves \Ax{D}$:
      \begin{derivation}
        1. & $\Log{KT} \Proves \Box !A \lif !A$ & \Ax{T} \\
        2. & $\Log{KT} \Proves !A \lif \Diamond!A$ & $\Ax{T_\Diamond}$ \\
        3. & $\Log{KT} \Proves \Box !A \lif \Diamond !A$ &  \PL, 1, 2
      \end{derivation}
  \end{enumerate}
\end{proof}

\begin{defn}
  પ્રચલિત પરંપરા અનુસાર આપણે \Log{S4}ને તંત્ર
  \Log{KT4} અને \Log{S5}ને તંત્ર \Log{KTB4} તરીકે વ્યાખ્યાયિત કરીએ છીએ.
\end{defn}

નીચેનો પ્રસ્તાવ બતાવે છે કે પરંપરાગત તંત્ર \Log{S5} માટે અનેક સમકક્ષ
સ્વયંસિદ્ધિકરણો છે. તેમાં આશ્ચર્ય નથી, કારણ કે સ્વયંસિદ્ધિઓનાં આ વિવિધ
સંયોજનો બધાં સામ્ય સંબંધોને લક્ષણિત કરે છે
(જુઓ \olref[frd][es5]{prop:equivalences}).

\begin{prop}\ollabel{prop:S5}
  $\Log{KTB4} = \Log{KT5} = \Log{KDB4} = \Log{KDB5}$.
\end{prop}

\begin{proof}
  અભ્યાસપ્રશ્ન.
\end{proof}

\begin{prob}
  \olref[nml][prf][prs]{prop:S5} સાબિત કરો.
\end{prob}
% END OLP-0435
\endgroup

\begingroup
\makeatletter\def\ol@sectioncs{section}\makeatother

% BEGIN OLP-0436
\olfileid{nml}{prf}{snd}

\olsection{યથાર્થતા}
\phantomsection\label{guunit:OLP-0436}


જો નિષ્પત્તિ તંત્રમાં જે કંઈ નિષ્પન્ન કરી શકાય તે બધું
પ્રામાણ્ય હોય, તો તે તંત્ર યથાર્થ કહેવાય છે. મોડલ તંત્રોમાં, એટલે કે
નિષ્પત્તિઓમાં જ્યાં \Ax{K} ઉપરાંત કેટલાંક સૂત્રો $!A_1$,
\dots,~$!A_n$ના દૃષ્ટાંતો વાપરી શકાય છે, આપણે ઇચ્છીએ છીએ કે દરેક
નિષ્પન્ન કરી શકાય એવું સૂત્ર એવા દરેક નિદર્શનમાં સત્ય હોય જેમાં $!A_1$,
\dots, $!A_n$ સત્ય હોય.

\begin{thm}[યથાર્થતાનું પ્રમેય]\ollabel{thm:soundness}
  જો $!A_1$, \dots, $!A_n$ના દરેક દૃષ્ટાંત અનુક્રમે નિદર્શનોના
  વર્ગો $\mClass{C}_1$, \dots, $\mClass{C}_n$માં પ્રામાણ્ય હોય,
  તો $\Log{K}!A_1\dots !A_n
  \Proves !B$ હોય ત્યારે $!B$ નિદર્શનોના વર્ગ
  $\mClass{C}_1 \cap \dots \cap \mClass{C}_n$માં પ્રામાણ્ય છે.
\end{thm}

\begin{proof}
  પ્રમાણોની લંબાઈ પર આગમનથી. સંક્ષેપમાં $\mClass{C} =
  \mClass{C}_1 \cap \dots \cap \mClass{C}_n$ મૂકો.
  \begin{enumerate}
  \item આગમનનો આધાર: જો $!B$નું પ્રમાણ લંબાઈ~$1$નું હોય, તો તે
    કાં તો પુનરુક્તિનો દૃષ્ટાંત, કાં~\Ax{K}નો દૃષ્ટાંત,
    કાં~\Dual{}નો દૃષ્ટાંત, અથવા $!A_1$, \dots,~$!A_n$ પૈકી કોઈ
    એકનો દૃષ્ટાંત છે. પ્રથમ કિસ્સામાં $!B$ એ $\mClass{C}$માં
    પ્રામાણ્ય છે, કારણ કે \olref[syn][tau]{prop:valid-taut} મુજબ
    પુનરુક્તિના દૃષ્ટાંતો નિદર્શનોના \emph{કોઈપણ} વર્ગમાં પ્રામાણ્ય છે.
    એ જ રીતે બીજા કિસ્સામાં
    \olref[syn][sch]{prop:Kvalid} અને
      \olref[syn][sch]{prop:Dual-valid} મુજબ પ્રામાણ્યતા મળે છે.
    અંતે ત્રીજા કિસ્સામાં, $!B$ એ $\mClass{C}_i$માં પ્રામાણ્ય છે અને
    $\mClass{C} \subseteq \mClass{C}_i$, તેથી
    \olref[syn][val]{prop:subset-class} મુજબ $!B$ એ $\mClass{C}$માં
    પણ પ્રામાણ્ય છે.
  \item આગમનનું પગલું: માનો કે $!B$નું પ્રમાણ લંબાઈ $k>1$નું છે. જો
    $!B$ પુનરુક્તિનો દૃષ્ટાંત, $!A_1$, \dots, $!A_n$ પૈકી કોઈનો
    દૃષ્ટાંત, અથવા ઉપર ગણાવેલી બીજી સ્વયંસિદ્ધિઓમાંથી કોઈનો દૃષ્ટાંત
    હોય, તો અગાઉના પગલાની જેમ જ ચાલીએ.\sourcecorrection{OLMD-027}{મૂળના આ આગમન-પગલામાં માત્ર પુનરુક્તિ અને વધારાના $!A_i$ દૃષ્ટાંતો ફરી ગણાવ્યા છે; K તથા શરતી Dualના દૃષ્ટાંતો પણ આધારકિસ્સાની જ રીતથી આવરે છે, તેથી બધા સ્વયંસિદ્ધિ-કિસ્સા સ્પષ્ટ કર્યા છે.}
    હવે માનો કે $!B$ અગાઉનાં સૂત્રો $!C \lif !B$ અને $!C$માંથી
    \MP{} વડે મળ્યું છે. ત્યારે $!C \lif !B$ અને $!C$નાં પ્રમાણોની
    લંબાઈ $<k$ છે અને આગમન-પરિકલ્પના મુજબ તે~$\mClass{C}$માં પ્રામાણ્ય
    છે. \olref[syn][sch]{prop:soundMP} મુજબ $!B$ પણ
    $\mClass{C}$માં પ્રામાણ્ય છે. અંતે માનો કે $!B$ એ $!C$માંથી
    \Nec{} વડે મળ્યું છે (એટલે $!B = \Box!C$). આગમન-પરિકલ્પના મુજબ
    $!C$ એ $\mClass{C}$માં પ્રામાણ્ય છે અને
    \olref[syn][val]{prop:Nec-rule} મુજબ~$!B$ પણ પ્રામાણ્ય છે.
  \end{enumerate}
\end{proof}
% END OLP-0436
\endgroup

\begingroup
\makeatletter\def\ol@sectioncs{section}\makeatother

% BEGIN OLP-0437
\olfileid{nml}{prf}{dis}

\olsection{તંત્રો ભિન્ન છે તે બતાવવું}
\phantomsection\label{guunit:OLP-0437}


\olref[prs]{sec}માં આપણે બે મોડલ તર્ક-તંત્રો હકીકતમાં એક જ તંત્ર છે તે
કેવી રીતે સાબિત કરવું જોયું. \olref[snd]{thm:soundness} વડે બે મોડલ
તંત્રો $\Sigma$ અને $\Sigma'$ ભિન્ન છે તે બતાવી શકાય: એવું
સૂત્ર~$!A$ શોધો કે $\Sigma' \Proves !A$, પરંતુ તે~$\Sigma$ના
કોઈ નિદર્શનમાં અસત્ય હોય.

\begin{prop}
  $\Log{KD} \subsetneq \Log{KT}$
\end{prop}

\begin{proof} આ બધા સ્વવાચક સંબંધો સિરિયલ હોય છે એ અર્થાનુસારી હકીકતનો
  વાક્યરચનાલક્ષી સમકક્ષ છે. $\Log{KD} \subseteq
  \Log{KT}$ બતાવવા $\Log{KD} \Proves !B$ હોય તો
  $\Log{KT} \Proves !B$ થાય છે તે જોવું પડે. તે $\Log{KT} \Proves
  \Ax{D}$ પરથી મળે છે, જેમ
  \olref[prs]{prop:S5facts}\olref[prs]{prop:S5facts-KT-D}માં બતાવ્યું છે.\sourcecorrection{OLMD-028}{મૂળમાં અહીં KTમાંથી 'Log D' નિષ્પન્ન થાય એમ છે; નિષ્પન્ન થતું તંત્ર નહિ પરંતુ Dનું સ્વયંસિદ્ધિ-સૂત્ર છે, તેથી Ax D મૂક્યું છે.}
  સમાવેશ ચુસ્ત છે તે બતાવવા યથાર્થતા
  (\olref[snd]{thm:soundness}) મુજબ \Log{KD}નું એવું નિદર્શન આપવું પૂરતું
  છે જેમાં \Ax{T}, એટલે કે $\Box p \lif p$, અસત્ય હોય (આ સરળ કામ
  અભ્યાસપ્રશ્ન તરીકે છોડ્યું છે). ત્યારે યથાર્થતા મુજબ $\Log{KD}
  \Proves/ \Box p \lif p$.
\end{proof}

\begin{prop}
  $\Log{KB} \neq \Log{K4}$. 
\end{prop}

\begin{proof}
  આપણે એવું સંમિત નિદર્શન રચીએ છીએ જેમાં \Ax{4}નો કોઈ દૃષ્ટાંત અસત્ય
  હોય. તે દૃષ્ટાંત \Log{K4}માં સ્પષ્ટપણે નિષ્પન્ન કરી શકાય એવું છે, પરંતુ
  \Log{KB}માં નથી; તેથી $\Log{K4} \nsubseteq \Log{KB}$ ફલિત થશે.
  \olref{fig:Bnot4}નું સંમિત નિદર્શન~$\mModel{M}$ લો. નિદર્શન સંમિત
  હોવાથી~$\mModel{M}$માં \Ax{K} અને~\Ax{B} સત્ય છે
  (અનુક્રમે \olref[syn][sch]{prop:Kvalid} અને
  \olref[frd][acc]{thm:soundschemas} મુજબ). પરંતુ
  $\mSat/{M}{\Box p \lif \Box\Box p}[w_1]$.
\end{proof}

\begin{figure}[htpb]
  \centering
  \begin{tikzpicture}[modal]
    \node[world] (w1) [label=above:\mFalse{p},
    label={below:$\mSat{{}}{\Box p}$\\ $\mSat/{{}}{\Box\Box p}$}] {$w_1$}; 
    \node[world] (w2) [label=above:\mTrue{p},
      label=below:$\mSat/{{}}{\Box p}$, right=of w1]{$w_2$}; 
    \draw[->, bend left] (w1) to (w2); 
    \draw[->, bend left] (w2) to (w1); 
  \end{tikzpicture}
  \caption{\Ax{4}ના દૃષ્ટાંતને અસત્ય બનાવતું સંમિત નિદર્શન.}
  \ollabel{fig:Bnot4}
\end{figure}
 
\begin{thm}\ollabel{thm:KTBnot45}
  $\Log{KTB} \Proves/ \Ax{4}$ અને $\Log{KTB} \Proves/ \Ax{5}$.
\end{thm}\sourcecorrection{OLMD-029}{મૂળ પ્રમેયમાં KTBમાંથી 'Log 4' અને 'Log 5' નિષ્પન્ન નથી એમ લખાયું છે; નિષ્પન્નતાના જમણે તંત્રના નામને બદલે અનુક્રમે 4 અને 5નાં સ્વયંસિદ્ધિ-સૂત્રો મૂક્યાં છે.}

\begin{proof}
  \olref[frd][acc]{thm:soundschemas} મુજબ \Ax{T} અને \Ax{B}ના બધા
  દૃષ્ટાંતો અનુક્રમે દરેક સ્વવાચક અને સંમિત નિદર્શનમાં સત્ય છે. તેથી
  યથાર્થતા મુજબ એવું સ્વવાચક-સંમિત નિદર્શન શોધવું પૂરતું છે જેના કોઈ
  જગતમાં~\Ax{4}નો દૃષ્ટાંત અસત્ય હોય, અને~\Ax{5} માટે પણ એમ જ. બંને
  દાવા માટે આપણે એક જ નિદર્શન વાપરીએ છીએ.
  \olref{fig:KTBnot45}નું સંમિત, સ્વવાચક નિદર્શન લો. ત્યારે
  $\mSat/{M}{\Box p \to \Box \Box
    p}[w_1]$, તેથી \Ax{4} એ $w_1$માં અસત્ય છે. એ જ રીતે,
  $\mSat/{M}{\Diamond
    \lnot p \lif \Box \Diamond \lnot p}[w_2]$, તેથી $!A = \lnot p$
  મૂકેલો~\Ax{5}નો દૃષ્ટાંત~$w_2$માં અસત્ય છે.
\end{proof}

\begin{figure}[htpb]
  \centering
  \begin{tikzpicture}[modal]
    \node[world] (w1) [label={right:\mTrue{p}},
      label={below: $\mSat{{}}{\Box p}$\\
        $\mSat/{{}}{\Box\Box p}$\\
        $\mSat/{{}}{\Diamond\lnot p}$}] {$w_1$};
    \draw[reflexive above] (w1) to (w1);
    \node[world] (w2) [label={right:\mTrue{p}}, label={below:
        $\mSat{{}}{\Diamond\lnot p}$\\
        $\mSat/{{}}{\Box\Diamond\lnot p}$}, right=of w1] {$w_2$} ;
    \draw[reflexive above] (w2) to (w2);
    \node[world] (w3) [label={right:\mFalse{p}},right=of w2] {$w_3$};
    \draw[reflexive above] (w3) to (w3);
    \draw[->, bend left] (w1) to (w2);  
    \draw[->, bend left] (w2) to (w3);  
    \draw[->, bend left] (w3) to (w2);  
    \draw[->, bend left] (w2) to (w1);  
  \end{tikzpicture}
  \caption{\olref{thm:KTBnot45} માટેનું નિદર્શન.}
  \ollabel{fig:KTBnot45}
\end{figure}

\begin{thm}\ollabel{thm:KD5not4}
  $\Log{KD5} \neq \Log{KT4} = \Log{S4}$.
\end{thm}

\begin{proof}
  \olref[frd][acc]{thm:soundschemas} મુજબ \Ax{D} અને~\Ax{5}ના બધા
  દૃષ્ટાંતો દરેક સિરિયલ યુક્લિડિયન નિદર્શનમાં સત્ય છે. તેથી એવું
  સિરિયલ યુક્લિડિયન નિદર્શન શોધવું પૂરતું છે જેના કોઈ જગતમાં~\Ax{4}નો
  દૃષ્ટાંત અસત્ય હોય. \olref{fig:KD5not4}નું નિદર્શન લો અને નોંધો કે
  $\mSat/{M}{\Box p
  \lif \Box\Box p}[w_1]$.
\end{proof}

\begin{figure}[t]
  \centering
  \begin{tikzpicture}[modal]
    \node[world] (w2) [label=north west:\mTrue{p}]{$w_2$} ;
    \draw[reflexive left] (w2) to (w2);
    \node[world] (w1) [label=right:\mFalse{p},
      label=below:{$\mSat{{}}{\Box p}, \mSat/{{}}{\Box\Box p}$},
      below right=of w2]{$w_1$};
    \node[world] (w3) [label=north east:\mTrue{p}, above right=of w1]{$w_3$} ;
    \draw[reflexive right] (w3) to (w3);
    \node[world] (w4) [label=right:\mFalse{p},above right=of w2]{$w_4$};
    \draw[reflexive above] (w4) to (w4);
    \draw[->] (w1) to (w2);  
    \draw[->] (w1) to (w3);  
    \draw[->, bend left=15] (w2) to (w3);
    \draw[->, bend left=15] (w3) to (w2);  
    \draw[->, bend left=15] (w2) to (w4);
    \draw[->, bend left=15] (w4) to (w2);  
    \draw[->, bend left=15] (w3) to (w4);
    \draw[->, bend left=15] (w4) to (w3);  
  \end{tikzpicture}
  \caption{\olref{thm:KD5not4} માટેનું નિદર્શન.}
  \ollabel{fig:KD5not4}
\end{figure}

\begin{prob}
  $3$ જગતોવાળું નિદર્શન વાપરીને \olref[nml][prf][dis]{thm:KD5not4}નું
  વૈકલ્પિક પ્રમાણ આપો.
\end{prob}

\begin{prob}
  $\Log{KT4} \Proves/ \Ax{B}$ અને $\Log{KT4} \Proves/ \Ax{5}$ બંને
  બતાવતું એક જ સ્વવાચક પરંપરિત નિદર્શન આપો.
\end{prob}
% END OLP-0437
\endgroup

\begingroup
\makeatletter\def\ol@sectioncs{section}\makeatother

% BEGIN OLP-0438
\olfileid{nml}{prf}{prg}

\olsection{સૂત્રોના ગણમાંથી નિષ્પન્નક્ષમતા}
\phantomsection\label{guunit:OLP-0438}


\olref[prf][prf]{sec}માં આપણે તંત્ર~$\Sigma$માં સૂત્રની
નિષ્પન્નતાનો ખ્યાલ વ્યાખ્યાયિત કર્યો હતો. હવે આ ખ્યાલને ગણ~$\Gamma$
માંનાં સૂત્રોમાંથી $\Sigma$માં નિષ્પન્નતા સુધી વિસ્તારીશું.\sourcecorrection{OLMD-030}{મૂળ સંદર્ભ prf/prs વિભાગ બતાવે છે, પરંતુ તંત્રમાં નિષ્પન્નતાની વ્યાખ્યા prf/prf વિભાગમાં છે; સંદર્ભને તે વિભાગ તરફ સુધાર્યો છે.}

\begin{defn}\ollabel{defn:Gammaproves}
  સૂત્ર~$!A$ એ ગણ~$\Gamma$માંના સૂત્રોમાંથી
  તંત્ર~$\Sigma$માં નિષ્પન્ન કરી શકાય એવું છે, જે $\Gamma \Proves[\Sigma] !A$
  લખાય છે, જો અને માત્ર જો $!B_1$, \dots, $!B_n \in \Gamma$ એવાં હોય કે
  $\Sigma \Proves !B_1 \lif (!B_2 \lif \cdots (!B_n \lif !A)
  \cdots)$.
\end{defn}
% END OLP-0438
\endgroup

\begingroup
\makeatletter\def\ol@sectioncs{section}\makeatother

% BEGIN OLP-0439
\olfileid{nml}{prf}{prp}

\olsection{નિષ્પન્નક્ષમતાના ગુણધર્મો}
\phantomsection\label{guunit:OLP-0439}


\begin{prop}\ollabel{prop:derivabilityfacts}
  માનો કે $\Sigma$ મોડલ તંત્ર અને $\Gamma$ મોડલ સૂત્રોનો ગણ છે.
  નીચેના ગુણધર્મો સાચા છે:
  \begin{enumerate}
  \item\ollabel{prop:derivabilityfacts-monotonicity} \emph{એકદિશવર્ધિતા}: જો
    $\Gamma \Proves[\Sigma] !A$ અને $\Gamma \subseteq \Delta$ હોય, તો
    $\Delta \Proves[\Sigma] !A$;
  \item\ollabel{prop:derivabilityfacts-reflexivity}
    \emph{સ્વવાચકતા}: જો $!A \in \Gamma$ હોય, તો $\Gamma
    \Proves[\Sigma] !A$;
  \item\ollabel{prop:derivabilityfacts-cut} \emph{કટ}: જો $\Gamma
    \Proves[\Sigma] !A$ અને $\Delta \cup \{!A\} \Proves[\Sigma] !B$
    હોય, તો $\Gamma \cup \Delta \Proves[\Sigma] !B$;
  \item\ollabel{prop:derivabilityfacts-deduction} \emph{નિષ્પત્તિ પ્રમેય}:
    $\Gamma \cup \{!B\}
    \Proves[\Sigma] !A$ જો અને માત્ર જો $\Gamma \Proves[\Sigma] !B \lif
      !A$;
  \item \ollabel{prop:derivabilityfacts-ruleT}\emph{નિયમ T}: જો
    $\Gamma \Proves[\Sigma] !A_1$ અને \dots અને $\Gamma
    \Proves[\Sigma] !A_n$ હોય અને $!A_1 \to (!A_2 \lif \cdots (!A_n \lif
    !B)\cdots)$ પુનરુક્તિનો દૃષ્ટાંત હોય, તો $\Gamma
    \Proves[\Sigma] !B$.\sourcecorrection{OLMD-031}{મૂળની છેલ્લી યાદી-પંક્તિમાં 'Rule T: If' ટિપ્પણીમાં હોવાથી છપાતું નથી; નિયમનું નામ અને 'જો'ની શરત દેખાય એમ પુનઃસ્થાપિત કર્યાં છે.}
  \end{enumerate}
\end{prop}

સાબિતી સરળ અભ્યાસપ્રશ્ન છે. \olref{prop:derivabilityfacts}ના
\olref{prop:derivabilityfacts-ruleT} ભાગથી, દાખલા તરીકે, જો
$\Gamma \Proves[\Sigma] !A \lor !B$ અને
$\Gamma \Proves[\Sigma] \lnot !A$ હોય, તો $\Gamma \Proves[\Sigma]
!B$ મળે છે. આગળ આપણે $\Gamma \cup \{ !A\} \Proves[\Sigma] !B$ના
સ્થાને $\Gamma, !A \Proves[\Sigma] !B$ લખીશું.

\begin{defn}
  ગણ $\Gamma$ એ તંત્ર~$\Sigma$ને અનુલક્ષીને \emph{નિષ્પત્તિ હેઠળ બંધ}
  છે જો અને માત્ર જો $\Gamma \Proves[\Sigma] !A$ હોય ત્યારે
  $!A \in \Gamma$ હોય.
\end{defn}
% END OLP-0439
\endgroup

\begingroup
\makeatletter\def\ol@sectioncs{section}\makeatother

% BEGIN OLP-0440
\olfileid{nml}{prf}{con}

\olsection{સુસંગતતા}
\phantomsection\label{guunit:OLP-0440}


સુસંગતતા એ સૂત્રોના ગણનો અગત્યનો ગુણધર્મ છે. જો સૂત્રોના કોઈ ગણમાંથી
વિસંવાદ, જેમ કે~$\lfalse$, નિષ્પન્ન કરી શકાય એવું હોય, તો તે ગણ અસુસંગત છે;
નહિ તો સુસંગત છે. ગણ અસુસંગત હોય, તો તેનાં બધાં સૂત્રો કોઈ
નિદર્શનના એક જગતમાં એકસાથે સત્ય હોઈ શકતાં નથી. પૂર્ણતાના પ્રમેય માટે
અમે વિપરીત વિધાન સાબિત કરીએ છીએ: દરેક સુસંગત ગણ કોઈ નિદર્શનના એક
જગતમાં સત્ય છે, ખાસ કરીને ``કૅનોનિકલ નિદર્શન''માં.

\begin{defn}
  ગણ $\Gamma$ એ તંત્ર~$\Sigma$ને અનુલક્ષીને \emph{સુસંગત} છે, અથવા
  આપણે કહીએ તેમ, $\Sigma$-સુસંગત છે, જો અને માત્ર જો $\Gamma
  \Proves/[\Sigma] \lfalse$.
\end{defn}

દાખલા તરીકે, ગણ $\{ \Box(p \lif q), \Box p, \lnot\Box q \}$
વિધાનાત્મક તર્કને અનુલક્ષીને સુસંગત છે, પરંતુ \Log{K}-સુસંગત નથી.
એ જ રીતે, ગણ $\{ \Diamond p, \Box\Diamond
p \lif q, \lnot q \}$ એ \Log{K5}-સુસંગત નથી.

\begin{prop}\ollabel{prop:consistencyfacts}
  માનો કે $\Gamma$ એ સૂત્રોનો ગણ છે. ત્યારે:
  \begin{enumerate}
  \item $\Gamma$ એ $\Sigma$-સુસંગત છે જો અને માત્ર જો એવું કોઈ
    સૂત્ર~$!A$ હોય કે $\Gamma \Proves/[\Sigma]
    !A$.
  \item \ollabel{prop:consistencyfacts-b}%
    $\Gamma \Proves[\Sigma] !A$ જો અને માત્ર જો $\Gamma \cup \{
    \lnot!A \}$ એ $\Sigma$-સુસંગત ન હોય.
  \item \ollabel{prop:consistencyfacts-c}%
    જો $\Gamma$ એ $\Sigma$-સુસંગત હોય, તો કોઈપણ સૂત્ર
    $!A$ માટે કાં $\Gamma \cup \{ !A \}$ એ $\Sigma$-સુસંગત છે,
    કાં $\Gamma \cup \{ \lnot!A \}$ એ $\Sigma$-સુસંગત છે.
  \end{enumerate}
\end{prop}

\begin{proof}
  આ હકીકતો પરંપરાગત વિધાનાત્મક તર્કથી સહેલાઈથી ફલિત થાય છે. આપણે
  \olref{prop:consistencyfacts-c}ની દલીલ આપીએ. પ્રતિપક્ષી વિધાનથી
  આગળ વધો અને માનો કે $\Gamma \cup \{ !A \}$ તથા
  $\Gamma \cup \{ \lnot!A \}$માંથી એકેય $\Sigma$-સુસંગત નથી.
  ત્યારે \olref{prop:consistencyfacts-b} મુજબ $\Gamma, !A \Proves[\Sigma]
  \lfalse$ અને $\Gamma, \lnot !A \Proves[\Sigma] \lfalse$ બંને.
  નિષ્પત્તિ પ્રમેયથી $\Gamma \Proves[\Sigma] !A \to \lfalse$ અને
  $\Gamma \Proves[\Sigma] \lnot!A \lif \lfalse$. પરંતુ $(!A \lif
  \lfalse) \lif ((\lnot!A \lif \lfalse) \lif \lfalse)$ પુનરુક્તિનો
  દૃષ્ટાંત છે; તેથી
  \olref[prp]{prop:derivabilityfacts}\olref[prp]{prop:derivabilityfacts-ruleT} મુજબ
  $\Gamma \Proves[\Sigma] \lfalse$.
\end{proof}
% END OLP-0440
\endgroup


\OLEndChapterHook
% END OLP-0427
\endgroup


\begingroup
\makeatletter\def\ol@sectioncs{section}\makeatother

% BEGIN OLP-0441
\olchapter{nml}{com}{પૂર્ણતા અને કૅનોનિકલ નિદર્શનો}
\olchapterid{nml}{com}
\phantomsection\label{guunit:OLP-0441}


\begingroup
\makeatletter\def\ol@sectioncs{section}\makeatother

% BEGIN OLP-0442
\olfileid{nml}{com}{int}

\olsection{પરિચય}
\phantomsection\label{guunit:OLP-0442}


જો $\Sigma$ મોડલ તંત્ર હોય, તો યથાર્થતાનો પ્રમેય સ્થાપિત કરે છે કે
$\Sigma \Proves !A$ હોય તો $\Sigma$નાં બધાં સૂત્રોના બધા
દૃષ્ટાંતો માન્ય હોય તેવા નિદર્શનોના કોઈપણ વર્ગ~$\mClass{C}$માં
$!A$ માન્ય છે. ખાસ કરીને, $\Log{K} \Proves !A$ હોય તો $!A$ બધાં
નિદર્શનોમાં સત્ય છે; $\Log{KT} \Proves !A$ હોય તો $!A$ બધાં
સ્વપ્રતિબિંબિત નિદર્શનોમાં સત્ય છે; $\Log{KD} \Proves !A$ હોય તો
$!A$ બધાં સિરિયલ નિદર્શનોમાં સત્ય છે, વગેરે.

પૂર્ણતા એ યથાર્થતાનું વિપરીત વિધાન છે: \Log{K} પૂર્ણ છે એટલે, દાખલા
તરીકે, સૂત્ર~$!A$ માન્ય હોય તો $\Proves !A$. પૂર્ણતા
સાબિત કરવી યથાર્થતા સાબિત કરવા કરતાં ઘણી વધુ અઘરી છે. પહેલાં તેનું
પ્રતિપક્ષી વિધાન ધ્યાનમાં લેવું ઉપયોગી છે: \Log{K} પૂર્ણ છે જો અને
માત્ર જો $\Proves/ !A$ હોય ત્યારે વિરોધી નિદર્શન, એટલે કે એવું
નિદર્શન~$\mModel{M}$ હોય કે $\mSat/{M}{!A}$. સમકક્ષ રીતે ($!A$નું
નિષેધન કરીને), $\Proves/ \lnot !A$ હોય ત્યારે $!A$નું નિદર્શન
હોય છે એવું આપણે સાબિત કરી શકીએ. આવું નિદર્શન રચતી વખતે આપણે
$!A$માં રહેલી માહિતી વાપરી શકીએ છીએ. ચોક્કસ સૂત્રો માટે
નિદર્શન શોધતી વખતે પણ આપણે ઘણી વાર એવું જ કરીએ છીએ: દાખલા તરીકે,
$p \lif \Box q$નું વિરોધી નિદર્શન શોધવું હોય તો તેમાં એવું જગત
હોવું જોઈએ જ્યાં $p$ સત્ય હોય અને $\Box q$ અસત્ય હોય. વળી જ્યાં
$\Box q$ અસત્ય હોય ત્યાંથી પ્રાપ્ય એવું જગત હોવું જોઈએ જ્યાં $q$
અસત્ય હોય. આપણને એટલું જ જાણવું જરૂરી છે: કયા જગતોમાં કયા
પ્રતિજ્ઞાપન ચલો સત્ય છે અને કયા જગતોમાંથી કયા જગતો
પ્રાપ્ય છે.

જોકે પૂર્ણતા સાબિત કરતી વખતે નિદર્શન રચવા માટે આપણી પાસે કોઈ
ચોક્કસ સૂત્ર~$!A$ નથી. આપણે સ્થાપિત કરવું છે કે
$\Proves/[\Sigma] \lnot !A$ હોય તેવા દરેક $!A$નું નિદર્શન હોય.
આ લઘુતમ આવશ્યકતા છે, કારણ કે \emph{જો} $\Proves[\Sigma] \lnot !A$
હોય, તો યથાર્થતા મુજબ $!A$નું ($\Sigma$ સત્ય હોય તેવું) કોઈ
નિદર્શન નથી. હવે નોંધો કે $\Proves/[\Sigma] \lnot !A$ જો અને માત્ર
જો $!A$ એ $\Sigma$-સુસંગત હોય. (યાદ કરો કે $\Sigma
\Proves/[\Sigma] \lnot !A$ અને $!A \Proves/[\Sigma] \lfalse$
સમકક્ષ છે.) તેથી આપણું કાર્ય દરેક $\Sigma$-સુસંગત સૂત્ર
માટે નિદર્શન રચવાનું છે.

આપણે એ યુક્તિ વાપરીશું કે $!A$ ધરાવતો એવો $\Sigma$-સુસંગત
સૂત્રોનો ગણ શોધવો જેમાં $!A$ને સત્ય બનાવતું જગત કેવું હોવું
જોઈએ તે જણાવતાં અન્ય સૂત્રો પણ હોય. આવા ગણો \emph{પૂર્ણ}
$\Sigma$-સુસંગત ગણો છે. $!A$ને સત્ય બનાવવા એક જ જગતવાળું નિદર્શન
રચવું પૂરતું નથી; તેમાં અનેક જગતો અને પ્રાપ્યતાનો સંબંધ હોવો પડશે.
$!A$ ધરાવતા પૂર્ણ $\Sigma$-સુસંગત ગણમાં $\Box !B$ અને
$\Diamond !C$ સ્વરૂપનાં અન્ય સૂત્રો પણ હશે. બધાં પ્રાપ્ય
જગતોમાં $!B$ સત્ય હોવું જોઈએ અને ઓછામાં ઓછા એકમાં $!C$ સત્ય
હોવું જોઈએ. આ માટે આપણે જગતોના ગણના આધાર તરીકે \emph{બધા}
શક્ય પૂર્ણ $\Sigma$-સુસંગત ગણો લઈશું. કઠિન ભાગ એ નક્કી કરવાનો
છે કે આપણા નિદર્શનમાં એક પૂર્ણ $\Sigma$-સુસંગત ગણ બીજામાંથી
ક્યારે પ્રાપ્ય ગણવો.

આ રીતે વ્યાખ્યાયિત નિદર્શનમાં $!A$ એક જગત પર સત્ય છે જો અને માત્ર
જો $!A$ તે ગણનો ઘટક છે—એ જગત પોતે પણ એક પૂર્ણ
$\Sigma$-સુસંગત ગણ છે—એવું આપણે બતાવીશું. $!A$ એ
$\Sigma$-સુસંગત હોય તો તે ઓછામાં ઓછા એક પૂર્ણ $\Sigma$-સુસંગત
ગણનો ઘટક હશે (આ હકીકત આપણે સાબિત કરીશું); તેથી એવું
જગત હશે જ્યાં $!A$ સત્ય છે. આમ આપણને એવું એક જ નિદર્શન મળશે જેમાં
દરેક $\Sigma$-સુસંગત સૂત્ર~$!A$ કોઈક જગત પર સત્ય છે.
આ એક જ નિદર્શન $\Sigma$ માટેનું \emph{કૅનોનિકલ} નિદર્શન છે.
% END OLP-0442
\endgroup

\begingroup
\makeatletter\def\ol@sectioncs{section}\makeatother

% BEGIN OLP-0443
\olfileid{nml}{com}{ccs}

\olsection{પૂર્ણ $\Sigma$-સુસંગત ગણો}
\phantomsection\label{guunit:OLP-0443}


ધારો કે $\Sigma$ મોડલ સૂત્રોનો ગણ છે—તેને સામાન્ય મોડલ
તર્કની સ્વયંસિદ્ધિઓ અથવા તેને વ્યાખ્યાયિત કરતા સિદ્ધાંતો માનો. ગણ
$\Gamma$ એ $\Sigma$-સુસંગત છે જો અને માત્ર જો $\Gamma
\Proves/[\Sigma] \lfalse$, એટલે કે $\Sigma$માંથી
$!A_1 \lif (!A_2 \lif \cdots (!A_n \lif \lfalse)\dots)$ની કોઈ
નિષ્પત્તિ નથી, જ્યાં દરેક $!A_i \in \Gamma$. આપણે એવું
``કૅનોનિકલ'' નિદર્શન રચીશું જેમાં દરેક જગત વિશિષ્ટ પ્રકારનો
$\Sigma$-સુસંગત ગણ હોય: જે માત્ર $\Sigma$-સુસંગત જ નહિ, પણ એ અર્થમાં
મહત્તમ હોય કે દરેક મોડલ સૂત્રનું સત્યમૂલ્ય નક્કી કરે—દરેક $!A$
માટે કાં તો $!A \in \Gamma$ અથવા $\lnot !A \in \Gamma$:

\begin{defn}
  ગણ $\Gamma$ \emph{પૂર્ણ $\Sigma$-સુસંગત} છે જો અને માત્ર જો તે
  $\Sigma$-સુસંગત છે અને દરેક~$!A$ માટે કાં $!A \in \Gamma$ અથવા
  $\lnot !A \in \Gamma$.
\end{defn}

પૂર્ણ $\Sigma$-સુસંગત ગણો $\Gamma$ના અનેક ઉપયોગી ગુણધર્મો છે.
ખાસ કરીને તેઓ નિષ્પત્તિ હેઠળ બંધ છે, એટલે કે $\Gamma
\Proves[\Sigma] !A$ હોય તો $!A \in \Gamma$. તેથી
સૂત્ર~$!A \in \Sigma$નો દરેક દૃષ્ટાંત પણ
$\in \Gamma$ છે. વધુમાં, $\Gamma$માં સભ્યતા વિધાનાત્મક સંયોજકોની
સત્યતાની શરતોને પ્રતિબિંબિત કરે છે. ``કૅનોનિકલ નિદર્શન'' વ્યાખ્યાયિત
કરીશું ત્યારે આ અગત્યનું બનશે.

\begin{prop}\ollabel{prop:ccs-properties}
  ધારો કે $\Gamma$ પૂર્ણ $\Sigma$-સુસંગત છે. ત્યારે:
  \begin{enumerate}
  \item \ollabel{prop:ccs-closed}%
    $\Gamma$ એ~$\Sigma$માં નિષ્પત્તિ હેઠળ બંધ છે.
  \item \ollabel{prop:ccs-sigma}%
    $\Sigma \subseteq \Gamma$.
  \item \ollabel{prop:ccs-lfalse}%
    $\lfalse \notin \Gamma$
  \item \ollabel{prop:ccs-ltrue}%
    $\ltrue \in \Gamma$
  \item \ollabel{prop:ccs-lnot}%
    $\lnot!A \in \Gamma$ જો અને માત્ર જો $!A \notin
    \Gamma$.
  \item \ollabel{prop:ccs-land}%
    $!A \land !B \in \Gamma$ જો અને માત્ર જો $!A \in \Gamma$ અને
    $!B \in \Gamma$
  \item \ollabel{prop:ccs-lor}%
    $!A \lor !B \in \Gamma$ જો અને માત્ર જો $!A \in \Gamma$ અથવા
    $!B \in \Gamma$
  \item \ollabel{prop:ccs-lif}%
    $!A \lif !B \in \Gamma$ જો અને માત્ર જો $!A \notin \Gamma$
    અથવા $!B \in \Gamma$
  \item \ollabel{prop:ccs-liff}%
    $!A \liff !B \in \Gamma$ જો અને માત્ર જો કાં તો $!A \in \Gamma$
    અને $!B \in \Gamma$ બંને, અથવા $!A \notin \Gamma$ અને
    $!B \notin \Gamma$ બંને
  \end{enumerate}
\end{prop}

\begin{proof}
  \begin{enumerate}
  \item ધારો કે $\Gamma \Proves[\Sigma] !A$, પણ $!A \notin
    \Gamma$. ત્યારે $\Gamma$ પૂર્ણ $\Sigma$-સુસંગત હોવાથી
    $\lnot!A \in \Gamma$. આથી $\Gamma$ અસુસંગત થઈ જાય, કારણ કે
    $!A, \lnot !A \Proves[\Sigma] \lfalse$.

  \item જો $!A \in \Sigma$ હોય તો $\Gamma \Proves[\Sigma] !A$,
    અને નિષ્પત્તિ હેઠળની બંધતાથી $!A \in \Gamma$, એટલે કે
    મુદ્દા~\olref{prop:ccs-closed} મુજબ.

  \item જો $\lfalse \in \Gamma$ હોય તો $\Gamma
    \Proves[\Sigma] \lfalse$; તેથી $\Gamma$ એ
    $\Sigma$-અસુસંગત બને.

  \item $\Gamma \Proves[\Sigma] \ltrue$, તેથી
    નિષ્પત્તિ હેઠળની બંધતા, એટલે કે મુદ્દા~\olref{prop:ccs-closed}
    મુજબ $\ltrue \in \Gamma$.

  \item જો $\lnot!A \in \Gamma$ હોય તો સુસંગતતાથી $!A \notin
    \Gamma$; અને જો $!A \notin \Gamma$ હોય તો $\Gamma$ પૂર્ણ
    $\Sigma$-સુસંગત હોવાથી $\lnot !A \in \Gamma$.
    \sourcecorrection{OLMD-032}{મૂળના આ બીજા ખંડમાં નિષેધચિહ્ન
    છૂટી ગયું છે; પૂર્ણતાની વ્યાખ્યા મુજબ અહીં નિષેધિત સૂત્ર ગણમાં
    આવવું જોઈએ.}

  \item ધારો કે $!A \land !B \in
      \Gamma$. $(!A \land !B) \lif !A$ પુનરુક્તિનો દૃષ્ટાંત હોવાથી
      નિષ્પત્તિ હેઠળની બંધતા, એટલે કે મુદ્દા~\olref{prop:ccs-closed}
      મુજબ $!A \in \Gamma$. $!B \in \Gamma$ માટે પણ તેવી જ દલીલ છે.
      ઊલટી દિશામાં, ધારો કે $!A \in \Gamma$ અને $!B \in
      \Gamma$ બંને. ત્યારે નિષ્પત્તિ હેઠળની બંધતાથી $(!A \land !B) \in
      \Gamma$, કારણ કે $!A \lif (!B \lif (!A \land !B))$
      પુનરુક્તિનો દૃષ્ટાંત છે.

  \item ધારો કે $!A \lor !B \in
      \Gamma$, પણ $!A \notin \Gamma$ અને $!B \notin \Gamma$.
      $\Gamma$ પૂર્ણ $\Sigma$-સુસંગત હોવાથી $\lnot!A \in \Gamma$
      અને $\lnot !B \in \Gamma$. ત્યારે $\lnot(!A \lor !B) \in
      \Gamma$, કારણ કે $\lnot !A \lif (\lnot !B \lif \lnot (!A \lor !B))$
      પુનરુક્તિનો દૃષ્ટાંત છે. આથી $\Gamma$ એ $\Sigma$-અસુસંગત
      બને, જે વિરોધાભાસ છે. ઊલટી દિશામાં, કોઈ એક વિયોજક ગણમાં હોય તો
      તેને વિયોજનમાં લઈ જતી વિધાનાત્મક પુનરુક્તિ અને નિષ્પત્તિ હેઠળની
      બંધતાથી વિયોજન પણ ગણમાં આવે છે.
      \sourcecorrection{OLMD-033}{મૂળની આ સાબિતીમાં વિયોજનના દ્વિમુખી
      વિધાનની માત્ર એક દિશા આપવામાં આવી છે; અહીં બીજી દિશાની દલીલ
      સ્પષ્ટ કરી છે.}

  \item ધારો કે $!A \lif !B \in
      \Gamma$ અને $!A \in \Gamma$; ત્યારે $\Gamma \Proves[\Sigma] !B$,
      અને નિષ્પત્તિ હેઠળની બંધતાથી $!B \in \Gamma$. ઊલટું, જો $!A
      \lif !B \notin \Gamma$ હોય તો $\Gamma$ પૂર્ણ
      $\Sigma$-સુસંગત હોવાથી $\lnot (!A \lif !B) \in \Gamma$.
      $\lnot(!A \lif !B) \lif !A$ પુનરુક્તિનો દૃષ્ટાંત હોવાથી
      નિષ્પત્તિ હેઠળની બંધતાથી $!A \in \Gamma$. એ જ રીતે
      $\lnot(!A \lif !B) \lif \lnot !B$ પુનરુક્તિનો દૃષ્ટાંત હોવાથી
      $\lnot !B \in \Gamma$. તેથી $\Gamma$ એ $\Sigma$-સુસંગત
      હોવાથી $!B \notin \Gamma$.

  \item ધારો કે $!A \liff !B \in
      \Gamma$. જો $!A \in \Gamma$ હોય તો $!B \in \Gamma$, કારણ કે
      $(!A \liff !B) \lif (!A \lif !B)$ પુનરુક્તિનો દૃષ્ટાંત છે.
      એ જ રીતે, $!B \in \Gamma$ હોય તો $!A \in \Gamma$. તેથી કાં
      તો $!A \in \Gamma$ અને $!B \in \Gamma$ બંને, અથવા
      $!A \in \Gamma$ પણ નહિ અને $!B \in \Gamma$ પણ નહિ.

      ઊલટું, ધારો કે $!A \liff !B \notin \Gamma$. $\Gamma$ પૂર્ણ
      $\Sigma$-સુસંગત હોવાથી $\lnot (!A \liff !B) \in \Gamma$.
      \sourcecorrection{OLMD-034}{મૂળની આ ઊલટી દિશાના આરંભમાં
      દ્વિમુખી પ્રેરણની જગ્યાએ એકમુખી પ્રેરણ લખાઈ છે; આગળનું નિષેધિત
      સૂત્ર દ્વિમુખી પ્રેરણનું જ છે.}
      $\lnot(!A \liff !B) \lif (!A \lif \lnot !B)$ પુનરુક્તિનો
      દૃષ્ટાંત હોવાથી $!A \in \Gamma$ હોય તો $\lnot !B \in
      \Gamma$; અને $\Gamma$ એ $\Sigma$-સુસંગત હોવાથી $!B \notin
      \Gamma$. એ જ રીતે, $!B \in \Gamma$ હોય તો $!A \notin \Gamma$.
      તેથી $!A \in \Gamma$ અને $!B \in \Gamma$ બંને નથી, તેમજ
      $!A \notin \Gamma$ અને $!B \notin \Gamma$ બંને પણ નથી.
  \end{enumerate}
\end{proof}

\begin{probtag}{probAnd,probOr,probIf,probIff}
  \olref[nml][com][ccs]{prop:ccs-properties}ની સાબિતી પૂરી કરો.
\end{probtag}
% END OLP-0443
\endgroup

\begingroup
\makeatletter\def\ol@sectioncs{section}\makeatother

% BEGIN OLP-0444
\olfileid{nml}{com}{lin}

\olsection{લિન્ડનબાઉમનું સહાયક પ્રમેય}
\phantomsection\label{guunit:OLP-0444}


લિન્ડનબાઉમનું સહાયક પ્રમેય સ્થાપિત કરે છે કે સૂત્રોનો દરેક
$\Sigma$-સુસંગત ગણ ઓછામાં ઓછા એક \emph{પૂર્ણ} $\Sigma$-સુસંગત ગણમાં
સમાયેલો છે. કૅનોનિકલ નિદર્શનની આપણી રચના બતાવશે કે દરેક પૂર્ણ
$\Sigma$-સુસંગત ગણ~$\Delta$ માટે કૅનોનિકલ નિદર્શનમાં એવું જગત છે
જ્યાં $\Delta$માંનાં બધાં અને માત્ર એ જ સૂત્રો સત્ય છે. તેથી
લિન્ડનબાઉમનું સહાયક પ્રમેય ખાતરી આપે છે કે દરેક $\Sigma$-સુસંગત ગણ
કૅનોનિકલ નિદર્શનના કોઈ જગત પર સત્ય છે.

\begin{thm}[લિન્ડનબાઉમનું સહાયક પ્રમેય]\ollabel{thm:lindenbaum}
  જો $\Gamma$ એ $\Sigma$-સુસંગત હોય, તો $\Gamma$ને વિસ્તારે એવો
  પૂર્ણ $\Sigma$-સુસંગત ગણ $\Delta$ છે.
\end{thm}

\begin{proof}
ભાષાનાં બધાં સૂત્રોની સંપૂર્ણ યાદી $!A_0$, $!A_1$, \dots{} લો
(પુનરાવર્તન માન્ય છે). દાખલા તરીકે, પહેલાં $\Obj p_0$ને યાદીમાં
મૂકો અને દરેક તબક્કે $n \ge 1$ માટે માત્ર $\Obj p_0$,
\dots,~$\Obj p_n$માંથી ચલ વાપરીને બનતાં લંબાઈ~$n$નાં સાન્ત
સંખ્યામાં સૂત્રોને યાદીમાં મૂકો. આપણે~$n$ પર અનુમાનથી સૂત્રોના
ગણો $\Delta_n$ વ્યાખ્યાયિત કરીએ છીએ અને પછી $\Delta = \bigcup_n
\Delta_n$ મૂકીએ છીએ. પહેલાં $\Delta_0 = \Gamma$ મૂકો. $\Delta_n$
વ્યાખ્યાયિત છે એમ ધારીને $\Delta_{n+1}$ આ રીતે વ્યાખ્યાયિત કરો:
\[
\Delta_{n+1} =
\begin{cases}
  \Delta_n \cup \{!A_n\}, & \text{જો $\Delta_n \cup \{ !A_n\}$
    એ $\Sigma$-સુસંગત હોય;} \\
  \Delta_n \cup \{ \lnot !A_n\}, & \text{અન્યથા.}
\end{cases}
\]
હવે $\Delta = \bigcup_{n=0}^\infty \Delta_n$ લો.

આ વ્યાખ્યાથી ખરેખર જરૂરી ગુણધર્મો ધરાવતો ગણ $\Delta$ મળે છે એવું
બતાવવું પડશે, એટલે કે $\Gamma \subseteq \Delta$ અને $\Delta$ એ
પૂર્ણ $\Sigma$-સુસંગત છે.

$\Gamma \subseteq \Delta$ સ્પષ્ટ છે, કારણ કે રચના મુજબ $\Delta_0
\subseteq \Delta$ અને $\Delta_0 = \Gamma$. હકીકતમાં, દરેક~$n$ માટે
$\Delta_n \subseteq \Delta$, કેમ કે $\Delta$ બધાં~$\Delta_n$નો સંઘ
છે. (રચનાના દરેક પગલે આપણે અગાઉના ગણમાં સૂત્ર ઉમેરીએ છીએ,
તેથી $\Delta_n \subseteq \Delta_{n+1}$; અને $\subseteq$ પરંપરિત
હોવાથી $n \le m$ હોય ત્યારે $\Delta_n \subseteq \Delta_{m}$.)
રચનાના દરેક તબક્કે આપણે કાં તો $!A_n$ અથવા $\lnot !A_n$ ઉમેરીએ
છીએ, અને દરેક સૂત્ર બધાં~$!A_n$ની યાદીમાં (ઓછામાં ઓછું એક વાર)
આવે છે. તેથી દરેક $!A$ માટે કાં તો $!A \in \Delta$ અથવા
$\lnot !A \in \Delta$; એટલે વ્યાખ્યા મુજબ $\Delta$ પૂર્ણ છે.

છેલ્લે, $\Delta$ એ $\Sigma$-સુસંગત છે તે બતાવવું પડશે. આ માટે
બતાવીએ છીએ કે (a) જો $\Delta$ એ $\Sigma$-અસુસંગત હોય, તો કોઈ
$\Delta_n$ એ $\Sigma$-અસુસંગત હશે; અને (b) બધાં $\Delta_n$ એ
$\Sigma$-સુસંગત છે.

ધારો કે $\Delta$ એ $\Sigma$-અસુસંગત છે. ત્યારે $\Delta
\Proves[\Sigma] \lfalse$, એટલે કે એવાં $!A_1$, \dots,~$!A_k \in
\Delta$ છે કે $\Sigma \Proves !A_1 \lif (!A_2 \lif \cdots (!A_k
\lif \lfalse)\dots)$. $\Delta = \bigcup_{n=0}^\infty \Delta_n$
હોવાથી દરેક $!A_i \in \Delta_{n_i}$ કોઈક~$n_i$ માટે સાચું છે.
આ સૂચકોમાં સૌથી મોટો $n$ લો. $n_i \le n$ હોવાથી $\Delta_{n_i}
\subseteq \Delta_n$. આમ, બધાં $!A_i$ કોઈ એક~$\Delta_n$માં છે.
આનો અર્થ $\Delta_n \Proves[\Sigma] \lfalse$, એટલે કે $\Delta_n$
એ $\Sigma$-અસુસંગત છે.

દરેક $\Delta_n$ એ $\Sigma$-સુસંગત છે તે બતાવવા~$n$ પર સરળ અનુમાન
વાપરીએ. $\Delta_0 = \Gamma$, અને આપણે $\Gamma$ એ
$\Sigma$-સુસંગત છે એમ ધાર્યું હતું. એટલે $n = 0$ માટે દાવો સાચો
છે. હવે $n$ માટે દાવો સાચો છે, એટલે કે $\Delta_n$ એ
$\Sigma$-સુસંગત છે, એમ ધારો. જો $\Delta_n \cup \{!A_n\}$ એ
$\Sigma$-સુસંગત હોય, તો $\Delta_{n+1}$ એ તે ગણ છે; નહિ તો તે
$\Delta_n \cup \{\lnot!A_n\}$ છે. પહેલા કિસ્સામાં
$\Delta_{n+1}$ એ સ્પષ્ટ રીતે $\Sigma$-સુસંગત છે. પરંતુ
\olref[prf][con]{prop:consistencyfacts}\olref[prf][con]{prop:consistencyfacts-c}
મુજબ $\Delta_n \cup \{!A_n\}$ અથવા $\Delta_n \cup \{\lnot!A_n\}$
બેમાંથી ઓછામાં ઓછો એક સુસંગત છે; તેથી બીજા કિસ્સામાં પણ
$\Delta_{n+1}$ સુસંગત છે.
\end{proof}

\begin{cor}\ollabel{cor:provability-characterization}
  $\Gamma \Proves[\Sigma] !A$ જો અને માત્ર જો $\Gamma$ને વિસ્તારે
  એવા દરેક પૂર્ણ $\Sigma$-સુસંગત ગણ $\Delta$માં $!A \in \Delta$
  હોય ($\Gamma = \emptyset$ હોય ત્યારે પણ; એ કિસ્સામાં મોડલ તંત્ર
  $\Sigma$નું બીજું લક્ષણચિત્રણ મળે છે).
\end{cor}

\begin{proof}
  ધારો કે $\Gamma \Proves[\Sigma] !A$, અને $\Gamma$ને વિસ્તારે
  એવો કોઈપણ પૂર્ણ $\Sigma$-સુસંગત ગણ $\Delta$ લો. જો $!A
  \notin \Delta$ હોય, તો મહત્તમતાથી $\lnot!A \in \Delta$;
  તેથી એકદિશવર્ધિતાથી $\Delta \Proves[\Sigma] !A$ અને સ્વવાચકતાથી
  $\Delta \Proves[\Sigma] \lnot!A$, એટલે $\Delta$ અસુસંગત છે.
  ઊલટું, જો $\Gamma \Proves/[\Sigma] !A$ હોય, તો $\Gamma \cup \{
  \lnot!A\}$ એ $\Sigma$-સુસંગત છે; લિન્ડનબાઉમના સહાયક પ્રમેયથી
  $\Gamma \cup \{ \lnot!A \}$ને વિસ્તારે એવો પૂર્ણ સુસંગત ગણ
  $\Delta$ છે. સુસંગતતાથી $!A \notin \Delta$.
\end{proof}
% END OLP-0444
\endgroup

\begingroup
\makeatletter\def\ol@sectioncs{section}\makeatother

% BEGIN OLP-0445
\olfileid{nml}{com}{mod}

\olsection{મોડલ સંચાલકો અને પૂર્ણ સુસંગત ગણો}
\phantomsection\label{guunit:OLP-0445}


\begin{explain}
જ્યારે આપણે કોઈ સામાન્ય મોડલ તર્ક~$\Sigma$ના પૂર્ણ
$\Sigma$-સુસંગત ગણો~$\Delta$ને જગતો તરીકે લઈને
નિદર્શન~$\mModel{M^\Sigma}$ રચીશું, ત્યારે આ ``જગતો'' વચ્ચે
પ્રાપ્યતા સંબંધ~$R^\Sigma$ પણ વ્યાખ્યાયિત કરવો પડશે. આપણે ઇચ્છીએ
છીએ કે પ્રાપ્યતા સંબંધ (અને નિયુક્તિ~$V^\Sigma$) એવી રીતે
વ્યાખ્યાયિત થાય કે $\mSat{M^\Sigma}{!A}[\Delta]$ જો અને માત્ર જો
$!A \in \Delta$. તે કેવી રીતે કરવું?
\sourcecorrection{OLMD-035}{મૂળમાં નિયુક્તિના ચિહ્નની પાછળનું બંધ
કૌંસ ગણિત-ચિહ્નની અંદર છે; અહીં કૌંસને ગણિત-ચિહ્નની બહાર રાખ્યું છે.}

પ્રાપ્યતા સંબંધ વ્યાખ્યાયિત કર્યા પછી જગત પર સત્યની વ્યાખ્યાથી
$\mSat{M^\Sigma}{\Box!A}[\Delta]$ જો અને માત્ર જો
  $R^\Sigma\Delta\Delta'$ હોય તેવા બધાં $\Delta'$ માટે
  $\mSat{M^\Sigma}{!A}[\Delta']$.
$\mSat{M^\Sigma}{!A}[\Delta]$ જો અને માત્ર જો $!A \in \Delta$
એ સાબિત કરવા ખાસ કરીને મોડલ સંચાલકથી શરૂ થતા સૂત્રો માટે પણ
આ સાચું હોવું જરૂરી છે, એટલે કે
$\mSat{M^\Sigma}{\Box!A}[\Delta]$ જો અને માત્ર જો
  $\Box !A \in \Delta$. આ આવશ્યકતાને
$\Box !A$ માટે જગત પર સત્યની
વ્યાખ્યા સાથે જોડવાથી મળે છે:
\[
%
  \Box!A \in \Delta \text{ જો અને માત્ર જો } !A \in \Delta' \text{
    એવા બધા $\Delta'$ માટે કે } R^\Sigma\Delta\Delta'
\]
ડાબેથી જમણી દિશા જુઓ: તે કહે છે કે જો
  $\Box !A \in \Delta$ હોય, તો કોઈપણ $!A$ અને
  $R^\Sigma\Delta\Delta'$ હોય તેવા કોઈપણ $\Delta'$ માટે
  $!A \in \Delta'$. જો આપણે $R^\Sigma\Delta\Delta'$ જો અને માત્ર
  જો $\Box!A \in \Delta$ હોય તેવા બધા સૂત્રો માટે $!A \in \Delta'$
  એવું નક્કી કરીએ, તો આ વિધાન સાચું રહેશે. ``જો અને માત્ર જો''ની
  જમણી બાજુની શરત વધુ સંક્ષેપમાં આમ લખી શકાય:
  $\Setabs{!A}{\Box!A \in \Delta} \subseteq \Delta'$.

તો પ્રશ્ન એ છે: $R^\Sigma$ની આ વ્યાખ્યા ખરેખર ખાતરી આપે છે કે
$\Box !A \in \Delta$ જો અને માત્ર જો
  $\mSat{M^\Sigma}{\Box !A}[\Delta]$?
  શું તે એ પણ ખાતરી આપે છે કે %
$\Diamond !A \in \Delta$ જો અને માત્ર જો
  $\mSat{M^\Sigma}{\Diamond !A}[\Delta]$?
આગળનાં થોડાં પરિણામો આ સ્થાપિત કરશે.
\end{explain}

\begin{defn}
  જો $\Gamma$ એ સૂત્રોનો ગણ હોય, તો મૂકો
  \begin{align*}
    \Box\Gamma & = \Setabs{\Box !B}{!B \in \Gamma}\\
    \Diamond\Gamma & = \Setabs{\Diamond !B}{!B \in \Gamma}\\
    \intertext{અને}
    \Box^{-1}\Gamma & = \Setabs{!B}{\Box !B \in \Gamma}\\
    \Diamond^{-1}\Gamma & = \Setabs{!B}{\Diamond !B \in \Gamma}\\
  \end{align*}
\end{defn}

બીજા શબ્દોમાં, $\Box\Gamma$ એટલે~$\Gamma$નો એવો ગણ જેમાં
દરેક સૂત્ર~$\Gamma$માંથી લઈને તેની આગળ $\Box$ મૂકેલો હોય;
અને $\Box^{-1}\Gamma$ એટલે~$\Gamma$ના
$\Box$થી શરૂ થતા બધાં સૂત્રોમાંથી શરૂઆતના $\Box$ દૂર કરીને
મળેલો ગણ. આ વ્યાખ્યા પોતે બહુ મહત્વની નથી, પરંતુ સંકેતન ઘણું સરળ
બનાવશે.

નોંધો કે $\Box\Box^{-1}\Gamma \subseteq \Gamma$:
\[
\Box\Box^{-1}\Gamma = \Setabs{\Box!B}{\Box!B \in \Gamma}
\]
એટલે કે આ~$\Gamma$નાં $\Box$થી શરૂ થતા બધાં સૂત્રોનો જ ગણ છે.

\begin{lem}\ollabel{lem:box1}
  જો $\Gamma \Proves[\Sigma] !A$ હોય, તો $\Box\Gamma
  \Proves[\Sigma] \Box !A$.
\end{lem}

\begin{proof}
  જો $\Gamma \Proves[\Sigma] !A$ હોય, તો એવાં $!B_1$, \dots,
  $!B_k \in \Gamma$ છે કે $\Sigma \Proves !B_1 \lif (!B_2 \lif
  \cdots (!B_k \lif !A)\cdots)$. $\Sigma$ સામાન્ય હોવાથી
  નિયમ~\RK{} મુજબ $\Sigma \Proves \Box!B_1 \lif (\Box!B_2 \lif
  \cdots (\Box!B_k \lif \Box!A)\cdots)$, જ્યાં સ્પષ્ટ રીતે
  $\Box!B_1$, \dots, $\Box!B_k \in \Box\Gamma$. તેથી વ્યાખ્યા મુજબ
  $\Box\Gamma \Proves[\Sigma] \Box !A$.
  \sourcecorrection{OLMD-036}{મૂળની બંને શૃંખલાના છેલ્લા પદમાં
  સૂચક બદલાઈ જાય છે; ગણમાં પસંદ કરાયેલા સૂત્રોની સંખ્યા સાથે મેળ
  રાખવા અહીં એ જ અંતિમ સૂચક વાપર્યો છે.}
\end{proof}

\begin{lem}\ollabel{lem:box2}
  જો $\Box^{-1} \Gamma \Proves[\Sigma] !A$ હોય, તો $\Gamma
  \Proves[\Sigma] \Box!A$.
\end{lem}

\begin{proof}
  ધારો કે $\Box^{-1}\Gamma \Proves[\Sigma] !A$; ત્યારે
  \olref{lem:box1} મુજબ $\Box\Box^{-1}\Gamma \Proves \Box!A$.
  પરંતુ $\Box\Box^{-1}\Gamma \subseteq \Gamma$ હોવાથી
  એકદિશવર્ધિતાથી $\Gamma \Proves[\Sigma] \Box!A$ પણ મળે છે.
\end{proof}

%

\begin{prop}\ollabel{prop:box}
  જો $\Gamma$ પૂર્ણ $\Sigma$-સુસંગત હોય, તો $\Box !A \in
  \Gamma$ જો અને માત્ર જો $\Box^{-1} \Gamma \subseteq \Delta$
  હોય તેવા દરેક પૂર્ણ $\Sigma$-સુસંગત ગણ $\Delta$ માટે
  $!A \in \Delta$.
\end{prop}

\begin{proof}
  ધારો કે $\Gamma$ પૂર્ણ $\Sigma$-સુસંગત છે. ``માત્ર જો'' દિશા
  સરળ છે: ધારો કે $\Box !A \in \Gamma$ અને $\Box^{-1}\Gamma
  \subseteq \Delta$. $\Box!A \in \Gamma$ હોવાથી
  $!A \in \Box^{-1}\Gamma \subseteq \Delta$; તેથી
  $!A \in \Delta$.

  ``જો'' દિશા માટે પ્રતિપક્ષી વિધાન સાબિત કરીએ: ધારો કે
  $\Box!A \notin \Gamma$. $\Gamma$ પૂર્ણ $\Sigma$-સુસંગત હોવાથી
  તે નિષ્પત્તિ હેઠળ બંધ છે; તેથી $\Gamma \Proves/[\Sigma] \Box !A$.
  \olref{lem:box2} મુજબ $\Box^{-1}\Gamma \Proves/[\Sigma] !A$.
  \olref[prf][con]{prop:consistencyfacts}\olref[prf][con]{prop:consistencyfacts-b}
  મુજબ $\Box^{-1}\Gamma \cup \{ \lnot!A \}$ એ
  $\Sigma$-સુસંગત છે. લિન્ડનબાઉમના સહાયક પ્રમેયથી એવો પૂર્ણ
  $\Sigma$-સુસંગત ગણ~$\Delta$ છે કે $\Box^{-1}\Gamma \cup
  \{ \lnot!A \} \subseteq \Delta$. સુસંગતતાથી $!A \notin \Delta$.
\end{proof}

\begin{lem}\ollabel{lem:box-iff-diamond}
  ધારો કે $\Gamma$ અને $\Delta$ પૂર્ણ $\Sigma$-સુસંગત છે.
  ત્યારે $\Box^{-1}\Gamma \subseteq \Delta$ જો અને માત્ર જો
  $\Diamond\Delta \subseteq \Gamma$.
\end{lem}

\begin{proof}
  ``માત્ર જો'' દિશા: $\Box^{-1}\Gamma \subseteq \Delta$ ધારો અને
  $\Diamond!A \in \Diamond\Delta$ (એટલે કે $!A \in \Delta$)
  માનો. $\Diamond!A \in \Gamma$ બતાવવા $\Box\lnot!A \notin
  \Gamma$ બતાવવું પૂરતું છે, કારણ કે ત્યારે મહત્તમતાથી
  $\lnot\Box\lnot !A \in \Gamma$. હવે જો $\Box\lnot!A \in
  \Gamma$ હોય, તો પૂર્વધારણાથી $\lnot!A \in \Delta$, જે
  $\Delta$ની સુસંગતતા વિરુદ્ધ છે (કારણ કે $!A \in \Delta$).
  તેથી જરૂરી મુજબ $\Box\lnot!A \notin \Gamma$.

  ``જો'' દિશા: $\Diamond\Delta \subseteq \Gamma$ ધારો.
  પ્રતિપક્ષી દલીલ કરીએ: $!A \notin \Delta$ માનીને
  $\Box!A \notin \Gamma$ બતાવીએ. જો $!A \notin \Delta$ હોય,
  તો મહત્તમતાથી $\lnot!A \in \Delta$ અને તેથી પૂર્વધારણાથી
  $\Diamond\lnot!A \in \Gamma$. પરંતુ સામાન્ય મોડલ તર્કમાં
  $\Diamond\lnot!A$ એ $\lnot\Box !A$ને સમકક્ષ છે; અને છેલ્લું
  સૂત્ર $\Gamma$માં હોય તો સુસંગતતાથી $\Box!A \notin\Gamma$,
  જેમ જોઈતું હતું.
\end{proof}

%
  %

\begin{prop}\ollabel{prop:diamond}
  જો $\Gamma$ પૂર્ણ $\Sigma$-સુસંગત હોય, તો $\Diamond !A \in
  \Gamma$ જો અને માત્ર જો એવો કોઈ પૂર્ણ $\Sigma$-સુસંગત ગણ
  $\Delta$ હોય કે $\Diamond\Delta \subseteq \Gamma$ અને
  $!A \in \Delta$.
\end{prop}

\begin{proof}
ધારો કે $\Gamma$ પૂર્ણ $\Sigma$-સુસંગત છે. \Dual{} અને નિષ્પત્તિ
હેઠળની બંધતા મુજબ $\Diamond!A \in \Gamma$ જો અને માત્ર જો
$\lnot\Box\lnot !A \in \Gamma$. $\Gamma$ પૂર્ણ
$\Sigma$-સુસંગત હોવાથી
\olref[ccs]{prop:ccs-properties}\olref[ccs]{prop:ccs-lnot} મુજબ
$\lnot\Box\lnot !A \in \Gamma$ જો અને માત્ર જો
$\Box\lnot !A \notin \Gamma$. \olref{prop:box} મુજબ
$\Box\lnot !A \notin \Gamma$ જો અને માત્ર જો એવો કોઈ પૂર્ણ
$\Sigma$-સુસંગત ગણ $\Delta$ હોય કે $\Box^{-1}\Gamma
\subseteq \Delta$ અને $\lnot!A \notin \Delta$. હવે આવો
કોઈપણ~$\Delta$ લો. \olref{lem:box-iff-diamond} મુજબ
$\Box^{-1}\Gamma \subseteq \Delta$ જો અને માત્ર જો
$\Diamond\Delta \subseteq \Gamma$. વધુમાં,
\olref[ccs]{prop:ccs-properties}\olref[ccs]{prop:ccs-lnot} મુજબ
$\lnot !A \notin \Delta$ જો અને માત્ર જો $!A \in \Delta$.
તેથી $\Diamond!A \in \Gamma$ જો અને માત્ર જો એવો કોઈ પૂર્ણ
$\Sigma$-સુસંગત ગણ $\Delta$ હોય કે $\Diamond\Delta
\subseteq \Gamma$ અને $!A \in \Delta$.
\end{proof}


\begin{prob}
  બતાવો કે જો $\Gamma$ પૂર્ણ $\Sigma$-સુસંગત હોય, તો
  $\Diamond !A \in \Gamma$ જો અને માત્ર જો એવો પૂર્ણ
    $\Sigma$-સુસંગત ગણ $\Delta$ છે કે $\Box^{-1}\Gamma
    \subseteq \Delta$ અને $!A \in \Delta$. આ સાબિતી
  \olref[nml][com][mod]{lem:box-iff-diamond} વાપર્યા \emph{વિના} કરો.
\end{prob}
% END OLP-0445
\endgroup

\begingroup
\makeatletter\def\ol@sectioncs{section}\makeatother

% BEGIN OLP-0446
\olfileid{nml}{com}{cmd}

\olsection{કૅનોનિકલ નિદર્શનો}
\phantomsection\label{guunit:OLP-0446}


મોડલ તંત્ર~$\Sigma$ માટેનું \emph{કૅનોનિકલ નિદર્શન} એ એવું ચોક્કસ
નિદર્શન~$\mModel{M^\Sigma}$ છે જેમાં જગતો બધા પૂર્ણ
$\Sigma$-સુસંગત ગણો છે. તેનો પ્રાપ્યતા સંબંધ~$R^\Sigma$ અને
સત્યમૂલ્ય-નિયુક્તિ $V^\Sigma$ એવી રીતે વ્યાખ્યાયિત થાય છે કે
જગત~$\Delta$ પર સત્ય સૂત્રો ચોક્કસપણે~$\Delta$માંનાં
સૂત્રો જ હોય.

\begin{defn}
  ધારો કે $\Sigma$ સામાન્ય મોડલ તર્ક છે. $\Sigma$ માટેનું
  \emph{કૅનોનિકલ નિદર્શન} $\mModel{M}^\Sigma =
  \tuple{W^\Sigma, R^\Sigma, V^\Sigma}$ છે, જ્યાં:
  \begin{enumerate}
  \item $W^\Sigma = \Setabs{ \Delta }{\Delta \text{
      એ પૂર્ણ $\Sigma$-સુસંગત છે} }$.
  \item $R^\Sigma \Delta\Delta'$ જો અને માત્ર જો
    $\Box^{-1}\Delta \subseteq
      \Delta'$.
  \item $V^\Sigma(p) = \Setabs{\Delta}{p \in \Delta}$.
  \end{enumerate}
\end{defn}
% END OLP-0446
\endgroup

\begingroup
\makeatletter\def\ol@sectioncs{section}\makeatother

% BEGIN OLP-0447
\olfileid{nml}{com}{tru}

\olsection{સત્યતા સહાયક પ્રમેય}
\phantomsection\label{guunit:OLP-0447}


કૅનોનિકલ નિદર્શન~$\mModel{M^\Sigma}$ એવી રીતે વ્યાખ્યાયિત છે કે
$\mSat{M^\Sigma}{!A}[\Delta]$ જો અને માત્ર જો $!A \in \Delta$.
પ્રતિજ્ઞાપન ચલો માટે $V^\Sigma$ની વ્યાખ્યાથી આ સીધું
મળે છે. જોકે બધાં સૂત્રો માટે આ સમકક્ષતા સાચી છે તે ચકાસવું
પડશે. તે આપણે અનુમાનથી કરીએ છીએ. અનુમાનના પગલામાં વિધાનાત્મક
સંચાલકો ધરાવતા સૂત્રો માટે (જ્યાં
\olref[ccs]{prop:ccs-properties} વાપરવું પડે) અને મોડલ સંચાલકો
ધરાવતા સૂત્રો માટે (જ્યાં \olref[mod]{sec}નાં પરિણામો વાપરીએ)
સમકક્ષતા સાબિત કરવી પડે છે.

\begin{prop}[સત્યતા સહાયક પ્રમેય]
  \ollabel{prop:truthlemma}
  દરેક સૂત્ર~$!A$ માટે $\mSat{M^\Sigma}{!A}[\Delta]$ જો અને
  માત્ર જો $!A \in \Delta$.
\end{prop}

\begin{proof}
  $!A$ પર અનુમાનથી.
  \begin{enumerate}
    \item \indcase{!A}{\lfalse}{$\mSat/{M^\Sigma}{\lfalse}[\Delta]$
        \olref[syn][trw]{defn:mmodels} મુજબ, અને $\lfalse \notin \Delta$
        \olref[ccs]{prop:ccs-properties}\olref[ccs]{prop:ccs-lfalse}
        મુજબ.}

    \item \indcase{!A}{\ltrue}{$\mSat{M^\Sigma}{\ltrue}[\Delta]$
        \olref[syn][trw]{defn:mmodels} મુજબ, અને $\ltrue \in \Delta$
        \olref[ccs]{prop:ccs-properties}\olref[ccs]{prop:ccs-ltrue}
        મુજબ.}

    \item \indcase{!A}{p}{$\mSat{M^\Sigma}{p}[\Delta]$ જો અને માત્ર જો
      $\Delta \in V^\Sigma(p)$, \olref[syn][trw]{defn:mmodels} મુજબ.
      વધુમાં, $V^\Sigma$ની વ્યાખ્યાથી $\Delta \in V^\Sigma(p)$ જો અને
      માત્ર જો $p \in \Delta$.}

    \item \indcase{!A}{\lnot !B}{%
        $\mSat{M^\Sigma}{\lnot !B}[\Delta]$
          જો અને માત્ર જો $\mSat/{M^\Sigma}{!B}[\Delta]$
          (\olref[syn][trw]{defn:mmodels}), જો અને માત્ર જો
          $!B \notin \Delta$ (અનુમાનની પૂર્વધારણાથી), જો અને માત્ર જો
          $\lnot !B \in \Delta$
          (\olref[ccs]{prop:ccs-properties}\olref[ccs]{prop:ccs-lnot}
          મુજબ).}

    \item \indcase{!A}{!B \land !C}{%
        $\mSat{M^\Sigma}{!B \land
            !C}[\Delta]$ જો અને માત્ર જો
          $\mSat{M^\Sigma}{!B}[\Delta]$ અને
          $\mSat{M^\Sigma}{!C}[\Delta]$ બંને
          (\olref[syn][trw]{defn:mmodels} મુજબ), જો અને માત્ર જો
          $!B \in \Delta$ અને $!C \in \Delta$ બંને (અનુમાનની
          પૂર્વધારણાથી), જો અને માત્ર જો $!B \land !C \in
          \Delta$ (\olref[ccs]{prop:ccs-properties}%
          \olref[ccs]{prop:ccs-land} મુજબ).}

    \item \indcase{!A}{!B \lor !C}{%
        $\mSat{M^\Sigma}{!B \lor
            !C}[\Delta]$ જો અને માત્ર જો
          $\mSat{M^\Sigma}{!B}[\Delta]$ અથવા
          $\mSat{M^\Sigma}{!C}[\Delta]$
          (\olref[syn][trw]{defn:mmodels} મુજબ), જો અને માત્ર જો
          $!B \in \Delta$ અથવા $!C \in \Delta$ (અનુમાનની
          પૂર્વધારણાથી), જો અને માત્ર જો $!B \lor !C \in
          \Delta$ (\olref[ccs]{prop:ccs-properties}%
          \olref[ccs]{prop:ccs-lor} મુજબ).}

    \item \indcase{!A}{!B \lif !C}{%
        $\mSat{M^\Sigma}{!B \lif
            !C}[\Delta]$ જો અને માત્ર જો
          $\mSat/{M^\Sigma}{!B}[\Delta]$ અથવા
          $\mSat{M^\Sigma}{!C}[\Delta]$
          (\olref[syn][trw]{defn:mmodels} મુજબ), જો અને માત્ર જો
          $!B \notin \Delta$ અથવા $!C \in \Delta$ (અનુમાનની
          પૂર્વધારણાથી), જો અને માત્ર જો $!B \lif !C \in
          \Delta$ (\olref[ccs]{prop:ccs-properties}%
          \olref[ccs]{prop:ccs-lif} મુજબ).}

    \item \indcase{!A}{!B \liff !C}{%
        $\mSat{M^\Sigma}{!B \liff
            !C}[\Delta]$ જો અને માત્ર જો કાં તો
          $\mSat{M^\Sigma}{!B}[\Delta]$ અને
          $\mSat{M^\Sigma}{!C}[\Delta]$ બંને, અથવા
          $\mSat/{M^\Sigma}{!B}[\Delta]$ અને
          $\mSat/{M^\Sigma}{!C}[\Delta]$ બંને
          (\olref[syn][trw]{defn:mmodels} મુજબ), જો અને માત્ર જો
          કાં તો $!B \in \Delta$ અને $!C \in \Delta$ બંને, અથવા
          $!B \notin \Delta$ અને $!C \notin \Delta$ બંને (અનુમાનની
          પૂર્વધારણાથી), જો અને માત્ર જો $!B \liff !C \in
          \Delta$ (\olref[ccs]{prop:ccs-properties}%
          \olref[ccs]{prop:ccs-liff} મુજબ).}

    \item \indcase{!A}{\Box !B}{%
        પહેલાં ધારો કે
          $\mSat{M^\Sigma}{\Box !B}[\Delta]$.
          \olref[syn][trw]{defn:mmodels} મુજબ $R^\Sigma
          \Delta\Delta'$ હોય તેવા દરેક $\Delta'$ માટે
          $\mSat{M^\Sigma}{!B}[\Delta']$. અનુમાનની પૂર્વધારણાથી
          $R^\Sigma \Delta\Delta'$ હોય તેવા દરેક $\Delta'$ માટે
          $!B \in \Delta'$. $R^\Sigma$ની વ્યાખ્યાથી
          $\Box^{-1} \Delta \subseteq \Delta'$ હોય તેવા દરેક
          $\Delta'$ માટે $!B \in \Delta'$.
          \olref[mod]{prop:box} મુજબ $\Box !B \in \Delta$.

          હવે $\Box !B \in \Delta$ ધારો. એવો $\Delta' \in
          W^\Sigma$ લો કે $R^\Sigma \Delta\Delta'$, એટલે કે
          $\Box^{-1}\Delta \subseteq \Delta'$. $\Box !B \in
          \Delta$ હોવાથી $!B \in \Box^{-1} \Delta$. તેથી
          $!B \in \Delta'$. અનુમાનની પૂર્વધારણાથી
          $\mSat{M^\Sigma}{!B}[\Delta']$. $R^\Sigma
          \Delta\Delta'$ હોય તેવો $\Delta'$ મનસ્વી હોવાથી,
          $R^\Sigma \Delta\Delta'$ હોય તેવા બધા $\Delta' \in
          W^\Sigma$ માટે $\mSat{M^\Sigma}{!B}[\Delta']$.
          \olref[syn][trw]{defn:mmodels} મુજબ
          $\mSat{M^\Sigma}{\Box !B}[\Delta]$.}

    \item \indcase{!A}{\Diamond !B}{%
        પહેલાં ધારો કે
          $\mSat{M^\Sigma}{\Diamond !B}[\Delta]$.
          \olref[syn][trw]{defn:mmodels} મુજબ એવો કોઈ $\Delta'$ છે કે
          $R^\Sigma \Delta\Delta'$ અને
          $\mSat{M^\Sigma}{!B}[\Delta']$. અનુમાનની પૂર્વધારણાથી
          એવો કોઈ $\Delta'$ છે કે $R^\Sigma \Delta\Delta'$ અને
          $!B \in \Delta'$. $R^\Sigma$ની વ્યાખ્યાથી,%
           એવો કોઈ $\Delta'$ છે કે
            $\Box^{-1} \Delta \subseteq \Delta'$ અને
            $!B \in \Delta'$. \olref[mod]{lem:box-iff-diamond}
            મુજબ, એવો કોઈ $\Delta'$ છે કે
          $\Diamond \Delta' \subseteq \Delta$ અને
          $!B \in \Delta'$.
          \sourcecorrection{OLMD-037}{મૂળમાં બૉક્સ-શાખામાં આ
          સમાવેશને ડાયમંડના ગુણધર્મથી મેળવ્યો છે; અહીં જરૂરી બે
          પ્રાપ્યતા-શરતોની સમકક્ષતાના સહાયક પ્રમેયનો સંદર્ભ આપ્યો છે.}
          $!B \in \Delta'$ હોવાથી $\Diamond !B \in \Diamond
          \Delta'$; તેથી $\Diamond !B \in \Delta$.

          હવે $\Diamond !B \in \Delta$ ધારો.
          \olref[mod]{prop:diamond} મુજબ એવો પૂર્ણ
          $\Sigma$-સુસંગત $\Delta' \in W^\Sigma$ છે કે
          $\Diamond \Delta' \subseteq \Delta$ અને $!B \in
          \Delta'$. \olref[mod]{lem:box-iff-diamond}
            મુજબ એવો $\Delta' \in W^\Sigma$ છે કે $\Box^{-1}
            \Delta \subseteq \Delta'$ અને $!B \in \Delta'$.
          $R^\Sigma$ની વ્યાખ્યાથી $R^\Sigma \Delta\Delta'$; એટલે
          એવો $\Delta' \in W^\Sigma$ છે કે $R^\Sigma
          \Delta\Delta'$ અને $!B \in \Delta'$. અનુમાનની
          પૂર્વધારણાથી તે સૂત્ર ત્યાં સત્ય છે.
          \sourcecorrection{OLMD-038}{મૂળની આ દિશામાં સભ્યતામાંથી
          જગત પર સત્યતા મેળવવા જરૂરી અનુમાનનો પગલો છોડાયો છે;
          તે અહીં સ્પષ્ટ કર્યો છે.}
          \olref[syn][trw]{defn:mmodels} મુજબ
          $\mSat{M^\Sigma}{\Diamond !B}[\Delta]$.}
  \end{enumerate}
\end{proof}

\begin{probtag}{probFalse,probTrue,probNot,probAnd,probOr,probIf,probIff,probBox,probDiamond}
  \sourcecorrection{OLMD-039}{મૂળની પ્રશ્ન-ટૅગ યાદીમાં સંયોજનના
  ટૅગનું એક અક્ષર નાનું છે, જ્યારે અનુમાન-કિસ્સામાં તેનું પહેલું
  અક્ષર મોટું છે; બંનેને મેળ ખાતાં કર્યા છે.}
  \olref[nml][com][tru]{prop:truthlemma}ની સાબિતી પૂરી કરો.
\end{probtag}
% END OLP-0447
\endgroup

\begingroup
\makeatletter\def\ol@sectioncs{section}\makeatother

% BEGIN OLP-0448
\olfileid{nml}{com}{cmk}

\olsection{\Log{K} માટે નિર્ધારણ અને પૂર્ણતા}
\phantomsection\label{guunit:OLP-0448}


હવે આપણે કૅનોનિકલ નિદર્શન વડે પૂર્ણતા સ્થાપિત કરવા તૈયાર છીએ.
પૂર્ણતા એ હકીકતમાંથી ફલિત થાય છે કે~$\Sigma$ના કૅનોનિકલ નિદર્શનમાં
સત્ય સૂત્રો ચોક્કસપણે $\Sigma$માંથી નિષ્પન્ન કરી શકાય એવું સૂત્રો જ છે.
આ ગુણધર્મવાળા નિદર્શનો $\Sigma$નું \emph{નિર્ધારણ કરે છે} એમ કહેવાય છે.

\begin{defn}
  નિદર્શન $\mModel{M}$ કોઈ સામાન્ય મોડલ તર્ક~$\Sigma$નું
  \emph{નિર્ધારણ કરે છે} જો અને માત્ર જો બધાં સૂત્રો~$!A$
  માટે $\mSat{M}{!A}$ જો અને માત્ર જો $\Sigma \Proves !A$.
\end{defn}

\begin{thm}[નિર્ધારણ]\ollabel{thm:determination}
   $\mSat{M^\Sigma}{!A}$ જો અને માત્ર જો $\Sigma \Proves !A$.
\end{thm}

\begin{proof}
  જો $\mSat{M^\Sigma}{!A}$ હોય, તો દરેક પૂર્ણ
  $\Sigma$-સુસંગત $\Delta$ માટે $\mSat{M^\Sigma}{!A}[\Delta]$.
  તેથી સત્યતા સહાયક પ્રમેય મુજબ દરેક પૂર્ણ $\Sigma$-સુસંગત
  $\Delta$માં $!A \in \Delta$; એટલે
  \olref[lin]{cor:provability-characterization} મુજબ ($\Gamma =
  \emptyset$ લઈને) $\Sigma \Proves !A$.

  ઊલટું, જો $\Sigma \Proves !A$ હોય, તો
  \olref[ccs]{prop:ccs-properties}\olref[ccs]{prop:ccs-closed}
  મુજબ દરેક પૂર્ણ $\Sigma$-સુસંગત $\Delta$માં~$!A$ છે; તેથી
  સત્યતા સહાયક પ્રમેય મુજબ દરેક~$\Delta \in W^\Sigma$ માટે
  $\mSat{M^\Sigma}{!A}[\Delta]$, એટલે કે
  $\mSat{M^\Sigma}{!A}$.
\end{proof}

\Log{K}નું કૅનોનિકલ નિદર્શન~\Log{K}નું નિર્ધારણ કરતું હોવાથી
પરિણામ તરીકે~\Log{K}ની પૂર્ણતા તરત મળે છે:

\begin{cor}\ollabel{cor:Kcomplete}
  પાયાનો મોડલ તર્ક \Log{K} બધાં નિદર્શનોના વર્ગને અનુલક્ષીને પૂર્ણ
  છે, એટલે કે $\Entails !A$ હોય તો $\Log{K} \Proves !A$.
\end{cor}

\begin{proof}
  પ્રતિપક્ષી રીતે, જો $\Log{K} \Proves/ !A$ હોય, તો નિર્ધારણ મુજબ
  $\mSat/{M^{\Log{K}}}{!A}$; તેથી $!A$ માન્ય નથી.
\end{proof}

નિદર્શનોના કોઈ વર્ગને અનુલક્ષીને તંત્ર $\Sigma$ની પૂર્ણતાનો સામાન્ય
કિસ્સો લો, જેમ કે સ્વપ્રતિબિંબિત, સંમિત અને પરંપરિત નિદર્શનોના વર્ગને
અનુલક્ષીને \Log{KTB4}ની પૂર્ણતા. અહીં માત્ર નિર્ધારણ પૂરતું નથી.
તંત્ર $\Sigma$નું કૅનોનિકલ નિદર્શન તે વર્ગનું સભ્ય છે એ પણ બતાવવું
પડશે. કૅનોનિકલ નિદર્શનની રચનામાંથી આ સ્પષ્ટ રીતે ફલિત થતું નથી—અને
તે હંમેશાં સાચું પણ નથી!
% END OLP-0448
\endgroup

\begingroup
\makeatletter\def\ol@sectioncs{section}\makeatother

% BEGIN OLP-0449
\olfileid{nml}{com}{fra}

\olsection{ફ્રેમ-પૂર્ણતા}
\phantomsection\label{guunit:OLP-0449}


જ્યારે આપણે બતાવીએ કે આપેલા તર્કનું કૅનોનિકલ નિદર્શન અનુરૂપ ફ્રેમ
ગુણધર્મ ધરાવે છે, ત્યારે \Log{K}નો પૂર્ણતા પ્રમેય બીજા મોડલ તંત્રો
સુધી વિસ્તારી શકાય છે.

\begin{thm}\ollabel{thm:completeframeprops}
  જો સામાન્ય મોડલ તર્ક $\Sigma$માં \olref{tab:correspondencetable}ની
  ડાબી બાજુનાં સૂત્રો પૈકી કોઈ એક હોય, તો~$\Sigma$નું કૅનોનિકલ
  નિદર્શન જમણી બાજુનો અનુરૂપ ગુણધર્મ ધરાવે છે.
\end{thm}

\begin{table}[htp]
  \centering
    \begin{tabular}{| l || l |}
      \hline
      \emph{જો $\Sigma$માં \dots{} હોય} & \emph{તો $\Sigma$નું કૅનોનિકલ
        નિદર્શન આ પ્રકારનું છે:} \\
      \hline \hline
      \Ax{D}:\;\; $\Box!A \lif \Diamond !A$ & \quad સિરિયલ; \\
      \hline
      \Ax{T}:\;\; $\Box!A \lif !A$ &\quad સ્વપ્રતિબિંબિત;\\
      \hline
      \Ax{B}: \;\;$!A \lif \Box\Diamond!A$ &\quad સંમિત; \\
      \hline
      \Ax{4}: \;\; $\Box!A \lif \Box\Box!A$ & \quad પરંપરિત; \\
      \hline
      \Ax{5}: \;\; $\Diamond !A \lif \Box\Diamond!A$& \quad યુક્લિડિયન.\\
      \hline
    \end{tabular}
    \caption{મૂળભૂત અનુરૂપતાની હકીકતો.}
    \ollabel{tab:correspondencetable}
\end{table}

\begin{proof}
  આ દરેક કિસ્સો ક્રમે લઈએ.

  ધારો કે $\Sigma$માં \Ax{D} છે, અને $\Delta \in W^\Sigma$ લો.
  એવું $\Delta'$ છે કે $R^\Sigma \Delta\Delta'$ તે બતાવવું છે.
  $\Box^{-1}\Delta$ એ $\Sigma$-સુસંગત છે તે બતાવવું પૂરતું છે:
  ત્યાર પછી લિન્ડનબાઉમના સહાયક પ્રમેયથી એવો પૂર્ણ
  $\Sigma$-સુસંગત ગણ $\Delta' \supseteq \Box^{-1}\Delta$ મળે છે,
  અને $R^\Sigma$ની વ્યાખ્યા
  મુજબ
  $R^\Sigma \Delta\Delta'$. હવે વિરોધ માટે ધારો કે
  $\Box^{-1}\Delta$ એ $\Sigma$-સુસંગત \emph{નથી}, એટલે કે
  $\Box^{-1}\Delta \Proves[\Sigma] \lfalse$.
  \olref[mod]{lem:box2} મુજબ $\Delta \Proves[\Sigma]
  \Box\lfalse$; અને $\Sigma$માં \Ax{D} હોવાથી
  $\Delta \Proves[\Sigma] \Diamond\lfalse$ પણ મળે છે.
  પરંતુ $\Sigma$ સામાન્ય હોવાથી $\Sigma \Proves
  \lnot\Diamond \lfalse$ (\olref[prf][nor]{prop:notDiamondBot});
  તેથી $\Delta \Proves[\Sigma] \lnot\Diamond \lfalse$ પણ,
  જે~$\Delta$ની સુસંગતતાનો વિરોધ કરે છે.

  હવે ધારો કે $\Sigma$માં \Ax{T} છે, અને $\Delta \in W^\Sigma$
  લો. $R^\Sigma \Delta\Delta$, એટલે કે
  %
  $\Box^{-1}\Delta \subseteq \Delta$, બતાવવું છે. પરંતુ
  $\Box!A \in \Delta$ હોય તો \Ax{T}થી $!A \in \Delta$ પણ,
  જેમ જોઈતું હતું.

  હવે ધારો કે $\Sigma$માં \Ax{B} છે, અને $\Delta$,
  $\Delta' \in W^\Sigma$ માટે $R^\Sigma \Delta\Delta'$.
  $R^\Sigma \Delta'\Delta$ બતાવવું છે, એટલે કે
  $\Box^{-1}\Delta' \subseteq \Delta$.
    \olref[mod]{lem:box-iff-diamond} મુજબ આ સમકક્ષ છે
  $\Diamond\Delta \subseteq \Delta'$. તેથી $!A \in \Delta$
  ધારો. \Ax{B}થી $\Box\Diamond!A \in \Delta$ પણ.
  $R^\Sigma \Delta\Delta'$ની પૂર્વધારણા
   મુજબ
  $\Box^{-1}\Delta \subseteq \Delta'$; તેથી જરૂરી મુજબ
  $\Diamond!A \in \Delta'$.

  હવે ધારો કે $\Sigma$માં \Ax{4} છે, અને $R^\Sigma
  \Delta_1\Delta_2$ તથા $R^\Sigma \Delta_2\Delta_3$.
  $R^\Sigma \Delta_1\Delta_3$ બતાવવું છે. પૂર્વધારણા
   મુજબ
  $\Box^{-1}\Delta_1 \subseteq \Delta_2$ અને
  $\Box^{-1}\Delta_2 \subseteq \Delta_3$ બંને મળે છે.
  $R^\Sigma \Delta_1\Delta_3$ માટે
  $\Box^{-1}\Delta_1 \subseteq \Delta_3$ બતાવવું પૂરતું છે.
  તેથી $!B \in \Box^{-1}\Delta_1$, એટલે કે
  $\Box!B \in \Delta_1$, લો. \Ax{4}થી
  $\Box\Box!B \in \Delta_1$ પણ; પછી પૂર્વધારણાથી પહેલાં
  $\Box!B \in \Delta_2$ અને ત્યાર બાદ $!B \in \Delta_3$,
  જેમ જોઈતું હતું.

  હવે ધારો કે $\Sigma$માં \Ax{5} છે, અને $R^\Sigma
  \Delta_1\Delta_2$ તથા $R^\Sigma \Delta_1\Delta_3$.
  $R^\Sigma \Delta_2\Delta_3$ બતાવવું છે. પહેલી પૂર્વધારણાથી
  $\Box^{-1}\Delta_1 \subseteq \Delta_2$
   મળે છે;
  બીજી પૂર્વધારણા $\Diamond\Delta_3 \subseteq \Delta_1$ને
  સમકક્ષ છે,
    \olref[mod]{lem:box-iff-diamond} મુજબ.
  $R^\Sigma \Delta_2\Delta_3$ બતાવવા
  \olref[mod]{lem:box-iff-diamond} મુજબ
    એટલું પૂરતું છે કે%
  $\Diamond\Delta_3 \subseteq \Delta_2$.
  તેથી $\Diamond!A \in \Diamond\Delta_3$, એટલે કે
  $!A \in \Delta_3$, લો. બીજી પૂર્વધારણાથી
  $\Diamond!A \in \Delta_1$ અને \Ax{5}થી
  $\Box\Diamond!A \in \Delta_1$ પણ. હવે પહેલી પૂર્વધારણાથી
  $\Diamond!A \in \Delta_2$, જેમ જોઈતું હતું.
  \end{proof}

પરિણામ તરીકે અનેક તંત્રોના પૂર્ણતા-પરિણામો મળે છે. દાખલા તરીકે,
$\Log{S5} = \Log{KT5} = \Log{KTB4}$ એ બધાં સ્વપ્રતિબિંબિત અને
યુક્લિડિયન નિદર્શનોના વર્ગને અનુલક્ષીને પૂર્ણ છે; આ વર્ગ બધાં
સ્વપ્રતિબિંબિત, સંમિત અને પરંપરિત નિદર્શનોના વર્ગ સમાન છે.

\begin{thm}\ollabel{thm:generaldet}
  $\mClass{C}_\Ax{D}$, $\mClass{C}_\Ax{T}$,
  $\mClass{C}_\Ax{B}$, $\mClass{C}_\Ax{4}$ અને
  $\mClass{C}_\Ax{5}$ અનુક્રમે બધાં સિરિયલ, સ્વપ્રતિબિંબિત,
  સંમિત, પરંપરિત અને યુક્લિડિયન નિદર્શનોના વર્ગો હોય. ત્યારે
  \Ax{D}, \Ax{T}, \Ax{B}, \Ax{4} અને \Ax{5}માંથી કોઈપણ
  સ્કીમા $!A_1$, \dots, $!A_n$ માટે તંત્ર
  $\Log{K}!A_1 \dots !A_n$ નિદર્શનોના વર્ગ
  $\mClass{C} = \mClass{C}_{!A_1} \cap \dots
  \cap \mClass{C}_{!A_n}$ વડે નિર્ધારિત થાય છે.
\end{thm}

\begin{prop}
  ધારો કે $\Sigma$ સામાન્ય મોડલ તર્ક છે. ત્યારે:
  \begin{enumerate}
  \item\ollabel{prop:anotherfive-a}%
    જો $\Sigma$માં સ્કીમા $\Diamond!A \lif \Box
    !A$ હોય, તો $\Sigma$નું કૅનોનિકલ નિદર્શન આંશિક વિધેયાત્મક છે.
  \item જો $\Sigma$માં સ્કીમા $\Diamond!A \liff \Box
    !A$ હોય, તો $\Sigma$નું કૅનોનિકલ નિદર્શન વિધેયાત્મક છે.
  \item જો $\Sigma$માં સ્કીમા $\Box\Box!A \lif \Box
    !A$ હોય, તો $\Sigma$નું કૅનોનિકલ નિદર્શન નબળે સઘન છે.
  \end{enumerate}
(આ ફ્રેમ ગુણધર્મોની વ્યાખ્યાઓ માટે \olref[frd][acc]{tab:anotherfive}
  જુઓ).
\end{prop}

\begin{proof}
  \begin{enumerate}
  \item ધારો કે $\Sigma$માં સ્કીમા $\Diamond !A \lif
    \Box !A$ છે. $R^\Sigma$ આંશિક વિધેયાત્મક છે તે બતાવવા
    કોઈપણ $\Delta_1$, $\Delta_2$, $\Delta_3 \in W^\Sigma$
    માટે $R^\Sigma \Delta_1\Delta_2$ અને $R^\Sigma
    \Delta_1\Delta_3$ હોય તો $\Delta_2=\Delta_3$ તે સાબિત કરવું
    પડશે. $R^\Sigma \Delta_1\Delta_2$ હોવાથી
    $\Box^{-1}\Delta_1 \subseteq \Delta_2$, અને
    $R^\Sigma \Delta_1\Delta_3$ હોવાથી
    $\Box^{-1}\Delta_1 \subseteq \Delta_3$ પણ
    .
    $\Delta_2=\Delta_3$ માટે બે સમાવેશ
    $\Delta_2 \subseteq \Delta_3$ અને $\Delta_3 \subseteq
    \Delta_2$ સ્થાપિત કરવાનું પૂરતું છે. પહેલા સમાવેશ માટે
    $!A \in \Delta_2$ લો; ત્યારે $\Diamond!A \in \Delta_1$,
    અને સ્કીમા તથા $\Delta_1$ની નિષ્પત્તિ હેઠળની બંધતાથી
    $\Box!A \in \Delta_1$ પણ. તેથી $R^\Sigma
    \Delta_1\Delta_3$ની પૂર્વધારણાથી $!A \in \Delta_3$.
    બીજો સમાવેશ એ જ રીતે મળે છે.

  \item આ મુદ્દા~\olref{prop:anotherfive-a} અને
    \olref{thm:completeframeprops}માં સિરિયલતાના પ્રમાણમાંથી
    તરત ફલિત થાય છે.

  \item ધારો કે $\Sigma$માં સ્કીમા $\Box\Box!A \lif \Box!A$
    છે. $R^\Sigma$ નબળે સઘન છે તે બતાવવા $R^\Sigma
    \Delta_1\Delta_2$ લો. એવું પૂર્ણ $\Sigma$-સુસંગત
    ગણ $\Delta_3$ છે કે $R^\Sigma \Delta_1\Delta_3$ અને
    $R^\Sigma \Delta_3\Delta_2$ તે બતાવવું છે. મૂકો:
    \[
    \Gamma = \Box^{-1}\Delta_1 \cup \Diamond\Delta_2.
    \]
    $\Gamma$ એ $\Sigma$-સુસંગત છે તે બતાવવું પૂરતું છે;
    ત્યારે લિન્ડનબાઉમના સહાયક પ્રમેયથી તેને એવા પૂર્ણ
    $\Sigma$-સુસંગત ગણ~$\Delta_3$ સુધી વિસ્તારી શકાય કે
    $\Box^{-1}\Delta_1 \subseteq \Delta_3$ અને
    $\Diamond\Delta_2 \subseteq \Delta_3$, એટલે કે
    $R^\Sigma \Delta_1\Delta_3$ અને $R^\Sigma
    \Delta_3\Delta_2$ (
      \olref[mod]{lem:box-iff-diamond} મુજબ).

    વિરોધ માટે ધારો કે $\Gamma$ સુસંગત નથી. ત્યારે એવાં
    સૂત્રો $\Box!A_1$, \dots, $\Box!A_n \in \Delta_1$
    અને $!B_1$, \dots,~$!B_m \in \Delta_2$ છે કે
    \[!A_1, \dots,
    !A_n, \Diamond!B_1, \dots, \Diamond!B_m \Proves[\Sigma]
    \lfalse.\]
    $\Diamond (!B_1 \land \dots \land !B_m) \to
    (\Diamond!B_1 \land \dots \land \Diamond!B_m)$ દરેક સામાન્ય
    મોડલ તર્કમાં નિષ્પન્ન કરી શકાય એવું હોવાથી આપણે નીચે મુજબ દલીલ કરીને
    $\Delta_2$ની સુસંગતતાનો વિરોધ મેળવીએ છીએ:
    \begin{align*}
      !A_1, \dots, !A_n, & \Diamond!B_1, \dots,
      \Diamond!B_m \Proves[\Sigma] \lfalse \\
      !A_1, \dots,!A_n & \Proves[\Sigma] (\Diamond!B_1 \land
      \dots \land \Diamond!B_m) \lif \lfalse\\
      & \qquad\text{નિષ્પત્તિ પ્રમેયથી}\\
      & \qquad\text{\olref[prf][prp]{prop:derivabilityfacts}\olref[prf][prp]{prop:derivabilityfacts-deduction}, અને \Taut} \\
      !A_1, \dots,!A_n & \Proves[\Sigma]
      \Diamond(!B_1 \land
      \dots \land !B_m) \lif \lfalse\\
      & \qquad \text{કારણ કે $\Sigma$ સામાન્ય છે} \\
      !A_1, \dots,!A_n & \Proves[\Sigma]
      \lnot\Diamond (!B_1 \land \dots \land !B_m)\\
      & \qquad\text{\PL\ મુજબ} \\
      !A_1, \dots,!A_n & \Proves[\Sigma]
      \Box\lnot (!B_1 \land \dots \land !B_m)\\
      & \qquad\text{$\lnot\Diamond$ને બદલે $\Box\lnot$} \\
      \Box!A_1, \dots,\Box!A_n
      & \Proves[\Sigma] \Box\Box \lnot (!B_1 \land \dots \land !B_m)\\
      & \qquad\text{\olref[mod]{lem:box1} મુજબ} \\
      \Box!A_1, \dots,\Box!A_n
      & \Proves[\Sigma]
      \Box\lnot (!B_1 \land \dots \land !B_m)\\
      &\qquad\text{સ્કીમા $\Box\Box!A \lif \Box!A$થી} \\
      \Delta_1 & \Proves[\Sigma]
      \Box\lnot (!B_1 \land \dots \land !B_m)\\
      &\qquad\text{એકદિશવર્ધિતાથી, \olref[prf][prp]{prop:derivabilityfacts}\olref[prf][prp]{prop:derivabilityfacts-monotonicity}} \\
      & \Box\lnot (!B_1 \land \dots \land !B_m) \in \Delta_1\\
      &\qquad\text{નિષ્પત્તિ હેઠળની બંધતાથી}; \\
      & \lnot (!B_1 \land \dots \land !B_m) \in \Delta_2\\
      &\qquad \text{કારણ કે }
      R^\Sigma \Delta_1\Delta_2.
    \end{align*}
    પરંતુ પસંદ કરેલાં બધાં સૂત્રો બીજા ગણમાં પણ છે, જેથી તેની
    સુસંગતતાનો વિરોધ મળે છે.
  \end{enumerate}
\end{proof}

આ ઉદાહરણોથી એવું લાગી શકે કે મોડલ તર્કનું દરેક તંત્ર~$\Sigma$
એ અર્થમાં \emph{પૂર્ણ} છે કે $\Sigma$ના દરેક પ્રમેય માન્ય હોય તેવી
દરેક ફ્રેમમાં માન્ય હોય તેવું દરેક સૂત્ર તે સાબિત કરે છે.
દુર્ભાગ્યે, આ અર્થમાં પૂર્ણ ન હોય એવાં ઘણાં તંત્રો છે.
% END OLP-0449
\endgroup


\OLEndChapterHook
% END OLP-0441
\endgroup


\begingroup
\makeatletter\def\ol@sectioncs{section}\makeatother

% BEGIN OLP-0450
\olchapter{nml}{fil}{ફિલ્ટ્રેશનો અને નિર્ણેયતા}
\olchapterid{nml}{fil}
\phantomsection\label{guunit:OLP-0450}


\begingroup
\makeatletter\def\ol@sectioncs{section}\makeatother

% BEGIN OLP-0451
\olfileid{nml}{fil}{int}

\olsection{પરિચય}
\phantomsection\label{guunit:OLP-0451}


કોઈ તર્ક વિશેનો એક અગત્યનો પ્રશ્ન હંમેશાં એ છે કે તે નિર્ણેય છે કે
નહિ, એટલે કે ``આ સૂત્ર માન્ય છે?'' એ પ્રશ્નનો ઉત્તર આપતી
કોઈ અસરકારક પ્રક્રિયા છે કે નહિ. વિધાનાત્મક તર્ક નિર્ણેય છે:
સત્યાર્થતા સારણી રચીને સૂત્ર પુનરુક્તિ છે કે નહિ તે
અસરકારક રીતે ચકાસી શકાય છે, અને આપેલા સૂત્ર માટે સારણી
સાન્ત છે. પરંતુ મોડલ સૂત્ર બધાં નિદર્શનોમાં સત્ય છે કે નહિ તે સીધું
ચકાસી શકાતું નથી, કારણ કે નિદર્શનો અનંત સંખ્યામાં છે. આપેલા સૂત્રને
સંબંધિત બધાં સાન્ત નિદર્શનોની યાદી બનાવી શકાય છે, કારણ કે માત્ર
સૂત્રમાં ખરેખર આવતાં પ્રતિજ્ઞાપન ચલોને જગતોના
ઉપગણો નિયુક્ત કરવાથી ફરક પડે છે. પ્રાપ્યતા સંબંધ નિશ્ચિત હોય, તો
$V(p)$ની જુદી જુદી શક્ય નિયુક્તિઓ એ~$W$ના બધા ઉપગણો જ છે;
$\card{W} = n$ હોય તો એવા $2^n$ ઉપગણો છે. આપણા
સૂત્ર~$!A$માં $m$ પ્રતિજ્ઞાપન ચલો હોય તો~$n$ જગતવાળાં $2^{nm}$ જુદાં નિદર્શનો મળે છે.
દરેક નિદર્શન માટે દરેક જગતમાં~$!A$નું સત્યમૂલ્ય ગણીને સૂત્ર~$!A$
બધાં જગતોમાં સત્ય છે કે નહિ તે ચકાસી શકાય છે. અલબત્ત, બધા શક્ય
પ્રાપ્યતા સંબંધો પણ ચકાસવા પડે; પરંતુ~$n$ જગતો પરના સંબંધો પણ
સાન્ત સંખ્યામાં જ છે (ચોક્કસપણે, $W \times W$ના ઉપગણોની સંખ્યા,
એટલે કે $2^{n^2}$).
\sourcecorrection{OLMD-040}{મૂળના આ ગણતરી-વાક્યમાં ખુલ્લા કૌંસનું
બંધ કૌંસ છૂટી ગયું છે; અહીં તે પુનઃસ્થાપિત કર્યું છે.}

જો આપણને તર્ક~\Log{K} નહિ, પરંતુ નિદર્શનોના કોઈ વર્ગથી
વ્યાખ્યાયિત તર્કમાં રસ હોય (દાખલા તરીકે, સ્વપ્રતિબિંબિત પરંપરિત
નિદર્શનોમાં), તો પ્રાપ્યતા સંબંધ યોગ્ય પ્રકારનો છે કે નહિ તે પણ
ચકાસી શકવું જોઈએ. આપણને રસ ધરાવતી ફ્રેમો મોડલ સૂત્રો વડે
વ્યાખ્યાયિત થઈ શકે ત્યારે એ કરી શકીએ છીએ (દાખલા તરીકે,
\Ax{T} અને \Ax{4} ફ્રેમમાં માન્ય છે કે નહિ તે ચકાસીને). તેથી વિચાર
એવો છે કે બધી સાન્ત ફ્રેમો ક્રમસર લઈએ, દરેક ફ્રેમ આપણને રસ ધરાવતા
વર્ગમાં છે કે નહિ તે ચકાસીએ, પછી તે ફ્રેમ પરનાં બધાં શક્ય
નિદર્શનોની યાદી બનાવી દરેકમાં $!A$ સત્ય છે કે નહિ તે તપાસીએ.
કોઈમાં સત્ય ન હોય તો અટકીએ: $!A$ રસના નિદર્શનોના વર્ગમાં માન્ય
નથી.

આ વિચારમાં સમસ્યા છે: શોધ ક્યારે, કે કદી, બંધ કરવી તે આપણે જાણતા
નથી. જો સૂત્રનું સાન્ત વિરોધી નિદર્શન હોય, તો આપણી પ્રક્રિયા
તે શોધી કાઢશે. પરંતુ સાન્ત વિરોધી નિદર્શન ન હોય તો ઉત્તર નહિ મળે.
સૂત્ર માન્ય હોઈ શકે (એકેય વિરોધી નિદર્શન ન હોય), અથવા તેનું
માત્ર અનંત વિરોધી નિદર્શન હોઈ શકે, જે આપણે કદી તપાસવાના નથી.
જો આપણે બતાવી શકીએ કે વિરોધી નિદર્શન ધરાવતું દરેક સૂત્ર
સાન્ત વિરોધી નિદર્શન પણ ધરાવે છે, તો આ સમસ્યા દૂર થાય. આવી
પરિસ્થિતિમાં તર્કમાં \emph{સાન્ત નિદર્શન ગુણધર્મ} છે એમ કહીએ છીએ.

પણ તર્કમાં સાન્ત નિદર્શન ગુણધર્મ છે તે કેવી રીતે બતાવવું? એક રસ્તો
એ છે કે~$!A$નું અનંત (વિરોધી) નિદર્શન સાન્ત નિદર્શનમાં ફેરવવાની
રીત શોધવી. એવું કરી શકાય તો જ્યાં $!A$ સત્ય ન હોય એવા કોઈ
નિદર્શનમાંથી મળેલું સાન્ત નિદર્શન પણ $!A$ને અસત્ય જ બનાવશે.
એ સાન્ત નિદર્શન બધાં સાન્ત નિદર્શનોની આપણી યાદીમાં આવશે, અને
આખરે દરેક અમાન્ય સૂત્ર માટે તેની અમાન્યતા નક્કી કરી શકીશું.
જો સૂત્ર માન્ય હોય, તો આ પ્રક્રિયા અટકશે નહિ. જો વધારામાં બતાવી
શકીએ કે પ્રક્રિયાથી મળતા સાન્ત નિદર્શનના કદની કોઈ મહત્તમ મર્યાદા
છે અને તે મર્યાદા માત્ર સૂત્ર~$!A$ પર આધારિત છે, તો
વિરોધી નિદર્શન શોધતી વખતે સાન્ત નિદર્શનો કયા કદ સુધી ચકાસવાં
તે જાણીશું. ત્યાં સુધીમાં વિરોધી નિદર્શન ન મળે તો એવું કોઈ
નથી. આમ આપણી પ્રક્રિયા દરેક સૂત્ર~$!A$ માટે ``$!A$ માન્ય
છે?'' એ પ્રશ્નનો ખરેખર નિર્ણય કરશે.

અનંત રચનાઓને સાન્ત રચનાઓમાં ફેરવવાની એક ઘણી વાર સફળ થતી
યુક્તિ એ છે કે સંબંધિત બાબતોમાં એકસરખું વર્તન કરતા રચનાના
ઘટકોને એક તરીકે ``ઓળખી'' લઈએ. જો~$\mModel{M}$માં
સંબંધિત બાબતોમાં એકસરખું વર્તન કરતાં અનંત જગતો હોય, તો કદાચ
\emph{બધાં} જગતોને એવા જગતોના સાન્ત સંખ્યામાં (સંભવતઃ અનંત)
``વર્ગો''માં મૂકી શકીએ. બીજા શબ્દોમાં, જગતોના ગણનું યોગ્ય
વિભાજન કરવું જોઈએ: દરેક વર્ગમાં સાન્ત અથવા અનંત જગતો હોઈ શકે,
પણ વર્ગોની સંખ્યા સાન્ત જ હોય.
\sourcecorrection{OLMD-041}{મૂળમાં દરેક વિભાજન-વર્ગમાં અનંત
જગતો હોવા જોઈએ એવું લખાયું છે; સાન્ત વર્ગો પણ માન્ય છે અને અહીં
માત્ર વર્ગોની સંખ્યા સાન્ત હોવી જરૂરી છે.}
પછી નવું નિદર્શન~$\mModel{M^*}$ વ્યાખ્યાયિત કરીએ જેમાં જગતો
આ વર્ગો હોય. જૂના નિદર્શનના સાન્ત સંખ્યામાં વર્ગો નવા નિદર્શનમાં
સાન્ત સંખ્યામાં જગતો આપે છે, એટલે સાન્ત નિદર્શન મળે છે.
જગત~$w$ જે વર્ગમાં હોય તેને $[w]$ કહીએ. આપણે ઇચ્છીએ છીએ કે
$\mSat{M}{!A}[w]$ જો અને માત્ર જો $\mSat{M^*}{!A}[{[w]}]$,
કારણ કે જૂનું નિદર્શન~$!A$નું વિરોધી નિદર્શન હોય તો નવું પણ
એવું જ જોઈએ. આ માટે વિભાજન તથા~$\mModel{M^*}$નો પ્રાપ્યતા
સંબંધ યોગ્ય રીતે વ્યાખ્યાયિત કરવો પડશે.

આ રીત કેવી રીતે ચાલે તે જોવા પહેલાં કલ્પો કે પ્રાપ્યતા સંબંધ
સર્વત્ર છે. $\mSat{M}{\Box !B}[w]$ જો અને માત્ર જો કોઈ
$v \in W$ માટે $\mSat{M}{\Box !B}[v]$;~$\mModel{M^*}$ માટે પણ
એ જ વાત, માત્ર $[w]$ અને $[v]$ સાથે.
\sourcecorrection{OLMD-042}{મૂળમાં અહીં પ્રાપ્યતા સંબંધ જ ન હોય
એમ કહે છે, પરંતુ આગળ બૉક્સના કિસ્સામાં દરેક જગતના સૂત્ર-સત્યની
શરત વાપરે છે. તે શરત માટે અહીં સર્વત્ર પ્રાપ્યતાનો કિસ્સો
સ્પષ્ટ કર્યો છે.}
પ્રથમ વિચાર તરીકે, બે જગતો $u$ અને $v$ને સમકક્ષ (એ જ વર્ગનાં)
કહીએ જો તેઓ~$\mModel{M}$નાં બધાં પ્રતિજ્ઞાપન ચલો
પર સહમત હોય, એટલે કે $\mSat{M}{p}[u]$ જો અને માત્ર જો
$\mSat{M}{p}[v]$. $V^*(p) = \Setabs{[w]}{\mSat{M}{p}[w]}$ લો.
આપણું લક્ષ્ય $\mSat{M}{!A}[w]$ જો અને માત્ર જો
$\mSat{M^*}{!A}[{[w]}]$ બતાવવાનું છે. સ્પષ્ટ છે કે આ અનુમાનથી
સાબિત કરીશું: આધાર-કિસ્સો $!A \equiv p$ હશે. પહેલાં ધારો કે
$\mSat{M}{p}[w]$. ત્યારે વ્યાખ્યા મુજબ $[w] \in V^*(p)$,
તેથી $\mSat{M^*}{p}[{[w]}]$.
\sourcecorrection{OLMD-043}{મૂળના આ આધાર-કિસ્સામાં નિયુક્તિના
ચિહ્ન પછી વિધાનચલનો આર્ગ્યુમેન્ટ છૂટી ગયો છે; અગાઉની વ્યાખ્યા
સાથે તેને મેળ ખાતો કર્યો છે.}
હવે ધારો કે $\mSat{M^*}{p}[{[w]}]$. તેનો અર્થ
$[w] \in V^*(p)$, એટલે કે~$w$ને સમકક્ષ કોઈ $v$ માટે
$\mSat{M}{p}[v]$. પરંતુ ``$w$ એ $v$ને સમકક્ષ છે'' એટલે
``$w$ અને $v$ એ બધા સરખા પ્રતિજ્ઞાપન ચલોને સત્ય
બનાવે છે''; તેથી $\mSat{M}{p}[w]$. હવે અનુમાનનું પગલું લો,
દાખલા તરીકે $!A \ident \lnot !B$. ત્યારે
$\mSat{M}{\lnot !B}[w]$ જો અને માત્ર જો $\mSat/{M}{!B}[w]$,
જો અને માત્ર જો $\mSat/{M^*}{!B}[{[w]}]$ (અનુમાનની
પૂર્વધારણાથી), જો અને માત્ર જો
$\mSat{M^*}{\lnot !B}[{[w]}]$. બીજા બિનમોડલ સંચાલકો માટે પણ
એ જ રીતે. $\Box$ માટે પણ ચાલે છે: ધારો કે
$\mSat{M^*}{\Box !B}[{[w]}]$. તેનો અર્થ કે દરેક $[u]$ માટે
$\mSat{M^*}{!B}[{[u]}]$. અનુમાનની પૂર્વધારણાથી દરેક $u$ માટે
$\mSat{M}{!B}[u]$. તેથી $\mSat{M}{\Box !B}[w]$.

સામાન્ય કિસ્સામાં~$\mModel{M^*}$નો પ્રાપ્યતા સંબંધ પણ વ્યાખ્યાયિત
કરવો પડે છે અને વાત વધુ જટિલ બને છે. ઉપરનું અનુમાન-પ્રમાણ ચાલે
તે માટે જરૂરી શરતો તેનો પ્રાપ્યતા સંબંધ~$R^*$ પૂરી કરતો હોય,
તો નિદર્શન~$\mModel{M^*}$ને \emph{ફિલ્ટ્રેશન} કહીશું.
ત્યારે કોઈપણ ફિલ્ટ્રેશન~$\mModel{M^*}$માં $[w]$ પર $!A$
સત્ય છે જો અને માત્ર જો~$\mModel{M}$માં~$w$ પર $!A$ સત્ય છે.
પરંતુ હવે ફિલ્ટ્રેશનો \emph{હોય છે} એ પણ બતાવવું પડે, એટલે કે
જરૂરી શરતો પૂરી કરે તેવો~$R^*$ વ્યાખ્યાયિત કરી શકાય છે.
આ શક્ય બને તે માટે જગતો $u$, $v$ બધાં પ્રતિજ્ઞાપન ચલો
પર જ નહિ, પણ~$!A$નાં બધાં ઉપ-સૂત્રો પર સહમત હોય ત્યારે જ
સમકક્ષ ગણવા જોઈએ. $!A$નાં ઉપ-સૂત્રો સાન્ત સંખ્યામાં હોવાથી
ફિલ્ટ્રેશન સાન્ત છે તેની ખાતરી હજી મળે છે. ફિલ્ટ્રેશનને વ્યાખ્યાયિત
કરવાનો એક જ રસ્તો નથી; અને ફિલ્ટ્રેશનનો પ્રાપ્યતા સંબંધ જરૂરી
ગુણધર્મો (જેમ કે સ્વપ્રતિબિંબિત, પરંપરિત વગેરે) ધરાવે તે માટે
~$R^*$ની વ્યાખ્યામાં સર્જનાત્મકતા વાપરવી પડે છે.
% END OLP-0451
\endgroup

\begingroup
\makeatletter\def\ol@sectioncs{section}\makeatother

% BEGIN OLP-0452
\olfileid{nml}{fil}{pre}

\olsection{પ્રારંભિક બાબતો}
\phantomsection\label{guunit:OLP-0452}


ફિલ્ટ્રેશનો બતાવે છે કે આપણાં મોડલ તર્કનાં તંત્રોમાં \emph{સાન્ત
નિદર્શન ગુણધર્મ} છે, અને એ વડે તેમની નિર્ણેયતા સ્થાપિત કરી શકાય
છે. તેનો અર્થ એ કે કોઈ નિદર્શનમાં સત્ય (અસત્ય) હોય તેવું કોઈપણ
સૂત્ર સાન્ત નિદર્શનમાં પણ સત્ય (અસત્ય) હોય છે.
ફિલ્ટ્રેશનો સૂત્રોના એવા ગણોને અનુલક્ષીને વ્યાખ્યાયિત થાય છે
જે ઉપસૂત્રો હેઠળ બંધ હોય.

\begin{defn}\ollabel{defn:modallyclosed}
  સૂત્રોનો ગણ $\Gamma$ \emph{ઉપસૂત્રો હેઠળ બંધ} છે જો તે
  $\Gamma$માંના દરેક સૂત્રનાં બધાં ઉપસૂત્રો ધરાવે.
  વધુમાં, $\Gamma$ \emph{મોડલ રીતે બંધ} છે જો તે ઉપસૂત્રો હેઠળ
  બંધ હોય અને $!A \in \Gamma$ હોય ત્યારે $\Box!A,
  \Diamond!A \in \Gamma$ પણ હોય.
\end{defn}

દાખલા તરીકે, આપેલું સૂત્ર~$!A$ હોય ત્યારે તેનાં બધાં
ઉપ-સૂત્રોનો ગણ ઉપ-સૂત્રો હેઠળ બંધ છે. $!A$નાં આ
ઉપ-સૂત્રોના ગણ દ્વારા કોઈ નિદર્શનનું ફિલ્ટ્રેશન વ્યાખ્યાયિત કરીએ ત્યારે
આપણને જોઈતો ગુણધર્મ તેને હશે: મૂળ નિદર્શન $!A$ને સત્ય (અસત્ય)
બનાવે જો અને માત્ર જો તે પણ એવું કરે.

~$\Gamma$ દ્વારા~$\mModel{M}$ના ફિલ્ટ્રેશનનાં જગતોનો ગણ નીચેના
સમકક્ષતા સંબંધના બધા સમકક્ષતા-વર્ગો તરીકે વ્યાખ્યાયિત થાય છે.

\begin{defn}
ધારો કે $\mModel{M} =\tuple{W, R, V}$ અને $\Gamma$ ઉપ-સૂત્રો
હેઠળ બંધ છે. $W$ પર એવો સંબંધ $\equiv$ વ્યાખ્યાયિત કરો કે
$\Gamma$માંથી સરખાં સૂત્રોને સત્ય બનાવતાં કોઈપણ બે જગતો
તેમાં સંબંધિત હોય, એટલે કે:
\[
u \equiv v \quad \text{જો અને માત્ર જો } \quad
\forall !A \in \Gamma : \mSat{M}{!A}[u] \Leftrightarrow \mSat{M}{!A}[v].
\]
જગત~$w$નો સમકક્ષતા-વર્ગ~$[w]_\equiv$, અથવા સંક્ષેપમાં $[w]$,
એ~$w$ને $\equiv$-સમકક્ષ બધાં જગતોનો ગણ છે:
\[
[w] = \Setabs{v}{v \equiv w}.
\]
\end{defn}

\begin{prop}
  $\mModel{M}$ અને $\Gamma$ આપેલાં હોય, તો ઉપર વ્યાખ્યાયિત
  $\equiv$ સમકક્ષતા સંબંધ છે, એટલે કે તે સ્વપ્રતિબિંબિત, સંમિત
  અને પરંપરિત છે.
\end{prop}

\begin{proof}
  સંબંધ $\equiv$ સ્વપ્રતિબિંબિત છે, કારણ કે $w$ પોતે જે
  સૂત્રોને~$\Gamma$માંથી સત્ય બનાવે છે, એ જ પોતે બનાવે છે.
  તે સંમિત છે: $u$ અને $v$~$\Gamma$માંથી સરખાં સૂત્રોને
  સત્ય બનાવે, તો $v$ અને~$u$ માટે પણ એ જ સાચું છે.
  તે પરંપરિત પણ છે: $u$ અને~$v$ તથા $v$ અને~$w$ જો~$\Gamma$માંથી
  સરખાં સૂત્રોને સત્ય બનાવે, તો $u$ અને~$w$ પણ~$\Gamma$માંથી
  સરખાં સૂત્રોને સત્ય બનાવે.
\end{proof}

દરેક સમકક્ષતા સંબંધની જેમ $\equiv$ પણ~$W$ને \emph{વિભાજન-વર્ગો}માં
વહેંચે છે: એ~$W$ના એવા ઉપગણો છે જે જોડીજોડીમાં અસંયુક્ત છે અને
સાથે મળીને આખા~$W$ને આવરે છે. દરેક $w \in W$ આ વર્ગોમાંના
એકનો ઘટક છે, એટલે કે~$[w]$નો, કારણ કે $w \equiv w$.
આમ વર્ગો $[w]$ આખા~$W$ને આવરે છે. તેઓ જોડીજોડીમાં અસંયુક્ત છે:
જો $u \in [w]$ અને $u \in [v]$ હોય, તો $u \equiv w$ અને
$u \equiv v$; તેથી સંમિતિ અને પરંપરિતતાથી $w \equiv v$,
એટલે $[w] = [v]$.
% END OLP-0452
\endgroup

\begingroup
\makeatletter\def\ol@sectioncs{section}\makeatother

% BEGIN OLP-0453
\olfileid{nml}{fil}{fil}

\olsection{ફિલ્ટ્રેશનો}
\phantomsection\label{guunit:OLP-0453}


$\Gamma$ દ્વારા~$\mModel{M}$નું ``તે'' ફિલ્ટ્રેશન વ્યાખ્યાયિત
કરવાને બદલે, નિદર્શન~$\mModel{M^*}$ ક્યારે~$\mModel{M}$નું
ફિલ્ટ્રેશન ગણાય તે વ્યાખ્યાયિત કરીએ છીએ. બધાં ફિલ્ટ્રેશનોમાં
જગતોનો ગણ~$W^*$ અને નિયુક્તિ~$V^*$ સરખાં હોય છે. પરંતુ જુદાં
ફિલ્ટ્રેશનોના પ્રાપ્યતા સંબંધો~$R^*$ જુદા હોઈ શકે. ફિલ્ટ્રેશન
ગણાવા માટે $R^*$એ કેટલીક શરતો પૂરી કરવી જ પડે. આ શરતો મુખ્ય
પરિણામ સાબિત કરવા બરાબર જરૂરી છે: $!A \in \Gamma$ હોય ત્યારે
$\mSat{M}{!A}[w]$ જો અને માત્ર જો
$\mSat{M^*}{!A}[{[w]}]$.

\begin{defn}\ollabel{defn:filtration}
  ધારો કે $\Gamma$ ઉપસૂત્રો હેઠળ બંધ છે અને
  $\mModel{M} = \tuple{W, R, V}$. $\Gamma$ દ્વારા
  $\mModel{M}$નું \emph{ફિલ્ટ્રેશન} એવું કોઈપણ નિદર્શન
  $\mModel{M^*} = \tuple{W^*,R^*,V^*}$ છે, જ્યાં:
  \begin{enumerate}
  \item $W^* = \Setabs{[w]}{w \in W}$;
  \item \ollabel{defn:filtration-R}%
    કોઈપણ $u,v \in W$ માટે:
    \begin{enumerate}
    \item \ollabel{defn:filtration-R1}%
      જો $Ruv$ હોય તો $R^*[u][v]$;
    \item \ollabel{defn:filtration-R2}%
      જો $R^*[u][v]$ હોય, તો કોઈપણ $\Box !A \in \Gamma$ માટે,
      $\mSat{M}{\Box!A}[u]$ હોય ત્યારે $\mSat{M}{!A}[v]$;
    \item \ollabel{defn:filtration-R3}%
      જો $R^*[u][v]$ હોય, તો કોઈપણ $\Diamond !A \in \Gamma$
      માટે, $\mSat{M}{!A}[v]$ હોય ત્યારે
      $\mSat{M}{\Diamond!A}[u]$.
    \end{enumerate}
  \item $V^*(p) = \Setabs{[u]}{u \in V(p)}$.
  \end{enumerate}
\end{defn}

$V^*(p)$ શું છે તે વિચારવા જેવું છે:~$\mModel{M}$માં જ્યાં $p$
સત્ય હોય એવા બધાં જગતો~$w$ના સમકક્ષતા-વર્ગો~$[w]$નો ગણ.
એક તરફ, જો $w \in V(p)$ હોય, તો એ વ્યાખ્યાથી $[w] \in
V^*(p)$. જોકે $[w] \in V^*(p)$ હોય તો $w \in V(p)$ પણ હોય જ એવું
જરૂરી નથી. $[w] \in V^*(p)$ પરથી માત્ર \emph{કોઈ}~$u \in V(p)$
માટે $[w] = [u]$ મળે છે. અલબત્ત, $[w] = [u]$ એટલે
$w \equiv u$. તેથી $[w] \in V^*(p)$ હોય ત્યારે આપણે
માત્ર એટલું તારવી શકીએ કે કોઈ $u \in V(p)$ માટે $w \equiv u$.

\begin{thm}\ollabel{thm:filtrations}
  જો $\mModel{M^*}$ એ $\Gamma$ દ્વારા $\mModel{M}$નું ફિલ્ટ્રેશન
  હોય, તો દરેક $!A \in \Gamma$ અને $w \in W$ માટે
  $\mSat{M}{!A}[w]$ જો અને માત્ર જો
  $\mSat{M^*}{!A}[{[w]}]$.
\end{thm}

\begin{proof}
  $!A$ પર અનુમાનથી; $\Gamma$ ઉપસૂત્રો હેઠળ બંધ છે એ વાપરીશું.
  $!A \in \Gamma$ અને $\Gamma$ ઉપ-સૂત્રો હેઠળ બંધ હોવાથી
  $!A$નાં બધાં ઉપ-સૂત્રો પણ $\in \Gamma$ છે. તેથી
  અનુમાનના દરેક પગલામાં~$!A$નાં ઉપ-સૂત્રો માટે અનુમાનની
  પૂર્વધારણા વાપરી શકાય છે.
  \begin{enumerate}
    \item \indcase{!A}{\lfalse}{ન તો
        $\mSat{M}{\indfrm}[w]$ અને ન તો
        $\mSat{M^*}{\indfrm}[{[w]}]$.}
    \item \indcase{!A}{\ltrue}{બંને
        $\mSat{M}{\indfrm}[w]$ અને
        $\mSat{M^*}{\indfrm}[{[w]}]$.}
    \item \indcase{!A}{p}{ડાબેથી જમણી દિશા સીધી છે:
      $\mSat{M}{\indfrm}[w]$ હોય તો જ $w \in V(p)$, જેથી
      $[w] \in V^*(p)$, એટલે કે
      $\mSat{M^*}{\indfrm}[{[w]}]$. ઊલટું, ધારો કે
      $\mSat{M^*}{\indfrm}[{[w]}]$, એટલે કે $[w] \in
      V^*(p)$. ત્યારે કોઈ $v \in V(p)$ માટે $w \equiv v$.
      તેથી $\mSat{M}{p}[v]$ પણ. $w \equiv v$ હોવાથી $w$ અને
      $v$ એ~$\Gamma$માંથી સરખાં સૂત્રોને સત્ય બનાવે છે.
      પૂર્વધારણા મુજબ $p \in \Gamma$ અને $\mSat{M}{p}[v]$
      હોવાથી $\mSat{M}{\indfrm}[w]$.}

    \item \indcase{!A}{\lnot
          !B}{$\mSat{M}{\indfrm}[w]$ જો અને માત્ર જો
          $\mSat/{M}{!B}[w]$. અનુમાનની પૂર્વધારણાથી
          $\mSat/{M}{!B}[w]$ જો અને માત્ર જો
          $\mSat/{M^*}{!B}[{[w]}]$. છેલ્લે,
          $\mSat/{M^*}{!B}[{[w]}]$ જો અને માત્ર જો
          $\mSat{M^*}{\indfrm}[{[w]}]$.}

    \item \indcase{!A}{(!B \land
          !C)}{$\mSat{M}{\indfrm}[w]$ જો અને માત્ર જો
          $\mSat{M}{!B}[w]$ અને $\mSat{M}{!C}[w]$ બંને.
          અનુમાનની પૂર્વધારણાથી $\mSat{M}{!B}[w]$ જો અને માત્ર જો
          $\mSat{M^*}{!B}[{[w]}]$, અને $\mSat{M}{!C}[w]$ જો અને
          માત્ર જો $\mSat{M^*}{!C}[{[w]}]$. વળી,
          $\mSat{M^*}{\indfrm}[{[w]}]$ જો અને માત્ર જો
          $\mSat{M^*}{!B}[{[w]}]$ અને
          $\mSat{M^*}{!C}[{[w]}]$ બંને.}

    \item \indcase{!A}{(!B \lor
          !C)}{$\mSat{M}{\indfrm}[w]$ જો અને માત્ર જો
          $\mSat{M}{!B}[w]$ અથવા $\mSat{M}{!C}[w]$.
          અનુમાનની પૂર્વધારણાથી $\mSat{M}{!B}[w]$ જો અને માત્ર જો
          $\mSat{M^*}{!B}[{[w]}]$, અને $\mSat{M}{!C}[w]$ જો અને
          માત્ર જો $\mSat{M^*}{!C}[{[w]}]$. વળી,
          $\mSat{M^*}{\indfrm}[{[w]}]$ જો અને માત્ર જો
          $\mSat{M^*}{!B}[{[w]}]$ અથવા
          $\mSat{M^*}{!C}[{[w]}]$.}

    \item \indcase{!A}{(!B \lif
          !C)}{$\mSat{M}{\indfrm}[w]$ જો અને માત્ર જો
          $\mSat/{M}{!B}[w]$ અથવા $\mSat{M}{!C}[w]$.
          અનુમાનની પૂર્વધારણાથી $\mSat{M}{!B}[w]$ જો અને માત્ર જો
          $\mSat{M^*}{!B}[{[w]}]$, અને $\mSat{M}{!C}[w]$ જો અને
          માત્ર જો $\mSat{M^*}{!C}[{[w]}]$. વળી,
          $\mSat{M^*}{\indfrm}[{[w]}]$ જો અને માત્ર જો
          $\mSat/{M^*}{!B}[{[w]}]$ અથવા
          $\mSat{M^*}{!C}[{[w]}]$.}

    \item \indcase{!A}{(!B \liff
          !C)}{$\mSat{M}{\indfrm}[w]$ જો અને માત્ર જો
          $\mSat{M}{!B}[w]$ અને $\mSat{M}{!C}[w]$ બંને, અથવા
          $\mSat/{M}{!B}[w]$ અને $\mSat/{M}{!C}[w]$ બંને.
          અનુમાનની પૂર્વધારણાથી $\mSat{M}{!B}[w]$ જો અને માત્ર જો
          $\mSat{M^*}{!B}[{[w]}]$, અને $\mSat{M}{!C}[w]$ જો અને
          માત્ર જો $\mSat{M^*}{!C}[{[w]}]$. વળી,
          $\mSat{M^*}{\indfrm}[{[w]}]$ જો અને માત્ર જો
          $\mSat{M^*}{!B}[{[w]}]$ અને
          $\mSat{M^*}{!C}[{[w]}]$ બંને, અથવા
          $\mSat/{M^*}{!B}[{[w]}]$ અને
          $\mSat/{M^*}{!C}[{[w]}]$ બંને.}

    \item \indcase{!A}{\Box
          !B}{ધારો કે $\mSat{M}{\indfrm}[w]$.
          $\mSat{M^*}{\indfrm}[{[w]}]$ બતાવવા એવો $v$ લો કે
          $R^*[w][v]$. \olref{defn:filtration}%
          \olref{defn:filtration-R2} મુજબ $\mSat{M}{!B}[v]$;
          અનુમાનની પૂર્વધારણાથી
          $\mSat{M^*}{!B}[{[v]}]$. $v$ મનસ્વી હોવાથી
          $\mSat{M^*}{\indfrm}[{[w]}]$ ફલિત થાય છે.

          ઊલટું, ધારો કે $\mSat{M^*}{\indfrm}[{[w]}]$ અને
          $Rwv$ હોય એવો મનસ્વી $v$ લો. \olref{defn:filtration}%
          \olref{defn:filtration-R1} મુજબ $R^*[w][v]$,
          તેથી $\mSat{M^*}{!B}[{[v]}]$; અનુમાનની પૂર્વધારણાથી
          $\mSat{M}{!B}[v]$. $v$ મનસ્વી હોવાથી
          $\mSat{M}{\indfrm}[w]$.}

    \item \indcase{!A}{\Diamond
          !B}{ધારો કે $\mSat{M}{\indfrm}[w]$.
          ત્યારે કોઈ $v \in W$ માટે $Rwv$ અને
          $\mSat{M}{!B}[v]$. અનુમાનની પૂર્વધારણાથી
          $\mSat{M^*}{!B}[{[v]}]$; અને
          \olref{defn:filtration}\olref{defn:filtration-R1} મુજબ
          $R^*[w][v]$. તેથી
          $\mSat{M^*}{\indfrm}[{[w]}]$.

          હવે ધારો કે $\mSat{M^*}{\indfrm}[{[w]}]$.
          ત્યારે એવો કોઈ $[v] \in W^*$ છે કે $R^*[w][v]$ અને
          $\mSat{M^*}{!B}[{[v]}]$. અનુમાનની પૂર્વધારણાથી
          $\mSat{M}{!B}[v]$.
          \olref{defn:filtration}\olref{defn:filtration-R3} મુજબ
          $\mSat{M}{\indfrm}[w]$.}
  \end{enumerate}
\end{proof}

\begin{prob}
  \olref[nml][fil][fil]{thm:filtrations}ની સાબિતી પૂરી કરો.
\end{prob}

નિદર્શનના જગતો પરની સત્યતા માટે જે સાચું છે તે નિદર્શનમાં સત્યતા
અને નિદર્શનોના વર્ગમાં માન્યતા માટે પણ સાચું છે.

\begin{cor}
  ધારો કે $\Gamma$ ઉપસૂત્રો હેઠળ બંધ છે. ત્યારે:
  \begin{enumerate}
  \item જો $\mModel{M^*}$ એ $\Gamma$ દ્વારા $\mModel{M}$નું
    ફિલ્ટ્રેશન હોય, તો કોઈપણ $!A \in \Gamma$ માટે:
    $\mSat{M}{!A}$ જો અને માત્ર જો $\mSat{M^*}{!A}$.
  \item જો $\mClass{C}$ નિદર્શનોનો વર્ગ હોય અને
    $\Gamma(\mClass{C})$ એ $\mClass{C}$માંનાં નિદર્શનોના
    $\Gamma$-ફિલ્ટ્રેશનોનો વર્ગ હોય, તો કોઈપણ
    સૂત્ર~$!A \in \Gamma$ એ $\mClass{C}$માં માન્ય છે
    જો અને માત્ર જો તે $\Gamma(\mClass{C})$માં માન્ય છે.
  \end{enumerate}
\end{cor}
% END OLP-0453
\endgroup

\begingroup
\makeatletter\def\ol@sectioncs{section}\makeatother

% BEGIN OLP-0454
\olfileid{nml}{fil}{exf}

\olsection{ફિલ્ટ્રેશનનાં ઉદાહરણો}
\phantomsection\label{guunit:OLP-0454}


ફિલ્ટ્રેશનો અસ્તિત્વમાં છે એવું હજી આપણે બતાવ્યું નથી. પરંતુ હકીકતમાં
દરેક નિદર્શન $\mModel{M}$ માટે $\Gamma$ દ્વારા~$\mModel{M}$નાં
ઘણાં ફિલ્ટ્રેશનો છે. તેમાંનાં બે ખાસ ઓળખીએ: સૌથી સૂક્ષ્મ અને સૌથી
સ્થૂલ ફિલ્ટ્રેશન. એ જ નિદર્શનનાં ફિલ્ટ્રેશનોનો પ્રાપ્યતા સંબંધ જુદો
હોય છે (કારણ કે \olref[fil]{defn:filtration} સીધું જ $W^*$ અને
~$V^*$ નક્કી કરે છે). સૌથી સૂક્ષ્મ ફિલ્ટ્રેશનમાં શક્ય તેટલાં ઓછાં
જગતો સંબંધિત હશે; સૌથી સ્થૂલમાં શક્ય તેટલાં વધુ.

\begin{defn}
  $\Gamma$ ઉપસૂત્રો હેઠળ બંધ હોય ત્યારે નિદર્શન $\mModel{M}$નું
  \emph{સૌથી સૂક્ષ્મ} ફિલ્ટ્રેશન $\mModel{M^*}$ આ રીતે વ્યાખ્યાયિત
  થાય છે:
  \[
  R^*[u][v] \quad \text{જો અને માત્ર જો} \quad \exists u'\in [u] \;
  \exists v' \in [v] : Ru'v'.
  \]
\end{defn}

\begin{prop}\ollabel{prop:finest}
  સૌથી સૂક્ષ્મ ફિલ્ટ્રેશન $\mModel{M^*}$ ખરેખર ફિલ્ટ્રેશન છે.
\end{prop}

\begin{proof}
  આ રીતે વ્યાખ્યાયિત $R^*$ એ
  \olref[fil]{defn:filtration}\olref[fil]{defn:filtration-R}ની
  શરતો પૂરી કરે છે તે ચકાસવું પડશે. ત્રણ શરતો ક્રમસર ચકાસીએ.

  જો $Ruv$ હોય તો $u \in [u]$ અને $v \in [v]$ હોવાથી
  $R^*[u][v]$ પણ; એટલે \olref[fil]{defn:filtration-R1} પૂરી થાય છે.

  %
      \olref[fil]{defn:filtration-R2} માટે $\Box!A \in \Gamma$,
      $R^*[u][v]$ અને $\mSat{M}{\Box!A}[u]$ ધારો.
      $R^*$ની વ્યાખ્યાથી એવાં $u' \equiv u$ અને
      $v' \equiv v$ છે કે $Ru'v'$. $u$ અને $u'$ એ $\Gamma$ પર
      સહમત હોવાથી $\mSat{M}{\Box!A}[u']$ પણ; તેથી
      $\mSat{M}{!A}[v']$. $\Gamma$ ઉપ-સૂત્રો હેઠળ બંધ હોવાથી
      $v$ અને $v'$ એ $!A$ પર સહમત છે; એટલે જરૂરી મુજબ
      $\mSat{M}{!A}[v]$.

  %
      \olref[fil]{defn:filtration-R3} ચકાસવા $\Diamond!A \in
      \Gamma$, $R^*[u][v]$ અને $\mSat{M}{!A}[v]$ ધારો.
      $R^*$ની વ્યાખ્યાથી એવાં $u' \equiv u$ અને
      $v' \equiv v$ છે કે $Ru'v'$. $v$ અને $v'$ એ~$\Gamma$ પર
      સહમત છે, અને $\Gamma$ ઉપ-સૂત્રો હેઠળ બંધ છે;
      તેથી $\mSat{M}{!A}[v']$ પણ અને પછી
      $\mSat{M}{\Diamond !A}[u']$. $u$ અને $u'$ પણ $\Gamma$ પર
      સહમત હોવાથી $\mSat{M}{\Diamond !A}[u]$.
\end{proof}


\begin{probtag}{probBox,probDiamond}
  \olref[nml][fil][exf]{prop:finest}ની સાબિતી પૂરી કરો.
\end{probtag}

\begin{defn}
  $\Gamma$ ઉપસૂત્રો હેઠળ બંધ હોય ત્યારે નિદર્શન~$\mModel{M}$નું
  \emph{સૌથી સ્થૂલ} ફિલ્ટ્રેશન~$\mModel{M^*}$માં
  $R^*[u][v]$ જો અને માત્ર જો
  નીચેની
    \emph{બંને} શરતો પૂરી થાય:
  \begin{enumerate}
  \item \ollabel{defn:coarsest-Box}જો $\Box!A \in \Gamma$
    અને $\mSat{M}{\Box!A}[u]$ હોય, તો
    $\mSat{M}{!A}[v]$;
  \item \ollabel{defn:coarsest-Diamond}જો $\Diamond!A
    \in \Gamma$ અને $\mSat{M}{!A}[v]$ હોય, તો
    $\mSat{M}{\Diamond!A}[u]$.
  \end{enumerate}
\end{defn}

\begin{prop}
  સૌથી સ્થૂલ ફિલ્ટ્રેશન $\mModel{M^*}$ ખરેખર ફિલ્ટ્રેશન છે.
\end{prop}

\begin{proof}
  $R^*$ની વ્યાખ્યા મુજબ માત્ર $Ruv$માંથી $R^*[u][v]$ મેળવવાની
  શરત ચકાસવાની બાકી છે. તેથી $Ruv$ ધારો.
  ધારો કે $\Box!A \in \Gamma$ અને
    $\mSat{M}{\Box!A}[u]$; ત્યારે સ્પષ્ટ રીતે
    $\mSat{M}{!A}[v]$, અને
      \olref{defn:coarsest-Box} પૂરી થાય છે. %
  ધારો કે $\Diamond!A \in
    \Gamma$ અને $\mSat{M}{!A}[v]$. $Ruv$ હોવાથી
    $\mSat{M}{\Diamond!A}[u]$, અને
      \olref{defn:coarsest-Diamond} પૂરી થાય છે.
\end{proof}

\begin{ex}
  ધારો કે $W = \PosInt$, $Rnm$ જો અને માત્ર જો $m = n + 1$,
  અને $V(p) = \Setabs{2n}{n \in \PosInt}$.
  \sourcecorrection{OLMD-044}{મૂળમાં $n\in\Nat$ હોવાથી મૂલ્યાંકનમાં
  $0$ પણ આવે છે, જ્યારે જગત-ગણ ધન પૂર્ણાંકોનો છે; $n\in\PosInt$
  રાખીને મૂલ્યાંકનને યોગ્ય જગતો સુધી મર્યાદિત કર્યું છે.}
  નિદર્શન $\mModel{M} = \tuple{W, R, V}$
  \olref{fig:ex-filtration}માં દર્શાવ્યું છે. જગતો $1$, $2$,
  વગેરે છે; દરેક જગતમાંથી બરાબર એક અન્ય જગત—તેનો અનુગામી—પ્રાપ્ય
  છે; અને $p$ બધાં અને માત્ર સમ સંખ્યાઓ પર સત્ય છે.

  \begin{figure}
    \centering
    \begin{tikzpicture}[modal]
      \node[world] (1) [label=below:\mFalse{p}] {$1$};
      \node[world] (2) [label=below:\mTrue{p}, right=of 1] {$2$};
      \node[world] (3) [label=below:\mFalse{p}, right=of 2] {$3$};
      \node[world] (4) [label=below:\mTrue{p}, right=of 3] {$4$};
      \node[phantom] (5) [right=of 4] {};
      \draw[->] (1) to (2);
      \draw[->] (2) to (3);
      \draw[->] (3) to (4);
      \draw[dotted] (4) to (5);
    \end{tikzpicture}

    \begin{tikzpicture}[modal]
      \node[world] (1) [label=below:\mFalse{p}] {$[1]$};
      \node[world] (2) [label=below:\mTrue{p}, right=of 1] {$[2]$};
      \draw[->,bend left] (1) to (2);
      \draw[->,bend left] (2) to (1);
      
      \node[world] (3) [label=below:\mFalse{p},right=of 2] {$[1]$};
      \node[world] (4) [label=below:\mTrue{p}, right=of 3] {$[2]$};
      \draw[->,bend left] (3) to (4);
      \draw[->,bend left] (4) to (3);
      \draw[->,reflexive above] (4) to (4);
    \end{tikzpicture}
    \caption{અનંત નિદર્શન અને તેનાં ફિલ્ટ્રેશનો.}
    \ollabel{fig:ex-filtration}
  \end{figure}
  
  હવે $\Gamma$ એ~$\Box p \lif p$નાં ઉપ-સૂત્રોનો ગણ,
  એટલે કે $\{p, \Box p, \Box p \lif p\}$, માનો. $p$ બધાં અને
  માત્ર સમ સંખ્યાઓ પર સત્ય છે; $\Box p$ બધાં અને માત્ર વિષમ
  સંખ્યાઓ પર સત્ય છે; તેથી $\Box p \lif p$ બધાં અને માત્ર સમ
  સંખ્યાઓ પર સત્ય છે. બીજા શબ્દોમાં, દરેક વિષમ સંખ્યા $\Box p$ને
  સત્ય, પરંતુ $p$ અને $\Box p \lif p$ને અસત્ય બનાવે છે; દરેક સમ
  સંખ્યા $p$ અને $\Box p \lif p$ને સત્ય, પરંતુ $\Box p$ને
  અસત્ય બનાવે છે. આમ $W^* = \{ [1], [2] \}$, જ્યાં
  $[1] = \{1, 3, 5, \dots\}$ અને
  $[2] = \{2, 4, 6, \dots\}$. $2 \in V(p)$ હોવાથી
  $[2] \in V^*(p)$; $1 \notin V(p)$ હોવાથી
  $[1] \notin V^*(p)$. તેથી $V^*(p) = \{[2]\}$.

  $W^*$ પર આધારિત કોઈપણ ફિલ્ટ્રેશનના પ્રાપ્યતા સંબંધમાં
  $\tuple{[1], [2]}, \tuple{[2],[1]}$ હોવાં જોઈએ: $R12$
  હોવાથી \olref[fil]{defn:filtration}%
  \olref[fil]{defn:filtration-R1} મુજબ $R^*[1][2]$ હોવું જોઈએ;
  અને $R23$ હોવાથી $R^*[2][3]$ હોવું જોઈએ, જ્યારે $[3]=[1]$.
  તેમાં $\tuple{[1],[1]}$ હોઈ શકે નહિ: એવું હોય તો
  $R^*[1][1]$, $\mSat{M}{\Box p}[1]$ પરંતુ
  $\mSat/{M}{p}[1]$ મળે, જે
  \olref[fil]{defn:filtration-R2}નો વિરોધ કરે છે.
  $R^*[2][2]$ હોવું જરૂરી પણ નથી અને પ્રતિબંધિત પણ નથી.
  તેથી~$\mModel{M}$નાં બે શક્ય ફિલ્ટ્રેશનો છે, જે નીચેના બે
  પ્રાપ્યતા સંબંધોને અનુરૂપ છે:
  \[
  \{\tuple{[1],[2]}, \tuple{[2],[1]}\} \text{ અને }
  \{\tuple{[1],[2]}, \tuple{[2],[1]}, \tuple{[2],[2]}\}.
  \]
  બંને કિસ્સામાં~$[1]$ પર $p$ અને $\Box p \lif p$ અસત્ય, જ્યારે
  $\Box p$ સત્ય છે;~$[2]$ પર $p$ અને $\Box p \lif p$ સત્ય,
  જ્યારે $\Box p$ અસત્ય છે.
\end{ex}

\begin{prob}
  નીચેનું નિદર્શન $\mModel{M} = \tuple{W, R, V}$ ધ્યાનમાં લો. અહીં
  $W = \Setabs{0\sigma}{\sigma \in \Bin^*}$, એટલે $0$થી શરૂ થતી
  $0$ અને~$1$ની શ્રેણીઓનો ગણ; $R\sigma\sigma'$ જો અને માત્ર જો
  $\sigma' = \sigma 0$ અથવા $\sigma' = \sigma 1$; અને
  $V(p) = \Setabs{\sigma 0}{\sigma \in \Bin^*}\cap W$ તથા
  $V(q) = \Setabs{\sigma 1}{\sigma \in
    \Bin^* \setminus \{1\}}\cap W$.
  મૂળમાં આપેલા મૂલ્યાંકન-ગણોને જગત-ગણ સાથે છેદવાથી તેઓ નિદર્શનના
  જગતોના ઉપગણ બને છે.\sourcecorrection{OLMD-045}{મૂળમાં $V(p)$ અને
  $V(q)$ને $W$ની બહારની શ્રેણીઓ પણ સમાવી શકે એવા ગણો તરીકે લખ્યાં છે;
  મૂલ્યાંકનને $W$ના ઉપગણો સુધી મર્યાદિત કરવા બંનેમાં $\cap W$ ઉમેર્યું છે.}
  તેનું ચિત્ર આ રહ્યું:
  \begin{center}
    \begin{tikzpicture}[modal,
        node distance=2cm
      ]

      \node[world] (0) [label={below:\mTrue{p}\\\mFalse{q}},font=\tiny] {$0$};
      \node[world] (00) [label={below:\mTrue{p}\\\mFalse{q}},
        above right=of 0,font=\tiny] {$00$};
      \node[world] (000) [label={below:\mTrue{p}\\\mFalse{q}},
        right=of 00, yshift=.8cm, font=\tiny] {$000$};
      \node[phantom] (0000) [right=of 000, yshift=.5cm] {};
      \node[phantom] (0001) [right=of 000, yshift=-.5cm] {};
      \node[world] (001) [label={below:\mFalse{p}\\\mTrue{q}},
        right=of 00, yshift=-.8cm, font=\tiny] {$001$};
      \node[phantom](0010) [right=of 001, yshift=.5cm] {};
      \node[phantom](0011) [right=of 001, yshift=-.5cm] {};
      \node[world] (01) [label={below:\mFalse{p}\\\mTrue{q}},
        below right=of 0,font=\tiny] {$01$};
      \node[world] (010) [label={below:\mTrue{p}\\\mFalse{q}},
        right=of 01, yshift=.8cm,font=\tiny] {$010$};
      \node[phantom] (0100) [right=of 010, yshift=.5cm] {};
      \node[phantom] (0101) [right=of 010, yshift=-.5cm] {};
      \node[world] (011) [label={below:\mFalse{p}\\\mTrue{q}},
        right=of 01, yshift=-.8cm,font=\tiny] {$011$};
      \node[phantom] (0110) [right=of 011, yshift=.5cm] {};
      \node[phantom] (0111) [right=of 011, yshift=-.5cm] {};

      \draw[->] (0) to (00);
      \draw[->] (0) to (01);
      \draw[->] (00) to (000);
      \draw[dotted] (000) to (0000);
      \draw[dotted] (000) to (0001);
      \draw[->] (00) to (001);
      \draw[dotted] (001) to (0010);
      \draw[dotted] (001) to (0011);
      \draw[->] (01) to (010);
      \draw[dotted] (010) to (0100);
      \draw[dotted] (010) to (0101);
      \draw[->] (01) to (011);
      \draw[dotted] (011) to (0110);
      \draw[dotted] (011) to (0111);
    \end{tikzpicture}
  \end{center}
  દરેક~$w$ માટે $\mSat/{M}{\Box(p \lor q) \lif (\Box p \lor \Box q)}[w]$
  છે.

  $\Gamma$ને~$\Box(p \lor q) \lif (\Box p \lor \Box q)$નાં
  ઉપ-સૂત્રોનો ગણ માનો. $W^*$ અને~$V^*$ શું છે?
  $\mModel{M}$ના સૌથી સૂક્ષ્મ ફિલ્ટ્રેશનનો પ્રાપ્યતા સંબંધ કયો છે?
  સૌથી સ્થૂલ ફિલ્ટ્રેશનનો કયો છે?
\end{prob}
% END OLP-0454
\endgroup

\begingroup
\makeatletter\def\ol@sectioncs{section}\makeatother

% BEGIN OLP-0455
\olfileid{nml}{fil}{fin}

\olsection{ફિલ્ટ્રેશનો સાંત છે}
\phantomsection\label{guunit:OLP-0455}


ઉપ-સૂત્રો હેઠળ બંધ એવા કોઈપણ ગણ~$\Gamma$ માટે ફિલ્ટ્રેશન
વ્યાખ્યાયિત કર્યું છે. વ્યાખ્યા પોતે ફિલ્ટ્રેશન સાંત હશે એની ખાતરી
આપતી નથી. ખરેખર, $\Gamma$ અનંત હોય (દાખલા તરીકે, બધાં
સૂત્રોનો ગણ હોય), તો ફિલ્ટ્રેશન પણ અનંત હોઈ શકે. પરંતુ
$\Gamma$ સાંત હોય (દાખલા તરીકે, આપેલા સૂત્ર~$!A$નાં
ઉપ-સૂત્રોનો ગણ હોય), તો~$\Gamma$ દ્વારા કોઈપણ ફિલ્ટ્રેશન
પણ સાંત હોય છે.

\begin{prop}\ollabel{prop:filt-are-finite}
  જો $\Gamma$ સાંત હોય, તો નિદર્શન $\mModel{M}$નું $\Gamma$ દ્વારા
  કોઈપણ ફિલ્ટ્રેશન $\mModel{M^*}$ પણ સાંત છે.
\end{prop}

\begin{proof}
  $W^*$નું કદ એ સમકક્ષતા સંબંધ~$\equiv$ હેઠળના જુદા વર્ગો~$[w]$ની
  સંખ્યા છે. આવા વર્ગનાં કોઈપણ બે જગત $u$, $v$—એટલે કે એવાં
  $u$ અને $v$ કે~$u \equiv v$—$\Gamma$નાં બધાં
  સૂત્રો~$!A$ પર સહમત હોય છે: $!A \in \Gamma$ હોય તો $!A$
  બંને $u$ અને $v$ પર $!A$ સત્ય હોય અથવા બંને પર અસત્ય હોય. તેથી દરેક
  વર્ગ~$[w]$ને~$\Gamma$નો એક ઉપગણ અનુરૂપ છે:~$[w]$નાં જગતો પર
  સત્ય હોય એવાં બધાં $!A \in \Gamma$નો ગણ. બે જુદા વર્ગ
  $[u]$ અને $[v]$ એ~$\Gamma$ના એ જ ઉપગણને અનુરૂપ હોઈ શકે નહિ. કારણ કે જો
  $u$ પર સત્ય સૂત્રોનો ગણ અને $v$ પર સત્ય સૂત્રોનો
  ગણ એક જ હોય, તો $u$ અને $v$ એ~$\Gamma$નાં બધાં સૂત્રો પર
  સહમત હોય, એટલે $u \equiv v$. પરંતુ ત્યારે $[u] = [v]$.
  આમ $W^*$માંથી $\Pow{\Gamma}$માં એક એક-એક વિધેય છે અને
  તેથી $\card{W^*} \le \card{\Pow{\Gamma}}$. જો $\Gamma$માં $n$
  વાક્યો હોય, તો $W^*$નું કદ વધુમાં વધુ~$2^n$ છે.
\end{proof}
% END OLP-0455
\endgroup

\begingroup
\makeatletter\def\ol@sectioncs{section}\makeatother

% BEGIN OLP-0456
\olfileid{nml}{fil}{fmp}

\olsection{\Log{K} અને \Log{S5}ને સાંત નિદર્શન ગુણધર્મ છે}
\phantomsection\label{guunit:OLP-0456}


\begin{defn}
  જો~$\Sigma$ના કોઈ નિદર્શનનાં એક જગત પર કોઈપણ સૂત્ર~$!A$ સત્ય હોય
  ત્યારે~$\Sigma$ના કોઈ \emph{સાંત} નિદર્શનનાં એક જગત પર પણ $!A$
  સત્ય હોય, તો નિદર્શ્ય તર્કની પદ્ધતિ~$\Sigma$ને \emph{સાંત
    નિદર્શન ગુણધર્મ} છે એમ કહેવાય.
\end{defn}

\begin{prop}\ollabel{prop:K-fmp}
  \Log{K}ને સાંત નિદર્શન ગુણધર્મ છે.
\end{prop}

\begin{proof}
  \Log{K} એ માન્ય સૂત્રોનો ગણ છે, એટલે દરેક નિદર્શન
  એ~\Log{K}નું નિદર્શન છે. \olref[fil]{thm:filtrations} મુજબ,
  જો $\mSat{M}{!A}[w]$ હોય, તો~$!A$નાં ઉપ-સૂત્રોના ગણ
  $\Gamma$ દ્વારા~$\mModel{M}$ના કોઈપણ ફિલ્ટ્રેશન માટે
  $\mSat{M^*}{!A}[{[w]}]$ હોય છે.\sourcecorrection{OLMD-046}{મૂળમાં
  ફિલ્ટ્રેશનની સત્યતા મૂળ જગત $w$ પર લખી છે; ફિલ્ટ્રેશનનું જગત
  તેનો સમકક્ષતા-વર્ગ $[w]$ હોવાથી અનુરૂપ સૂચક બદલ્યો છે.}
  કોઈપણ સૂત્રનાં સાંત સંખ્યામાં જ ઉપ-સૂત્રો હોય છે;
  તેથી $\Gamma$ સાંત છે. \olref[fin]{prop:filt-are-finite} મુજબ
  $\card{W^*} \le 2^n$, જ્યાં $n$ એ~$\Gamma$નાં
  સૂત્રોની સંખ્યા છે. અને \Log{K} નિદર્શનો પર કોઈ નિયંત્રણ
  મૂકતું ન હોવાથી $\mModel{M^*}$ પણ \Log{K}-નિદર્શન છે.
\end{proof}

ફિલ્ટ્રેશન દ્વારા તર્કશાસ્ત્ર~\Log{L}ને સાંત નિદર્શન ગુણધર્મ છે એવું
બતાવવા, \Log{L}-નિદર્શનનું ફિલ્ટ્રેશન પોતે પણ \Log{L}-નિદર્શન
હોવું આવશ્યક છે. ઘણી વાર આ માટે નોંધપાત્ર કાર્ય કરવું પડે છે, અને
દરેક ફિલ્ટ્રેશન \Log{L}-નિદર્શન આપતું નથી. જોકે સાર્વત્રિક
નિદર્શનો માટે આ ગુણધર્મ જળવાય છે.

\begin{prop}\ollabel{prop:univ-fin}
  $\mClass{U}$ને સાર્વત્રિક નિદર્શનોનો વર્ગ માનો
  (જુઓ \olref[frd][es5]{prop:S5=univ}) અને
  $\mClass{U}_\mathrm{Fin}$ને બધાં સાંત સાર્વત્રિક નિદર્શનોનો વર્ગ
  માનો. ત્યારે કોઈપણ સૂત્ર~$!A$ એ $\mClass{U}$માં માન્ય છે
  જો અને માત્ર જો તે $\mClass{U}_\mathrm{Fin}$માં માન્ય છે.
\end{prop}

\begin{proof}
  સાંત સાર્વત્રિક નિદર્શનો પણ સાર્વત્રિક નિદર્શનો છે; તેથી
  ડાબેથી જમણી દિશા સ્પષ્ટ છે. જમણેથી ડાબી દિશા માટે ધારો કે કોઈ
  સાર્વત્રિક નિદર્શન $\mModel{M}$ના જગત $w$ પર~$!A$ અસત્ય છે.
  $\Gamma$માં $!A$ તથા તેનાં બધાં ઉપસૂત્રો લો; સ્પષ્ટ છે કે
  $\Gamma$ સાંત છે. $\mModel{M}$નું એક ફિલ્ટ્રેશન
  $\mModel{M^*}$ લો. \olref[fin]{prop:filt-are-finite} મુજબ તે
  $\mModel{M^*}$ સાંત છે અને \olref[fil]{thm:filtrations} મુજબ~$\mModel{M^*}$માં
  $[w]$ પર $!A$ અસત્ય છે. હવે $\mModel{M^*}$ પણ સાર્વત્રિક છે
  તે જોવાનું રહે છે: કોઈપણ $u$ અને $v$ માટે પૂર્વધારણાથી $Ruv$,
  અને \olref[fil]{defn:filtration}%
  \olref[fil]{defn:filtration-R} મુજબ $R^*[u][v]$ પણ.
\end{proof}

\begin{cor}\ollabel{cor:S5fmp}
  \Log{S5}ને સાંત નિદર્શન ગુણધર્મ છે.
\end{cor}

\begin{proof}
  \olref[frd][es5]{prop:S5=univ} મુજબ, જો $!A$ કોઈ સ્વવાચક
  અને યુક્લિડિયન નિદર્શનનાં એક જગત પર સત્ય હોય, તો તે કોઈ
  સાર્વત્રિક નિદર્શનનાં એક જગત પર સત્ય છે.
  \olref{prop:univ-fin} મુજબ તે કોઈ સાંત સાર્વત્રિક નિદર્શનનાં એક
  જગત પર સત્ય છે (એટલે કે~$!A$નાં ઉપ-સૂત્રોના ગણ દ્વારા
  તે નિદર્શનના ફિલ્ટ્રેશનમાં). દરેક સાર્વત્રિક નિદર્શન
  સ્વવાચક અને યુક્લિડિયન પણ છે; તેથી $!A$ સાંત,
  સ્વવાચક તથા યુક્લિડિયન નિદર્શનનાં એક જગત પર સત્ય છે.
\end{proof}

\begin{prob}
  ક્રમશઃ સિરિયલ અથવા સ્વવાચક નિદર્શનનું કોઈપણ
  ફિલ્ટ્રેશન પણ સિરિયલ અથવા સ્વવાચક હોય છે તે બતાવો.
\end{prob}

\begin{prob}
  સંમિત (અથવા પરંપરિત, અથવા યુક્લિડિયન) નિદર્શનનું એવું
  ફિલ્ટ્રેશન શોધો જે સંમિત નથી (અથવા પરંપરિત નથી, અથવા
  યુક્લિડિયન નથી).
\end{prob}
% END OLP-0456
\endgroup

\begingroup
\makeatletter\def\ol@sectioncs{section}\makeatother

% BEGIN OLP-0457
\olfileid{nml}{fil}{dec}

\olsection{\Log{S5} નિર્ણેય છે}
\phantomsection\label{guunit:OLP-0457}

સાંત નિદર્શન ગુણધર્મથી સૂત્ર-ઢાંચાઓ દ્વારા અપાયેલી નિદર્શ્ય-તર્ક
પદ્ધતિઓ \emph{નિર્ણેય} છે તે બતાવવાનો સરળ માર્ગ મળે છે (એટલે કે
કોઈ સૂત્ર પદ્ધતિમાં નિષ્પન્ન કરી શકાય એવું છે કે નહિ તે નક્કી કરવાની
ગણનક્ષમ પદ્ધતિ છે).

\begin{thm}
  \Log{S5} નિર્ણેય છે.
\end{thm}

\begin{proof}
  $!A$ આપેલું હોય અને $!A$માં આવતાં પ્રતિજ્ઞાપન ચલો
  $p_1$, \dots, $p_k$માંનાં હોય એમ ધારો. દરેક $n$ માટે $n$ જગતવાળાં
  અને $p_1$, \dots, $p_k$ને સત્યમૂલ્ય આપતાં સાંત સંખ્યામાં જ નિદર્શનો
  હોવાથી આપણે બે પ્રક્રિયાઓ \emph{સમાંતરે} ચલાવી શકીએ: કોઈ વ્યવસ્થિત
  રીતે પ્રમાણો રચીને \Log{S5}નાં બધાં પ્રમેયોની ગણના કરવી; અને
  $1$, $2$, \dots જગતવાળાં બધાં સાર્વત્રિક નિદર્શનોની ગણના કરીને
  તેમાંનાં કોઈ જગત પર $!A$ અસત્ય છે કે નહિ તે ચકાસવું.
  \sourcecorrection{OLMD-047}{મૂળમાં બીજી પ્રક્રિયા માટે બધાં સાંત
  નિદર્શનોની ગણના લખી છે; તે S5નું નિર્ણેયન આપતું નથી, કારણ કે
  અસર્વત્રિક નિદર્શન S5-પ્રમેયને પણ ખોટું પાડી શકે. S5ના
  સાર્વત્રિક નિદર્શનો સુધી ગણના મર્યાદિત કરી છે.}
  અંતે બેમાંથી એક પ્રક્રિયા ઉત્તર આપશે: કારણ કે
  \olref[com][fra]{thm:generaldet} અને \olref[fmp]{cor:S5fmp} મુજબ,
  કાં તો $!A$ નિષ્પન્ન કરી શકાય એવું છે, અથવા તે કોઈ સાંત સાર્વત્રિક
  નિદર્શનમાં અસત્ય છે.
\end{proof}

ઉપરની સાબિતી \Log{S5} માટે કામ કરે છે કારણ કે સાર્વત્રિક
નિદર્શનોનાં ફિલ્ટ્રેશનો આપમેળે સાર્વત્રિક હોય છે. સ્વવાચકતા અને
સિરિયલતા માટે પણ એવું જ છે, પરંતુ અન્ય ગુણધર્મો માટે વધુ કામ
જરૂરી છે.
% END OLP-0457
\endgroup

\begingroup
\makeatletter\def\ol@sectioncs{section}\makeatother

% BEGIN OLP-0458
\olfileid{nml}{fil}{acc}

\olsection{ફિલ્ટ્રેશનો અને પ્રાપ્યતાના ગુણધર્મો}
\phantomsection\label{guunit:OLP-0458}


જણાવ્યા મુજબ, સાર્વત્રિક, સિરિયલ અને સ્વવાચક નિદર્શનોનાં
ફિલ્ટ્રેશનો પણ હંમેશા અનુક્રમે સાર્વત્રિક, સિરિયલ અથવા સ્વવાચક
હોય છે. પરંતુ સંમિત અથવા પરંપરિત નિદર્શનનું દરેક ફિલ્ટ્રેશન
અનુક્રમે સંમિત અથવા પરંપરિત હોતું નથી. છતાં, કેટલાક કિસ્સામાં
આ ગુણધર્મ જળવાય તે રીતે ફિલ્ટ્રેશન વ્યાખ્યાયિત કરી શકાય છે.
એ માટે આપણે સૌથી સ્થૂલ ફિલ્ટ્રેશનની વ્યાખ્યા જેવી રીત લઈએ છીએ,
પણ~$R^*$ની વ્યાખ્યામાં વધારાની શરતો ઉમેરીએ છીએ. $\Gamma$
ઉપ-સૂત્રો હેઠળ બંધ માનો. નિદર્શન
$\mModel{M} =\tuple{W, R, V}$નાં જગતો $u$, $v$ વચ્ચેના
\olref{tab:Cn-filtrations}માંના સંબંધો~$C_i(u,v)$ જુઓ.
આ શરતોના સંયોજનો પરથી $R^*[u][v]$ વ્યાખ્યાયિત કરી શકાય છે.
દાખલા તરીકે, $R^*[u][v]$ જો અને માત્ર જો $C_1(u,v)$ હોય એવી
શરત રાખીએ, તો બરાબર સૌથી સ્થૂલ ફિલ્ટ્રેશન મળે છે. જો
$R^*[u][v]$ જો અને માત્ર જો $C_1(u,v)$ અને $C_2(u, v)$ બંને
હોય એમ રાખીએ, તો જુદું ફિલ્ટ્રેશન મળે છે. તે સૌથી સ્થૂલ
ફિલ્ટ્રેશન કરતાં ``વધુ સૂક્ષ્મ'' છે, કારણ કે એકલા $C_1(u,v)$
કરતાં $C_1(u,v)$ અને $C_2(u,v)$ બંનેને ઓછાં જગત-યુગ્મો
સંતોષે છે.

\begin{table}[ht]
  \centering
  \begin{tabular}{|ll|}
    \hline
    \multirow{2}{*}{$C_1(u,v)$:}
    &
      જો $\Box!A \in \Gamma$ અને $\mSat{M}{\Box!A}[u]$ હોય,
      તો $\mSat{M}{!A}[v]$; અને\\
    &
      જો $\Diamond!A \in \Gamma$ અને $\mSat{M}{!A}[v]$ હોય,
      તો $\mSat{M}{\Diamond!A}[u]$; \\
    \hline
    \multirow{2}{*}{$C_2(u,v)$:}
    &
      જો $\Box!A \in \Gamma$ અને $\mSat{M}{\Box!A}[v]$ હોય,
      તો $\mSat{M}{!A}[u]$; અને\\
    &
      જો $\Diamond!A \in \Gamma$ અને $\mSat{M}{!A}[u]$ હોય,
      તો $\mSat{M}{\Diamond!A}[v]$; \\
    \hline
    \multirow{2}{*}{$C_3(u,v)$:}
    &
      જો $\Box!A \in \Gamma$ અને $\mSat{M}{\Box!A}[u]$ હોય,
      તો $\mSat{M}{\Box!A}[v]$; અને\\
    &
      જો $\Diamond!A \in \Gamma$ અને $\mSat{M}{\Diamond!A}[v]$ હોય,
      તો $\mSat{M}{\Diamond!A}[u]$; \\
    \hline
    \multirow{2}{*}{$C_4(u,v)$:}
    &
      જો $\Box!A \in \Gamma$ અને $\mSat{M}{\Box!A}[v]$ હોય,
      તો $\mSat{M}{\Box!A}[u]$; અને\\
    &
      જો $\Diamond!A \in \Gamma$ અને $\mSat{M}{\Diamond!A}[u]$ હોય,
      તો $\mSat{M}{\Diamond!A}[v]$; \\
    \hline
    \end{tabular}
  \caption{ફિલ્ટ્રેશનો વ્યાખ્યાયિત કરવા માટે જગતો પરની શરતો.}
  \ollabel{tab:Cn-filtrations}
\end{table}

\begin{thm}\ollabel{thm:more-filtrations}
  $\mModel{M} =\tuple{W,R,V}$ નિદર્શન હોય, $\Gamma$ ઉપ-સૂત્રો
  હેઠળ બંધ હોય અને $W^*$ તથા $V^*$ની વ્યાખ્યા
  \olref[fil]{defn:filtration} મુજબ હોય. ત્યારે:
  \begin{enumerate}
  \item ધારો કે $R^*[u][v]$ જો અને માત્ર જો $C_1(u, v) \land
    C_2(u,v)$. ત્યારે $R^*$ સંમિત છે; અને $\mModel{M}$ સંમિત હોય
    તો $\mModel{M^*} = \tuple{W^*,R^*,V^*}$ ફિલ્ટ્રેશન છે.
  \item ધારો કે $R^*[u][v]$ જો અને માત્ર જો $C_1(u, v)
    \land C_3(u,v)$. ત્યારે $R^*$ પરંપરિત છે; અને
    $\mModel{M}$ પરંપરિત હોય તો
    $\mModel{M^*}=\tuple{W^*,R^*,V^*}$ ફિલ્ટ્રેશન છે.
  \item ધારો કે $R^*[u][v]$ જો અને માત્ર જો $C_1(u, v) \land C_2(u,v)
    \land C_3(u,v) \land C_4(u,v)$. ત્યારે $R^*$ સંમિત અને
    પરંપરિત છે; અને $\mModel{M}$ સંમિત તથા પરંપરિત હોય તો
    $\mModel{M^*}=\tuple{W^*,R^*,V^*}$ ફિલ્ટ્રેશન છે.
  \item ધારો કે $R^*$ એ રીતે વ્યાખ્યાયિત છે કે $R^*[u][v]$ જો
    અને માત્ર જો $C_1(u,
    v) \land C_3(u,v) \land C_4(u,v)$. ત્યારે $R^*$ પરંપરિત
    અને યુક્લિડિયન છે; અને $\mModel{M}$ પરંપરિત તથા યુક્લિડિયન
    હોય તો $\mModel{M^*}=\tuple{W^*,R^*,V^*}$ ફિલ્ટ્રેશન છે.
  \end{enumerate}
\end{thm}

\begin{proof}
  \begin{enumerate}
    \item $R^*$ સંમિત છે તે તરત જ દેખાય છે, કારણ કે $C_1(u,v)
      \Leftrightarrow C_2(v,u)$ અને $C_2(u,v) \Leftrightarrow
      C_1(v,u)$. તેથી $\mModel{M}$ સંમિત હોય તો
      $\mModel{M^*}$ એ~$\Gamma$ દ્વારા ફિલ્ટ્રેશન છે તે
      બતાવવાનું રહે છે. $C_1(u,v)$ શરતથી
      %
        %
          \olref[fil]{defn:filtration-R2} અને
          \olref[fil]{defn:filtration-R3}, એટલે
          \olref[fil]{defn:filtration}ની શરતો
      પૂરી થાય છે. આમ માત્ર
      \olref[fil]{defn:filtration}%
      \olref[fil]{defn:filtration-R1} ચકાસવી રહે છે, એટલે કે
      $Ruv$ હોય તો $R^*[u][v]$.

      $Ruv$ ધારો. $R^*[u][v]$ બતાવવા $C_1(u,v)$ અને
      $C_2(u,v)$ સ્થાપિત કરવાં પડે. $C_1$ માટે:
      જો $\Box!A
        \in\Gamma$ અને $\mSat{M}{\Box!A}[u]$ હોય, તો
        $\mSat{M}{!A}[v]$ પણ (કારણ કે $Ruv$).
        તેવી જ રીતે, %
      જો $\Diamond!A \in \Gamma$ અને
        $\mSat{M}{!A}[v]$ હોય, તો $Ruv$ હોવાથી
        $\mSat{M}{\Diamond!A}[u]$.
      $C_2$ માટે:
      જો $\Box!A \in \Gamma$ અને
        $\mSat{M}{\Box!A}[v]$ હોય, તો સંમિતતાથી $Ruv$માંથી
        $Rvu$ મળે છે; એટલે $\mSat{M}{!A}[u]$.
        તેવી જ રીતે, %
      જો $\Diamond!A \in\Gamma$ અને
        $\mSat{M}{!A}[u]$ હોય, તો $\mSat{M}{\Diamond!A}[v]$
        (કારણ કે સંમિતતાથી $Rvu$).
    \item અભ્યાસ માટે છોડ્યું છે.
    \item અભ્યાસ માટે છોડ્યું છે.
    \item અભ્યાસ માટે છોડ્યું છે.
  \end{enumerate}
\end{proof}

\begin{prob}
  \olref[nml][fil][acc]{thm:more-filtrations}ની સાબિતી પૂરી કરો.
\end{prob}
% END OLP-0458
\endgroup

\begingroup
\makeatletter\def\ol@sectioncs{section}\makeatother

% BEGIN OLP-0459
\olfileid{nml}{fil}{euc}

\olsection{યુક્લિડિયન નિદર્શનોનાં ફિલ્ટ્રેશનો}
\phantomsection\label{guunit:OLP-0459}


\olref[acc]{sec}ની રીત યુક્લિડિયન અથવા સિરિયલ અને યુક્લિડિયન
નિદર્શનોના કિસ્સામાં કામ કરતી નથી. \olref{fig:ser-eucl}ની ઉપરની
આકૃતિનું નિદર્શન જુઓ, જે યુક્લિડિયન અને સિરિયલ બંને છે. ધારો કે
$\Gamma = \{p, \Box p \}$. $\Gamma$ દ્વારા ફિલ્ટ્રેશન લેતાં
$[w_1] = [w_3]$ થાય છે, કારણ કે $w_1$ અને $w_3$ જ
એવાં જગતો છે જે~$\Gamma$ પર સહમત છે. દરેક ફિલ્ટ્રેશનમાં
$\mModel{M}$માંથી મળતી કિનારીઓ પણ રહેશે, જેમ
\olref{fig:ser-eucl2}માં દર્શાવ્યું છે. તે નિદર્શન યુક્લિડિયન નથી.
વળી તેને યુક્લિડિયન બનાવવા કિનારીઓ ઉમેરી પણ શકાતી નથી:
$[w_2]$ અને $[w_4]$ વચ્ચે બંને દિશામાં, અને પછી $[w_2]$
અને~$[w_5]$ વચ્ચે પણ બંને દિશામાં કિનારીઓ ઉમેરવી પડે.
\sourcecorrection{OLMD-050}{મૂળમાં બીજા કિનારી-યુગ્મ માટે
મૂળ જગતો $w_2,w_5$ લખ્યાં છે; અહીં ચર્ચા ફિલ્ટ્રેશનની હોવાથી
સમકક્ષતા-વર્ગો $[w_2],[w_5]$ લખ્યા છે.}
પરંતુ
~$w_2$ પર $\Box p$ સત્ય હોવું જોઈએ, જ્યારે~$w_5$ પર $p$
અસત્ય છે.

\begin{figure}[htpb]
  \centering
  \begin{tikzpicture}[modal]
    \node[world] (w1)
          [label={left: \mFalse{p}},
            label={below: $\mSat{{}}{\Box p}$}] {$w_1$};
    \node[world] (w2)
          [label={right: \mTrue{p}},
            label = {below: $\mSat{{}}{\Box p}$},
            right=of w1] {$w_2$};
    \draw[->] (w1) to (w2);
    \draw[reflexive above] (w2) to (w2);
    \node[world] (w3)
          [label={left: \mFalse{p}},
            label={below: $\mSat{{}}{\Box p}$},
            below=of w1] {$w_3$};
    \node[world] (w4)
          [label={right:\mTrue{p}},
            label={below: $\mSat/{{}}{\Box p}$},
            right=of w3] {$w_4$};
    \draw[->] (w3) to (w4);
    \draw[reflexive above] (w4) to (w4);
    \node[world] (w5)
          [label={right: \mFalse{p}},
            label=below:$\mSat/{{}}{\Box p}$,
            right=of w4] {$w_5$};
    \draw[->, bend left] (w4) to (w5);
    \draw[->, bend left] (w5) to (w4);
    \draw[reflexive above] (w5) to (w5);
  \end{tikzpicture}
  \caption{સિરિયલ અને યુક્લિડિયન નિદર્શન.}
  \ollabel{fig:ser-eucl}
\end{figure}
\sourcecorrection{OLMD-048}{મૂળની બંને આકૃતિમાં બીજા જગતનું
સ્વ-લૂપ છૂટી ગયું છે. તે વિના ઉપરનું નિદર્શન સિરિયલ પણ નથી અને
યુક્લિડિયન પણ નથી; મૂળ દાવાને સાચો બનાવવા બંને આકૃતિમાં તે
લૂપ ઉમેર્યો છે.}

\begin{figure}[ht]
  \centering
  \begin{tikzpicture}[modal,every text node part/.style={align=center}]
    \node[world] (w1)
          [label={left: \mFalse{p}},
            label={right:$[w_1]=[w_3]$},
            label={below: $\mSat{{}}{\Box p}$}] {$[w_1]$};
    \node[world] (w2)
         [label={right: \mTrue{p}},
           label = {below: $\mSat{{}}{\Box p}$},
           right=of w1, yshift=1.5cm] {$[w_2]$};
    \draw[->] (w1) to (w2);
    \draw[reflexive above] (w2) to (w2);
    \node[world] (w4)
         [label={right:\mTrue{p}},
           label={below: $\mSat/{{}}{\Box p}$},
           right=of w1,yshift=-1.5cm] {$[w_4]$};
    \draw[->] (w1) to (w4);
    \draw[reflexive above] (w4) to (w4);
    \node[world] (w5)
         [label={right: \mFalse{p}},
             label=below:$\mSat/{{}}{\Box p}$,
             right=of w4] {$[w_5]$};
    \draw[->, bend left] (w4) to (w5);
    \draw[->, bend left] (w5) to (w4);
    \draw[reflexive above] (w5) to (w5);
  \end{tikzpicture}
  \caption{\olref{fig:ser-eucl}ના નિદર્શનનું ફિલ્ટ્રેશન.}
  \ollabel{fig:ser-eucl2}
\end{figure}

ખાસ કરીને યુક્લિડિયન ફિલ્ટ્રેશન મેળવવા માત્ર મનસ્વી
ઉપ-સૂત્રો-બંધ ગણો $\Gamma$ દ્વારા ફિલ્ટ્રેશન લેવું પૂરતું
નથી. તેના બદલે \emph{મોડલ રીતે બંધ} ગણો $\Gamma$
લેવા પડે છે (જુઓ \olref[pre]{defn:modallyclosed}). વાક્યોના
આવા ગણો અનંત હોય છે; તેથી તેઓ તાત્કાલિક સાંત નિદર્શન ગુણધર્મ
અથવા અનુરૂપ પદ્ધતિની નિર્ણેયતા આપતા નથી.

\begin{thm}\ollabel{thm:modal-closed-filt}
  $\Gamma$ મોડલ રીતે બંધ હોય,
  $\mModel{M}=\tuple{W,R,V}$ હોય, અને
  $\mModel{M^*} = \tuple{W^*,R^*,V^*}$ એ
  $\mModel{M}$નું સૌથી સ્થૂલ ફિલ્ટ્રેશન હોય.
  \begin{enumerate}
  \item જો $\mModel{M}$ સંમિત હોય, તો $\mModel{M^*}$ પણ સંમિત છે.
  \item જો $\mModel{M}$ પરંપરિત હોય, તો $\mModel{M^*}$ પણ પરંપરિત છે.
  \item જો $\mModel{M}$ યુક્લિડિયન હોય, તો $\mModel{M^*}$ પણ યુક્લિડિયન છે.
  \end{enumerate}
\end{thm}

\begin{proof}
\sourcecorrection{OLMD-049}{મૂળના પ્રમાણમાં પ્રથમ કિસ્સો
પરંપરિતતા માટે છે, જ્યારે પ્રમેયમાં પહેલું વિધાન સંમિતતાનું છે;
બાકીના બે પ્રશ્નો પણ સરકેલા છે. પ્રમાણના ત્રણ કિસ્સાઓ પ્રમેયના
ક્રમ સાથે ગોઠવ્યા છે.}
\begin{enumerate}
  \item અભ્યાસ માટે છોડ્યું છે. બધાં સંમિત નિદર્શનોમાં
    \Ax{B} અને $\Ax{B_\Diamond}$ માન્ય છે તેનો ઉપયોગ કરો.
  \item જો $\mModel{M^*}$ સૌથી સ્થૂલ ફિલ્ટ્રેશન હોય, તો વ્યાખ્યા
    મુજબ $R^*[u][v]$ જો અને માત્ર જો $C_1(u,v)$. પરંપરિતતા
    માટે $C_1(u,v)$ અને $C_1(v,w)$ ધારો; $C_1(u,w)$ બતાવવાનું છે.
    ધારો કે $\mSat{M}{\Box !A}[u]$; ત્યારે
      બધાં પરંપરિત નિદર્શનોમાં \Ax{4} માન્ય હોવાથી
      $\mSat{M}{\Box\Box!A}[u]$. મોડલ બંધતાથી
      $\Box\Box!A \in \Gamma$ હોવાથી $C_1(u,v)$ દ્વારા
      $\mSat{M}{\Box!A}[v]$ અને $C_1(v,w)$ દ્વારા
      $\mSat{M}{!A}[w]$. %
    ધારો કે $\mSat{M}{!A}[w]$; મોડલ
      બંધતાથી $\Diamond !A \in \Gamma$ હોવાથી $C_1(v,w)$
      દ્વારા $\mSat{M}{\Diamond !A}[v]$. મોડલ બંધતાથી
      $\Diamond\Diamond!A \in \Gamma$ હોવાથી $C_1(u,v)$
      દ્વારા $\mSat{M}{\Diamond\Diamond !A}[u]$. બધાં
      પરંપરિત નિદર્શનોમાં $\Ax{4}_\Diamond$ માન્ય હોવાથી
      $\mSat{M}{\Diamond!A}[u]$.
  \item અભ્યાસ માટે છોડ્યું છે. બધાં યુક્લિડિયન નિદર્શનોમાં
    \Ax{5} અને $\Ax{5_\Diamond}$ માન્ય છે તેનો ઉપયોગ કરો.
\end{enumerate}
\end{proof}

\begin{prob}
  \olref[nml][fil][euc]{thm:modal-closed-filt}ની સાબિતી પૂરી કરો.
\end{prob}
% END OLP-0459
\endgroup


\OLEndChapterHook
% END OLP-0450
\endgroup


\begingroup
\makeatletter\def\ol@sectioncs{section}\makeatother

% BEGIN OLP-0460
\olchapter{nml}{tab}{મોડલ ટેબ્લો}
\olchapterid{nml}{tab}
\phantomsection\label{guunit:OLP-0460}


\begin{editorial}
  નિદર્શ્ય તર્કશાસ્ત્ર માટે ઉપસર્ગવાળા ટેબ્લો પરનું આ પ્રારૂપ અધ્યાય છે.
  તેમાં વધુ ઉદાહરણો, પૂર્ણતાનાં પ્રમાણો અને બંધ ટેબ્લોની નિષ્ફળ
  શોધમાંથી પ્રતિનિદર્શનો કેવી રીતે મેળવવાં તેની ચર્ચા જરૂરી છે.
\end{editorial}

\begingroup
\makeatletter\def\ol@sectioncs{section}\makeatother

% BEGIN OLP-0461
\olfileid{nml}{tab}{int}

\olsection{પરિચય}
\phantomsection\label{guunit:OLP-0461}


ટેબ્લો એ ચિહ્નિત સૂત્રોનાં ચોક્કસ પ્રકારનાં (નીચે તરફ શાખા પાડતાં) વૃક્ષો છે.
ચિહ્નિત સૂત્ર એટલે સત્યમૂલ્ય-ચિહ્ન ($\True$ અથવા
$\False$) અને વાક્યની જોડી:
\[
\sFmla{\True}{!A} \text{ અથવા } \sFmla{\False}{!A}.
\]
ટેબ્લો કેટલીક \emph{ધારણાઓ}થી શરૂ થાય છે. પછીનું દરેક
ચિહ્નિત સૂત્ર અનુમાનનો કોઈ નિયમ લાગુ કરીને મળે છે.
કેટલાક નિયમો વૃક્ષની શાખાના છેડે એક અથવા વધુ
ચિહ્નિત સૂત્રો ઉમેરે છે; બીજા નિયમો બે નવા છેડા ઉમેરી
બે શાખાઓ બનાવે છે. કોઈ નિયમથી મળતા ચિહ્નિત સૂત્રોમાંનું
સૂત્ર જે ચિહ્નિત સૂત્ર પર નિયમ લાગુ થયો તેના
સૂત્ર કરતાં ઓછું સંકુલ હોય છે. કોઈ શાખામાં
$\sFmla{\True}{!A}$ અને $\sFmla{\False}{!A}$ બંને હોય,
તો તે શાખાને \emph{બંધ} કહીએ છીએ. ટેબ્લોની દરેક શાખા
બંધ હોય તો આખું ટેબ્લો બંધ છે. બંધ ટેબ્લો એક
નિષ્પત્તિ છે, જે બતાવે છે કે ટેબ્લોની શરૂઆતમાં લીધેલા
ચિહ્નિત સૂત્રોનો ગણ અસંતોષ્ય છે. એના આધારે
$\Proves$ સંબંધ વ્યાખ્યાયિત કરી શકાય:
$\Gamma \Proves !A$ જો અને માત્ર જો એવો કોઈ સાંત ગણ
~$\Gamma_0 = \{!B_1,
\dots, !B_n\} \subseteq \Gamma$ હોય કે નીચેની ધારણાઓ માટે
બંધ ટેબ્લો મળે:
\[
\{\sFmla{\False}{!A}, \sFmla{\True}{!B_1}, \dots, \sFmla{\True}{!B_n}\}.
\]

નિદર્શ્ય તર્કશાસ્ત્રો માટે ચિહ્નિત સૂત્રની સંકલ્પના વિસ્તૃત કરવી પડે છે અને
~$\Box$ અને
    $\Diamond$ને આવરી લેતા નિયમો ઉમેરવા પડે છે.
\sourcecorrection{OLMD-051}{મૂળમાં અંદરના
\texttt{\string\iftag\{prvDiamond\}} માટે જરૂરી ત્રીજું, ખાલી
શાખા-આર્ગ્યુમેન્ટ ખૂટે છે; પસંદગી-મેક્રો યોગ્ય રીતે ચાલે તે માટે
તે ઉમેર્યું છે.}
ચિહ્ન ($\True$ અથવા $\False$) ઉપરાંત, મોડલ ટેબ્લોનાં
સૂત્રોને \emph{ઉપસર્ગો}~$\sigma$ પણ હોય છે. ઉપસર્ગો ધન
પૂર્ણાંકોની ખાલી ન હોય તેવી શ્રેણીઓ છે, એટલે
$\sigma \in (\PosInt)^* \setminus
\{\emptyseq\}$. આવા ઉપસર્ગો લખતી વખતે આપણે આસપાસના
$\tuple{\ }$ છોડીએ છીએ અને જુદા ઘટકો વચ્ચે
$,$ને બદલે~$.$ મૂકીએ છીએ. જો $\sigma$ ઉપસર્ગ હોય, તો
$\sigma.n$ એ $\sigma
\concat \tuple{n}$ છે; દાખલા તરીકે, $\sigma = 1.2.1$ હોય
તો $\sigma.3$ એ $1.2.1.3$ છે. આમ, દાખલા તરીકે,
\[
\sFmla{\True}{\Box !A \lif !A}[1.2]
\]
એ \emph{ઉપસર્ગવાળું ચિહ્નિત સૂત્ર} છે (અથવા ટૂંકમાં
\emph{ઉપસર્ગવાળું સૂત્ર}).

સહજ રીતે, ઉપસર્ગ એ નિદર્શનમાં એવા જગતનું નામ આપે છે જે
ટેબ્લોની શાખા પરનાં સૂત્રોને સંતોષી શકે; અને જો
$\sigma$ કોઈ જગતનું નામ હોય, તો $\sigma.n$ એ~$\sigma$ નામના
જગતમાંથી પ્રાપ્ય જગતનું નામ આપે છે.
% END OLP-0461
\endgroup

\begingroup
\makeatletter\def\ol@sectioncs{section}\makeatother

% BEGIN OLP-0462
\olfileid{nml}{tab}{rul}

\olsection{\Ax{K} માટેના નિયમો}
\phantomsection\label{guunit:OLP-0462}


સામાન્ય વિધાનસંબંધી જોડાણીઓ માટેના નિયમો સામાન્ય વિધાનસંબંધી
ચિહ્નિત ટેબ્લો જેવા જ છે; ફક્ત ઉપસર્ગો ઉમેરાય છે.
દરેક કિસ્સામાં, ચિહ્નિત સૂત્ર
$\sFmla{S}{!A}[\sigma]$ પર લાગુ થયેલો નિયમ નવાં સૂત્રો
આપે છે, જેમને પણ~$\sigma$ ઉપસર્ગ હોય છે. આ સહજ છે: દાખલા તરીકે,
જો $!A \land !B$ એ~$\sigma$ નામના જગત પર સત્ય હોય, તો
$!A$ અને $!B$ પણ~$\sigma$ પર જ સત્ય છે (અન્ય જગત પર જરૂરી
નથી). વિધાનસંબંધી નિયમો \olref{tab:prop-rules}માં એકત્ર કર્યા છે.

\begin{table}
  \[\def\arraystretch{3}\begin{array}{|c|c|}
    \hline
    \AxiomC{\sFmla{\True}{\lnot !A}[\sigma]}
    \RightLabel{\TRule{\True}{\lnot}}
    \UnaryInfC{\sFmla{\False}{!A}[\sigma]}
    \DisplayProof
    &
    \AxiomC{\sFmla{\False}{\lnot !A}[\sigma]}
    \RightLabel{\TRule{\False}{\lnot}}
    \UnaryInfC{\sFmla{\True}{!A}[\sigma]}
    \DisplayProof
    \\[1ex]
    \hline
    \AxiomC{\sFmla{\True}{!A \land !B}[\sigma]}
    \RightLabel{\TRule{\True}{\land}}
    \UnaryInfC{\sFmla{\True}{!A}[\sigma]}
    \noLine
    \UnaryInfC{\sFmla{\True}{!B}[\sigma]}
    \DisplayProof
    &
    \AxiomC{\sFmla{\False}{!A \land !B}[\sigma]}
    \RightLabel{\TRule{\False}{\land}}
    \UnaryInfC{$\sFmla{\False}{!A}[\sigma] \quad \mid \quad
      \sFmla{\False}{!B}[\sigma]$}
    \DisplayProof
    \\[2ex]
    \hline
    \AxiomC{\sFmla{\True}{!A \lor !B}[\sigma]}
    \RightLabel{\TRule{\True}{\lor}}
    \UnaryInfC{$\sFmla{\True}{!A}[\sigma] \quad \mid \quad
      \sFmla{\True}{!B}[\sigma]$}
    \DisplayProof
    &
    \AxiomC{\sFmla{\False}{!A \lor !B}[\sigma]}
    \RightLabel{\TRule{\False}{\lor}}
    \UnaryInfC{\sFmla{\False}{!A}[\sigma]}
    \noLine
    \UnaryInfC{\sFmla{\False}{!B}[\sigma]}
    \DisplayProof
    \\[2ex]
    \hline
    \AxiomC{\sFmla{\True}{!A \lif !B}[\sigma]}
    \RightLabel{\TRule{\True}{\lif}}
    \UnaryInfC{$\sFmla{\False}{!A}[\sigma] \quad \mid
      \quad \sFmla{\True}{!B}[\sigma]$}
    \DisplayProof
    &
    \AxiomC{\sFmla{\False}{!A \lif !B}[\sigma]}
    \RightLabel{\TRule{\False}{\lif}}
    \UnaryInfC{\sFmla{\True}{!A}[\sigma]}
    \noLine
    \UnaryInfC{\sFmla{\False}{!B}[\sigma]}
    \DisplayProof
    \\[2ex]
    \hline
  \end{array}\]
  \caption{વિધાનસંબંધી જોડાણીઓ માટે ઉપસર્ગવાળા ટેબ્લો નિયમો}
  \ollabel{tab:prop-rules}
\end{table}

બંધતાની શરત સામાન્ય ટેબ્લો જેવી જ છે, પરંતુ અહીં માત્ર
સૂત્રો જ નહિ, ઉપસર્ગો પણ સરખા હોવા જોઈએ. તેથી કોઈ શાખામાં
નીચેના બંને હોય તો તે બંધ છે:
\[
\sFmla{\True}{!A}[\sigma] \quad\text{અને}\quad \sFmla{\False}{!A}[\sigma]
\]
કોઈ ઉપસર્ગ $\sigma$ અને સૂત્ર~$!A$ માટે.

ધારણાઓ ગોઠવવાના નિયમો પણ સામાન્ય ટેબ્લો જેવા જ છે;
ફક્ત ધારણાઓ માટે આપણે હંમેશા ઉપસર્ગ~$1$ લઈએ છીએ.
(કયો ઉપસર્ગ લઈએ તે મહત્ત્વનું નથી, જો બધી ધારણાઓ માટે તે એક જ હોય.)
દાખલા તરીકે, આપણે કહીએ છીએ કે
\[
!B_1, \dots, !B_n \Proves !A
\]
જો અને માત્ર જો નીચેની ધારણાઓ માટે બંધ ટેબ્લો હોય:
\[
\sFmla{\True}{!B_1}[1], \dots, \sFmla{\True}{!B_n}[1],
\sFmla{\False}{!A}[1].
\]

$\Box$ અને %
$\Diamond$ જેવા મોડલ સંક્રિયકો માટે,
ઉપસર્ગ~$\sigma$વાળા સૂત્ર પર લાગુ થયેલા નિયમના
નિષ્કર્ષનો ઉપસર્ગ $\sigma.n$ હોય છે. જોકે કયું $n$ લઈ શકાય
તે ચિહ્ન~$\True$ છે કે~$\False$ તેના પર આધાર રાખે છે.

$\TRule{\True}{\Box}$ નિયમ
  $\sFmla{\True}{\Box !A}[\sigma]$ ધરાવતી શાખાને
  $\sFmla{\True}{!A}[\sigma.n]$થી વિસ્તારે છે.
  તેવી જ રીતે, %
$\TRule{\False}{\Diamond}$ નિયમ
  $\sFmla{\False}{\Diamond !A}[\sigma]$ ધરાવતી શાખાને
  $\sFmla{\False}{!A}[\sigma.n]$થી વિસ્તારે છે.
આ નિયમો ફક્ત એવા
ઉપસર્ગ~$\sigma.n$ માટે લાગુ કરી શકાય છે જે સંબંધિત શાખા પર
\emph{પહેલેથી} આવે છે. આવા ઉપસર્ગને તે શાખા પર ``વપરાયેલો'' કહીએ.

$\TRule{\False}{\Box}$ નિયમ
  $\sFmla{\False}{\Box !A}[\sigma]$ ધરાવતી શાખાને
  $\sFmla{\False}{!A}[\sigma.n]$થી વિસ્તારે છે.
  તેવી જ રીતે, %
$\TRule{\True}{\Diamond}$ નિયમ
  $\sFmla{\True}{\Diamond !A}[\sigma]$ ધરાવતી શાખાને
  $\sFmla{\True}{!A}[\sigma.n]$થી વિસ્તારે છે.
જોકે આ નિયમો ફક્ત એવા
ઉપસર્ગ~$\sigma.n$ માટે લાગુ કરી શકાય છે જે સંબંધિત શાખા પર
\emph{પહેલેથી આવતો નથી}. આવા ઉપસર્ગોને તે શાખા માટે ``નવા'' કહીએ.

નિયમો \olref{tab:rules-K}માં આપ્યા છે.

\begin{table}
  \begin{center}
    \def\arraystretch{3}\def\fCenter{}
    \begin{tabular}{|c|c|}
    \hline
          \AxiomC{\sFmla{\True}{\Box !A}[\sigma]}
      \RightLabel{\TRule{\True}{\Box}}
      \UnaryInfC{\sFmla{\True}{!A}[\sigma.n]}
      \DisplayProof
      &
      \AxiomC{\sFmla{\False}{\Box !A}[\sigma]}
      \RightLabel{\TRule{\False}{\Box}}
      \UnaryInfC{\sFmla{\False}{!A}[\sigma.n]}
      \DisplayProof\\
      $\sigma.n$ વપરાયેલો છે & $\sigma.n$ નવો છે
      \\[1ex]
      \hline
          \AxiomC{\sFmla{\True}{\Diamond !A}[\sigma]}
      \RightLabel{\TRule{\True}{\Diamond}}
      \UnaryInfC{\sFmla{\True}{!A}[\sigma.n]}
      \DisplayProof
      &
      \AxiomC{\sFmla{\False}{\Diamond !A}[\sigma]}
      \RightLabel{\TRule{\False}{\Diamond}}
      \UnaryInfC{\sFmla{\False}{!A}[\sigma.n]}
      \DisplayProof\\
      $\sigma.n$ નવો છે & $\sigma.n$ વપરાયેલો છે
      \\[1ex]
      \hline
    \end{tabular}
  \end{center}
  \caption{\Ax{K} માટેના મોડલ નિયમો.}
  \ollabel{tab:rules-K}
\end{table}

\TRule{\True}{\Box}
માટેનો ઉપસર્ગ વપરાયેલો હોવો જોઈએ એવી મર્યાદા જરૂરી છે;
નહિ તો નીચેનું પણ બંધ ટેબ્લો ગણવું પડે:
  \begin{oltableau}
    [\pFmla{\True}{\Box \formula{A}}{1}, just = \TAss
      [\pFmla{\False}{\Diamond \formula{A}}{1}, just = \TAss
        [\pFmla{\True}{\formula{A}}{1.1}, just = {\TRule{\True}{\Box}[1]}
          [\pFmla{\False}{\formula{A}}{1.1},
            just = {\TRule{\False}{\Diamond}[2]}, close]
        ]
      ]
    ]
\end{oltableau}

પરંતુ $\Box \formula{A} \Entails/ \Diamond \formula{A}$,
તેથી આપણી પ્રમાણપદ્ધતિ અવિશ્વસનીય બની જાય. તેવી જ રીતે,
$\Diamond \formula{A} \Entails/
\Box \formula{A}$ પણ નથી; છતાં
\TRule{\False}{\Box}
માટેનો ઉપસર્ગ નવો હોવાની મર્યાદા ન હોય, તો નીચેનું બંધ ટેબ્લો થાય:
  \begin{oltableau}
    [\pFmla{\True}{\Diamond \formula{A}}{1}, just = \TAss,name=one
      [\pFmla{\False}{\Box \formula{A}}{1}, just = \TAss, name=two
        [\pFmla{\True}{\formula{A}}{1.1}, just = {\TRule{\True}{\Diamond}[1]}
          [\pFmla{\False}{\formula{A}}{1.1},
            just = {\TRule{\False}{\Box}[2]}, close]
        ]
      ]
    ]
  \end{oltableau}
% END OLP-0462
\endgroup

\begingroup
\makeatletter\def\ol@sectioncs{section}\makeatother

% BEGIN OLP-0463
\olfileid{nml}{tab}{prk}

\olsection{\Log{K} માટેના ટેબ્લો}
\phantomsection\label{guunit:OLP-0463}



\begin{ex}
  $\Proves (\Box!A \land \Box!B)
  \lif \Box (!A \land !B)$ દર્શાવતું બંધ ટેબ્લો આપીએ છીએ.
  \begin{oltableau}
    [\pFmla{\False}{(\Box\formula{A} \land \Box\formula{B}) \lif
        \Box (\formula{A} \land \formula{B})}{1},
      just =\TAss
      [\pFmla{\True}{\Box\formula{A} \land \Box\formula{B}}{1},
        just = {\TRule{\False}{\lif}[1]}
        [\pFmla{\False}{\Box(\formula{A} \land \formula{B})}{1},
          just = {\TRule{\False}{\lif}[1]}
          [\pFmla{\True}{\Box\formula{A}}{1},
            just = {\TRule{\True}{\land}[2]}
            [\pFmla{\True}{\Box\formula{B}}{1},
              just = {\TRule{\True}{\land}[2]}
              [\pFmla{\False}{\formula{A} \land \formula{B}}{1.1},
                just = {\TRule{\False}{\Box}[3]}
                [\pFmla{\False}{\formula{A}}{1.1},
                  just = {\TRule{\False}{\land}[6]}
                  [\pFmla{\True}{\formula{A}}{1.1},
                    just= {\TRule{\True}{\Box}[4]}, close]]
                [\pFmla{\False}{\formula{B}}{1.1},
                  just = {\TRule{\False}{\land}[6]}
                  [\pFmla{\True}{\formula{B}}{1.1},
                    just= {\TRule{\True}{\Box}[5]}, close]]
              ]
            ]
          ]
        ]
      ]
    ]
  \end{oltableau}
\end{ex}

\begin{ex}
  $\Proves \Diamond(!A \lor !B)
  \lif (\Diamond !A \lor \Diamond !B)$ દર્શાવતું બંધ ટેબ્લો આપીએ છીએ:
  \begin{oltableau}
    [\pFmla{\False}{\Diamond(\formula{A} \lor \formula{B}) \lif
        (\Diamond \formula{A} \lor \Diamond \formula{B})}{1},
      just =\TAss
      [\pFmla{\True}{\Diamond(\formula{A} \lor \formula{B})}{1},
        just = {\TRule{\False}{\lif}[1]}
        [\pFmla{\False}{\Diamond\formula{A} \lor \Diamond\formula{B}}{1},
          just = {\TRule{\False}{\lif}[1]}
          [\pFmla{\False}{\Diamond\formula{A}}{1},
            just = {\TRule{\False}{\lor}[3]}
            [\pFmla{\False}{\Diamond\formula{B}}{1},
              just = {\TRule{\False}{\lor}[3]}
              [\pFmla{\True}{\formula{A} \lor \formula{B}}{1.1},
                just = {\TRule{\True}{\Diamond}[2]},
                [\pFmla{\True}{\formula{A}}{1.1},
                  just = {\TRule{\True}{\lor}[6]}
                  [\pFmla{\False}{\formula{A}}{1.1},
                    just= {\TRule{\False}{\Diamond}[4]}, close]]
                [\pFmla{\True}{\formula{B}}{1.1},
                  just = {\TRule{\True}{\lor}[6]}
                  [\pFmla{\False}{\formula{B}}{1.1},
                    just= {\TRule{\False}{\Diamond}[5]}, close]]
              ]
            ]
          ]
        ]
      ]
    ]
  \end{oltableau}
\end{ex}


\begin{prob}
  નીચેના સૂત્રો માટે~$\Log{K}$માં બંધ ટેબ્લો શોધો:
  \begin{enumerate}
    \item $\Box \lnot p \lif \Box(p \lif q)$
    \item $(\Box p \lor \Box q) \lif \Box(p \lor q)$
    \item $\Diamond p \lif \Diamond(p \lor q)$
    \item $\Box(p \land q) \lif \Box p$
  \end{enumerate}
\end{prob}
% END OLP-0463
\endgroup

\begingroup
\makeatletter\def\ol@sectioncs{section}\makeatother

% BEGIN OLP-0464
\olfileid{nml}{tab}{sou}

\olsection{\Log{K} માટેની યથાર્થતા}
\phantomsection\label{guunit:OLP-0464}


\begin{editorial}
  યથાર્થતાનું આ પ્રમાણ શાસ્ત્રીય વિધાનતર્કનું યથાર્થતા-પ્રમાણ ફરી
  વાપરે છે: એટલે તે બધું શરૂઆતથી સાબિત કરે છે. સ્વતંત્ર રીતે
  વાંચી શકાય એવું યથાર્થતા-પ્રમાણ જોઈએ તો આ યોગ્ય છે. સામાન્ય
  ટેબ્લો માટેની યથાર્થતા પહેલાં જોઈ હોય તો આ પુનરાવર્તન લાગશે.
  આગળ જતાં સ્વતંત્ર આવૃત્તિ અને અમોડલ કિસ્સા પર આધારિત આવૃત્તિ
  વચ્ચે પસંદગી શક્ય બનાવવાની યોજના છે.
\end{editorial}

\begin{explain}
  ઉપસર્ગવાળા ટેબ્લો યથાર્થ છે તે બતાવવા આપણે સાબિત કરવું પડે કે
  \[
  \sFmla{\True}{!B_1}[1], \dots, \sFmla{\True}{!B_n}[1], \sFmla{\False}{!A}[1]
  \]
  માટે બંધ ટેબ્લો હોય, તો $!B_1, \dots, !B_n \Entails !A$.
  પ્રતિવિધાન સાબિત કરવું વધુ સરળ છે: કોઈ $\mModel{M}$ અને
  જગત~$w$ માટે બધા $i=1$, \dots,~$n$ પર
  $\mSat{M}{!B_i}[w]$ હોય, પરંતુ
  $\mSat/{M}{!A}[w]$ હોય, તો કોઈ ટેબ્લો બંધ થઈ શકે નહિ.
  \sourcecorrection{OLMD-052}{મૂળમાં પ્રતિનિદર્શનમાં $!A$ પણ સત્ય
  લખાયું છે; પ્રતિવિધાન અને આગળની દલીલ માટે તે અસત્ય હોવું જરૂરી
  છે, તેથી અસંતોષનું ચિહ્ન પુનઃસ્થાપિત કર્યું છે.}
  આવું પ્રતિનિદર્શન બતાવે છે કે ટેબ્લોની પ્રારંભિક ધારણાઓ
  સંતોષ્ય છે. પ્રમાણની રીત એ બતાવવાની છે કે ટેબ્લોની કોઈ
  શાખા પરનાં બધાં ઉપસર્ગવાળાં સૂત્રો સંતોષ્ય હોય ત્યારે
  કોઈપણ નિયમ લાગુ કરવાથી ઓછામાં ઓછી એક વિસ્તૃત શાખા પણ
  સંતોષ્ય રહે છે. બંધ શાખાઓ અસંતોષ્ય હોવાથી, ઉપસર્ગવાળાં
  સૂત્રોના સંતોષ્ય ગણ માટેના દરેક ટેબ્લોમાં
  ઓછામાં ઓછી એક ખુલ્લી શાખા હશે.

  આ રીત મોડલ કિસ્સામાં લાગુ કરવા, ``સંતોષ્ય''ની વ્યાખ્યા
  મોડલ નિદર્શનો અને ઉપસર્ગો સુધી વિસ્તૃત કરવી પડે. તે થઈ જાય
  પછી પ્રમાણ સીધું છે.
\end{explain}

\begin{defn}
  $P$ ઉપસર્ગોનો ગણ હોય, એટલે $P \subseteq (\PosInt)^*
  \setminus \{\emptyseq\}$, અને $\mModel{M}$ નિદર્શન હોય.
  વિધેય~$f\colon P \to W$ એ~$\mModel{M}$માં~$P$નું
  \emph{અર્થઘટન} ત્યારે કહેવાય જ્યારે $\sigma$ અને $\sigma.n$
  બંને~$P$માં હોય ત્યારે $Rf(\sigma)f(\sigma.n)$ હોય.

  ઉપસર્ગો~$P$ના અર્થઘટનને અનુલક્ષીને વ્યાખ્યાયિત કરીએ છીએ:
  \begin{enumerate}
  \item $\mModel{M}$ એ $\sFmla{\True}{!A}[\sigma]$ને સંતોષે છે
    જો અને માત્ર જો
    $\mSat{M}{!A}[f(\sigma)]$.
  \item $\mModel{M}$ એ $\sFmla{\False}{!A}[\sigma]$ને સંતોષે છે
    જો અને માત્ર જો
    $\mSat/{M}{!A}[f(\sigma)]$.
  \end{enumerate}
\end{defn}

\begin{defn}
  $\Gamma$ ઉપસર્ગવાળાં સૂત્રોનો ગણ અને $P(\Gamma)$ તેમાં
  આવતા ઉપસર્ગોનો ગણ હોય. $f$ એ~$\mModel{M}$માં
  $P(\Gamma)$નું અર્થઘટન હોય, તો $\mModel{M}$ એ~$f$ને અનુલક્ષીને
  $\Gamma$ને સંતોષે છે, $\mSat{M}{\Gamma}[f]$, એમ ત્યારે કહીએ
  જ્યારે $\mModel{M}$ એ~$\Gamma$ના દરેક ઉપસર્ગવાળા
  સૂત્રને~$f$ને અનુલક્ષીને સંતોષે. $\Gamma$
  \emph{સંતોષ્ય} છે જો અને માત્ર જો એવું નિદર્શન~$\mModel{M}$
  અને $P(\Gamma)$નું અર્થઘટન~$f$ હોય કે
  $\mSat{M}{\Gamma}[f]$.
\end{defn}

\begin{prop}
  જો $\Gamma$માં $\sFmla{\True}{!A}[\sigma]$ અને
  $\sFmla{\False}{!A}[\sigma]$ બંને હોય, કોઈ સૂત્ર~$!A$
  અને ઉપસર્ગ~$\sigma$ માટે, તો $\Gamma$ અસંતોષ્ય છે.
\end{prop}

\begin{proof}
  એવું નિદર્શન~$\mModel{M}$ અને $P(\Gamma)$નું અર્થઘટન~$f$
  હોઈ શકે નહિ કે $\mSat{M}{!A}[f(\sigma)]$ અને
  $\mSat/{M}{!A}[f(\sigma)]$ બંને હોય.
\end{proof}

\begin{thm}[યથાર્થતા]
  \ollabel{thm:tableau-soundness}
  જો $\Gamma$ માટે બંધ ટેબ્લો હોય, તો $\Gamma$ અસંતોષ્ય છે.
\end{thm}

\begin{proof}
ટેબ્લોની શાખા પરનાં ચિહ્નિત સૂત્રોનો ગણ સંતોષ્ય હોય
તો અને માત્ર તો શાખાને સંતોષ્ય કહીએ; અને ઓછામાં ઓછી એક સંતોષ્ય
શાખા ધરાવતું ટેબ્લો સંતોષ્ય કહેવાય.

નીચેનું બતાવીએ છીએ: સંતોષ્ય ટેબ્લોને અનુમાનના કોઈ એક નિયમથી
વિસ્તારતાં હંમેશા સંતોષ્ય ટેબ્લો મળે છે. તેનાથી પ્રમેય
સાબિત થશે: કોઈપણ બંધ ટેબ્લો, માત્ર~$\Gamma$માંથી લીધેલી
ધારણાઓવાળા ટેબ્લોને અનુમાનના નિયમો લાગુ કરવાથી મળે છે.
તેથી $\Gamma$ સંતોષ્ય હોય, તો તેના માટેનું દરેક ટેબ્લો
સંતોષ્ય હશે. પરંતુ બંધ ટેબ્લો દેખીતી રીતે સંતોષ્ય નથી:
તેની બધી શાખાઓ બંધ છે અને બંધ શાખાઓ અસંતોષ્ય છે.

ધારો કે આપણી પાસે સંતોષ્ય ટેબ્લો છે, એટલે ઓછામાં ઓછી એક
સંતોષ્ય શાખાવાળું ટેબ્લો. અનુમાનનો નિયમ લાગુ કરવાથી
શાખામાં ચિહ્નિત સૂત્રો ઉમેરાય છે અથવા શાખા બે ભાગમાં
વહેંચાય છે. જો ટેબ્લોમાં એવી સંતોષ્ય શાખા હોય જેને સંબંધિત
નિયમથી વિસ્તારી ન હોય, તો વિસ્તૃત ટેબ્લોમાં પણ તે
સંતોષ્ય જ રહે છે; તેથી વિસ્તૃત ટેબ્લો સંતોષ્ય છે. આમ માત્ર એવો
કિસ્સો ધ્યાનમાં લેવો પડે જેમાં નિયમ સંતોષ્ય શાખાને લાગુ થાય.

તે શાખા પરનાં ચિહ્નિત સૂત્રોનો ગણ $\Gamma$ લો, અને નિયમ
જેને લાગુ પડે તે ચિહ્નિત સૂત્ર
$\sFmla{S}{!A}[\sigma] \in \Gamma$ લો. જો નિયમથી શાખા
વહેંચાતી ન હોય, તો વિસ્તૃત શાખા, એટલે નિયમના નિષ્કર્ષો સાથેનો
$\Gamma$, હજી સંતોષ્ય છે તે બતાવવું પડે. જો શાખા વહેંચાતી
હોય, તો મળતી બે શાખામાંથી ઓછામાં ઓછી એક સંતોષ્ય છે તે બતાવવું
પડે.
\sourcecorrection{OLMD-053}{મૂળમાં અહીં
\texttt{\string\tagfalse\{prvDiamond\}} લખીને ડાયમંડ-શાખાઓ
બંધ કરી દેવાય છે, જ્યારે આગળની સાબિતીમાં એ જ શાખાઓ જરૂરી છે.
આ અણધાર્યો ટૅગ-ફેરફાર દૂર કર્યો છે.}
પહેલાં શાખા ન વહેંચતા સંભવિત અનુમાનો ધ્યાનમાં લઈએ.
\begin{enumerate}
\item $\sFmla{\True}{\lnot !B}[\sigma] \in \Gamma$ને
  $\TRule{\True}{\lnot}$ લાગુ કરીને શાખા વિસ્તારીએ. ત્યારે
  વિસ્તૃત શાખામાં ચિહ્નિત સૂત્રોનો ગણ $\Gamma \cup
  \{\sFmla{\False}{!B}[\sigma]\}$ છે. ધારો કે
  $\mSat{M}{\Gamma}[f]$. ખાસ કરીને,
  $\mSat{M}{\lnot !B}[f(\sigma)]$. તેથી
  $\mSat/{M}{!B}[f(\sigma)]$, એટલે~$\mModel{M}$ એ
  $\sFmla{\False}{!B}[\sigma]$ને~$f$ને અનુલક્ષીને સંતોષે છે.
\item $\sFmla{\False}{\lnot !B}[\sigma] \in \Gamma$ને
  $\TRule{\False}{\lnot}$ લાગુ કરીને શાખા વિસ્તારીએ:
  અભ્યાસપ્રશ્ન.
\item $\sFmla{\True}{!B \land !C}[\sigma] \in \Gamma$ને
  $\TRule{\True}{\land}$ લાગુ કરીને શાખા વિસ્તારીએ. તેનાથી
  શાખા પર બે નવાં ચિહ્નિત સૂત્રો મળે છે:
  $\sFmla{\True}{!B}[\sigma]$ અને $\sFmla{\True}{!C}[\sigma]$.
  ધારો કે $\mSat{M}{\Gamma}[f]$, ખાસ કરીને
  $\mSat{M}{!B \land !C}[f(\sigma)]$. ત્યારે
  $\mSat{M}{!B}[f(\sigma)]$ અને $\mSat{M}{!C}[f(\sigma)]$.
  એટલે $\mModel{M}$ એ $\sFmla{\True}{!B}[\sigma]$ અને
  $\sFmla{\True}{!C}[\sigma]$ બંનેને~$f$ને અનુલક્ષીને સંતોષે છે.
\item $\sFmla{\False}{!B \lor !C}[\sigma] \in \Gamma$ને
  $\TRule{\False}{\lor}$ લાગુ કરીને શાખા વિસ્તારીએ:
  અભ્યાસપ્રશ્ન.\sourcecorrection{OLMD-054}{મૂળમાં આ ચિહ્નિત
  સૂત્રનો ઉપસર્ગ છૂટી ગયો છે; શાખાના અન્ય નિયમો સાથે સુસંગત
  રહે તે માટે $[\sigma]$ ઉમેર્યું છે.}
\item $\sFmla{\False}{!B \lif !C}[\sigma] \in \Gamma$ને
  $\TRule{\False}{\lif}$ લાગુ કરીને શાખા વિસ્તારીએ. તેનાથી
  શાખા પર બે નવાં ચિહ્નિત સૂત્રો મળે છે:
  $\sFmla{\True}{!B}[\sigma]$ અને
  $\sFmla{\False}{!C}[\sigma]$. ધારો કે
  $\mSat{M}{\Gamma}[f]$, ખાસ કરીને
  $\mSat/{M}{!B \lif !C}[f(\sigma)]$. ત્યારે
  $\mSat{M}{!B}[f(\sigma)]$ અને
  $\mSat/{M}{!C}[f(\sigma)]$. એટલે
  $\mModel{M}, f$ એ $\sFmla{\True}{!B}[\sigma]$ અને
  $\sFmla{\False}{!C}[\sigma]$ બંનેને સંતોષે છે.

%
\item $\sFmla{\True}{\Box !B}[\sigma] \in \Gamma$ને
  $\TRule{\True}{\Box}$ લાગુ કરીને શાખા વિસ્તારીએ:
  તેનાથી કોઈ
    $\sigma.n \in P(\Gamma)$ માટે નવું ચિહ્નિત સૂત્ર~$\sFmla{\True}{!B}[\sigma.n]$ શાખા પર મળે છે
    (કારણ કે $\sigma.n$ વપરાયેલો હોવો જોઈએ). ધારો કે
    $\mSat{M}{\Gamma}[f]$, ખાસ કરીને
    $\mSat{M}{\Box !B}[f(\sigma)]$. $f$ ઉપસર્ગોનું અર્થઘટન છે
    અને $\sigma$, $\sigma.n \in P(\Gamma)$ બંને હોવાથી
    $Rf(\sigma)f(\sigma.n)$ છે. તેથી
    $\mSat{M}{!B}[f(\sigma.n)]$; એટલે $\mModel{M}, f$ એ
    $\sFmla{\True}{!B}[\sigma.n]$ને સંતોષે છે.

\item $\sFmla{\False}{\Box !B}[\sigma] \in \Gamma$ને
  $\TRule{\False}{\Box}$ લાગુ કરીને શાખા વિસ્તારીએ:
  તેનાથી નવું ચિહ્નિત સૂત્ર~$\sFmla{\False}{!B}[\sigma.n]$ મળે છે, જ્યાં
    $\sigma.n$ શાખા પર નવો ઉપસર્ગ છે, એટલે
    $\sigma.n \notin P(\Gamma)$.
    \sourcecorrection{OLMD-055}{મૂળમાં આ નિયમનું નવું સૂત્ર $!A$
    સાથે લખાયું છે; પૂર્વપક્ષ $\Box !B$ અને આગળની સાબિતી મુજબ
    અહીં $!B$ હોવું જોઈએ.}
    $\Gamma$ સંતોષ્ય હોવાથી એવી $\Struct{M}$ અને
    $P(\Gamma)$નું અર્થઘટન~$f$ છે કે
    $\mSat{M}{\Gamma}[f]$, ખાસ કરીને
    $\mSat/{M}{\Box !B}[f(\sigma)]$.
    $\Gamma \cup \{\sFmla{\False}{!B}[\sigma.n]\}$ સંતોષ્ય છે
    તે બતાવવાનું છે. તે માટે $P(\Gamma)
    \cup \{\sigma.n\}$નું અર્થઘટન નીચે મુજબ વ્યાખ્યાયિત કરીએ:

  $\mSat/{M}{\Box !B}[f(\sigma)]$ હોવાથી એવું $w \in W$ છે કે
  $Rf(\sigma)w$ અને $\mSat/{M}{!B}[w]$. $f'$ એ $f$ જેવું જ
  રાખો, માત્ર $f'(\sigma.n) = w$. $f'(\sigma) = f(\sigma)$
  અને $Rf(\sigma)w$ હોવાથી $Rf'(\sigma)f'(\sigma.n)$;
  તેથી $f'$ એ $P(\Gamma) \cup \{\sigma.n\}$નું અર્થઘટન છે.
  સ્પષ્ટ છે કે $\mSat/{M}{!B}[f'(\sigma.n)]$.
  બધા ઉપસર્ગ $\sigma' \in P(\Gamma)$ માટે
  $f(\sigma') = f'(\sigma')$ હોવાથી $\mSat{M}{\Gamma}[f']$.
  આમ $\mModel{M}, f'$ એ $\Gamma \cup
  \{\sFmla{\False}{!B}[\sigma.n]\}$ને સંતોષે છે.

%
\item $\sFmla{\True}{\Diamond !B}[\sigma] \in \Gamma$ને
  $\TRule{\True}{\Diamond}$ લાગુ કરીને શાખા વિસ્તારીએ:
  તેનાથી નવું ચિહ્નિત સૂત્ર~$\sFmla{\True}{!B}[\sigma.n]$ મળે છે, જ્યાં
    $\sigma.n$ શાખા પર નવો ઉપસર્ગ છે, એટલે
    $\sigma.n \notin P(\Gamma)$.
    \sourcecorrection{OLMD-056}{મૂળમાં આ નિયમનું નવું સૂત્ર $!A$
    સાથે લખાયું છે; પૂર્વપક્ષ $\Diamond !B$ અને આગળની સાબિતી
    મુજબ અહીં $!B$ હોવું જોઈએ.}
    $\Gamma$ સંતોષ્ય હોવાથી એવી $\Struct{M}$ અને
    $P(\Gamma)$નું અર્થઘટન~$f$ છે કે
    $\mSat{M}{\Gamma}[f]$, ખાસ કરીને
    $\mSat{M}{\Diamond !B}[f(\sigma)]$.
    $\Gamma \cup \{\sFmla{\True}{!B}[\sigma.n]\}$ સંતોષ્ય છે
    તે બતાવવાનું છે. તે માટે $P(\Gamma)
    \cup \{\sigma.n\}$નું અર્થઘટન નીચે મુજબ વ્યાખ્યાયિત કરીએ:

  $\mSat{M}{\Diamond !B}[f(\sigma)]$ હોવાથી એવું $w \in W$
  છે કે $Rf(\sigma)w$ અને $\mSat{M}{!B}[w]$. $f'$ એ $f$
  જેવું જ રાખો, માત્ર $f'(\sigma.n) = w$.
  $f'(\sigma) = f(\sigma)$ અને $Rf(\sigma)w$ હોવાથી
  $Rf'(\sigma)f'(\sigma.n)$; તેથી $f'$ એ
  $P(\Gamma) \cup \{\sigma.n\}$નું અર્થઘટન છે. સ્પષ્ટ છે કે
  $\mSat{M}{!B}[f'(\sigma.n)]$. બધા ઉપસર્ગ
  $\sigma' \in P(\Gamma)$ માટે $f(\sigma') = f'(\sigma')$
  હોવાથી $\mSat{M}{\Gamma}[f']$. આમ $\mModel{M}, f'$ એ
  $\Gamma \cup \{\sFmla{\True}{!B}[\sigma.n]\}$ને સંતોષે છે.
\item $\sFmla{\False}{\Diamond !B}[\sigma] \in \Gamma$ને
  $\TRule{\False}{\Diamond}$ લાગુ કરીને શાખા વિસ્તારીએ:
  તેનાથી કોઈ
    $\sigma.n \in P(\Gamma)$ માટે નવું ચિહ્નિત સૂત્ર~$\sFmla{\False}{!B}[\sigma.n]$ શાખા પર મળે છે
    (કારણ કે $\sigma.n$ વપરાયેલો હોવો જોઈએ). ધારો કે
    $\mSat{M}{\Gamma}[f]$, ખાસ કરીને
    $\mSat/{M}{\Diamond !B}[f(\sigma)]$. $f$ ઉપસર્ગોનું
    અર્થઘટન છે અને $\sigma$, $\sigma.n \in P(\Gamma)$ બંને
    હોવાથી $Rf(\sigma)f(\sigma.n)$ છે. તેથી
    $\mSat/{M}{!B}[f(\sigma.n)]$; એટલે $\mModel{M}, f$ એ
    $\sFmla{\False}{!B}[\sigma.n]$ને સંતોષે છે.
\end{enumerate}
\sourcecorrection{OLMD-057}{મૂળમાં બૉક્સ અને ડાયમંડના
નવા-ઉપસર્ગ કિસ્સાઓમાં $\Sat{M}{\Gamma}[f]$ વપરાયું છે;
અહીં વ્યાખ્યાયિત મોડલ-સંતોષ $\mSat{M}{\Gamma}[f]$ રાખ્યું છે.}
હવે બે શાખા આપતા સંભવિત અનુમાનો ધ્યાનમાં લઈએ.
\sourcecorrection{OLMD-058}{મૂળમાં આ નિયમોને ``બે પૂર્વપક્ષવાળા''
કહ્યા છે; દરેક નિયમને એક પૂર્વપક્ષ છે અને તે શાખાને બે ભાગમાં
વહેંચે છે, તેથી શાખાઓની સંખ્યા સ્પષ્ટ કરી છે.}
\begin{enumerate}
\item $\sFmla{\False}{!B \land !C}[\sigma] \in \Gamma$ને
  $\TRule{\False}{\land}$ લાગુ કરીને શાખા વિસ્તારીએ. તેનાથી
  બે શાખા મળે છે: ડાબી શાખા $\sFmla{\False}{!B}[\sigma]$માંથી
  અને જમણી શાખા $\sFmla{\False}{!C}[\sigma]$માંથી આગળ વધે છે.
  ધારો કે $\mSat{M}{\Gamma}[f]$, ખાસ કરીને
  $\mSat/{M}{!B \land !C}[f(\sigma)]$. ત્યારે
  $\mSat/{M}{!B}[f(\sigma)]$ અથવા
  $\mSat/{M}{!C}[f(\sigma)]$. પહેલા કિસ્સામાં
  $\mModel{M}, f$ એ $\sFmla{\False}{!B}[\sigma]$ને સંતોષે છે;
  એટલે ડાબી શાખા સંતોષ્ય છે. બીજા કિસ્સામાં
  $\mModel{M}, f$ એ $\sFmla{\False}{!C}[\sigma]$ને સંતોષે છે;
  એટલે જમણી શાખા સંતોષ્ય છે.
\item $\sFmla{\True}{!B \lor !C}[\sigma] \in \Gamma$ને
    $\TRule{\True}{\lor}$ લાગુ કરીને શાખા વિસ્તારીએ:
    અભ્યાસપ્રશ્ન.
\item $\sFmla{\True}{!B \lif !C}[\sigma] \in \Gamma$ને
    $\TRule{\True}{\lif}$ લાગુ કરીને શાખા વિસ્તારીએ:
    અભ્યાસપ્રશ્ન.
\end{enumerate}
\end{proof}

\begin{prob}
\olref[nml][tab][sou]{thm:tableau-soundness}ની સાબિતી પૂરી કરો.
\end{prob}

\begin{cor}
\ollabel{cor:entailment-soundness}
જો $\Gamma \Proves !A$ હોય, તો $\Gamma \Entails !A$.
\end{cor}

\begin{proof}
  જો $\Gamma \Proves !A$ હોય, તો કેટલાંક $!B_1$, \dots,
  $!B_n \in \Gamma$ માટે
  $\Delta = \{\sFmla{\False}{!A}[1], \sFmla{\True}{!B_1}[1],
  \dots, \sFmla{\True}{!B_n}[1]\}$નું બંધ ટેબ્લો છે.
  બતાવવું છે કે $\Gamma \Entails !A$. એવું ન હોય તો કોઈ
  $\mModel{M}$ અને $w$ માટે બધા $i=1$, \dots,~$n$ પર
  $\mSat{M}{!B_i}[w]$ હોય, પરંતુ
  $\mSat/{M}{!A}[w]$. $f(1) = w$ રાખો. ત્યારે $f$ એ
  $P(\Delta)$નું~$\mModel{M}$માં અર્થઘટન છે અને
  $\mModel{M}$ એ~$f$ને અનુલક્ષીને~$\Delta$ને સંતોષે છે.
  પરંતુ \olref{thm:tableau-soundness} મુજબ બંધ ટેબ્લો
  હોવાથી $\Delta$ અસંતોષ્ય છે—વિરોધાભાસ. તેથી
  $\Gamma \Entails !A$.
  \sourcecorrection{OLMD-059}{મૂળની છેલ્લી પંક્તિમાં ફરી
  $\Gamma \Proves !A$ નિષ્કર્ષ લખ્યો છે; આ ઉપપ્રમેય માટે જરૂરી
  નિષ્કર્ષ $\Gamma \Entails !A$ પુનઃસ્થાપિત કર્યો છે.}
\end{proof}

\begin{cor}
\ollabel{cor:weak-soundness}
જો $\Proves !A$ હોય, તો $!A$ બધાં નિદર્શનોમાં સત્ય છે.
\end{cor}
% END OLP-0464
\endgroup

\begingroup
\makeatletter\def\ol@sectioncs{section}\makeatother

% BEGIN OLP-0465
\olfileid{nml}{tab}{mru}

\olsection{અન્ય સુલભતા-સંબંધો માટેના નિયમો}
\phantomsection\label{guunit:OLP-0465}


વિશિષ્ટ સુલભતા-સંબંધોથી નિર્ધારિત તર્કપદ્ધતિઓ માટે
\olref{tab:more-rules}ના વધારાના નિયમો લઈએ છીએ.

\begin{table}
  \begin{center}
    \def\arraystretch{3}\def\fCenter{}
    \begin{tabular}{|c|c|}
    \hline
          \AxiomC{\sFmla{\True}{\Box !A}[\sigma]}
      \RightLabel{T$\Box$}
      \UnaryInfC{\sFmla{\True}{!A}[\sigma]}
      \DisplayProof
    &
          \AxiomC{\sFmla{\False}{\Diamond !A}[\sigma]}
      \RightLabel{T$\Diamond$}
      \UnaryInfC{\sFmla{\False}{!A}[\sigma]}
      \DisplayProof
    \\[1ex]
    \hline
          \AxiomC{\sFmla{\True}{\Box !A}[\sigma]}
      \RightLabel{D$\Box$}
      \UnaryInfC{\sFmla{\True}{\Diamond!A}[\sigma]}
      \DisplayProof
    &
          \AxiomC{\sFmla{\False}{\Diamond !A}[\sigma]}
      \RightLabel{D$\Diamond$}
      \UnaryInfC{\sFmla{\False}{\Box!A}[\sigma]}
      \DisplayProof
    \\[1ex]
    \hline
          \AxiomC{\sFmla{\True}{\Box !A}[\sigma.n]}
      \RightLabel{B$\Box$}
      \UnaryInfC{\sFmla{\True}{!A}[\sigma]}
      \DisplayProof
    &
          \AxiomC{\sFmla{\False}{\Diamond !A}[\sigma.n]}
      \RightLabel{B$\Diamond$}
      \UnaryInfC{\sFmla{\False}{!A}[\sigma]}
      \DisplayProof
    \\[1ex]
    \hline
          \AxiomC{\sFmla{\True}{\Box !A}[\sigma]}
      \RightLabel{4$\Box$}
      \UnaryInfC{\sFmla{\True}{\Box!A}[\sigma.n]}
      \DisplayProof
    &
          \AxiomC{\sFmla{\False}{\Diamond !A}[\sigma]}
      \RightLabel{4$\Diamond$}
      \UnaryInfC{\sFmla{\False}{\Diamond!A}[\sigma.n]}
      \DisplayProof
    \\
    $\sigma.n$ વપરાયેલો છે & $\sigma.n$ વપરાયેલો છે
    \\[1ex]
    \hline
          \AxiomC{\sFmla{\True}{\Box !A}[\sigma.n]}
      \RightLabel{4r$\Box$}
      \UnaryInfC{\sFmla{\True}{\Box!A}[\sigma]}
      \DisplayProof
    &
          \AxiomC{\sFmla{\False}{\Diamond !A}[\sigma.n]}
      \RightLabel{4r$\Diamond$}
      \UnaryInfC{\sFmla{\False}{\Diamond!A}[\sigma]}
      \DisplayProof
    \\[1ex]
    \hline
    \end{tabular}
  \end{center}
  \caption{વધારાના મોડલ નિયમો.}
  \ollabel{tab:more-rules}
\end{table}

આ નિયમો ઉમેરવાથી \olref{tab:logics-rules}માં આપેલી
તર્કપદ્ધતિઓ માટે યથાર્થ અને પૂર્ણ એવી પ્રણાલીઓ મળે છે.

\begin{table}
  \begin{center}
    \begin{tabular}{lll}
      \hline
      તર્કપદ્ધતિ & $R$નો ગુણધર્મ & નિયમો\\
      \hline
      $\Log{T} = \Log{KT}$ & સ્વવાચક &
      T$\Box$%
      , %
      T$\Diamond$
      \\ \hline
      $\Log{D} = \Log{KD}$ & સિરિયલ &
      D$\Box$%
      , %
      D$\Diamond$
      \\ \hline
      $\Log{K4}$ & પરંપરિત &
      4$\Box$%
      , %
      4$\Diamond$
      \\ \hline
      $\Log{B} = \Log{KTB}$ & સ્વવાચક, &
      T$\Box$%
      , %
      T$\Diamond$\\
      & સંમિત &
      B$\Box$%
      , %
      B$\Diamond$
      \\ \hline
      $\Log{S4} = \Log{KT4}$ & સ્વવાચક, &
      T$\Box$%
      , %
      T$\Diamond$,\\
      & પરંપરિત &
      4$\Box$%
      , %
      4$\Diamond$
      \\ \hline
      $\Log{S5} = \Log{KT4B}$ & સ્વવાચક, &
      T$\Box$%
      , %
      T$\Diamond$,\\
      & પરંપરિત, &
      4$\Box$%
      , %
      4$\Diamond$,\\
      &  યુક્લિડિયન &
      4r$\Box$%
      , %
      4r$\Diamond$
      \\ \hline
    \end{tabular}
  \end{center}
  \caption{વિવિધ મોડલ તર્કપદ્ધતિઓ માટેના ટેબ્લો નિયમો.}
  \ollabel{tab:logics-rules}
\end{table}


\begin{ex}
  $\Log{S5} \Proves \Ax{5}$નો દાવો સાચો છે; નીચેનું બંધ ટેબ્લો
  $\Box!A \lif \Box\Diamond!A$ બતાવે છે.
  \sourcecorrection{OLMD-060}{મૂળમાં બીજા સૂત્રને જ \Ax{5} તરીકે
  ઓળખાવ્યું છે. પરંતુ \Ax{5}નું સૂત્ર
  $\Diamond!A \lif \Box\Diamond!A$ છે; પ્રદર્શિત વૃક્ષ
  $\Box!A \lif \Box\Diamond!A$ સાબિત કરે છે. તેથી બંને દાવા
  અલગ કરીને વૃક્ષનો ખરેખર નિષ્કર્ષ સ્પષ્ટ કર્યો છે.}
  \begin{oltableau}
    [\pFmla{\False}{\Box\formula{A} \lif \Box\Diamond \formula{A}}{1},
      just = \TAss
      [\pFmla{\True}{\Box \formula{A}}{1}, just = {\TRule{\False}{\lif}[1]}
        [\pFmla{\False}{\Box\Diamond \formula{A}}{1},
          just = {\TRule{\False}{\lif}[1]}
          [\pFmla{\False}{\Diamond \formula{A}}{1.1},
            just = {\TRule{\False}{\Box}[3]}
            [\pFmla{\False}{\Diamond \formula{A}}{1},
              just = {4r$\Diamond$ 4}
              [\pFmla{\False}{\formula{A}}{1.1},
                just = {\TRule{\False}{\Diamond}[5]}
                [\pFmla{\True}{\formula{A}}{1.1},
                  just = {\TRule{\True}{\Box}[2]}, close]
              ]
            ]
          ]
        ]
      ]
    ]
  \end{oltableau}
\end{ex}

\begin{prob}
નીચેના દાવા બતાવતાં બંધ ટેબ્લો આપો:
  \begin{enumerate}
  \item $\Log{KT5} \Proves \Ax{B}$;
  \item $\Log{KT5} \Proves \Ax{4}$;
  \item $\Log{KDB4} \Proves \Ax{T}$;
  \item $\Log{KB4} \Proves \Ax{5}$;
  \item $\Log{KB5} \Proves \Ax{4}$;
  \item $\Log{KT} \Proves \Ax{D}$.
  \end{enumerate}
\end{prob}
% END OLP-0465
\endgroup

\begingroup
\makeatletter\def\ol@sectioncs{section}\makeatother

% BEGIN OLP-0466
\olfileid{nml}{tab}{msn}
      
\olsection{વધારાના નિયમોની યથાર્થતા}
\phantomsection\label{guunit:OLP-0466}


નિયમને નિદર્શનોના કોઈ વર્ગ માટે યથાર્થ ત્યારે કહીએ જ્યારે
ટેબ્લોની કોઈ શાખા તે વર્ગના નિદર્શનમાં સંતોષ્ય હોય ત્યારે
નિયમ લાગુ કરવાથી મળતી શાખા પણ તે વર્ગના નિદર્શનમાં સંતોષ્ય હોય.

\begin{prop}
  \ollabel{prop:soundness-T} 
  \Ax{T}$\Box$%
   અને %
  \Ax{T}$\Diamond$
  નિયમો
  સ્વવાચક નિદર્શનો માટે યથાર્થ છે.
\end{prop}

\begin{proof}
\begin{enumerate}
\item $\sFmla{\True}{\Box !B}[\sigma] \in \Gamma$ને
  T$\Box$ લાગુ કરીને શાખા વિસ્તારીએ:
  તેનાથી શાખા પર નવું ચિહ્નિત સૂત્ર~$\sFmla{\True}{!B}[\sigma]$ મળે છે. ધારો કે
    $\mSat{M}{\Gamma}[f]$, ખાસ કરીને $\mSat{M}{\Box
      !B}[f(\sigma)]$. $R$ સ્વવાચક હોવાથી
    $Rf(\sigma)f(\sigma)$. તેથી $\mSat{M}{!B}[f(\sigma)]$;
    એટલે $\mModel{M}, f$ એ $\sFmla{\True}{!B}[\sigma]$ને
    સંતોષે છે.
\item $\sFmla{\False}{\Diamond !B}[\sigma]
  \in \Gamma$ને T$\Diamond$ લાગુ કરીને શાખા વિસ્તારીએ:
  તેનાથી શાખા પર નવું ચિહ્નિત સૂત્ર~$\sFmla{\False}{!B}[\sigma]$ મળે છે. ધારો કે
    $\mSat{M}{\Gamma}[f]$, ખાસ કરીને $\mSat/{M}{\Diamond
      !B}[f(\sigma)]$. $R$ સ્વવાચક હોવાથી
    $Rf(\sigma)f(\sigma)$. તેથી $\mSat/{M}{!B}[f(\sigma)]$;
    એટલે $\mModel{M}, f$ એ $\sFmla{\False}{!B}[\sigma]$ને
    સંતોષે છે.
\end{enumerate}
\end{proof}



\begin{prop}
  \ollabel{prop:soundness-D} 
  \Ax{D}$\Box$%
   અને %
  \Ax{D}$\Diamond$
  નિયમો
  સિરિયલ નિદર્શનો માટે યથાર્થ છે.
\end{prop}

\begin{proof}
\begin{enumerate}
\item D$\Box$ને
  $\sFmla{\True}{\Box !B}[\sigma] \in \Gamma$ પર લાગુ કરીને
  શાખા વિસ્તારીએ:
  તેનાથી શાખા પર નવું ચિહ્નિત સૂત્ર~$\sFmla{\True}{\Diamond !B}[\sigma]$ મળે છે.
    ધારો કે $\mSat{M}{\Gamma}[f]$, ખાસ કરીને
    $\mSat{M}{\Box !B}[f(\sigma)]$. $R$ સિરિયલ હોવાથી એવું
    $w \in W$ છે કે $Rf(\sigma)w$. ત્યારે $\mSat{M}{!B}[w]$;
    તેથી $\mSat{M}{\Diamond !B}[f(\sigma)]$.
    આમ $\mModel{M}, f$ એ
    $\sFmla{\True}{\Diamond!B}[\sigma]$ને સંતોષે છે.
\item D$\Diamond$ને
  $\sFmla{\False}{\Diamond !B}[\sigma] \in \Gamma$ પર લાગુ
  કરીને શાખા વિસ્તારીએ:
  તેનાથી શાખા પર નવું ચિહ્નિત સૂત્ર~$\sFmla{\False}{\Box !B}[\sigma]$ મળે છે.
    ધારો કે $\mSat{M}{\Gamma}[f]$, ખાસ કરીને
    $\mSat/{M}{\Diamond !B}[f(\sigma)]$. $R$ સિરિયલ હોવાથી
    એવું $w \in W$ છે કે $Rf(\sigma)w$. ત્યારે
    $\mSat/{M}{!B}[w]$; તેથી
    $\mSat/{M}{\Box !B}[f(\sigma)]$. આમ $\mModel{M}, f$ એ
    $\sFmla{\False}{\Box!B}[\sigma]$ને સંતોષે છે.
\end{enumerate}
\end{proof}



\begin{prop}
  \ollabel{prop:soundness-B} 
  \Ax{B}$\Box$%
   અને %
  \Ax{B}$\Diamond$
  નિયમો
  સંમિત નિદર્શનો માટે યથાર્થ છે.
\end{prop}

\begin{proof}
\begin{enumerate}
\item B$\Box$ને
  $\sFmla{\True}{\Box !B}[\sigma.n] \in \Gamma$ પર લાગુ કરીને
  શાખા વિસ્તારીએ:
  તેનાથી શાખા પર નવું ચિહ્નિત સૂત્ર~$\sFmla{\True}{!B}[\sigma]$ મળે છે. ધારો કે
    $\mSat{M}{\Gamma}[f]$, ખાસ કરીને $\mSat{M}{\Box
      !B}[f(\sigma.n)]$. $f$ એ શાખાના ઉપસર્ગોનું
    $\mModel{M}$માં અર્થઘટન હોવાથી $Rf(\sigma)f(\sigma.n)$.
    $R$ સંમિત હોવાથી $Rf(\sigma.n)f(\sigma)$ પણ છે.
    $\mSat{M}{\Box !B}[f(\sigma.n)]$ હોવાથી
    $\mSat{M}{!B}[f(\sigma)]$. તેથી $\mModel{M}, f$ એ
    $\sFmla{\True}{!B}[\sigma]$ને સંતોષે છે.
\item B$\Diamond$ને
  $\sFmla{\False}{\Diamond !B}[\sigma.n] \in \Gamma$ પર લાગુ
  કરીને શાખા વિસ્તારીએ:
  તેનાથી શાખા પર નવું ચિહ્નિત સૂત્ર~$\sFmla{\False}{!B}[\sigma]$ મળે છે. ધારો કે
    $\mSat{M}{\Gamma}[f]$, ખાસ કરીને $\mSat/{M}{\Diamond
      !B}[f(\sigma.n)]$. $f$ એ શાખાના ઉપસર્ગોનું
    $\mModel{M}$માં અર્થઘટન હોવાથી $Rf(\sigma)f(\sigma.n)$.
    $R$ સંમિત હોવાથી $Rf(\sigma.n)f(\sigma)$ પણ છે.
    $\mSat/{M}{\Diamond !B}[f(\sigma.n)]$ હોવાથી
    $\mSat/{M}{!B}[f(\sigma)]$. તેથી $\mModel{M}, f$ એ
    $\sFmla{\False}{!B}[\sigma]$ને સંતોષે છે.
\end{enumerate}
\end{proof}



\begin{prop}
  \ollabel{prop:soundness-4} 
  \Ax{4}$\Box$%
   અને %
  \Ax{4}$\Diamond$
  નિયમો
  પરંપરિત નિદર્શનો માટે યથાર્થ છે.
\end{prop}

\begin{proof}
\begin{enumerate}
\item 4$\Box$ને
  $\sFmla{\True}{\Box !B}[\sigma] \in \Gamma$ પર લાગુ કરીને
  શાખા વિસ્તારીએ:
  તેનાથી શાખા પર નવું ચિહ્નિત સૂત્ર~$\sFmla{\True}{\Box!B}[\sigma.n]$ મળે છે.
    ધારો કે $\mSat{M}{\Gamma}[f]$, ખાસ કરીને
    $\mSat{M}{\Box !B}[f(\sigma)]$. $f$ એ શાખાના
    ઉપસર્ગોનું~$\mModel{M}$માં અર્થઘટન છે અને $\sigma.n$
    વપરાયેલો હોવો જોઈએ, તેથી $Rf(\sigma)f(\sigma.n)$.
    હવે $Rf(\sigma.n)w$ હોય એવું કોઈપણ જગત~$w$ લો.
    $R$ પરંપરિત હોવાથી $Rf(\sigma)w$.
    $\mSat{M}{\Box !B}[f(\sigma)]$ હોવાથી
    $\mSat{M}{!B}[w]$. તેથી
    $\mSat{M}{\Box !B}[f(\sigma.n)]$ અને $\mModel{M}, f$ એ
    $\sFmla{\True}{\Box !B}[\sigma.n]$ને સંતોષે છે.
\item 4$\Diamond$ને
  $\sFmla{\False}{\Diamond !B}[\sigma] \in \Gamma$ પર લાગુ
  કરીને શાખા વિસ્તારીએ:
  તેનાથી શાખા પર નવું ચિહ્નિત સૂત્ર~$\sFmla{\False}{\Diamond!B}[\sigma.n]$ મળે છે.
    ધારો કે $\mSat{M}{\Gamma}[f]$, ખાસ કરીને
    $\mSat/{M}{\Diamond !B}[f(\sigma)]$. $f$ એ શાખાના
    ઉપસર્ગોનું~$\mModel{M}$માં અર્થઘટન છે અને $\sigma.n$
    વપરાયેલો હોવો જોઈએ, તેથી $Rf(\sigma)f(\sigma.n)$.
    હવે $Rf(\sigma.n)w$ હોય એવું કોઈપણ જગત~$w$ લો.
    $R$ પરંપરિત હોવાથી $Rf(\sigma)w$.
    $\mSat/{M}{\Diamond !B}[f(\sigma)]$ હોવાથી
    $\mSat/{M}{!B}[w]$. તેથી
    $\mSat/{M}{\Diamond !B}[f(\sigma.n)]$ અને $\mModel{M}, f$ એ
    $\sFmla{\False}{\Diamond !B}[\sigma.n]$ને સંતોષે છે.
\end{enumerate}
\end{proof}



\begin{prop}
  \ollabel{prop:soundness-4r} 
  \Ax{4r}$\Box$%
   અને %
  \Ax{4r}$\Diamond$
  નિયમો
  યુક્લિડિયન નિદર્શનો માટે યથાર્થ છે.
\end{prop}

\begin{proof}
\sourcecorrection{OLMD-061}{મૂળમાં બંને યુક્લિડિયન કિસ્સામાં
$f(\sigma).n$ લખાયું છે, જે વિધેયના મૂલ્યને ઉપસર્ગ જેવું ગણે છે;
પૂર્વપક્ષ મુજબ અહીં $f(\sigma.n)$ જરૂરી છે.}
\begin{enumerate}
\item 4r$\Box$ને
  $\sFmla{\True}{\Box !B}[\sigma.n] \in \Gamma$ પર લાગુ કરીને
  શાખા વિસ્તારીએ:
  તેનાથી શાખા પર નવું ચિહ્નિત સૂત્ર~$\sFmla{\True}{\Box!B}[\sigma]$ મળે છે. ધારો કે
    $\mSat{M}{\Gamma}[f]$, ખાસ કરીને $\mSat{M}{\Box
      !B}[f(\sigma.n)]$. $f$ એ શાખાના ઉપસર્ગોનું
    $\mModel{M}$માં અર્થઘટન હોવાથી $Rf(\sigma)f(\sigma.n)$.
    હવે $Rf(\sigma)w$ હોય એવું કોઈપણ જગત~$w$ લો.
    $R$ યુક્લિડિયન હોવાથી $Rf(\sigma.n)w$.
    $\mSat{M}{\Box !B}[f(\sigma.n)]$ હોવાથી
    $\mSat{M}{!B}[w]$. તેથી $\mSat{M}{\Box !B}[f(\sigma)]$;
    એટલે $\mModel{M}, f$ એ
    $\sFmla{\True}{\Box !B}[\sigma]$ને સંતોષે છે.
\item 4r$\Diamond$ને
  $\sFmla{\False}{\Diamond !B}[\sigma.n] \in \Gamma$ પર લાગુ
  કરીને શાખા વિસ્તારીએ:
  તેનાથી શાખા પર નવું ચિહ્નિત સૂત્ર~$\sFmla{\False}{\Diamond!B}[\sigma]$ મળે છે.
    \sourcecorrection{OLMD-062}{મૂળમાં ડાયમંડ-નિયમનો નિષ્કર્ષ
    $\sFmla{\True}{\Box!B}[\sigma]$ લખાયો છે; નિયમ-કોષ્ટક અને
    આ જ પ્રમાણના અંતિમ વાક્ય મુજબ તે
    $\sFmla{\False}{\Diamond!B}[\sigma]$ હોવો જોઈએ.}
    ધારો કે $\mSat{M}{\Gamma}[f]$, ખાસ કરીને
    $\mSat/{M}{\Diamond !B}[f(\sigma.n)]$.
    $f$ એ શાખાના ઉપસર્ગોનું $\mModel{M}$માં અર્થઘટન હોવાથી
    $Rf(\sigma)f(\sigma.n)$. હવે $Rf(\sigma)w$ હોય એવું
    કોઈપણ જગત~$w$ લો. $R$ યુક્લિડિયન હોવાથી
    $Rf(\sigma.n)w$. $\mSat/{M}{\Diamond !B}[f(\sigma.n)]$
    હોવાથી $\mSat/{M}{!B}[w]$. તેથી
    $\mSat/{M}{\Diamond !B}[f(\sigma)]$; એટલે
    $\mModel{M}, f$ એ
    $\sFmla{\False}{\Diamond !B}[\sigma]$ને સંતોષે છે.
\end{enumerate}
\end{proof}



\begin{cor}
\ollabel{cor:soundness-logics} \olref[mru]{tab:logics-rules}માં
આપેલી ટેબ્લો પ્રણાલીઓ નિદર્શનોના અનુરૂપ વર્ગો માટે
યથાર્થ છે.
\end{cor}
% END OLP-0466
\endgroup

\begingroup
\makeatletter\def\ol@sectioncs{section}\makeatother

% BEGIN OLP-0467
\olfileid{nml}{tab}{s5}

\olsection{\Log{S5} માટે સરળ ટેબ્લો}
\phantomsection\label{guunit:OLP-0467}


\Log{S5} સાર્વત્રિક નિદર્શનોના વર્ગને અનુલક્ષીને યથાર્થ અને પૂર્ણ
છે, એટલે એવા નિદર્શનો જેમાં દરેક જગત દરેક જગતમાંથી સુલભ
હોય છે. સાર્વત્રિક નિદર્શનમાં સુલભતા-સંબંધથી કોઈ વધારાનો
ભેદ પડતો નથી: ``એવું જગત~$w$ છે જ્યાં
$\mSat{M}{!A}[w]$'' ત્યારે અને માત્ર ત્યારે સાચું છે જ્યારે
એવું જગત~$w$~$u$થી સુલભ હોય. તેથી \Log{S5}માં નિદર્શનને
જગતોના ગણ અને નિમણૂક~$V$ વડે જ વ્યાખ્યાયિત કરી શકીએ.
આથી ટેબ્લોના નિયમો પણ સરળ કરી શકાય એવું સૂચન મળે છે.
સામાન્ય કિસ્સામાં ઉપસર્ગ તરીકે ધન પૂર્ણાંકોની શ્રેણીઓ લઈએ
છીએ, જેથી કયા ઉપસર્ગ બીજા કયા ઉપસર્ગથી સુલભ જગતોને નામ આપે
છે તેનો હિસાબ રાખી શકાય: $\sigma.n$ એ~$\sigma$થી સુલભ
જગતનું નામ આપે છે. પરંતુ \Log{S5}માં દરેક જગત બીજા દરેક
જગતથી સુલભ છે, તેથી આવો હિસાબ રાખવો જરૂરી નથી. તેના
બદલે ધન પૂર્ણાંકોને જ ઉપસર્ગ તરીકે વાપરી શકીએ. સરળ બનાવેલા
નિયમો \olref{tab:rules-S5}માં આપ્યા છે.

\begin{table}
  \begin{center}
    \def\arraystretch{3}\def\fCenter{}
    \begin{tabular}{|c|c|}
    \hline
          \AxiomC{\sFmla{\True}{\Box !A}[n]}
      \RightLabel{\TRule{\True}{\Box}}
      \UnaryInfC{\sFmla{\True}{!A}[m]}
      \DisplayProof
      &
      \AxiomC{\sFmla{\False}{\Box !A}[n]}
      \RightLabel{\TRule{\False}{\Box}}
      \UnaryInfC{\sFmla{\False}{!A}[m]}
      \DisplayProof\\
      $m$ વપરાયેલો છે & $m$ નવો છે
      \\[1ex]
      \hline
          \AxiomC{\sFmla{\True}{\Diamond !A}[n]}
      \RightLabel{\TRule{\True}{\Diamond}}
      \UnaryInfC{\sFmla{\True}{!A}[m]}
      \DisplayProof
      &
      \AxiomC{\sFmla{\False}{\Diamond !A}[n]}
      \RightLabel{\TRule{\False}{\Diamond}}
      \UnaryInfC{\sFmla{\False}{!A}[m]}
      \DisplayProof\\
      $m$ નવો છે & $m$ વપરાયેલો છે
      \\[1ex]
      \hline
    \end{tabular}
  \end{center}
  \caption{\Log{S5} માટેના સરળ બનાવેલા નિયમો.}
  \ollabel{tab:rules-S5}
\end{table}

\begin{ex}
  $\Log{S5} \Proves \Ax{5}$, એટલે કે,
  $\Diamond!A \lif \Box\Diamond!A$ બતાવતું સરળ બનાવેલું બંધ
  ટેબ્લો આપીએ છીએ.
  \begin{oltableau}
    [\pFmla{\False}{\Diamond\formula{A} \lif \Box\Diamond \formula{A}}{1},
      just = \TAss
      [\pFmla{\True}{\Diamond \formula{A}}{1}, just = {\TRule{\False}{\lif}[1]}
        [\pFmla{\False}{\Box\Diamond \formula{A}}{1},
          just = {\TRule{\False}{\lif}[1]}
          [\pFmla{\False}{\Diamond \formula{A}}{2},
            just = {\TRule{\False}{\Box}[3]}
            [\pFmla{\True}{\formula{A}}{3},
              just = {\TRule{\True}{\Diamond}[2]}
                [\pFmla{\False}{\formula{A}}{3},
                  just = {\TRule{\False}{\Diamond}[4]}, close]
            ]
          ]
        ]
      ]
    ]
  \end{oltableau}
\end{ex}
% END OLP-0467
\endgroup

\begingroup
\makeatletter\def\ol@sectioncs{section}\makeatother

% BEGIN OLP-0468
\olfileid{nml}{tab}{cpl}

\olsection{\Log{K} માટેની પૂર્ણતા}
\phantomsection\label{guunit:OLP-0468}


\begin{explain}
  ટેબ્લોની રીત પૂર્ણ છે તે બતાવવા સાબિત કરવું પડે કે
  $\Gamma \Proves !A$ બતાવતું કોઈ બંધ ટેબ્લો ન હોય ત્યારે
  $\Gamma \Entails/ !A$, એટલે પ્રતિવિધાન હોય. પરંતુ ``બંધ
  ટેબ્લો નથી''નો અર્થ એવો છે કે તેને બનાવવાના દરેક પ્રયાસમાં
  શાખા બંધ થવામાં નિષ્ફળતા મળે. યુક્તિ એ છે કે બધા પ્રયાસ નિષ્ફળ
  જાય તો એક નિશ્ચિત \emph{પદ્ધતિસર અને સર્વગ્રાહી} પ્રયાસ પણ
  નિષ્ફળ જાય; અને બંધ ટેબ્લો અસ્તિત્વમાં હોય તો આ
  પદ્ધતિસરનો પ્રયાસ શાખાઓ બંધ કરી દે. આ એક ટેબ્લોની ખુલ્લી
  શાખાઓમાં પ્રતિવિધાન વ્યાખ્યાયિત કરવા જરૂરી બધી માહિતી હશે.
  તેની કોઈ ખુલ્લી શાખાથી મળતા પ્રતિવિધાનનાં જગતો તે શાખા પર
  વપરાયેલા બધા ઉપસર્ગો હશે; અને પ્રતિજ્ઞાપન ચલ~$p$ એ~$\sigma$ પર ત્યારે અને માત્ર ત્યારે સત્ય હશે
  જ્યારે $\sFmla{\True}{p}[\sigma]$ એ શાખા પર આવે.
\end{explain}

\begin{defn}
  ટેબ્લોની શાખાને \emph{પૂર્ણ} ત્યારે કહીએ જ્યારે તેમાં
  નિયમ લાગુ કરી શકાય એવું ઉપસર્ગવાળું સૂત્ર
  $\sFmla{S}{!A}[\sigma]$ હોય ત્યારે તેમાં આગળનું પણ હોય:
  \begin{enumerate}
    \item વિધાનસંબંધી એક-શાખી નિયમોમાં તેના અનુરૂપ નિષ્કર્ષરૂપ
      ઉપસર્ગવાળાં સૂત્રો;
    \item વિધાનસંબંધી શાખા-વિભાજક નિયમોમાં અનુરૂપ નિષ્કર્ષના
      સૂત્રોમાંથી ઓછામાં ઓછું એક;
    \item નવો ઉપસર્ગ માગતા મોડલ નિયમોમાં શક્ય નિષ્કર્ષોમાંથી
      ઓછામાં ઓછું એક;
    \item વપરાયેલો ઉપસર્ગ માગતા મોડલ નિયમોમાં શાખા પર આવતા
      દરેક અનુરૂપ ઉપસર્ગ માટેનો નિષ્કર્ષ.
  \end{enumerate}
\end{defn}

\begin{explain}
દાખલા તરીકે, પૂર્ણ શાખામાં $\sFmla{\True}{!B \land !C}[\sigma]$
હોય ત્યારે તેમાં $\sFmla{\True}{!B}[\sigma]$ અને
$\sFmla{\True}{!C}[\sigma]$ બંને હોય. તેમાં
$\sFmla{\True}{!B \lor !C}[\sigma]$ હોય તો
$\sFmla{\True}{!B}[\sigma]$ અને
$\sFmla{\True}{!C}[\sigma]$માંથી ઓછામાં ઓછું એક હોય.
\sourcecorrection{OLMD-063}{મૂળના આ બે વિધાનસંબંધી ઉદાહરણમાં
પ્રથમ સૂત્રનો ઉપસર્ગ છૂટી ગયો છે અને વિયોજનના પહેલા
નિષ્કર્ષને અસત્ય ચિહ્ન અપાયું છે; નિયમો મુજબ ઉપસર્ગ અને સત્ય
ચિહ્ન પુનઃસ્થાપિત કર્યાં છે.}
તેમાં $\sFmla{\False}{\Box !B}[\sigma]$ હોય તો કોઈ એક~$n$
માટે તેમાં $\sFmla{\False}{!B}[\sigma.n]$ પણ હોય. અને તેમાં
$\sFmla{\True}{\Box !B}[\sigma]$ હોય ત્યારે,
$\sigma.n$ શાખા પર વપરાયેલો હોય એવા દરેક~$n$ માટે તેમાં
$\sFmla{\True}{!B}[\sigma.n]$ પણ હોય.
\sourcecorrection{OLMD-064}{મૂળના ચાર મોડલ ઉદાહરણમાં
$\Box/\Diamond$નો સૂત્ર-આર્ગ્યુમેન્ટ ખૂટે છે; અનુગામી ઉપસર્ગે
નિયમનો નિષ્કર્ષ મૂળના મોડલ સૂત્રને બદલે $!B$ માટે હોવો જોઈએ.
અહીં પૂર્વપક્ષમાં $!B$ અને અનુગામી નિષ્કર્ષમાં ચિહ્નિત $!B$
રાખ્યાં છે.}
\end{explain}

\begin{prop}\ollabel{prop:complete-tableau}
  દરેક સાંત $\Gamma$ માટે એવું ટેબ્લો છે કે તેની દરેક
  ખુલ્લી શાખા પૂર્ણ છે.
\end{prop}

\begin{proof}
  $\Gamma$ માટેના ટેબ્લોમાં કોઈ ખુલ્લી શાખા લો. તેના
  દરેક સાંત તબક્કે નિયમ લાગુ કરી શકાય એવાં ઉપસર્ગવાળાં
  સૂત્રો સાંત જ હોય છે. નિશ્ચિત ક્રમમાં (જેમ કે ઉપરથી
  નીચે), જેના માટે (1)--(4)ની શરતો હજી પૂરી ન થતી હોય એવા
  દરેક ઉપસર્ગવાળા સૂત્ર પર લાગુ પડતા નિયમો ચલાવી
  શાખા વિસ્તારો. ક્યારેક શાખા વહેંચાશે; ત્યારે બાકીનાં
  ઉપસર્ગવાળાં સૂત્રો માટે મળતી દરેક શાખાના છેડે
  નિયમો ચલાવો. સૂત્રો અને વપરાયેલા ઉપસર્ગો બંને સાંત હોવાથી
  આ ફેરી અંતે મૂળ શાખાને વિસ્તરતી સાંત સંખ્યામાં શાખાઓ
  મળે છે. દરેક ખુલ્લી શાખા પર પ્રક્રિયા ફરી કરો અને ફેરીઓ
  ચાલુ રાખો. આ તબક્કાઓના સંઘરૂપ સીમા-ટેબ્લોમાં દરેક
  ખુલ્લી શાખા પૂર્ણ છે: શાખા પર આવેલું દરેક સૂત્ર અને દરેક
  વપરાયેલો અનુગામી ઉપસર્ગ કોઈ સાંત તબક્કે હાજર હોય છે,
  તેથી પછીની ફેરીમાં તેની બાકી શરત પૂરી થાય છે.
  \sourcecorrection{OLMD-065}{મૂળની પ્રતિજ્ઞા બધી શાખાઓને
  પૂર્ણ કહે છે, પરંતુ આપેલું નિર્માણ માત્ર ખુલ્લી શાખાઓ
  વિસ્તારે છે; અંતે ``દરેક શાખા બંધ છે'' પણ લખાયું છે,
  જે પૂર્ણતાના હેતુથી વિરુદ્ધ છે. દાવો ખુલ્લી શાખાઓ માટે
  સ્પષ્ટ કર્યો અને ન્યાયી પુનરાવર્તનનું સીમા-વૃક્ષ ઉમેર્યું છે.}
\end{proof}

\begin{thm}[પૂર્ણતા]
  \ollabel{thm:tableau-completeness}
  $\Gamma$ માટે કોઈ બંધ ટેબ્લો ન હોય, તો $\Gamma$ સંતોષ્ય છે.
\end{thm}

\begin{proof}
ઉપરની પ્રતિજ્ઞા સાંત ધારણાઓ માટે છે. અહીં ધારણાઓનો ગણ
અનંત પણ હોઈ શકે. એવો ટેબ્લો મેળવવા, ભાષા ગણનીય
હોવાથી ધારણાઓને ક્રમમાં
ગોઠવો. દરેક સાંત તબક્કે આગળની ધારણા દરેક ખુલ્લી શાખામાં
ઉમેરો અને ત્યાં સુધી ઊભી થયેલી નિયમ-શરતોને ન્યાયી ક્રમમાં
પૂરી કરો. બંધ ટેબ્લો અસ્તિત્વમાં ન હોવાથી કોઈ સાંત
તબક્કે બધી શાખાઓ બંધ થઈ શકતી નથી. દરેક તબક્કે શાખાઓની
સંખ્યા સાંત હોવાથી, સતત ખુલ્લા વિસ્તરણો ધરાવતી શાખા
પસંદ કરતાં એક ખુલ્લી પૂર્ણ શાખા મળે છે.
\sourcecorrection{OLMD-066}{મૂળની પ્રતિજ્ઞા માત્ર સાંત
ધારણાઓ માટે છે, પરંતુ પ્રમેયના પ્રમાણમાં તેને મનસ્વી
ધારણાગણ પર સીધી લાગુ કરી છે. ગણનીય ભાષામાં ધારણાઓનું
ન્યાયી ક્રમમાં ઉમેરણ અને સાંત-શાખાવાળા તબક્કાઓમાંથી
અનંત ખુલ્લી શાખાની પસંદગી સ્પષ્ટ કરી છે.}
ધારણાઓનો ગણ ખાલી હોય તો એક-જગતનું કોઈપણ નિદર્શન તેને
સંતોષે છે; તેથી આગળ અખાલી કિસ્સો લો.
\sourcecorrection{OLMD-067}{મૂળના નિદર્શન-નિર્માણમાં
જગતોનો ગણ ઉપસર્ગોનો ગણ છે; ખાલી ધારણાગણ માટે તે ખાલી
પડી શકે, જ્યારે નિદર્શનનો જગત-ગણ અખાલી હોવો જોઈએ.
ખાલી કિસ્સો અલગથી નિકાલ કર્યો છે.}
આ રીતે $\Gamma$ માટે એવું ટેબ્લો મળે છે જેમાં
ઓછામાં ઓછી એક શાખા ખુલ્લી અને પૂર્ણ છે. તે શાખા પરનાં
ઉપસર્ગવાળાં સૂત્રોનો ગણ $\Delta$ લો, અને તેમાં
આવતા ઉપસર્ગોનો ગણ $P(\Delta)$ લો.

નિદર્શન~$\mModel{M(\Delta)} = \tuple{P(\Delta), R, V}$
વ્યાખ્યાયિત કરીએ છીએ. તેમાં~$\Delta$માં આવતા ઉપસર્ગો
જગતો છે અને સુલભતા-સંબંધ આ રીતે અપાય છે:
\[
R\sigma\sigma' \quad \text{જો અને માત્ર જો} \quad
\sigma'=\sigma.n \quad \text{કોઈ~$n$ માટે}
\]
અને
\[
V(p) = \Setabs{\sigma}{\sFmla{\True}{p}[\sigma] \in \Delta}.
\]
હવે~$!A$ પર આગમનથી બતાવીએ છીએ કે
$\sFmla{\True}{!A}[\sigma] \in \Delta$ હોય તો
$\mSat{M(\Delta)}{!A}[\sigma]$, અને
$\sFmla{\False}{!A}[\sigma] \in \Delta$ હોય તો
$\mSat/{M(\Delta)}{!A}[\sigma]$.

\begin{enumerate}
  \item \indcase{!A}{p}{$\sFmla{\True}{\indfrm}[\sigma] \in \Delta$
    હોય તો~$V$ની વ્યાખ્યા મુજબ $\sigma \in V(p)$, તેથી
    $\mSat{M(\Delta)}{\indfrm}[\sigma]$.

    $\sFmla{\False}{\indfrm}[\sigma] \in \Delta$ હોય તો
    $\sFmla{\True}{\indfrm}[\sigma] \notin \Delta$, કારણ કે
    અન્યથા શાખા બંધ હોત. તેથી $\sigma \notin V(p)$ અને
    $\mSat/{M(\Delta)}{\indfrm}[\sigma]$.}
  \item \indcase{!A}{\lnot !B}{
      $\sFmla{\True}{\indfrm}[\sigma] \in \Delta$ હોય તો
      શાખા પૂર્ણ હોવાથી $\sFmla{\False}{!B}[\sigma] \in \Delta$.
      આગમન-ધારણા મુજબ
      $\mSat/{M(\Delta)}{!B}[\sigma]$, તેથી
      $\mSat{M(\Delta)}{\indfrm}[\sigma]$.

      $\sFmla{\False}{\indfrm}[\sigma] \in \Delta$ હોય તો
      શાખા પૂર્ણ હોવાથી $\sFmla{\True}{!B}[\sigma] \in \Delta$.
      આગમન-ધારણા મુજબ
      $\mSat{M(\Delta)}{!B}[\sigma]$, તેથી
      $\mSat/{M(\Delta)}{\indfrm}[\sigma]$.}
  \item \indcase{!A}{!B \land !C}{
      $\sFmla{\True}{\indfrm}[\sigma] \in \Delta$ હોય તો
      શાખા પૂર્ણ હોવાથી $\sFmla{\True}{!B}[\sigma] \in \Delta$
      અને $\sFmla{\True}{!C}[\sigma] \in \Delta$ બંને હોય.
      આગમન-ધારણા મુજબ $\mSat{M(\Delta)}{!B}[\sigma]$ અને
      $\mSat{M(\Delta)}{!C}[\sigma]$; તેથી
      $\mSat{M(\Delta)}{\indfrm}[\sigma]$.

      $\sFmla{\False}{\indfrm}[\sigma] \in \Delta$ હોય તો
      શાખા પૂર્ણ હોવાથી $\sFmla{\False}{!B}[\sigma] \in \Delta$
      અથવા $\sFmla{\False}{!C}[\sigma] \in \Delta$ હોય.
      આગમન-ધારણા મુજબ
      $\mSat/{M(\Delta)}{!B}[\sigma]$ અથવા
      $\mSat/{M(\Delta)}{!C}[\sigma]$; તેથી
      $\mSat/{M(\Delta)}{\indfrm}[\sigma]$.}
  \sourcecorrection{OLMD-068}{મૂળમાં અસત્ય સંયોજનના
    આગમન-નિષ્કર્ષમાં બીજો ભાગ ફરી $!B$ લખાયો છે;
    પૂર્વપક્ષ અને શાખા-નિયમ મુજબ ત્યાં $!C$ રાખ્યું છે.}
  \item \indcase{!A}{!B \lor !C}{
      $\sFmla{\True}{\indfrm}[\sigma] \in \Delta$ હોય તો
      શાખા પૂર્ણ હોવાથી $\sFmla{\True}{!B}[\sigma] \in \Delta$
      અથવા $\sFmla{\True}{!C}[\sigma] \in \Delta$ હોય.
      આગમન-ધારણા મુજબ
      $\mSat{M(\Delta)}{!B}[\sigma]$ અથવા
      $\mSat{M(\Delta)}{!C}[\sigma]$; તેથી
      $\mSat{M(\Delta)}{\indfrm}[\sigma]$.

      $\sFmla{\False}{\indfrm}[\sigma] \in \Delta$ હોય તો
      શાખા પૂર્ણ હોવાથી $\sFmla{\False}{!B}[\sigma] \in \Delta$
      અને $\sFmla{\False}{!C}[\sigma] \in \Delta$ બંને હોય.
      આગમન-ધારણા મુજબ
      $\mSat/{M(\Delta)}{!B}[\sigma]$ અને
      $\mSat/{M(\Delta)}{!C}[\sigma]$; તેથી
      $\mSat/{M(\Delta)}{\indfrm}[\sigma]$.}
  \sourcecorrection{OLMD-069}{મૂળમાં અસત્ય વિયોજનના
    આગમન-નિષ્કર્ષમાં બીજો ભાગ ફરી $!B$ લખાયો છે;
    બંને અસત્ય શાખા-સૂત્રો મુજબ ત્યાં $!C$ રાખ્યું છે.}
  \item \indcase{!A}{!B \lif !C}{
      $\sFmla{\True}{\indfrm}[\sigma] \in \Delta$ હોય તો
      શાખા પૂર્ણ હોવાથી $\sFmla{\False}{!B}[\sigma] \in \Delta$
      અથવા $\sFmla{\True}{!C}[\sigma] \in \Delta$ હોય.
      આગમન-ધારણા મુજબ
      $\mSat/{M(\Delta)}{!B}[\sigma]$ અથવા
      $\mSat{M(\Delta)}{!C}[\sigma]$; તેથી
      $\mSat{M(\Delta)}{\indfrm}[\sigma]$.

      $\sFmla{\False}{\indfrm}[\sigma] \in \Delta$ હોય તો
      શાખા પૂર્ણ હોવાથી $\sFmla{\True}{!B}[\sigma] \in \Delta$
      અને $\sFmla{\False}{!C}[\sigma] \in \Delta$ બંને હોય.
      આગમન-ધારણા મુજબ
      $\mSat{M(\Delta)}{!B}[\sigma]$ અને
      $\mSat/{M(\Delta)}{!C}[\sigma]$; તેથી
      $\mSat/{M(\Delta)}{\indfrm}[\sigma]$.}
  \sourcecorrection{OLMD-070}{મૂળમાં અસત્ય અનુબંધના
    આગમન-નિષ્કર્ષમાં બીજો ભાગ ફરી $!B$ લખાયો છે;
    અનુબંધનો ઉત્તરપક્ષ અસત્ય હોવાથી ત્યાં $!C$ રાખ્યું છે.}
  \item \indcase{!A}{\Box !B}{
      $\sFmla{\True}{\indfrm}[\sigma] \in \Delta$ હોય તો
      શાખા પૂર્ણ હોવાથી તેમાં વપરાયેલા દરેક $\sigma.n$ માટે
      $\sFmla{\True}{!B}[\sigma.n] \in \Delta$; એટલે
      $R\sigma\sigma'$ ધરાવતા દરેક $\sigma' \in P(\Delta)$
      માટે આ વાત સાચી છે. આગમન-ધારણા મુજબ
      $R\sigma\sigma'$ ધરાવતા એવા દરેક $\sigma'$ માટે
      $\mSat{M(\Delta)}{!B}[\sigma']$.
      તેથી $\mSat{M(\Delta)}{\indfrm}[\sigma]$.

      $\sFmla{\False}{\indfrm}[\sigma] \in \Delta$ હોય તો
      શાખા પૂર્ણ હોવાથી કોઈ $\sigma.n$ માટે
      $\sFmla{\False}{!B}[\sigma.n] \in \Delta$.
      આગમન-ધારણા મુજબ $\mSat/{M(\Delta)}{!B}[\sigma.n]$.
      $R\sigma(\sigma.n)$ હોવાથી એવું~$\sigma'$ છે કે
      $\mSat/{M(\Delta)}{!B}[\sigma']$. આમ
      $\mSat/{M(\Delta)}{\indfrm}[\sigma]$.}
  \item \indcase{!A}{\Diamond !B}{
      $\sFmla{\True}{\indfrm}[\sigma] \in \Delta$ હોય તો
      શાખા પૂર્ણ હોવાથી કોઈ $\sigma.n$ માટે
      $\sFmla{\True}{!B}[\sigma.n] \in \Delta$.
      આગમન-ધારણા મુજબ $\mSat{M(\Delta)}{!B}[\sigma.n]$.
      $R\sigma(\sigma.n)$ હોવાથી એવું~$\sigma'$ છે કે
      $\mSat{M(\Delta)}{!B}[\sigma']$. આમ
      $\mSat{M(\Delta)}{\indfrm}[\sigma]$.

      $\sFmla{\False}{\indfrm}[\sigma] \in \Delta$ હોય તો
      શાખા પૂર્ણ હોવાથી તેમાં વપરાયેલા દરેક $\sigma.n$ માટે
      $\sFmla{\False}{!B}[\sigma.n] \in \Delta$; એટલે
      $R\sigma\sigma'$ ધરાવતા દરેક
      $\sigma' \in P(\Delta)$ માટે આ વાત સાચી છે.
      આગમન-ધારણા મુજબ $R\sigma\sigma'$ ધરાવતા એવા દરેક
      $\sigma'$ માટે $\mSat/{M(\Delta)}{!B}[\sigma']$.
      તેથી $\mSat/{M(\Delta)}{\indfrm}[\sigma]$.}
\end{enumerate}
$\Gamma \subseteq \Delta$ હોવાથી $\mSat{M(\Delta)}{\Gamma}$.
\end{proof}

\begin{prob}
\olref[nml][tab][cpl]{thm:tableau-completeness}નું પ્રમાણ પૂર્ણ કરો.
\end{prob}

\begin{cor}
\ollabel{cor:entailment-completeness}
$\Gamma \Entails !A$ હોય તો $\Gamma \Proves !A$.
\end{cor}

\begin{cor}
\ollabel{cor:weak-completeness}
બધાં નિદર્શનોમાં $!A$ સત્ય હોય તો $\Proves !A$.
\end{cor}
% END OLP-0468
\endgroup

\begingroup
\makeatletter\def\ol@sectioncs{section}\makeatother

% BEGIN OLP-0469
\olfileid{nml}{tab}{cou}
      
\olsection{ટેબ્લોમાંથી પ્રતિવિધાનો}
\phantomsection\label{guunit:OLP-0469}


\begin{explain}
  પૂર્ણતા પ્રમેયનું પ્રમાણ $\Entails !A$ હોય તો $\Proves !A$
  એટલું જ નથી બતાવતું; $\Entails/ !A$ હોય ત્યારે~$!A$નું
  પ્રતિવિધાન બનાવવાની રીત પણ આપે છે.
  \sourcecorrection{OLMD-071}{મૂળમાં આ પ્રતિવિધાનની શરત
  $\Entails/ A$ લખાઈ છે; બાકીના પરિચ્છેદ અને પૂર્ણતા-દાવા
  મુજબ $!A$નું અનુસરણ ન થવું જોઈએ, તેથી $!A$ રાખ્યું છે.}
  $\Log{K}$ના કિસ્સામાં આ રીત \emph{નિર્ણય-પ્રક્રિયા} છે.
  ધારો કે $\Entails/ !A$. ત્યારે
  \olref[cpl]{prop:complete-tableau}નું પ્રમાણ પૂર્ણ
  ટેબ્લો બનાવવાની રીત આપે છે. આ એક-સૂત્રના કિસ્સામાં
  રીત હંમેશા સમાપ્ત થાય છે. $\Log{K}$ના વિધાનસંબંધી નિયમો
  ઓછી જટિલતાવાળાં ઉપસર્ગવાળાં સૂત્રો જ ઉમેરે છે;
  એટલે દરેક ચિહ્નિત સૂત્ર $\sFmla{S}{!A}[\sigma]$ પર
  શાખામાં કોઈ વિધાનસંબંધી નિયમ એક વાર લાગુ કરવો પૂરતો છે.
  નવા ઉપસર્ગો માત્ર
  $\TRule{\False}{\Box}$
    અને $\TRule{\True}{\Diamond}$
  નિયમોથી બને છે; તેમને
  પણ એક જ વાર લાગુ કરવાના હોય છે (અને એક જ નવો ઉપસર્ગ મળે છે).
  $\TRule{\True}{\Box}$
    અને $\TRule{\False}{\Diamond}$
  નિયમો કદાચ અનેક વાર
  લાગુ કરવા પડે, પરંતુ દરેક ઉપસર્ગ માટે એક જ વાર; અને નવા
  ઉપસર્ગો સાંત સંખ્યામાં જ બને છે. તેથી સાંત તબક્કાઓ પછી
  રચનામાં બંધ અથવા પૂર્ણ શાખા મળે છે.

  ખુલ્લી પૂર્ણ શાખાવાળું ટેબ્લો બની જાય પછી
  \olref[cpl]{thm:tableau-completeness}નું પ્રમાણ મૂળના
  ઉપસર્ગવાળાં સૂત્રોના ગણને સંતોષતું સ્પષ્ટ નિદર્શન
  આપે છે. આમ $\Gamma \Entails !A$ હોય તો બંધ ટેબ્લો
  મળે છે અને $\Gamma \Proves !A$ થાય છે એટલું જ નહિ;
  બંધ ટેબ્લોને યોગ્ય રીતે શોધતાં જો ``પૂર્ણ''
  ટેબ્લો મળે, તો $\Gamma \Entails/ !A$ છે તે જાણીએ
  છીએ અને પ્રતિવિધાન પણ બનાવી શકીએ છીએ.
\end{explain}

\begin{ex}
  આપણે જાણીએ છીએ કે
  $\Proves/ \Box(p \lor q) \lif (\Box p \lor \Box q)$.
  ટેબ્લોનું નિર્માણ આ રીતે શરૂ થાય છે:
  \begin{oltableau}
    [\pFmla{\False}{\Box(p \lor q) \lif (\Box p \lor \Box q)}{1},
      just = \TAss, checked
      [\pFmla{\True}{\Box(p \lor q)}{1},
        just = {\TRule{\False}{\lif}[1]}, 
        [\pFmla{\False}{\Box p \lor \Box q}{1},
          just = {\TRule{\False}{\lif}[1]}, checked
          [\pFmla{\False}{\Box p}{1},
            just = {\TRule{\False}{\lor}[3]}, checked
            [\pFmla{\False}{\Box q}{1},
              just = {\TRule{\False}{\lor}[3]}, checked
              [\pFmla{\False}{p}{1.1}, 
                just = {\TRule{\False}{\Box}[4]}, checked
                [\pFmla{\False}{q}{1.2}, 
                  just = {\TRule{\False}{\Box}[5]}, checked
                ]
              ]
            ]
          ]
        ]
      ]
    ]
  \end{oltableau}
  આ ટેબ્લો હજી પૂરું થયું નથી. આગળ ટીકચિહ્ન વિનાની
  એકમાત્ર પંક્તિ જોઈએ: પંક્તિ~$2$ પરનું ઉપસર્ગવાળું
  સૂત્ર $\sFmla{\True}{\Box(p \lor q)}[1]$.
  પદ્ધતિસર રચનામાં $\TRule{\True}{\Box}$ નિયમ શાખા પર
  વપરાયેલા દરેક ઉપસર્ગ, એટલે $1.1$ અને $1.2$ બંને માટે
  લાગુ કરવો પડે છે:
  \begin{oltableau}
    [\pFmla{\False}{\Box(p \lor q) \lif (\Box p \lor \Box q)}{1},
      just = \TAss, checked
      [\pFmla{\True}{\Box(p \lor q)}{1},
        just = {\TRule{\False}{\lif}[1]}, 
        [\pFmla{\False}{\Box p \lor \Box q}{1},
          just = {\TRule{\False}{\lif}[1]}, checked
          [\pFmla{\False}{\Box p}{1},
            just = {\TRule{\False}{\lor}[3]}, checked
            [\pFmla{\False}{\Box q}{1},
              just = {\TRule{\False}{\lor}[3]}, checked
              [\pFmla{\False}{p}{1.1}, 
                just = {\TRule{\False}{\Box}[4]}, checked
                [\pFmla{\False}{q}{1.2}, 
                  just = {\TRule{\False}{\Box}[5]}, checked
                  [\pFmla{\True}{p \lor q}{1.1}, 
                    just = {\TRule{\True}{\Box}[2]}
                    [\pFmla{\True}{p \lor q}{1.2}, 
                      just = {\TRule{\True}{\Box}[2]}
                    ]
                  ]
                ]
              ]
            ]
          ]
        ]
      ]
    ]
  \end{oltableau}
  હવે પંક્તિઓ 2, 8 અને 9 પર ટીકચિહ્ન નથી. પરંતુ નવો
  ઉપસર્ગ ઉમેરાયો નથી; તેથી મળતી બધી શાખાઓ પર, તે બંધ ન
  થાય ત્યાં સુધી, પંક્તિઓ~8 અને~9ને
  $\TRule{\True}{\lor}$ લાગુ કરીએ છીએ:
  \begin{oltableau}
    [\pFmla{\False}{\Box(p \lor q) \lif (\Box p \lor \Box q)}{1},
      just = \TAss, checked
      [\pFmla{\True}{\Box(p \lor q)}{1},
        just = {\TRule{\False}{\lif}[1]}, checked
        [\pFmla{\False}{\Box p \lor \Box q}{1},
          just = {\TRule{\False}{\lif}[1]}, checked
          [\pFmla{\False}{\Box p}{1},
            just = {\TRule{\False}{\lor}[3]}, checked
            [\pFmla{\False}{\Box q}{1},
              just = {\TRule{\False}{\lor}[3]}, checked
              [\pFmla{\False}{p}{1.1}, 
                just = {\TRule{\False}{\Box}[4]}, checked
                [\pFmla{\False}{q}{1.2}, 
                  just = {\TRule{\False}{\Box}[5]}, checked
                  [\pFmla{\True}{p \lor q}{1.1}, 
                    just = {\TRule{\True}{\Box}[2]}, checked
                    [\pFmla{\True}{p \lor q}{1.2}, 
                      just = {\TRule{\True}{\Box}[2]}, checked
                      [\pFmla{\True}{p}{1.1},
                        just = {\TRule{\True}{\lor}[8]}, checked, close
                      ]
                      [\pFmla{\True}{q}{1.1},
                        just = {\TRule{\True}{\lor}[8]}, checked
                        [\pFmla{\True}{p}{1.2},
                          just = {\TRule{\True}{\lor}[9]}, checked]
                        [\pFmla{\True}{q}{1.2},
                          just = {\TRule{\True}{\lor}[9]}, checked, close]
                      ]
                    ]
                  ]
                ]
              ]
            ]
          ]
        ]
      ]
    ]
  \end{oltableau}
  એક ખુલ્લી શાખા બાકી રહે છે અને તે પૂર્ણ છે. તેમાંથી
  જગતો $W = \{1, 1.1, 1.2\}$ (ખુલ્લી શાખા પરના ઉપસર્ગો),
  સુલભતા-સંબંધ $R = \{\tuple{1, 1.1}, \tuple{1, 1.2}\}$,
  અને નિમણૂક $V(p) = \{1.2\}$ (કારણ કે પંક્તિ~11 પર
  $\sFmla{\True}{p}[1.2]$ છે) તથા $V(q) = \{1.1\}$
  (કારણ કે પંક્તિ~10 પર $\sFmla{\True}{q}[1.1]$ છે)
  ધરાવતું નિદર્શન વ્યાખ્યાયિત કરીએ છીએ. તે
  \olref{fig:counter-Box}માં દર્શાવ્યું છે; તે
  $\Box(p \lor q) \lif (\Box p \lor \Box q)$નું પ્રતિવિધાન
  છે તે ચકાસી શકો છો.
  \begin{figure}
  \begin{center}
    \begin{tikzpicture}[modal]
      \node[world] (w1) [label={[align=right]right:\mFalse{p}\\ \mFalse{q}}]{$1$}; 
      \node[world] (w2) [label={[align=right]right:\mFalse{p}\\ \mTrue{q}},
        above left=of w1]{$1.1$}; 
      \node[world] (w3) [label={[align=right]right:\mTrue{p}\\ \mFalse{q}},
        above right=of w1] {$1.2$};
      \draw[->] (w1) to (w2);
      \draw[->] (w1) to (w3);
    \end{tikzpicture}
  \end{center}
  \caption{$\Box(p \lor q) \lif (\Box p \lor \Box q)$નું પ્રતિવિધાન.}
  \ollabel{fig:counter-Box}
\end{figure}
\end{ex}          
% END OLP-0469
\endgroup


\OLEndChapterHook
% END OLP-0460
\endgroup


\begingroup
\makeatletter\def\ol@sectioncs{section}\makeatother

% BEGIN OLP-0470
\olchapter{nml}{seq}{મોડલ સિક્વન્ટ કલન}
\olchapterid{nml}{seq}
\phantomsection\label{guunit:OLP-0470}

\begin{editorial}
  મોડલ તર્ક માટેનાં સિક્વન્ટ કલનોનું આ મુસદ્દા-પ્રકરણ છે.
  તેમાં વધુ ઉદાહરણો તથા યથાર્થતા અને પૂર્ણતાનાં પ્રમાણો
  જરૂરી છે.
\end{editorial}

\begingroup
\makeatletter\def\ol@sectioncs{section}\makeatother

% BEGIN OLP-0471
\olfileid{nml}{seq}{int}

\olsection{પરિચય}
\phantomsection\label{guunit:OLP-0471}


વિધાનતર્કના સિક્વન્ટ કલનમાં
$\Box$ અને~%
$\Diamond$ માટેના વધારાના નિયમો
ઉમેરી શકાય છે. દાખલા તરીકે,~$\Log{K}$ માટે
\Log{LK} ઉપરાંત આ નિયમો છે:
\[
  % <> and [] primitive
    \Axiom$\Gamma \fCenter \Delta, !A$
    \RightLabel{$\Box$}
    \UnaryInf$\Box\Gamma \fCenter \Diamond\Delta, \Box!A$
    \DisplayProof  
    \qquad
    \Axiom$!A, \Gamma \fCenter \Delta$
    \RightLabel{$\Diamond$}
    \UnaryInf$\Diamond!A, \Box\Gamma \fCenter \Diamond\Delta$
    \DisplayProof
\]
$\Log{K}$ના વિસ્તરણો માટે બીજા નિયમો પણ ઉમેરવા પડે છે.

દરેક મોડલ તર્ક માટે આવું સિક્વન્ટ કલન જાણીતું નથી.
$\Log{S5}$ અર્થવિજ્ઞાનની દૃષ્ટિએ સરળ છે (તેને સુલભતા
સંબંધ વાપર્યા વિના પણ વ્યાખ્યાયિત કરી શકાય છે), છતાં
$\Log{LK}$માંથી મળતું અને~\Cut નિયમ વિના પૂર્ણ હોય એવું
સિક્વન્ટ કલન તેના માટે જાણીતું નથી. જોકે તેના માટે કટ વિના
પૂર્ણ \emph{હાઇપર-સિક્વન્ટ} કલન છે.
% END OLP-0471
\endgroup

\begingroup
\makeatletter\def\ol@sectioncs{section}\makeatother

% BEGIN OLP-0472
\olfileid{nml}{seq}{rul}

\olsection{\Log{K} માટેના નિયમો}
\phantomsection\label{guunit:OLP-0472}


સામાન્ય વિધાનસંબંધી સંયોજકો માટેના નિયમો સામાન્ય સિક્વન્ટ
કલન~$\Log{LK}$ના નિયમો જેવા જ છે. સ્વયંસિદ્ધિઓ પણ એ જ છે:
$!A \Sequent !A$ સ્વરૂપનું કોઈપણ સિક્વન્ટ સ્વયંસિદ્ધિ ગણાય છે.

મોડલ સંક્રિયકો~$\Box$
અને~$\Diamond$ માટે નીચેના
વધારાના નિયમો છે:
\[
  % <> and [] primitive
    \Axiom$\Gamma \fCenter \Delta, !A$
    \RightLabel{$\Box$}
    \UnaryInf$\Box\Gamma \fCenter \Diamond\Delta, \Box!A$
    \DisplayProof  
    \qquad
    \Axiom$!A, \Gamma \fCenter \Delta$
    \RightLabel{$\Diamond$}
    \UnaryInf$\Diamond!A, \Box\Gamma \fCenter \Diamond\Delta$
    \DisplayProof
\]
અહીં $\Box\Gamma$ એટલે~$\Gamma$ના દરેક
સૂત્ર પહેલાં $\Box$ મૂકવાથી~$\Gamma$માંથી મળતી
સૂત્રોની
શ્રેણી અને
$\Diamond\Delta$ એટલે~$\Delta$ના દરેક
સૂત્ર પહેલાં $\Diamond$ મૂકવાથી~$\Delta$માંથી મળતી
સૂત્રોની
શ્રેણી. %
%
$\Gamma$~અને~$\Delta$
ખાલી હોઈ શકે; ત્યારે નિષ્કર્ષ સિક્વન્ટનો અનુરૂપ ભાગ
$\Box\Gamma$~અને~$\Diamond\Delta$
પણ ખાલી હોય છે.

જમણી બાજુ $\Box$ અને
ડાબી બાજુ $\Diamond$ માત્ર એક જ
સૂત્ર~$!A$ પર મૂકવાની મર્યાદા જરૂરી છે. જો
જમણી બાજુ ગમે તેટલાં સૂત્રો પર $\Box$
મૂકવાની અથવા ડાબી
બાજુ ગમે તેટલાં સૂત્રો પર $\Diamond$ મૂકવાની
છૂટ આપીએ, તો નીચેનું નિષ્પન્ન કરી શકીએ:
  \[
          \Axiom$!A \fCenter !A$
      \RightLabel{\RightR{\lnot}}
      \UnaryInf$\fCenter !A, \lnot !A$
      \RightLabel{$\Box*$}
      \UnaryInf$\fCenter \Box!A, \Box\lnot !A$
      \doubleLine
      \RightLabel{\RightR{\lor}}
      \UnaryInf$\fCenter \Box!A \lor \Box \lnot !A$
      \DisplayProof
    \qquad
          \Axiom$!A \fCenter !A$
      \RightLabel{\LeftR{\lnot}}
      \UnaryInf$\lnot !A, !A \fCenter$
      \RightLabel{$\Diamond*$}
      \UnaryInf$\Diamond\lnot !A,\Diamond!A \fCenter $
      \RightLabel{\RightR{\lnot}}
      \UnaryInf$\Diamond!A \fCenter \lnot\Diamond\lnot !A$
      \RightLabel{\RightR{\lif}}
      \UnaryInf$ \fCenter \Diamond!A \lif \lnot\Diamond\lnot !A$
      \DisplayProof
  \]
પરંતુ $\Box!A \lor \Box \lnot !A$ અને
$\Diamond!A \lif \lnot\Diamond\lnot !A$ બંને~$\Log{K}$માં માન્ય નથી.

જો પૂર્વપક્ષમાં~$!A$ ઉપરાંત બાજુનાં સૂત્રો રાખીએ અને
$\Box$ નિયમને જમણી બાજુ માત્ર~$!A$ પર
$\Box$ મૂકવા દઈએ અથવા $\Diamond$
નિયમને ડાબી બાજુ માત્ર~$!A$ પર $\Diamond$ મૂકવા દઈએ
(પણ બાજુનાં સૂત્રોને યથાવત્ રાખીએ), તો નીચેનું
નિષ્પન્ન કરી શકીએ:
\[
      \Axiom$!A \fCenter !A$
    \RightLabel{\RightR{\lnot}}
    \UnaryInf$\fCenter !A, \lnot !A$
    \RightLabel{\RightR{\Exchange}}
    \UnaryInf$\fCenter \lnot !A, !A$
    \RightLabel{$\Box*$}
    \UnaryInf$\fCenter \lnot!A, \Box !A$
    \doubleLine
    \RightLabel{\RightR{\lor}}
    \UnaryInf$\fCenter \lnot!A \lor \Box !A$
    \DisplayProof
  \qquad
      \Axiom$!A \fCenter !A$
    \RightLabel{\LeftR{\lnot}}
    \UnaryInf$\lnot !A, !A \fCenter$
    \RightLabel{$\Diamond*$}
    \UnaryInf$\Diamond\lnot !A, !A \fCenter $
    \RightLabel{\RightR{\lnot}}
    \UnaryInf$!A \fCenter \lnot\Diamond\lnot !A$
    \RightLabel{\RightR{\lif}}
    \UnaryInf$ \fCenter !A \lif \lnot\Diamond\lnot !A$
    \DisplayProof
\]
પરંતુ $\lnot!A \lor \Box !A$ (જે $!A
\lif \Box!A$ને સમકક્ષ છે) અને
$!A \lif \lnot\Diamond\lnot !A$ બંને~$\Log{K}$માં માન્ય નથી.
% END OLP-0472
\endgroup

\begingroup
\makeatletter\def\ol@sectioncs{section}\makeatother

% BEGIN OLP-0473
\olfileid{nml}{seq}{prk}

\olsection{\Log{K} માટે સિક્વન્ટ નિષ્પત્તિઓ}
\phantomsection\label{guunit:OLP-0473}



\begin{ex}
  $\Proves (\Box!A \land \Box!B) \lif \Box (!A \land !B)$
  બતાવતું સિક્વન્ટ કલનનું નિષ્પત્તિ આપીએ છીએ.
  \begin{prooftree}
    \Axiom$!A \fCenter !A$
    \doubleLine
    \UnaryInf$!B, !A \fCenter !A$   
    \Axiom$!B \fCenter !B$
    \doubleLine
    \UnaryInf$!B, !A \fCenter !B$   
    \RightLabel{\RightR{\land}}
    \BinaryInf$!B, !A \fCenter !A \land !B$    
    \RightLabel{$\Box$}
    \UnaryInf$\Box!B, \Box!A \fCenter \Box
    (!A \land !B)$
    \RightLabel{\LeftR{\land}}
    \UnaryInf$\Box!A \land \Box!B, \Box!A \fCenter \Box
    (!A \land !B)$
    \RightLabel{\LeftR{\Exchange}}
    \UnaryInf$\Box!A, \Box!A \land \Box!B \fCenter \Box
    (!A \land !B)$
    \RightLabel{\LeftR{\land}}
    \UnaryInf$\Box!A \land \Box!B, \Box!A \land \Box!B \fCenter \Box (!A \land !B)$
    \RightLabel{\LeftR{\Contraction}}
    \UnaryInf$\Box!A \land \Box!B \fCenter \Box (!A \land !B)$
    \RightLabel{\RightR{\lif}}
    \UnaryInf$\fCenter (\Box!A \land \Box!B)
    \lif \Box (!A \land !B)$
  \end{prooftree}
\end{ex}

  \begin{ex}
    $\Proves \Diamond(!A \lor !B) \lif
    (\Diamond !A \lor \Diamond!B)$ બતાવતું સિક્વન્ટ કલનનું
    નિષ્પત્તિ આપીએ છીએ.
    \begin{prooftree}
      \Axiom$!A \fCenter !A$
      \doubleLine
      \UnaryInf$!A \fCenter !A, !B$   
      \Axiom$!B \fCenter !B$
      \doubleLine
      \UnaryInf$!B \fCenter !A, !B$   
      \RightLabel{\LeftR{\lor}}
      \BinaryInf$!A \lor !B \fCenter !A, !B$    
      \RightLabel{$\Diamond$}
      \UnaryInf$\Diamond(!A \lor !B) \fCenter 
      \Diamond !A, \Diamond!B$
      \RightLabel{\RightR{\lor}}
      \UnaryInf$\Diamond(!A \lor !B) \fCenter 
      \Diamond !A, \Diamond !A \lor \Diamond!B$
      \RightLabel{\RightR{\Exchange}}
      \UnaryInf$\Diamond(!A \lor !B) \fCenter \Diamond !A \lor \Diamond!B, \Diamond !A$
      \RightLabel{\RightR{\lor}}
      \UnaryInf$\Diamond(!A \lor !B) \fCenter 
      \Diamond !A \lor \Diamond!B, \Diamond !A \lor \Diamond!B$
      \RightLabel{\RightR{\Contraction}}
      \UnaryInf$\Diamond(!A
      \lor !B) \fCenter \Diamond !A \lor \Diamond!B$
      \RightLabel{\RightR{\lif}}
      \UnaryInf$\fCenter \Diamond(!A
      \lor !B) \lif (\Diamond !A \lor \Diamond!B)$
    \end{prooftree}
  \end{ex}

  અહીં~\Dual નું નિષ્પત્તિ છે.
  \sourcecorrection{OLMD-075}{નીચેના વૃક્ષની પ્રથમ
  $!A \fCenter !A$ પંક્તિમાંથી $\lnot !A,!A \fCenter$
  મેળવવા ડાબી નિષેધનો નિયમ જોઈએ; મૂળમાં તેને જમણી
  નિષેધનો નિયમ કહ્યો છે. નિયમનું લેબલ સુધાર્યું છે.}
  \begin{prooftree}
    \Axiom$!A \fCenter !A$
    \RightLabel{\LeftR{\lnot}}
    \UnaryInf$\lnot !A, !A \fCenter $
    \RightLabel{$\Diamond$}
    \UnaryInf$\Diamond\lnot !A, \Box !A \fCenter $
    \RightLabel{\RightR{\lnot}}
    \UnaryInf$\Box !A \fCenter \lnot\Diamond\lnot !A$
    \RightLabel{\RightR{\lif}}
    \UnaryInf$\fCenter \Box !A \lif \lnot\Diamond\lnot !A$
    \Axiom$!A \fCenter !A$
    \RightLabel{\RightR{\lnot}}
    \UnaryInf$\fCenter !A, \lnot !A $
    \RightLabel{\RightR{\Exchange}}
    \UnaryInf$\fCenter \lnot !A, !A $
    \RightLabel{$\Box$}
    \UnaryInf$\fCenter \Diamond\lnot !A, \Box !A$
    \RightLabel{\RightR{\Exchange}}
    \UnaryInf$\fCenter \Box !A, \Diamond\lnot !A$
    \RightLabel{\RightR{\lnot}}
    \UnaryInf$\lnot\Diamond\lnot !A \fCenter \Box !A$
    \RightLabel{\RightR{\lif}}
    \UnaryInf$\fCenter \lnot \Diamond \lnot !A \lif \Box !A$
    \RightLabel{\RightR{\land}}
    \BinaryInf$\fCenter \Box !A \liff \lnot\Diamond\lnot !A$ 
  \end{prooftree} 


\begin{prob}
  $\Log{K}$માં નીચેના સૂત્રો માટે સિક્વન્ટ કલનનાં
  પ્રમાણો શોધો:
  \begin{enumerate}
    \item $\Box \lnot p \lif \Box(p \lif q)$
    \item $(\Box p \lor \Box q) \lif \Box(p \lor q)$
    \item $\Diamond p \lif \Diamond(p \lor q)$
    \item $\Box(p \land q) \lif \Box p$
  \end{enumerate}
\end{prob}
% END OLP-0473
\endgroup

\begingroup
\makeatletter\def\ol@sectioncs{section}\makeatother

% BEGIN OLP-0474
\olfileid{nml}{seq}{mru}

\olsection{અન્ય સુલભતા-સંબંધો માટેના નિયમો}
\phantomsection\label{guunit:OLP-0474}


વિશિષ્ટ સુલભતા-સંબંધોથી નિર્ધારિત તર્કપદ્ધતિઓ માટે
\olref{tab:more-rules}ના વધારાના નિયમો લઈએ છીએ.

% [] and <> primitive
\begin{table}
  \begin{center}
    \def\arraystretch{3}
    \begin{tabular}{cc}
    \hline
      \Axiom$!A, \Gamma \fCenter \Delta$
      \RightLabel{T$\Box$}
      \UnaryInf$\Box!A, \Gamma \fCenter \Delta$
      \DisplayProof
    &
      \Axiom$\Gamma \fCenter \Delta, !A$
      \RightLabel{T$\Diamond$}
      \UnaryInf$\Gamma \fCenter \Delta, \Diamond!A$
      \DisplayProof
    \\[1ex]
    \multicolumn{2}{c}{
      \Axiom$\Gamma \fCenter \Delta$
      \RightLabel{D}
      \UnaryInf$\Box\Gamma \fCenter \Diamond\Delta$
      \DisplayProof}
    \\[1ex]
      \Axiom$\Gamma, \Diamond\Pi \fCenter \Box\Delta, \Lambda, !A$
      \RightLabel{B$\Box$}
      \UnaryInf$\Box\Gamma, \Pi \fCenter \Delta, \Diamond\Lambda, \Box!A$
      \DisplayProof
      &
      \Axiom$!A, \Diamond\Gamma, \Pi \fCenter \Box\Lambda, \Delta$
      \RightLabel{B$\Diamond$}
      \UnaryInf$\Diamond!A, \Gamma, \Box\Pi \fCenter \Lambda, \Diamond\Delta$
      \DisplayProof
    \\[1ex]
      \Axiom$\Box\Gamma \fCenter \Diamond\Delta, !A$
      \RightLabel{4$\Box$}
      \UnaryInf$\Box\Gamma \fCenter \Diamond\Delta, \Box!A$
      \DisplayProof
      &
      \Axiom$!A, \Box\Gamma \fCenter \Diamond\Delta$
      \RightLabel{4$\Diamond$}
      \UnaryInf$\Diamond!A, \Box\Gamma \fCenter \Diamond\Delta$
      \DisplayProof
    \\[1ex]
      \Axiom$\Box\Gamma, \Diamond\Pi \fCenter \Box\Delta,
      \Diamond\Lambda, !A$
      \RightLabel{5$\Box$}
      \UnaryInf$\Box\Gamma, \Diamond\Pi \fCenter \Box\Delta,
      \Diamond\Lambda, \Box!A$
      \DisplayProof
&
      \Axiom$!A, \Diamond\Gamma, \Box\Pi \fCenter \Diamond\Delta, \Box\Lambda$
      \RightLabel{5$\Diamond$}
      \UnaryInf$\Diamond!A,\Diamond\Gamma,\Box\Pi \fCenter \Diamond\Delta,\Box\Lambda$
      \DisplayProof
    \\[1ex]
    \hline
    \end{tabular}
  \end{center}
  \caption{વધારાના મોડલ નિયમો.}
  \ollabel{tab:more-rules}
\end{table}

આ નિયમો ઉમેરવાથી \olref{tab:logics-rules}માં આપેલી
તર્કપદ્ધતિઓ માટે યથાર્થ અને પૂર્ણ એવી પ્રણાલીઓ મળે છે.

\begin{table}
  \begin{center}
    \begin{tabular}{lll}
      \hline
      તર્કપદ્ધતિ & $R$નો ગુણધર્મ & નિયમો\\
      \hline
      $\Log{T} = \Log{KT}$ & સ્વવાચક & $\Box$,
      T$\Box$%
      , %
      T$\Diamond$
      \\ \hline
      $\Log{D} = \Log{KD}$ & સિરિયલ & $\Box$,
      %
        D
      \\ \hline
      $\Log{K4}$ & પરંપરિત & $\Box$,
      4$\Box$% 
      , %
      4$\Diamond$
      \\ \hline
      $\Log{B} = \Log{KTB}$ & સ્વવાચક, & $\Box$,
      T$\Box$%
      , %
      T$\Diamond$\\
      & સંમિત &
      B$\Box$%
      , %
      B$\Diamond$
      \\ \hline
      $\Log{S4} = \Log{KT4}$ & સ્વવાચક, & $\Box$,
      T$\Box$%
      , %
      T$\Diamond$\\
      & પરંપરિત &
      4$\Box$%
      , %
      4$\Diamond$
      \\ \hline
      $\Log{S5} = \Log{KT5}$ & સ્વવાચક, & $\Box$,
      T$\Box$%
      , %
      T$\Diamond$\\
      & પરંપરિત, &
      5$\Box$%
      , %
      5$\Diamond$\\
      &  યુક્લિડિયન &
      \\ \hline
    \end{tabular}
  \end{center}
  \caption{વિવિધ મોડલ તર્કપદ્ધતિઓ માટેના સિક્વન્ટ નિયમો.}
  \ollabel{tab:logics-rules}
\end{table}

\begin{ex}
  $\Log{K4} \Proves \Ax{4}$, એટલે કે,
  $\Box!A \lif \Box\Box!A$ બતાવતું સિક્વન્ટ
  નિષ્પત્તિ આપીએ છીએ.
  \begin{prooftree}
    \Axiom$\Box!A \fCenter \Box!A$
    \RightLabel{4$\Box$}
    \UnaryInf$\Box!A \fCenter \Box\Box!A$
    \RightLabel{\RightR{\lif}}
    \UnaryInf$\fCenter \Box!A \lif \Box\Box!A$
  \end{prooftree}
\end{ex}

% <> and [] primitive
\begin{ex}
  $\Log{S5} \Proves \Ax{5}$, એટલે કે,
  $\Diamond!A \lif \Box\Diamond!A$ બતાવતું સિક્વન્ટ
  નિષ્પત્તિ આપીએ છીએ.
  \begin{prooftree}
    \Axiom$\Diamond!A \fCenter \Diamond!A$
    \RightLabel{5$\Box$}
    \UnaryInf$\Diamond!A \fCenter \Box\Diamond!A$
    \RightLabel{\RightR{\lif}}
    \UnaryInf$\fCenter \Diamond!A \lif \Box\Diamond!A$
  \end{prooftree}
\end{ex}
\begin{ex}
\Log{S5} માટેનું સિક્વન્ટ કલન \Cut{} નિયમ વિના પૂર્ણ નથી;
દાખલા તરીકે, $\Diamond\Box !A \lif !A$ એ~$\Log{S5}$માં માન્ય છે, પણ તેનું
\Cut વિના પ્રમાણ નથી. અહીં~\Cut વાપરતું નિષ્પત્તિ છે:
\begin{prooftree}
% both [] and <> primitive
\Axiom$\Box !A \fCenter \Box!A$
\RightLabel{5$\Diamond$}
\UnaryInf$\Diamond\Box!A \fCenter \Box!A$

\Axiom$!A \fCenter !A$
\RightLabel{T$\Box$}
\UnaryInf$\Box!A \fCenter !A$
\RightLabel{\Cut}
\BinaryInf$\Diamond\Box!A \fCenter !A$
\RightLabel{\RightR{\lif}}
\UnaryInf$\fCenter \Diamond\Box!A \lif !A$
\end{prooftree}
\end{ex}

\begin{prob}
નીચેના દાવા બતાવતાં સિક્વન્ટ નિષ્પત્તિઓ આપો:
  \begin{enumerate}
  \item $\Log{KT5} \Proves \Ax{B}$;
  \item $\Log{KT5} \Proves \Ax{4}$;
  \item $\Log{KDB4} \Proves \Ax{T}$;
  \item $\Log{KB4} \Proves \Ax{5}$;
  \item $\Log{KB5} \Proves \Ax{4}$;
  \item $\Log{KT} \Proves \Ax{D}$.
  \end{enumerate}
\end{prob}
% END OLP-0474
\endgroup


\OLEndChapterHook
% END OLP-0470
\endgroup


\OLEndPartHook
% END OLP-0407
